% !TeX program = xelatex
% 使用 XeLaTeX 编译本文件两次，以生成并更新目录。
% 仅需本文件；采用通用 TeX Live 宏包及 Noto CJK 或 Fandol 中文字体。
\documentclass[11pt,a4paper]{article}
\usepackage[left=21mm,right=21mm,top=20mm,bottom=21mm,headheight=15pt,headsep=6mm,footskip=10mm]{geometry}
\usepackage{fontspec,xcolor,fancyhdr,needspace,amsmath,amssymb}
\usepackage[unicode,colorlinks=true,linkcolor=ink,bookmarksdepth=3,bookmarksnumbered]{hyperref}
\usepackage{bookmark}
\IfFontExistsTF{Noto Serif CJK SC}{
  \setmainfont{Noto Serif CJK SC}
  \setsansfont{Noto Sans CJK SC}
}{
  \setmainfont{FandolSong-Regular.otf}[BoldFont=FandolSong-Bold.otf]
  \setsansfont{FandolHei-Regular.otf}[BoldFont=FandolHei-Bold.otf]
}
\XeTeXlinebreaklocale "zh"
\XeTeXlinebreakskip=0pt plus 1pt
\definecolor{ink}{HTML}{172F3A}
\definecolor{accent}{HTML}{176A65}
\definecolor{muted}{HTML}{415D64}
\definecolor{linecol}{HTML}{B6C9CA}
\hypersetup{pdftitle={数据结构算法手册},pdfsubject={算法C 风格伪代码与复杂度}}
\setcounter{tocdepth}{3}
\setcounter{secnumdepth}{3}
\setlength{\parindent}{0pt}
\renewcommand{\contentsname}{目录}
\newcommand{\running}{目录}
\pagestyle{fancy}\fancyhf{}
\fancyhead[L]{\footnotesize\sffamily\color{muted}数据结构算法手册}
\fancyhead[R]{\footnotesize\sffamily\color{muted}\running}
\fancyfoot[C]{\small\sffamily\thepage}
\renewcommand{\headrulewidth}{0.3pt}
\raggedbottom
\setlength{\emergencystretch}{2em}
\makeatletter
\renewcommand{\@pnumwidth}{2.25em}\renewcommand{\@tocrmarg}{3em}
\renewcommand*\l@subsubsection{\@dottedtocline{3}{4em}{4.6em}}
\renewcommand*\l@subsection{\@dottedtocline{2}{1.6em}{3.1em}}
\makeatother
\newcommand{\Part}[1]{\clearpage\renewcommand{\running}{#1}\section{#1}}
\newcommand{\Group}[1]{\Needspace{10\baselineskip}\subsection{#1}}
\newenvironment{entry}[1]{\par\addvspace{6pt}\noindent\begin{minipage}{\linewidth}
  \subsubsection{#1}\vspace{1pt}\sffamily\fontsize{10.3}{15.9}\selectfont
}{\end{minipage}\par\vspace{4pt}}
\newcommand{\Step}[2]{\par\begingroup\advance\leftskip by #1em\relax\advance\leftskip by 0.7em\relax\hangindent=0.8em\hangafter=1\noindent #2\par\endgroup}
\newcommand{\Note}[1]{{\color{muted}\quad // #1}}
\newcommand{\Cost}[2]{\par\vspace{4pt}{\color{accent}\small 时间\enspace #1\qquad 辅助空间\enspace #2}\par\vspace{3pt}{\color{linecol}\hrule}}

\IfFontExistsTF{DejaVu Sans Mono}{\setmonofont{DejaVu Sans Mono}[Scale=0.91]}{}
\newcommand{\Code}[2]{\par\begingroup\leftskip=#1em\relax\advance\leftskip by .35em\hangindent=1.2em\hangafter=1\rightskip=0pt plus 1fil\ttfamily\fontsize{9.8}{13.7}\selectfont\noindent #2\strut\par\endgroup}
\newcommand{\Hint}[2]{\par\begingroup\leftskip=#1em\relax\advance\leftskip by .35em\rightskip=0pt plus 1fil\sffamily\fontsize{9.4}{13.8}\selectfont\color{accent}\noindent // #2\strut\par\endgroup}
\begin{document}\color{ink}
\begin{titlepage}\thispagestyle{empty}
\vspace*{0.27\textheight}
{\sffamily\bfseries\fontsize{30}{43}\selectfont 数据结构算法手册\par}
\vspace{9pt}{\color{accent}\rule{43mm}{1.4pt}}\par
\vspace{9pt}{\large C 风格伪代码与复杂度}\par
\end{titlepage}
\pagenumbering{roman}
\pdfbookmark[0]{目录}{toc}\tableofcontents\bookmarksetup{startatroot}
\clearpage\pagenumbering{arabic}
\renewcommand{\running}{阅读约定}
\section*{阅读约定}
\textbf{怎样读这本手册}\quad 每个算法按函数入口、执行步骤、时间与辅助空间排列。绿色的双斜线注释解释代码中的变量、边界或操作顺序。多种实现列在同一个条目下。\par\medskip
\textbf{指针}\quad \texttt{p->next} 是结点 p 保存后继地址的域。写 \texttt{p = p->next} 会移动游标；写 \texttt{p->next = q} 会改动链。\texttt{\&p} 是指针变量自身的地址，将它传给函数，就能让函数修改外部的 p。H 为哨兵，H 的后继才是首元。\par\medskip
\textbf{共同约定}\quad 数组下标从 0 起，线性表的位序从 1 起。区间 $[l,r)$ 包含 l、排除 r。整数相除取整数商；\texttt{require} 表示调用须满足的前提。\texttt{NULL} 表示空指针，\texttt{NONE} 表示没有有效编号，\texttt{INF} 表示正无穷。逻辑运算采用 C 的短路求值。\par\medskip
\textbf{伪码简写}\quad 标量的类型可由上下文推断。数组和容器形参默认可原地修改。\texttt{for (x in A)} 顺次枚举 A 中的项，\texttt{repeat(k)} 重复 k 次。元组表示同时取回多个结果。\texttt{NewNode} 等分配函数初始化空指针域；一般条目假设内存足够。\par\medskip
\textbf{栈与队列}\quad \texttt{Push(\&S,x)} 入栈。\texttt{x = Pop(\&S)} 是弹出并取值的伪码简写，也可按输出参数写成 \texttt{Pop(\&S,\&x)}。\texttt{EnQueue} 在队尾加入，\texttt{DeQueue} 在队头取出。局部工作栈、队列、路径和输出缓冲缺省为空；访问标记缺省为假。\par\medskip
\textbf{树与图}\quad 根为第 1 层，空树高度为 0，路径长度按边数计。图默认用邻接表，V、E 分别表示顶点数与边数。无向边在邻接表的两条记录共享边编号。\texttt{visit} 表示处理一个元素，\texttt{Emit} 表示立即输出或复制一份当前结果。\par\medskip
\textbf{复杂度}\quad 默认最坏情形，均摊与期望值另标。辅助空间包含递归栈，排除输入和必要输出。n、m 表示输入长度，h 为树高，w 为最大层宽，L 为输入编码或输出总长度。集合与散列表默认期望常数操作；算术、比较和基本结点操作按单位代价计。

\Part{顺序表}

\Group{基本操作}

\Needspace{16\baselineskip}
\subsubsection{顺序表的初始化与创建}
\Code{0.0}{CreateList(input, n) \{}
\Hint{1.3}{空表也预留一个槽，第一次插入即可写入。}\nopagebreak[4]
\Code{1.3}{L.cap = max(1, n);}
\Code{1.3}{L.data = AllocElems(L.cap);}
\Code{1.3}{for (int i = 0; i<n; i++) \{}
\Code{2.6}{L.data[i] = input[i];}
\Code{1.3}{\}}
\Hint{1.3}{cap 记录能装多少，n 记录已经写入多少。空槽不参与后续查找。}\nopagebreak[4]
\Code{1.3}{L.n = n;}
\Code{1.3}{return L;}
\Code{0.0}{\}}
\nopagebreak[4]
\Cost{\ensuremath{O\!\left(n\right)}}{\ensuremath{O\!\left(1\right)}}
\par\medskip

\Needspace{11\baselineskip}
\subsubsection{按位查找}
\Code{0.0}{GetAt(L, i) \{}
\Hint{1.3}{位序从 1 开始，数组下标从 0 开始。}\nopagebreak[4]
\Code{1.3}{require(1 <= i \&\& i <= L.n);}
\Hint{1.3}{i 等于 n 时读取最后一项，下标是 n−1。}\nopagebreak[4]
\Code{1.3}{return L.data[i-1];}
\Code{0.0}{\}}
\nopagebreak[4]
\Cost{\ensuremath{O\!\left(1\right)}}{\ensuremath{O\!\left(1\right)}}
\par\medskip

\Needspace{13\baselineskip}
\subsubsection{按值查找}
\Code{0.0}{Locate(L, x) \{}
\Hint{1.3}{首次命中就返回位序，0 表示没有匹配。}\nopagebreak[4]
\Code{1.3}{for (int i = 0; i<L.n; i++) \{}
\Code{2.6}{if (L.data[i] == x) return i+1;}
\Code{1.3}{\}}
\Hint{1.3}{循环走完还没有返回，才可以断定 x 不在表内。}\nopagebreak[4]
\Code{1.3}{return 0;}
\Code{0.0}{\}}
\nopagebreak[4]
\Cost{\ensuremath{O\!\left(n\right)}}{\ensuremath{O\!\left(1\right)}}
\par\medskip

\Needspace{16\baselineskip}
\subsubsection{指定位置插入元素}
\Code{0.0}{InsertAt(L, i, x) \{}
\Hint{1.3}{从末端搬起，避免覆盖还没搬走的值。}\nopagebreak[4]
\Code{1.3}{require(1 <= i \&\& i <= L.n+1);}
\Code{1.3}{if (L.n == L.cap) Grow(L);}
\Code{1.3}{for (int j = L.n; j >= i; j--) \{}
\Code{2.6}{L.data[j] = L.data[j-1];}
\Code{1.3}{\}}
\Hint{1.3}{当 i=n+1 时移动循环执行零次，直接写入尾部。}\nopagebreak[4]
\Code{1.3}{L.data[i-1] = x;}
\Code{1.3}{L.n++;}
\Code{0.0}{\}}
\nopagebreak[4]
\Cost{\ensuremath{O\!\left(n\right)}}{\ensuremath{O\!\left(1\right)}，扩容时 \ensuremath{O\!\left(n\right)}}
\par\medskip

\Needspace{16\baselineskip}
\subsubsection{指定位置删除元素}
\Code{0.0}{EraseAt(L, i) \{}
\Hint{1.3}{用后继填补空位，最后缩短有效区。}\nopagebreak[4]
\Code{1.3}{require(1 <= i \&\& i <= L.n);}
\Code{1.3}{x = L.data[i-1];}
\Code{1.3}{for (int j = i; j<L.n; j++) \{}
\Code{2.6}{L.data[j-1] = L.data[j];}
\Code{1.3}{\}}
\Hint{1.3}{被删值先存入 x，左移后原槽里的值就变了。}\nopagebreak[4]
\Code{1.3}{L.n--;}
\Code{1.3}{return x;}
\Code{0.0}{\}}
\nopagebreak[4]
\Cost{\ensuremath{O\!\left(n\right)}}{\ensuremath{O\!\left(1\right)}}
\par\medskip

\Needspace{12\baselineskip}
\subsubsection{顺序表的遍历}
\Code{0.0}{VisitList(L) \{}
\Hint{1.3}{只读已经写入的前 n 个槽。}\nopagebreak[4]
\Code{1.3}{for (int i = 0; i<L.n; i++) \{}
\Hint{2.6}{循环上界是 i<n；下标 n 已经位于有效区之外。}\nopagebreak[4]
\Code{2.6}{visit(L.data[i]);}
\Code{1.3}{\}}
\Code{0.0}{\}}
\nopagebreak[4]
\Cost{\ensuremath{O\!\left(n\right)}}{\ensuremath{O\!\left(1\right)}}
\par\medskip

\Needspace{17\baselineskip}
\subsubsection{顺序表的动态扩容}
\Code{0.0}{Grow(L) \{}
\Hint{1.3}{拷贝完成后才释放旧数组。}\nopagebreak[4]
\Code{1.3}{c = max(1, 2*L.cap);}
\Code{1.3}{B = AllocElems(c);}
\Code{1.3}{for (int i = 0; i<L.n; i++) \{}
\Code{2.6}{B[i] = L.data[i];}
\Code{1.3}{\}}
\Hint{1.3}{cap 翻倍后有效长度 n 不变，新数组前 n 项仍是原顺序。}\nopagebreak[4]
\Code{1.3}{free(L.data);}
\Code{1.3}{L.data = B;}
\Code{1.3}{L.cap = c;}
\Code{0.0}{\}}
\nopagebreak[4]
\Cost{\ensuremath{O\!\left(n\right)}}{\ensuremath{O\!\left(n\right)}}
\par\medskip

\Group{综合算法}

\Needspace{16\baselineskip}
\subsubsection{顺序表的原地逆置}
\Code{0.0}{Reverse(A, n) \{}
\Hint{1.3}{l、r 分别指向尚未交换区间的两端。}\nopagebreak[4]
\Code{1.3}{l = 0;}
\Code{1.3}{r = n-1;}
\Code{1.3}{while (l<r) \{}
\Code{2.6}{swap(A[l], A[r]);}
\Hint{2.6}{一旦两个指针相遇或交错，中间已经没有需要交换的一对。}\nopagebreak[4]
\Code{2.6}{l++;}
\Code{2.6}{r--;}
\Code{1.3}{\}}
\Code{0.0}{\}}
\nopagebreak[4]
\Cost{\ensuremath{O\!\left(n\right)}}{\ensuremath{O\!\left(1\right)}}
\par\medskip

\Needspace{14\baselineskip}
\subsubsection{删除所有值为x的元素}
\Code{0.0}{RemoveValue(L, x) \{}
\Hint{1.3}{w 是下一个保留元素的写入位置。}\nopagebreak[4]
\Code{1.3}{w = 0;}
\Code{1.3}{for (v in L.data[0:L.n]) \{}
\Hint{2.6}{读到 x 时只推进读取位置，留出的位置会被后面的保留项覆盖。}\nopagebreak[4]
\Code{2.6}{if (v != x) L.data[w++] = v;}
\Code{1.3}{\}}
\Code{1.3}{L.n = w;}
\Code{0.0}{\}}
\nopagebreak[4]
\Cost{\ensuremath{O\!\left(n\right)}}{\ensuremath{O\!\left(1\right)}}
\par\medskip

\Needspace{14\baselineskip}
\subsubsection{删除指定值域内的元素}
\Code{0.0}{RemoveRange(L, lo, hi) \{}
\Hint{1.3}{闭区间外的元素按原顺序写回。}\nopagebreak[4]
\Code{1.3}{w = 0;}
\Code{1.3}{for (v in L.data[0:L.n]) \{}
\Hint{2.6}{w 最终就是新长度；lo 和 hi 本身也属于删除范围。}\nopagebreak[4]
\Code{2.6}{if (v<lo || v>hi) L.data[w++] = v;}
\Code{1.3}{\}}
\Code{1.3}{L.n = w;}
\Code{0.0}{\}}
\nopagebreak[4]
\Cost{\ensuremath{O\!\left(n\right)}}{\ensuremath{O\!\left(1\right)}}
\par\medskip

\Needspace{14\baselineskip}
\subsubsection{删除有序顺序表中的重复元素}
\Code{0.0}{Unique(L) \{}
\Hint{1.3}{有序表的相等元素连续，只需和最后保留项比较。}\nopagebreak[4]
\Code{1.3}{w = 0;}
\Code{1.3}{for (v in L.data[0:L.n]) \{}
\Hint{2.6}{w 为零时还没有上一项，必须先判 w，才能读 A[w−1]。}\nopagebreak[4]
\Code{2.6}{if (w == 0 || v != L.data[w-1]) L.data[w++] = v;}
\Code{1.3}{\}}
\Code{1.3}{L.n = w;}
\Code{0.0}{\}}
\nopagebreak[4]
\Cost{\ensuremath{O\!\left(n\right)}}{\ensuremath{O\!\left(1\right)}}
\par\medskip

\Needspace{18\baselineskip}
\subsubsection{合并两个有序顺序表}
\Code{0.0}{Merge(A, B) \{}
\Hint{1.3}{相等时取左表，保留相等元素的原有顺序。}\nopagebreak[4]
\Code{1.3}{i = j = 0;}
\Code{1.3}{InitList(\&C);}
\Code{1.3}{while (i<len(A) \&\& j<len(B)) \{}
\Code{2.6}{if (A[i] <= B[j]) Append(\&C, A[i++]);}
\Code{2.6}{else Append(\&C, B[j++]);}
\Code{1.3}{\}}
\Hint{1.3}{一侧耗尽后，另一侧已经排好序，可以逐项接到末尾。}\nopagebreak[4]
\Code{1.3}{while (i<len(A)) Append(\&C, A[i++]);}
\Code{1.3}{while (j<len(B)) Append(\&C, B[j++]);}
\Code{1.3}{return C;}
\Code{0.0}{\}}
\nopagebreak[4]
\Cost{\ensuremath{O\!\left(m+n\right)}}{\ensuremath{O\!\left(1\right)}}
\par\medskip

\Needspace{16\baselineskip}
\subsubsection{顺序表的循环左移与右移}
\Code{0.0}{Rotate(A, k, right) \{}
\Hint{1.3}{两段各自翻转，再翻转整体，就能交换两段位置。}\nopagebreak[4]
\Code{1.3}{n = len(A);}
\Code{1.3}{if (n == 0) return;}
\Code{1.3}{k = k\%n;}
\Code{1.3}{if (right) k = (n-k)\%n;}
\Hint{1.3}{k 先对 n 取余。空数组须提前返回，以免除以零。}\nopagebreak[4]
\Code{1.3}{reverse(A[0:k]);}
\Code{1.3}{reverse(A[k:n]);}
\Code{1.3}{reverse(A[0:n]);}
\Code{0.0}{\}}
\nopagebreak[4]
\Cost{\ensuremath{O\!\left(n\right)}}{\ensuremath{O\!\left(1\right)}}
\par\medskip

\Needspace{17\baselineskip}
\subsubsection{将顺序表按指定条件划分}
\Code{0.0}{Partition(A, P) \{}
\Hint{1.3}{w 左侧都已满足谓词 P。}\nopagebreak[4]
\Code{1.3}{w = 0;}
\Code{1.3}{for (i = 0; i<len(A); i++) \{}
\Code{2.6}{if (P(A[i])) \{}
\Code{3.9}{swap(A[w], A[i]);}
\Hint{3.9}{每找到一个符合 P 的元素，就交换到 w 所指的首个未归类位置。}\nopagebreak[4]
\Code{3.9}{w++;}
\Code{2.6}{\}}
\Code{1.3}{\}}
\Code{1.3}{return w;}
\Code{0.0}{\}}
\nopagebreak[4]
\Cost{\ensuremath{O\!\left(n\right)}}{\ensuremath{O\!\left(1\right)}}
\par\medskip

\Needspace{27\baselineskip}
\subsubsection{求两个有序顺序表的中位数}
\Code{0.0}{Median(A, B) \{}
\Hint{1.3}{切分线左边合计一半元素，两侧交叉比较决定移动方向。}\nopagebreak[4]
\Code{1.3}{if (len(A)>len(B)) return Median(B, A);}
\Code{1.3}{m = len(A);}
\Code{1.3}{n = len(B);}
\Code{1.3}{require(m+n>0);}
\Code{1.3}{lo = 0;}
\Code{1.3}{hi = m;}
\Code{1.3}{while (lo <= hi) \{}
\Code{2.6}{i = (lo+hi)/2;}
\Code{2.6}{j = (m+n+1)/2-i;}
\Code{2.6}{a = (i == 0 ? -INF : A[i-1]);}
\Code{2.6}{b = (i == m ? INF : A[i]);}
\Code{2.6}{c = (j == 0 ? -INF : B[j-1]);}
\Code{2.6}{d = (j == n ? INF : B[j]);}
\Hint{2.6}{a、b 是 A 的切口两侧，c、d 是 B 的切口两侧。若 a>d，A 在左半边放多了。}\nopagebreak[4]
\Code{2.6}{if (a>d) hi = i-1;}
\Code{2.6}{else if (c>b) lo = i+1;}
\Code{2.6}{else if ((m+n)\%2) return max(a, c);}
\Code{2.6}{else return (max(a, c)+min(b, d))/2.0;}
\Code{1.3}{\}}
\Code{0.0}{\}}
\nopagebreak[4]
\Cost{\ensuremath{O\!\left(\log (\min (m,n)+2)\right)}}{\ensuremath{O\!\left(1\right)}}
\par\medskip

\Needspace{18\baselineskip}
\subsubsection{有序顺序表的二分查找}
\Code{0.0}{BinarySearch(A, x) \{}
\Hint{1.3}{闭区间搜索，每次排除已检查的中点。}\nopagebreak[4]
\Code{1.3}{lo = 0;}
\Code{1.3}{hi = len(A)-1;}
\Code{1.3}{while (lo <= hi) \{}
\Code{2.6}{mid = lo+(hi-lo)/2;}
\Code{2.6}{if (A[mid] == x) return mid;}
\Hint{2.6}{当 lo>hi 时搜索区间已空；此时返回 −1。}\nopagebreak[4]
\Code{2.6}{if (A[mid]<x) lo = mid+1;}
\Code{2.6}{else hi = mid-1;}
\Code{1.3}{\}}
\Code{1.3}{return -1;}
\Code{0.0}{\}}
\nopagebreak[4]
\Cost{\ensuremath{O\!\left(\log (n+1)\right)}}{\ensuremath{O\!\left(1\right)}}
\par\medskip

\Needspace{17\baselineskip}
\subsubsection{二分查找的左右边界确定}
\Code{0.0}{Bound(A, x, upper) \{}
\Hint{1.3}{upper 为真时跳过所有等于 x 的项。两种返回值组成半开匹配区间。}\nopagebreak[4]
\Code{1.3}{lo = 0;}
\Code{1.3}{hi = len(A);}
\Code{1.3}{while (lo<hi) \{}
\Code{2.6}{mid = (lo+hi)/2;}
\Hint{2.6}{循环结束时 lo 是首个满足边界条件的下标。若所有值都更小，则 lo=n。}\nopagebreak[4]
\Code{2.6}{if (A[mid]<x || (upper \&\& A[mid] == x)) lo = mid+1;}
\Code{2.6}{else hi = mid;}
\Code{1.3}{\}}
\Code{1.3}{return lo;}
\Code{0.0}{\}}
\nopagebreak[4]
\Cost{\ensuremath{O\!\left(\log (n+1)\right)}}{\ensuremath{O\!\left(1\right)}}
\par\medskip

\Needspace{19\baselineskip}
\subsubsection{寻找数组中缺失的最小正整数}
\Code{0.0}{FirstMissing(A, n) \{}
\Hint{1.3}{正数 v 应归位到 v−1，已有相同值时停止交换。}\nopagebreak[4]
\Code{1.3}{for (i = 0; i<n; i++) \{}
\Code{2.6}{while (1 <= A[i] \&\& A[i] <= n \&\& A[A[i]-1] != A[i]) \{}
\Code{3.9}{j = A[i]-1;}
\Code{3.9}{swap(A[i], A[j]);}
\Code{2.6}{\}}
\Code{1.3}{\}}
\Hint{1.3}{交换前把目标下标存到 j，避免数组值交换后再用它计算位置。}\nopagebreak[4]
\Code{1.3}{for (i = 0; i<n; i++) \{}
\Code{2.6}{if (A[i] != i+1) return i+1;}
\Code{1.3}{\}}
\Code{1.3}{return n+1;}
\Code{0.0}{\}}
\nopagebreak[4]
\Cost{\ensuremath{O\!\left(n\right)}}{\ensuremath{O\!\left(1\right)}}
\par\medskip

\Needspace{16\baselineskip}
\subsubsection{求最大连续子数组和}
\Code{0.0}{MaxSubarray(A, n) \{}
\Hint{1.3}{ending 只记录恰在当前位置结束的非空段。}\nopagebreak[4]
\Code{1.3}{require(n>0);}
\Code{1.3}{best = ending = A[0];}
\Code{1.3}{for (i = 1; i<n; i++) \{}
\Code{2.6}{ending = max(A[i], ending+A[i]);}
\Hint{2.6}{ending 可以丢弃负贡献的旧段，从当前值重新开始；best 保留历史最优。}\nopagebreak[4]
\Code{2.6}{best = max(best, ending);}
\Code{1.3}{\}}
\Code{1.3}{return best;}
\Code{0.0}{\}}
\nopagebreak[4]
\Cost{\ensuremath{O\!\left(n\right)}}{\ensuremath{O\!\left(1\right)}}
\par\medskip

\Part{链表}

\Group{单链表的基本操作}

\Needspace{16\baselineskip}
\subsubsection{头插法建立单链表}
\Code{0.0}{BuildFront(input) \{}
\Hint{1.3}{新结点总是接在头结点之后，结果顺序与输入相反。}\nopagebreak[4]
\Code{1.3}{H = NewNode();}
\Code{1.3}{for (x in input) \{}
\Code{2.6}{p = NewNode(x);}
\Hint{2.6}{先把旧首元地址存进新结点，下一行才能安全改写头指针。}\nopagebreak[4]
\Code{2.6}{p->next = H->next;}
\Hint{2.6}{先令 p 接到旧首元，再令 H 接到 p。若交换这两句，容易产生自环。}\nopagebreak[4]
\Code{2.6}{H->next = p;}
\Code{1.3}{\}}
\Code{1.3}{return H;}
\Code{0.0}{\}}
\nopagebreak[4]
\Cost{\ensuremath{O\!\left(n\right)}}{\ensuremath{O\!\left(1\right)}}
\par\medskip

\Needspace{17\baselineskip}
\subsubsection{尾插法建立单链表}
\Code{0.0}{BuildBack(input) \{}
\Hint{1.3}{tail 保存最后一个结点，空表时它就是 H。}\nopagebreak[4]
\Code{1.3}{H = NewNode();}
\Code{1.3}{tail = H;}
\Code{1.3}{for (x in input) \{}
\Code{2.6}{p = NewNode(x);}
\Hint{2.6}{把 tail->next 指给 p 是连边，把 tail 改成 p 是移动游标。}\nopagebreak[4]
\Code{2.6}{tail->next = p;}
\Code{2.6}{tail = p;}
\Code{1.3}{\}}
\Code{1.3}{return H;}
\Code{0.0}{\}}
\nopagebreak[4]
\Cost{\ensuremath{O\!\left(n\right)}}{\ensuremath{O\!\left(1\right)}}
\par\medskip

\Needspace{13\baselineskip}
\subsubsection{按位查找}
\Code{0.0}{NodeAt(H, i) \{}
\Hint{1.3}{首元是第 1 个数据结点。}\nopagebreak[4]
\Code{1.3}{if (i<1) return NULL;}
\Code{1.3}{p = H->next;}
\Hint{1.3}{每走一次后继，距离目标位序还差一步的计数就减少一。}\nopagebreak[4]
\Code{1.3}{while (p \&\& --i>0) p = p->next;}
\Code{1.3}{return p;}
\Code{0.0}{\}}
\nopagebreak[4]
\Cost{\ensuremath{O\!\left(n\right)}}{\ensuremath{O\!\left(1\right)}}
\par\medskip

\Needspace{12\baselineskip}
\subsubsection{按值查找}
\Code{0.0}{Find(H, x) \{}
\Hint{1.3}{找不到时游标自然走到 NULL。}\nopagebreak[4]
\Code{1.3}{p = H->next;}
\Hint{1.3}{判断 p 非空后才能取 p->data；这就是循环条件从左往右写的原因。}\nopagebreak[4]
\Code{1.3}{while (p \&\& p->data != x) p = p->next;}
\Code{1.3}{return p;}
\Code{0.0}{\}}
\nopagebreak[4]
\Cost{\ensuremath{O\!\left(n\right)}}{\ensuremath{O\!\left(1\right)}}
\par\medskip

\Needspace{19\baselineskip}
\subsubsection{指定位置插入结点}
\Code{0.0}{InsertNode(H, i, x) \{}
\Hint{1.3}{前驱走 i−1 步，插首位时前驱就是 H。}\nopagebreak[4]
\Code{1.3}{require(i >= 1);}
\Code{1.3}{p = H;}
\Code{1.3}{repeat (i-1) \{}
\Code{2.6}{if (!p->next) return false;}
\Code{2.6}{p = p->next;}
\Code{1.3}{\}}
\Hint{1.3}{p 要停在插入位置之前，这样无需单独修改首元情形。}\nopagebreak[4]
\Code{1.3}{s = NewNode(x);}
\Code{1.3}{s->next = p->next;}
\Code{1.3}{p->next = s;}
\Code{1.3}{return true;}
\Code{0.0}{\}}
\nopagebreak[4]
\Cost{\ensuremath{O\!\left(n\right)}}{\ensuremath{O\!\left(1\right)}}
\par\medskip

\Needspace{21\baselineskip}
\subsubsection{指定位置删除结点}
\Code{0.0}{DeleteNode(H, i) \{}
\Hint{1.3}{先接通后半链，再释放被删除结点。}\nopagebreak[4]
\Code{1.3}{require(i >= 1);}
\Code{1.3}{p = H;}
\Code{1.3}{repeat (i-1) \{}
\Code{2.6}{if (!p->next) return FAIL;}
\Code{2.6}{p = p->next;}
\Code{1.3}{\}}
\Hint{1.3}{p 存在仍不足以删除，p->next 还必须是一个数据结点。}\nopagebreak[4]
\Code{1.3}{q = p->next;}
\Code{1.3}{if (!q) return FAIL;}
\Code{1.3}{p->next = q->next;}
\Code{1.3}{x = q->data;}
\Code{1.3}{free(q);}
\Code{1.3}{return x;}
\Code{0.0}{\}}
\nopagebreak[4]
\Cost{\ensuremath{O\!\left(n\right)}}{\ensuremath{O\!\left(1\right)}}
\par\medskip

\Needspace{18\baselineskip}
\subsubsection{删除指定值的结点}
\Code{0.0}{DeleteValue(H, x) \{}
\Hint{1.3}{删除后前驱不动，才能检查连续出现的目标值。}\nopagebreak[4]
\Code{1.3}{p = H;}
\Code{1.3}{while (p->next) \{}
\Code{2.6}{q = p->next;}
\Code{2.6}{if (q->data == x) \{}
\Code{3.9}{p->next = q->next;}
\Hint{3.9}{只有当前项保留时才让 p 前进；删掉当前项后新的后继尚未检查。}\nopagebreak[4]
\Code{3.9}{free(q);}
\Code{2.6}{\}}
\Code{2.6}{else p = q;}
\Code{1.3}{\}}
\Code{0.0}{\}}
\nopagebreak[4]
\Cost{\ensuremath{O\!\left(n\right)}}{\ensuremath{O\!\left(1\right)}}
\par\medskip

\Needspace{12\baselineskip}
\subsubsection{单链表的遍历}
\Code{0.0}{VisitNodes(H) \{}
\Hint{1.3}{访问从首元开始，哨兵不存有效数据。}\nopagebreak[4]
\Code{1.3}{for (p = H->next; p != NULL; p = p->next) \{}
\Hint{2.6}{p=p->next 只移动局部指针，不改变任何结点的后继。}\nopagebreak[4]
\Code{2.6}{visit(p->data);}
\Code{1.3}{\}}
\Code{0.0}{\}}
\nopagebreak[4]
\Cost{\ensuremath{O\!\left(n\right)}}{\ensuremath{O\!\left(1\right)}}
\par\medskip

\Needspace{14\baselineskip}
\subsubsection{单链表的销毁}
\Code{0.0}{DestroyList(Node **H) \{}
\Hint{1.3}{free 后不能再取后继，必须先保存。}\nopagebreak[4]
\Code{1.3}{while ((*H)) \{}
\Code{2.6}{q = (*H)->next;}
\Hint{2.6}{最后一次迭代把 H 变为空，调用方不再持有已释放的地址。}\nopagebreak[4]
\Code{2.6}{free((*H));}
\Code{2.6}{(*H) = q;}
\Code{1.3}{\}}
\Code{0.0}{\}}
\nopagebreak[4]
\Cost{\ensuremath{O\!\left(n\right)}}{\ensuremath{O\!\left(1\right)}}
\par\medskip

\Group{双向链表}

\Needspace{18\baselineskip}
\subsubsection{双向链表的建立}
\Code{0.0}{BuildDouble(input) \{}
\Hint{1.3}{向右接新结点时同时填它的左向指针。}\nopagebreak[4]
\Code{1.3}{H = NewDNode();}
\Code{1.3}{t = H;}
\Code{1.3}{for (x in input) \{}
\Code{2.6}{p = NewDNode(x);}
\Code{2.6}{p->prior = t;}
\Hint{2.6}{每次循环结束，t 都是新尾结点，下一次通过它连接新项。}\nopagebreak[4]
\Code{2.6}{t->next = p;}
\Code{2.6}{t = p;}
\Code{1.3}{\}}
\Code{1.3}{return H;}
\Code{0.0}{\}}
\nopagebreak[4]
\Cost{\ensuremath{O\!\left(n\right)}}{\ensuremath{O\!\left(1\right)}}
\par\medskip

\Needspace{15\baselineskip}
\subsubsection{双向链表的插入}
\Code{0.0}{InsertAfter(p, x) \{}
\Hint{1.3}{q 可能为空，此时插入点在表尾。}\nopagebreak[4]
\Code{1.3}{q = p->next;}
\Code{1.3}{s = NewDNode(x);}
\Code{1.3}{s->prior = p;}
\Hint{1.3}{s 的左右域先接好，最后让原来的左邻指向 s。}\nopagebreak[4]
\Code{1.3}{s->next = q;}
\Code{1.3}{if (q) q->prior = s;}
\Code{1.3}{p->next = s;}
\Code{0.0}{\}}
\nopagebreak[4]
\Cost{\ensuremath{O\!\left(1\right)}}{\ensuremath{O\!\left(1\right)}}
\par\medskip

\Needspace{14\baselineskip}
\subsubsection{双向链表的删除}
\Code{0.0}{EraseDouble(s) \{}
\Hint{1.3}{s 是数据结点，其左邻可能是哨兵。}\nopagebreak[4]
\Code{1.3}{a = s->prior;}
\Code{1.3}{b = s->next;}
\Hint{1.3}{若 b 为空，说明删的是尾结点，不可访问 b 的 prior 域。}\nopagebreak[4]
\Code{1.3}{a->next = b;}
\Code{1.3}{if (b) b->prior = a;}
\Code{1.3}{free(s);}
\Code{0.0}{\}}
\nopagebreak[4]
\Cost{\ensuremath{O\!\left(1\right)}}{\ensuremath{O\!\left(1\right)}}
\par\medskip

\Needspace{19\baselineskip}
\subsubsection{双向链表的逆置}
\Code{0.0}{ReverseDouble(H) \{}
\Hint{1.3}{摘下旧首元，再插入新链的最前端。}\nopagebreak[4]
\Code{1.3}{p = H->next;}
\Code{1.3}{H->next = NULL;}
\Code{1.3}{while (p) \{}
\Code{2.6}{q = p->next;}
\Code{2.6}{p->next = H->next;}
\Hint{2.6}{q 保存旧后继；移动 p 时必须使用 q，不能沿已经改向的 next。}\nopagebreak[4]
\Code{2.6}{p->prior = H;}
\Code{2.6}{if (H->next) H->next->prior = p;}
\Code{2.6}{H->next = p;}
\Code{2.6}{p = q;}
\Code{1.3}{\}}
\Code{0.0}{\}}
\nopagebreak[4]
\Cost{\ensuremath{O\!\left(n\right)}}{\ensuremath{O\!\left(1\right)}}
\par\medskip

\Group{循环链表}

\Needspace{19\baselineskip}
\subsubsection{循环单链表的建立}
\Code{0.0}{BuildRing(input) \{}
\Hint{1.3}{空环的头和尾是同一哨兵。}\nopagebreak[4]
\Code{1.3}{H = NewNode();}
\Code{1.3}{t = H;}
\Code{1.3}{H->next = H;}
\Code{1.3}{for (x in input) \{}
\Code{2.6}{p = NewNode(x);}
\Hint{2.6}{新结点的 next 指回 H，尾插之后仍然是一个封闭的环。}\nopagebreak[4]
\Code{2.6}{p->next = H;}
\Code{2.6}{t->next = p;}
\Code{2.6}{t = p;}
\Code{1.3}{\}}
\Code{1.3}{return (H, t);}
\Code{0.0}{\}}
\nopagebreak[4]
\Cost{\ensuremath{O\!\left(n\right)}}{\ensuremath{O\!\left(1\right)}}
\par\medskip

\Needspace{22\baselineskip}
\subsubsection{循环单链表的插入与删除}
\Code{0.0}{UpdateRing(H, Node **tail, p, x, insert) \{}
\Hint{1.3}{尾部插删必须同步修改 tail。}\nopagebreak[4]
\Code{1.3}{if (insert) \{}
\Code{2.6}{s = NewNode(x);}
\Code{2.6}{s->next = p->next;}
\Code{2.6}{p->next = s;}
\Code{2.6}{if ((*tail) == p) (*tail) = s;}
\Code{1.3}{\}}
\Code{1.3}{else \{}
\Hint{2.6}{不允许删除哨兵；删去唯一数据项时 tail 会重新指向 H。}\nopagebreak[4]
\Code{2.6}{q = p->next;}
\Code{2.6}{require(q != H);}
\Code{2.6}{p->next = q->next;}
\Code{2.6}{if ((*tail) == q) (*tail) = p;}
\Code{2.6}{free(q);}
\Code{1.3}{\}}
\Code{0.0}{\}}
\nopagebreak[4]
\Cost{\ensuremath{O\!\left(1\right)}}{\ensuremath{O\!\left(1\right)}}
\par\medskip

\Needspace{23\baselineskip}
\subsubsection{循环双向链表的插入与删除}
\Code{0.0}{UpdateDoubleRing(H, p, x, insert) \{}
\Hint{1.3}{哨兵参与环接，首尾边界也有两个邻居。}\nopagebreak[4]
\Code{1.3}{if (insert) \{}
\Code{2.6}{q = p->next;}
\Code{2.6}{s = NewDNode(x);}
\Code{2.6}{s->prior = p;}
\Code{2.6}{s->next = q;}
\Code{2.6}{q->prior = s;}
\Code{2.6}{p->next = s;}
\Code{1.3}{\}}
\Hint{1.3}{删除 p 只需让两侧邻居互指，随后再释放 p。}\nopagebreak[4]
\Code{1.3}{else \{}
\Code{2.6}{require(p != H);}
\Code{2.6}{p->prior->next = p->next;}
\Code{2.6}{p->next->prior = p->prior;}
\Code{2.6}{free(p);}
\Code{1.3}{\}}
\Code{0.0}{\}}
\nopagebreak[4]
\Cost{\ensuremath{O\!\left(1\right)}}{\ensuremath{O\!\left(1\right)}}
\par\medskip

\Needspace{15\baselineskip}
\subsubsection{循环链表的合并}
\Code{0.0}{JoinRings(A, B) \{}
\Hint{1.3}{B 的结点转移给 A，B 的哨兵恢复为空环。}\nopagebreak[4]
\Code{1.3}{if (B.tail == B.head) return;}
\Code{1.3}{A.tail->next = B.head->next;}
\Code{1.3}{B.tail->next = A.head;}
\Hint{1.3}{A 为空时它的尾就是自己的哨兵，同一组赋值也能接入 B。}\nopagebreak[4]
\Code{1.3}{A.tail = B.tail;}
\Code{1.3}{B.head->next = B.head;}
\Code{1.3}{B.tail = B.head;}
\Code{0.0}{\}}
\nopagebreak[4]
\Cost{\ensuremath{O\!\left(1\right)}}{\ensuremath{O\!\left(1\right)}}
\par\medskip

\Needspace{26\baselineskip}
\subsubsection{约瑟夫环问题}
\Code{0.0}{Josephus(first, tail, k) \{}
\Hint{1.3}{无哨兵环中 prev 始终是 cur 的前驱。}\nopagebreak[4]
\Code{1.3}{require(k>0);}
\Code{1.3}{cur = first;}
\Code{1.3}{prev = tail;}
\Code{1.3}{while (cur) \{}
\Code{2.6}{repeat (k-1) \{}
\Code{3.9}{prev = cur;}
\Code{3.9}{cur = cur->next;}
\Code{2.6}{\}}
\Code{2.6}{emit(cur->data);}
\Hint{2.6}{当前人计作第 1 人，所以前进 k−1 次；最后一人须单独终止。}\nopagebreak[4]
\Code{2.6}{if (cur->next == cur) \{}
\Code{3.9}{free(cur);}
\Code{3.9}{break;}
\Code{2.6}{\}}
\Code{2.6}{prev->next = cur->next;}
\Code{2.6}{free(cur);}
\Code{2.6}{cur = prev->next;}
\Code{1.3}{\}}
\Code{0.0}{\}}
\nopagebreak[4]
\Cost{\ensuremath{O\!\left(nk\right)}}{\ensuremath{O\!\left(1\right)}}
\par\medskip

\Group{静态链表}

\Needspace{14\baselineskip}
\subsubsection{静态链表的初始化}
\Code{0.0}{InitStatic(A, C) \{}
\Hint{1.3}{槽 0 为哨兵，−1 表示没有后继。}\nopagebreak[4]
\Code{1.3}{A[0].next = -1;}
\Code{1.3}{freeHead = (C>0 ? 1 : -1);}
\Hint{1.3}{最后一个可用槽的 next 设为 −1，防止分配时越界。}\nopagebreak[4]
\Code{1.3}{for (i = 1; i <= C; i++) \{}
\Code{2.6}{A[i].next = (i<C ? i+1 : -1);}
\Code{1.3}{\}}
\Code{0.0}{\}}
\nopagebreak[4]
\Cost{\ensuremath{O\!\left(C\right)}}{\ensuremath{O\!\left(1\right)}}
\par\medskip

\Needspace{17\baselineskip}
\subsubsection{空闲结点的分配与回收}
\Code{0.0}{Slot(op, i) \{}
\Hint{1.3}{分配和归还都只操作空闲链的表头。}\nopagebreak[4]
\Code{1.3}{if (op == ALLOC) \{}
\Code{2.6}{require(freeHead != -1);}
\Code{2.6}{i = freeHead;}
\Code{2.6}{freeHead = A[i].next;}
\Hint{2.6}{alloc 保存原空闲表头作为结果，再把空闲表头移到下一槽。}\nopagebreak[4]
\Code{2.6}{return i;}
\Code{1.3}{\}}
\Code{1.3}{A[i].next = freeHead;}
\Code{1.3}{freeHead = i;}
\Code{0.0}{\}}
\nopagebreak[4]
\Cost{\ensuremath{O\!\left(1\right)}}{\ensuremath{O\!\left(1\right)}}
\par\medskip

\Needspace{12\baselineskip}
\subsubsection{静态链表的查找}
\Code{0.0}{StaticFind(A, x) \{}
\Hint{1.3}{沿 next 游标走，物理下标无需连续。}\nopagebreak[4]
\Code{1.3}{i = A[0].next;}
\Hint{1.3}{i 为 −1 时不能用作数组下标，应先结束查找。}\nopagebreak[4]
\Code{1.3}{while (i != -1 \&\& A[i].data != x) i = A[i].next;}
\Code{1.3}{return i;}
\Code{0.0}{\}}
\nopagebreak[4]
\Cost{\ensuremath{O\!\left(n\right)}}{\ensuremath{O\!\left(1\right)}}
\par\medskip

\Needspace{21\baselineskip}
\subsubsection{静态链表的插入与删除}
\Code{0.0}{StaticUpdate(A, p, x, insert) \{}
\Hint{1.3}{回收槽之前先从数据链解开。}\nopagebreak[4]
\Code{1.3}{if (insert) \{}
\Code{2.6}{i = Slot(ALLOC, 0);}
\Code{2.6}{A[i].data = x;}
\Code{2.6}{A[i].next = A[p].next;}
\Code{2.6}{A[p].next = i;}
\Code{1.3}{\}}
\Hint{1.3}{插入时先保存前驱的旧后继；删除时先接通后继，再把槽归还。}\nopagebreak[4]
\Code{1.3}{else \{}
\Code{2.6}{i = A[p].next;}
\Code{2.6}{require(i != -1);}
\Code{2.6}{A[p].next = A[i].next;}
\Code{2.6}{Slot(FREE, i);}
\Code{1.3}{\}}
\Code{0.0}{\}}
\nopagebreak[4]
\Cost{\ensuremath{O\!\left(1\right)}}{\ensuremath{O\!\left(1\right)}}
\par\medskip

\Group{链表的综合算法}

\Needspace{18\baselineskip}
\subsubsection{单链表的原地逆置}
\Code{0.0}{ReverseNodes(first) \{}
\Hint{1.3}{prev 是已经逆好的链首。改写 cur 的后继之前，先保存剩余链。}\nopagebreak[4]
\Code{1.3}{prev = NULL;}
\Code{1.3}{cur = first;}
\Code{1.3}{while (cur) \{}
\Code{2.6}{after = cur->next;}
\Hint{2.6}{这一步改变链的方向。cur 本身仍指向原结点。}\nopagebreak[4]
\Code{2.6}{cur->next = prev;}
\Hint{2.6}{after 指向未处理后缀，prev 指向已逆置前缀；两者合起来保住整条链。}\nopagebreak[4]
\Code{2.6}{prev = cur;}
\Code{2.6}{cur = after;}
\Code{1.3}{\}}
\Code{1.3}{return prev;}
\Code{0.0}{\}}
\nopagebreak[4]
\Cost{\ensuremath{O\!\left(n\right)}}{\ensuremath{O\!\left(1\right)}}
\par\medskip

\Needspace{14\baselineskip}
\subsubsection{单链表的递归逆置}
\Code{0.0}{ReverseRec(p) \{}
\Hint{1.3}{子调用返回新首元；原后继现在是后半链的末尾。}\nopagebreak[4]
\Code{1.3}{if (!p || !p->next) return p;}
\Code{1.3}{head = ReverseRec(p->next);}
\Hint{1.3}{p 的原后继现在接回 p，使 p 排在新链最后。}\nopagebreak[4]
\Code{1.3}{p->next->next = p;}
\Hint{1.3}{把 p 的 next 清空，使原链首变成新尾，避免两个结点互指。}\nopagebreak[4]
\Code{1.3}{p->next = NULL;}
\Code{1.3}{return head;}
\Code{0.0}{\}}
\nopagebreak[4]
\Cost{\ensuremath{O\!\left(n\right)}}{\ensuremath{O\!\left(n\right)}}
\par\medskip

\Needspace{19\baselineskip}
\subsubsection{单链表的局部反转（第m至第n个结点）}
\Code{0.0}{ReverseRange(H, m, n) \{}
\Hint{1.3}{tail 不移动，每次把它后面的一个结点搬到 pre 后。}\nopagebreak[4]
\Code{1.3}{require(1 <= m \&\& m <= n);}
\Code{1.3}{pre = H;}
\Code{1.3}{repeat (m-1) pre = pre->next;}
\Code{1.3}{tail = pre->next;}
\Code{1.3}{repeat (n-m) \{}
\Hint{2.6}{s 原来紧跟 tail，摘下它后先补好缺口，再让它成为区间首元。}\nopagebreak[4]
\Code{2.6}{s = tail->next;}
\Code{2.6}{tail->next = s->next;}
\Code{2.6}{s->next = pre->next;}
\Code{2.6}{pre->next = s;}
\Code{1.3}{\}}
\Code{0.0}{\}}
\nopagebreak[4]
\Cost{\ensuremath{O\!\left(n\right)}}{\ensuremath{O\!\left(1\right)}}
\par\medskip

\Needspace{26\baselineskip}
\subsubsection{单链表每k个结点一组翻转}
\Code{0.0}{ReverseGroups(H, k) \{}
\Hint{1.3}{先确认有完整一组，最后不足 k 项的后缀保持原序。}\nopagebreak[4]
\Code{1.3}{require(k>0);}
\Code{1.3}{pre = H;}
\Code{1.3}{while (true) \{}
\Code{2.6}{end = pre;}
\Code{2.6}{repeat (k) \{}
\Code{3.9}{end = end->next;}
\Code{3.9}{if (!end) return;}
\Code{2.6}{\}}
\Code{2.6}{tail = pre->next;}
\Hint{2.6}{一组翻完后原组首 tail 变成组尾，也就是下一组的前驱。}\nopagebreak[4]
\Code{2.6}{repeat (k-1) \{}
\Code{3.9}{s = tail->next;}
\Code{3.9}{tail->next = s->next;}
\Code{3.9}{s->next = pre->next;}
\Code{3.9}{pre->next = s;}
\Code{2.6}{\}}
\Hint{2.6}{原来的组首已成组尾，从这里准备下一组。}\nopagebreak[4]
\Code{2.6}{pre = tail;}
\Code{1.3}{\}}
\Code{0.0}{\}}
\nopagebreak[4]
\Cost{\ensuremath{O\!\left(n\right)}}{\ensuremath{O\!\left(1\right)}}
\par\medskip

\Needspace{18\baselineskip}
\subsubsection{单链表相邻结点的两两交换}
\Code{0.0}{SwapPairs(H) \{}
\Hint{1.3}{a 交换后变成这一对的末尾。}\nopagebreak[4]
\Code{1.3}{p = H;}
\Code{1.3}{while (p->next \&\& p->next->next) \{}
\Code{2.6}{a = p->next;}
\Code{2.6}{b = a->next;}
\Code{2.6}{a->next = b->next;}
\Hint{2.6}{先让 a 跨过 b 指向下一对，再让 b 指回 a。}\nopagebreak[4]
\Code{2.6}{b->next = a;}
\Code{2.6}{p->next = b;}
\Code{2.6}{p = a;}
\Code{1.3}{\}}
\Code{0.0}{\}}
\nopagebreak[4]
\Cost{\ensuremath{O\!\left(n\right)}}{\ensuremath{O\!\left(1\right)}}
\par\medskip

\Needspace{22\baselineskip}
\subsubsection{单链表的循环右移}
\Code{0.0}{RotateRight(first, k) \{}
\Hint{1.3}{临时成环，再在新尾处断开。}\nopagebreak[4]
\Code{1.3}{if (!first) return NULL;}
\Code{1.3}{t = first;}
\Code{1.3}{n = 1;}
\Code{1.3}{while (t->next) \{}
\Code{2.6}{t = t->next;}
\Code{2.6}{n++;}
\Code{1.3}{\}}
\Hint{1.3}{新尾是从原首元出发走 n−k−1 步的位置，新首是它的后继。}\nopagebreak[4]
\Code{1.3}{t->next = first;}
\Code{1.3}{p = first;}
\Code{1.3}{repeat (n-1-k\%n) p = p->next;}
\Code{1.3}{head = p->next;}
\Code{1.3}{p->next = NULL;}
\Code{1.3}{return head;}
\Code{0.0}{\}}
\nopagebreak[4]
\Cost{\ensuremath{O\!\left(n\right)}}{\ensuremath{O\!\left(1\right)}}
\par\medskip

\Needspace{22\baselineskip}
\subsubsection{删除单链表中的重复结点}
\Code{0.0}{RemoveDuplicates(H) \{}
\Hint{1.3}{哈希集合记录已经保留的值。}\nopagebreak[4]
\Code{1.3}{InitSet(\&seen);}
\Code{1.3}{p = H;}
\Code{1.3}{while (p->next) \{}
\Code{2.6}{q = p->next;}
\Code{2.6}{if (SetContains(\&seen, q->data)) \{}
\Code{3.9}{p->next = q->next;}
\Code{3.9}{free(q);}
\Code{2.6}{\}}
\Hint{2.6}{先检查是否见过该值；否则每项都会被自己刚插入的记录误判为重复。}\nopagebreak[4]
\Code{2.6}{else \{}
\Code{3.9}{SetInsert(\&seen, q->data);}
\Code{3.9}{p = q;}
\Code{2.6}{\}}
\Code{1.3}{\}}
\Code{0.0}{\}}
\nopagebreak[4]
\Cost{\ensuremath{O\!\left(n\right)}（期望）}{\ensuremath{O\!\left(n\right)}}
\par\medskip

\Needspace{18\baselineskip}
\subsubsection{删除单链表中满足指定条件的结点}
\Code{0.0}{RemoveIf(H, P) \{}
\Hint{1.3}{p 始终位于待检查结点之前。}\nopagebreak[4]
\Code{1.3}{p = H;}
\Code{1.3}{while (p->next) \{}
\Code{2.6}{q = p->next;}
\Code{2.6}{if (P(q->data)) \{}
\Code{3.9}{p->next = q->next;}
\Hint{3.9}{P 返回真表示删除；返回假才把当前数据结点作为下一轮的前驱。}\nopagebreak[4]
\Code{3.9}{free(q);}
\Code{2.6}{\}}
\Code{2.6}{else p = q;}
\Code{1.3}{\}}
\Code{0.0}{\}}
\nopagebreak[4]
\Cost{\ensuremath{O\!\left(n\right)}}{\ensuremath{O\!\left(1\right)}}
\par\medskip

\Needspace{24\baselineskip}
\subsubsection{合并两个有序单链表}
\Code{0.0}{MergeNodes(a, b) \{}
\Hint{1.3}{复用原结点，较小链头接到 t 后面。}\nopagebreak[4]
\Code{1.3}{D = ZeroNode();}
\Code{1.3}{t = \&D;}
\Code{1.3}{while (a \&\& b) \{}
\Code{2.6}{if (a->data <= b->data) \{}
\Code{3.9}{t->next = a;}
\Code{3.9}{a = a->next;}
\Code{2.6}{\}}
\Code{2.6}{else \{}
\Hint{3.9}{归并结束可直接接上剩余链，不必把剩余结点逐个拆下。}\nopagebreak[4]
\Code{3.9}{t->next = b;}
\Code{3.9}{b = b->next;}
\Code{2.6}{\}}
\Code{2.6}{t = t->next;}
\Code{1.3}{\}}
\Hint{1.3}{只有一侧还有余项，直接接整段。}\nopagebreak[4]
\Code{1.3}{t->next = (a ? a : b);}
\Code{1.3}{return D.next;}
\Code{0.0}{\}}
\nopagebreak[4]
\Cost{\ensuremath{O\!\left(m+n\right)}}{\ensuremath{O\!\left(1\right)}}
\par\medskip

\Needspace{25\baselineskip}
\subsubsection{两个单链表的交替合并}
\Code{0.0}{Alternate(a, b) \{}
\Hint{1.3}{两条链各取一项，短链先耗尽也可以继续。}\nopagebreak[4]
\Code{1.3}{D = ZeroNode();}
\Code{1.3}{t = \&D;}
\Code{1.3}{while (a || b) \{}
\Code{2.6}{if (a) \{}
\Code{3.9}{t->next = a;}
\Code{3.9}{t = a;}
\Code{3.9}{a = a->next;}
\Code{2.6}{\}}
\Code{2.6}{if (b) \{}
\Hint{3.9}{每次接入的是原结点，结果与输入共享这些地址，原表头不再单独拥有链。}\nopagebreak[4]
\Code{3.9}{t->next = b;}
\Code{3.9}{t = b;}
\Code{3.9}{b = b->next;}
\Code{2.6}{\}}
\Code{1.3}{\}}
\Code{1.3}{t->next = NULL;}
\Code{1.3}{return D.next;}
\Code{0.0}{\}}
\nopagebreak[4]
\Cost{\ensuremath{O\!\left(m+n\right)}}{\ensuremath{O\!\left(1\right)}}
\par\medskip

\Needspace{18\baselineskip}
\subsubsection{单链表的拆分}
\Code{0.0}{SplitAt(H, k) \{}
\Hint{1.3}{切口前驱保留在前表，后表另建哨兵。}\nopagebreak[4]
\Code{1.3}{p = H;}
\Code{1.3}{repeat (k) \{}
\Code{2.6}{require(p->next);}
\Code{2.6}{p = p->next;}
\Code{1.3}{\}}
\Hint{1.3}{先把 p 的后继保存为新表的首元，再切断，两段都可重新找到。}\nopagebreak[4]
\Code{1.3}{B = NewNode();}
\Code{1.3}{B->next = p->next;}
\Code{1.3}{p->next = NULL;}
\Code{1.3}{return (H, B);}
\Code{0.0}{\}}
\nopagebreak[4]
\Cost{\ensuremath{O\!\left(n\right)}}{\ensuremath{O\!\left(1\right)}}
\par\medskip

\Needspace{19\baselineskip}
\subsubsection{按结点位置进行奇偶拆分}
\Code{0.0}{SplitParity(first) \{}
\Hint{1.3}{奇偶按原位序判定，两条子链都用尾插。}\nopagebreak[4]
\Code{1.3}{T[0..1] = twoEmptyLists();}
\Code{1.3}{i = 1;}
\Code{1.3}{while (first) \{}
\Code{2.6}{p = first;}
\Code{2.6}{first = p->next;}
\Hint{2.6}{尾插前将 p->next 清空，防止结果表尾还拖着原链后缀。}\nopagebreak[4]
\Code{2.6}{p->next = NULL;}
\Code{2.6}{AppendNode(\&T[i\%2], p);}
\Code{2.6}{i++;}
\Code{1.3}{\}}
\Code{1.3}{return (T[1], T[0]);}
\Code{0.0}{\}}
\nopagebreak[4]
\Cost{\ensuremath{O\!\left(n\right)}}{\ensuremath{O\!\left(1\right)}}
\par\medskip

\Needspace{20\baselineskip}
\subsubsection{按给定值对单链表进行稳定划分}
\Code{0.0}{StablePartition(first, x) \{}
\Hint{1.3}{两条临时链的尾插顺序就是原链次序。}\nopagebreak[4]
\Code{1.3}{small = emptyList();}
\Code{1.3}{big = emptyList();}
\Code{1.3}{while (first) \{}
\Code{2.6}{p = first;}
\Code{2.6}{first = p->next;}
\Code{2.6}{p->next = NULL;}
\Hint{2.6}{两个结果都带哨兵，因此 small 为空时也可通过它的尾部接入 big。}\nopagebreak[4]
\Code{2.6}{if (p->data<x) AppendNode(\&small, p);}
\Code{2.6}{else AppendNode(\&big, p);}
\Code{1.3}{\}}
\Code{1.3}{small.tail->next = big.first;}
\Code{1.3}{return small.head->next;}
\Code{0.0}{\}}
\nopagebreak[4]
\Cost{\ensuremath{O\!\left(n\right)}}{\ensuremath{O\!\left(1\right)}}
\par\medskip

\Needspace{15\baselineskip}
\subsubsection{查找单链表的中间结点}
\Code{0.0}{Middle(first) \{}
\Hint{1.3}{偶数长度返回靠右的中点。}\nopagebreak[4]
\Code{1.3}{slow = fast = first;}
\Code{1.3}{while (fast \&\& fast->next) \{}
\Code{2.6}{slow = slow->next;}
\Hint{2.6}{快指针若为空就结束；它若没有后继，也不能再连续走两步。}\nopagebreak[4]
\Code{2.6}{fast = fast->next->next;}
\Code{1.3}{\}}
\Code{1.3}{return slow;}
\Code{0.0}{\}}
\nopagebreak[4]
\Cost{\ensuremath{O\!\left(n\right)}}{\ensuremath{O\!\left(1\right)}}
\par\medskip

\Needspace{20\baselineskip}
\subsubsection{查找单链表的倒数第k个结点}
\Code{0.0}{KthLast(first, k) \{}
\Hint{1.3}{两指针拉开 k 步后保持间距。}\nopagebreak[4]
\Code{1.3}{require(k>0);}
\Code{1.3}{a = b = first;}
\Code{1.3}{repeat (k) \{}
\Code{2.6}{if (!a) return NULL;}
\Code{2.6}{a = a->next;}
\Code{1.3}{\}}
\Hint{1.3}{先走 k 步时若提前遇空，说明链长不足 k；正好走到空仍然合法。}\nopagebreak[4]
\Code{1.3}{while (a) \{}
\Code{2.6}{a = a->next;}
\Code{2.6}{b = b->next;}
\Code{1.3}{\}}
\Code{1.3}{return b;}
\Code{0.0}{\}}
\nopagebreak[4]
\Cost{\ensuremath{O\!\left(n\right)}}{\ensuremath{O\!\left(1\right)}}
\par\medskip

\Needspace{23\baselineskip}
\subsubsection{删除单链表的倒数第k个结点}
\Code{0.0}{EraseKthLast(H, k) \{}
\Hint{1.3}{慢指针从哨兵出发，停在待删项前一格。}\nopagebreak[4]
\Code{1.3}{require(k>0);}
\Code{1.3}{a = H->next;}
\Code{1.3}{p = H;}
\Code{1.3}{repeat (k) \{}
\Code{2.6}{require(a);}
\Code{2.6}{a = a->next;}
\Code{1.3}{\}}
\Code{1.3}{while (a) \{}
\Hint{2.6}{a 到空时 p 停在目标前驱，所以目标通过 p->next 取得。}\nopagebreak[4]
\Code{2.6}{a = a->next;}
\Code{2.6}{p = p->next;}
\Code{1.3}{\}}
\Code{1.3}{q = p->next;}
\Code{1.3}{p->next = q->next;}
\Code{1.3}{free(q);}
\Code{0.0}{\}}
\nopagebreak[4]
\Cost{\ensuremath{O\!\left(n\right)}}{\ensuremath{O\!\left(1\right)}}
\par\medskip

\Needspace{16\baselineskip}
\subsubsection{判断单链表是否存在环}
\Code{0.0}{Meet(first) \{}
\Hint{1.3}{先走再比较；初始重合不能证明有环。}\nopagebreak[4]
\Code{1.3}{slow = fast = first;}
\Code{1.3}{while (fast \&\& fast->next) \{}
\Code{2.6}{slow = slow->next;}
\Code{2.6}{fast = fast->next->next;}
\Hint{2.6}{fast 一次要走两次 next，循环条件必须同时保护 fast 和它的后继。}\nopagebreak[4]
\Code{2.6}{if (slow == fast) return slow;}
\Code{1.3}{\}}
\Code{1.3}{return NULL;}
\Code{0.0}{\}}
\nopagebreak[4]
\Cost{\ensuremath{O\!\left(n\right)}}{\ensuremath{O\!\left(1\right)}}
\par\medskip

\Needspace{17\baselineskip}
\subsubsection{查找单链表的环入口}
\Code{0.0}{CycleEntry(first) \{}
\Hint{1.3}{相遇点与首元各走一步，再次相遇处是入口。}\nopagebreak[4]
\Code{1.3}{q = Meet(first);}
\Code{1.3}{if (!q) return NULL;}
\Code{1.3}{p = first;}
\Code{1.3}{while (p != q) \{}
\Hint{2.6}{两个指针比较的是结点地址。入口恰在首元时，循环可以一次也不走。}\nopagebreak[4]
\Code{2.6}{p = p->next;}
\Code{2.6}{q = q->next;}
\Code{1.3}{\}}
\Code{1.3}{return p;}
\Code{0.0}{\}}
\nopagebreak[4]
\Cost{\ensuremath{O\!\left(n\right)}}{\ensuremath{O\!\left(1\right)}}
\par\medskip

\Needspace{18\baselineskip}
\subsubsection{求单链表中环的长度}
\Code{0.0}{CycleLength(first) \{}
\Hint{1.3}{固定相遇点，数回到它所需的步数。}\nopagebreak[4]
\Code{1.3}{p = Meet(first);}
\Code{1.3}{if (!p) return 0;}
\Code{1.3}{k = 1;}
\Code{1.3}{q = p->next;}
\Code{1.3}{while (q != p) \{}
\Hint{2.6}{长度从 1 起算，因为 q 已经从 p 沿环走了一步。}\nopagebreak[4]
\Code{2.6}{k++;}
\Code{2.6}{q = q->next;}
\Code{1.3}{\}}
\Code{1.3}{return k;}
\Code{0.0}{\}}
\nopagebreak[4]
\Cost{\ensuremath{O\!\left(n\right)}}{\ensuremath{O\!\left(1\right)}}
\par\medskip

\Needspace{13\baselineskip}
\subsubsection{判断两个单链表是否相交}
\Code{0.0}{Intersect(a, b) \{}
\Hint{1.3}{无环链表相交后共享尾部，比较地址。}\nopagebreak[4]
\Code{1.3}{if (!a || !b) return false;}
\Code{1.3}{while (a->next) a = a->next;}
\Hint{1.3}{尾结点的值相同不足以说明共享结点，需要比较尾指针本身。}\nopagebreak[4]
\Code{1.3}{while (b->next) b = b->next;}
\Code{1.3}{return a == b;}
\Code{0.0}{\}}
\nopagebreak[4]
\Cost{\ensuremath{O\!\left(m+n\right)}}{\ensuremath{O\!\left(1\right)}}
\par\medskip

\Needspace{16\baselineskip}
\subsubsection{寻找两个单链表的第一个公共结点}
\Code{0.0}{FirstCommon(a, b) \{}
\Hint{1.3}{各走完两条链，长度差被抵消。}\nopagebreak[4]
\Code{1.3}{p = a;}
\Code{1.3}{q = b;}
\Code{1.3}{while (p != q) \{}
\Code{2.6}{p = (p ? p->next : b);}
\Hint{2.6}{一侧走到 NULL 后从另一条链首继续；不相交时最终同时为空。}\nopagebreak[4]
\Code{2.6}{q = (q ? q->next : a);}
\Code{1.3}{\}}
\Code{1.3}{return p;}
\Code{0.0}{\}}
\nopagebreak[4]
\Cost{\ensuremath{O\!\left(m+n\right)}}{\ensuremath{O\!\left(1\right)}}
\par\medskip

\Needspace{26\baselineskip}
\subsubsection{判断单链表是否为回文链表}
\Code{0.0}{Palindrome(first) \{}
\Hint{1.3}{只逆右半；比较结束后接回原来的链。}\nopagebreak[4]
\Code{1.3}{if (!first || !first->next) return true;}
\Code{1.3}{s = f = first;}
\Code{1.3}{while (f->next \&\& f->next->next) \{}
\Code{2.6}{s = s->next;}
\Code{2.6}{f = f->next->next;}
\Code{1.3}{\}}
\Code{1.3}{r = ReverseNodes(s->next);}
\Code{1.3}{a = first;}
\Code{1.3}{b = r;}
\Hint{1.3}{b 遍历较短的右半链；奇数长度的中央结点不需要和自己比较。}\nopagebreak[4]
\Code{1.3}{ok = true;}
\Code{1.3}{while (b) \{}
\Code{2.6}{if (a->data != b->data) ok = false;}
\Code{2.6}{a = a->next;}
\Code{2.6}{b = b->next;}
\Code{1.3}{\}}
\Code{1.3}{s->next = ReverseNodes(r);}
\Code{1.3}{return ok;}
\Code{0.0}{\}}
\nopagebreak[4]
\Cost{\ensuremath{O\!\left(n\right)}}{\ensuremath{O\!\left(1\right)}}
\par\medskip

\Needspace{26\baselineskip}
\subsubsection{单链表的重排}
\Code{0.0}{Reorder(first) \{}
\Hint{1.3}{从左中点断开，将逆置的右半链交错插回。}\nopagebreak[4]
\Code{1.3}{if (!first) return;}
\Code{1.3}{s = f = first;}
\Code{1.3}{while (f->next \&\& f->next->next) \{}
\Code{2.6}{s = s->next;}
\Code{2.6}{f = f->next->next;}
\Code{1.3}{\}}
\Code{1.3}{b = ReverseNodes(s->next);}
\Code{1.3}{s->next = NULL;}
\Code{1.3}{a = first;}
\Hint{1.3}{x、y 各保存一侧未用后缀，重新连边后靠它们继续向后走。}\nopagebreak[4]
\Code{1.3}{while (b) \{}
\Code{2.6}{x = a->next;}
\Code{2.6}{y = b->next;}
\Code{2.6}{a->next = b;}
\Code{2.6}{b->next = x;}
\Code{2.6}{a = x;}
\Code{2.6}{b = y;}
\Code{1.3}{\}}
\Code{0.0}{\}}
\nopagebreak[4]
\Cost{\ensuremath{O\!\left(n\right)}}{\ensuremath{O\!\left(1\right)}}
\par\medskip

\Needspace{19\baselineskip}
\subsubsection{单链表的排序}
\Code{0.0}{InsertionSortNodes(first) \{}
\Hint{1.3}{在已排序链中越过相等项再插入，排序稳定。}\nopagebreak[4]
\Code{1.3}{D = ZeroNode();}
\Code{1.3}{while (first) \{}
\Code{2.6}{x = first;}
\Code{2.6}{first = x->next;}
\Code{2.6}{p = \&D;}
\Hint{2.6}{寻找插入前驱时跳过等值项，保证后读入的相等元素排在已有元素后面。}\nopagebreak[4]
\Code{2.6}{while (p->next \&\& p->next->data <= x->data) p = p->next;}
\Code{2.6}{x->next = p->next;}
\Code{2.6}{p->next = x;}
\Code{1.3}{\}}
\Code{1.3}{return D.next;}
\Code{0.0}{\}}
\nopagebreak[4]
\Cost{\ensuremath{O\!\left(n^{2}\right)}}{\ensuremath{O\!\left(1\right)}}
\par\medskip

\Needspace{19\baselineskip}
\subsubsection{单链表的归并排序}
\Code{0.0}{MergeSortNodes(a) \{}
\Hint{1.3}{使用左中点保证长度为 2 的链也能拆成两段。}\nopagebreak[4]
\Code{1.3}{if (!a || !a->next) return a;}
\Code{1.3}{s = a;}
\Code{1.3}{f = a->next;}
\Code{1.3}{while (f \&\& f->next) \{}
\Code{2.6}{s = s->next;}
\Hint{2.6}{切断 s 的后继后左右递归才真正变小，否则可能反复处理原链。}\nopagebreak[4]
\Code{2.6}{f = f->next->next;}
\Code{1.3}{\}}
\Code{1.3}{b = s->next;}
\Code{1.3}{s->next = NULL;}
\Code{1.3}{return MergeNodes(MergeSortNodes(a), MergeSortNodes(b));}
\Code{0.0}{\}}
\nopagebreak[4]
\Cost{\ensuremath{O\!\left(n \log (n+1)\right)}}{\ensuremath{O\!\left(\log (n+1)\right)}}
\par\medskip

\Needspace{27\baselineskip}
\subsubsection{带随机指针的复杂链表复制}
\Code{0.0}{CopyRandom(first) \{}
\Hint{1.3}{先把副本夹在原结点后面，随机指针目标的后继就是相应副本。}\nopagebreak[4]
\Code{1.3}{if (!first) return NULL;}
\Code{1.3}{for (p = first; p; ) \{}
\Code{2.6}{s = NewNode(p->data);}
\Code{2.6}{s->next = p->next;}
\Code{2.6}{p->next = s;}
\Code{2.6}{p = s->next;}
\Code{1.3}{\}}
\Code{1.3}{for (p = first; p; p = p->next->next) \{}
\Code{2.6}{p->next->random = (p->random ? p->random->next : NULL);}
\Code{1.3}{\}}
\Hint{1.3}{拆链时先还原原后继，再从下一原结点找到下一副本；两个尾部都要处理空值。}\nopagebreak[4]
\Hint{1.3}{此时原首元的后继就是整条副本链的首元，先保存结果。}\nopagebreak[4]
\Code{1.3}{copy = first->next;}
\Code{1.3}{for (p = first; p; ) \{}
\Code{2.6}{s = p->next;}
\Code{2.6}{p->next = s->next;}
\Code{2.6}{s->next = (p->next ? p->next->next : NULL);}
\Code{2.6}{p = p->next;}
\Code{1.3}{\}}
\Code{1.3}{return copy;}
\Code{0.0}{\}}
\nopagebreak[4]
\Cost{\ensuremath{O\!\left(n\right)}}{\ensuremath{O\!\left(1\right)}}
\par\medskip

\Part{栈}

\Group{顺序栈}

\Needspace{12\baselineskip}
\subsubsection{顺序栈的初始化}
\Code{0.0}{InitStack(S, A, C) \{}
\Hint{1.3}{top 为最后一个已用下标，空时为 −1。}\nopagebreak[4]
\Code{1.3}{S->a = A;}
\Hint{1.3}{空栈第一次入栈把 top 从 −1 变为 0，正好写入首槽。}\nopagebreak[4]
\Code{1.3}{S->cap = C;}
\Code{1.3}{S->top = -1;}
\Code{0.0}{\}}
\nopagebreak[4]
\Cost{\ensuremath{O\!\left(1\right)}}{\ensuremath{O\!\left(1\right)}}
\par\medskip

\Needspace{12\baselineskip}
\subsubsection{入栈操作}
\Code{0.0}{Push(S, x) \{}
\Hint{1.3}{先判满再增加 top，避免写到容量之外。}\nopagebreak[4]
\Code{1.3}{if (S->top+1 == S->cap) return false;}
\Hint{1.3}{top+1 等于容量时再加一就会越过数组末端。}\nopagebreak[4]
\Code{1.3}{S->a[++S->top] = x;}
\Code{1.3}{return true;}
\Code{0.0}{\}}
\nopagebreak[4]
\Cost{\ensuremath{O\!\left(1\right)}}{\ensuremath{O\!\left(1\right)}}
\par\medskip

\Needspace{12\baselineskip}
\subsubsection{出栈操作}
\Code{0.0}{Pop(S, out) \{}
\Hint{1.3}{输出参数接收原栈顶，随后栈顶回退。}\nopagebreak[4]
\Code{1.3}{if (S->top<0) return false;}
\Hint{1.3}{这里的后置减一先取原 top 对应元素，再缩短栈。}\nopagebreak[4]
\Code{1.3}{*out = S->a[S->top--];}
\Code{1.3}{return true;}
\Code{0.0}{\}}
\nopagebreak[4]
\Cost{\ensuremath{O\!\left(1\right)}}{\ensuremath{O\!\left(1\right)}}
\par\medskip

\Needspace{12\baselineskip}
\subsubsection{获取栈顶元素}
\Code{0.0}{Peek(S, out) \{}
\Hint{1.3}{读取数据时保留 top。}\nopagebreak[4]
\Code{1.3}{if (S->top<0) return false;}
\Hint{1.3}{out 是调用方的输出地址，写 *out 就能把值带回去。}\nopagebreak[4]
\Code{1.3}{*out = S->a[S->top];}
\Code{1.3}{return true;}
\Code{0.0}{\}}
\nopagebreak[4]
\Cost{\ensuremath{O\!\left(1\right)}}{\ensuremath{O\!\left(1\right)}}
\par\medskip

\Needspace{10\baselineskip}
\subsubsection{判断栈空与栈满}
\Code{0.0}{StackState(S) \{}
\Hint{1.3}{满栈的最后一个槽位是 cap−1。}\nopagebreak[4]
\Hint{1.3}{空和满是两个独立条件，需要分别判断。}\nopagebreak[4]
\Code{1.3}{return (S->top == -1, S->top == S->cap-1);}
\Code{0.0}{\}}
\nopagebreak[4]
\Cost{\ensuremath{O\!\left(1\right)}}{\ensuremath{O\!\left(1\right)}}
\par\medskip

\Needspace{22\baselineskip}
\subsubsection{双栈共享存储空间}
\Code{0.0}{SharedStack(S, side, op, x) \{}
\Hint{1.3}{两栈向中间生长，左右栈顶相邻就没有空槽。}\nopagebreak[4]
\Code{1.3}{if (op == PUSH) \{}
\Code{2.6}{require(S.l+1<S.r);}
\Code{2.6}{if (side == LEFT) S.a[++S.l] = x;}
\Code{2.6}{else S.a[--S.r] = x;}
\Code{1.3}{\}}
\Code{1.3}{else if (side == LEFT) \{}
\Code{2.6}{require(S.l >= 0);}
\Hint{2.6}{左栈向高下标增长，右栈向低下标增长；各自出栈条件不同。}\nopagebreak[4]
\Code{2.6}{return S.a[S.l--];}
\Code{1.3}{\}}
\Code{1.3}{else \{}
\Code{2.6}{require(S.r<S.cap);}
\Code{2.6}{return S.a[S.r++];}
\Code{1.3}{\}}
\Code{0.0}{\}}
\nopagebreak[4]
\Cost{\ensuremath{O\!\left(1\right)}/次}{\ensuremath{O\!\left(1\right)}}
\par\medskip

\Group{链栈}

\Needspace{10\baselineskip}
\subsubsection{链栈的初始化}
\Code{0.0}{InitLinkedStack(Node **top) \{}
\Hint{1.3}{调用方传入栈顶指针的地址。}\nopagebreak[4]
\Hint{1.3}{修改 *top 会改变调用方保存的栈顶；普通局部指针赋值没有这个效果。}\nopagebreak[4]
\Code{1.3}{*top = NULL;}
\Code{0.0}{\}}
\nopagebreak[4]
\Cost{\ensuremath{O\!\left(1\right)}}{\ensuremath{O\!\left(1\right)}}
\par\medskip

\Needspace{22\baselineskip}
\subsubsection{链栈的入栈与出栈}
\Code{0.0}{LinkedStack(Node **top, op, x) \{}
\Hint{1.3}{头插与摘除首元都只改一条链。}\nopagebreak[4]
\Code{1.3}{if (op == PUSH) \{}
\Code{2.6}{p = NewNode(x);}
\Code{2.6}{p->next = *top;}
\Code{2.6}{*top = p;}
\Code{1.3}{\}}
\Code{1.3}{else \{}
\Code{2.6}{require(*top);}
\Hint{2.6}{出栈先把 *top 改为后继，原结点释放后仍能找到其余元素。}\nopagebreak[4]
\Code{2.6}{p = *top;}
\Code{2.6}{*top = p->next;}
\Code{2.6}{x = p->data;}
\Code{2.6}{free(p);}
\Code{2.6}{return x;}
\Code{1.3}{\}}
\Code{0.0}{\}}
\nopagebreak[4]
\Cost{\ensuremath{O\!\left(1\right)}/次}{\ensuremath{O\!\left(1\right)}}
\par\medskip

\Needspace{14\baselineskip}
\subsubsection{链栈的销毁}
\Code{0.0}{DestroyStack(Node **top) \{}
\Hint{1.3}{先推进栈顶，再释放旧首元。}\nopagebreak[4]
\Code{1.3}{while (*top) \{}
\Code{2.6}{p = *top;}
\Hint{2.6}{只要顶部还不是 NULL，就仍有结点需要释放。}\nopagebreak[4]
\Code{2.6}{*top = p->next;}
\Code{2.6}{free(p);}
\Code{1.3}{\}}
\Code{0.0}{\}}
\nopagebreak[4]
\Cost{\ensuremath{O\!\left(n\right)}}{\ensuremath{O\!\left(1\right)}}
\par\medskip

\Group{栈的应用}

\Needspace{17\baselineskip}
\subsubsection{括号匹配}
\Code{0.0}{MatchBrackets(s) \{}
\Hint{1.3}{栈顶必须与当前右括号配对。}\nopagebreak[4]
\Code{1.3}{InitStack(\&S);}
\Code{1.3}{for (c in s) \{}
\Code{2.6}{if (IsLeft(c)) Push(\&S, c);}
\Code{2.6}{else if (IsRight(c)) \{}
\Hint{3.9}{遇右括号时若栈空，说明它没有可配对的左括号。}\nopagebreak[4]
\Code{3.9}{if (Empty(\&S) || !Pair(Pop(\&S), c)) return false;}
\Code{2.6}{\}}
\Code{1.3}{\}}
\Code{1.3}{return Empty(\&S);}
\Code{0.0}{\}}
\nopagebreak[4]
\Cost{\ensuremath{O\!\left(n\right)}}{\ensuremath{O\!\left(n\right)}}
\par\medskip

\Needspace{17\baselineskip}
\subsubsection{进制转换}
\Code{0.0}{ConvertBase(x, b) \{}
\Hint{1.3}{余数从低位向高位产生，出栈正好颠倒顺序。}\nopagebreak[4]
\Code{1.3}{require(x >= 0 \&\& b >= 2);}
\Code{1.3}{InitStack(\&S);}
\Code{1.3}{do \{}
\Code{2.6}{Push(\&S, x\%b);}
\Hint{2.6}{即使 x=0 也要输出一位，因此用至少执行一次的循环。}\nopagebreak[4]
\Code{2.6}{x = x/b;}
\Code{1.3}{\}}
\Code{1.3}{while (x>0);}
\Code{1.3}{while (!Empty(\&S)) EmitDigit(Pop(\&S));}
\Code{0.0}{\}}
\nopagebreak[4]
\Cost{\ensuremath{O\!\left(\log (x+2)\right)}}{\ensuremath{O\!\left(\log (x+2)\right)}}
\par\medskip

\Needspace{19\baselineskip}
\subsubsection{判断合法的出栈序列}
\Code{0.0}{ValidPop(in, out, n) \{}
\Hint{1.3}{假设入栈元素互不相同，两个序列均长 n。栈顶匹配下一个目标时立即出栈。}\nopagebreak[4]
\Code{1.3}{InitStack(\&S);}
\Code{1.3}{j = 0;}
\Code{1.3}{for (i = 0; i<n; i++) \{}
\Code{2.6}{Push(\&S, in[i]);}
\Code{2.6}{while (j<n \&\& !Empty(\&S) \&\& Top(\&S) == out[j]) \{}
\Hint{3.9}{只在匹配到 out[j] 时推进 j；目标序列中的次序不能跳过。}\nopagebreak[4]
\Code{3.9}{Pop(\&S);}
\Code{3.9}{j++;}
\Code{2.6}{\}}
\Code{1.3}{\}}
\Code{1.3}{return j == n;}
\Code{0.0}{\}}
\nopagebreak[4]
\Cost{\ensuremath{O\!\left(n\right)}}{\ensuremath{O\!\left(n\right)}}
\par\medskip

\Needspace{28\baselineskip}
\subsubsection{中缀表达式转换为后缀表达式}
\Code{0.0}{ToPostfix(tokens) \{}
\Hint{1.3}{输入为合法的中缀记号序列，运算符均为二元。}\nopagebreak[4]
\Code{1.3}{InitStack(\&S);}
\Code{1.3}{out = EmptyList();}
\Code{1.3}{for (t in tokens) \{}
\Code{2.6}{if (IsValue(t)) Append(\&out, t);}
\Code{2.6}{else if (t == '(') Push(\&S, t);}
\Code{2.6}{else if (t == ')') \{}
\Code{3.9}{while (Top(\&S) != '(') Append(\&out, Pop(\&S));}
\Code{3.9}{Pop(\&S);}
\Hint{3.9}{匹配的左括号直接丢弃。}\nopagebreak[4]
\Code{2.6}{\}}
\Code{2.6}{else \{}
\Code{3.9}{while (!Empty(\&S) \&\& Top(\&S) != '(' \&\& ShouldPop(Top(\&S), t)) \{}
\Code{5.2}{Append(\&out, Pop(\&S));}
\Code{3.9}{\}}
\Code{3.9}{Push(\&S, t);}
\Code{2.6}{\}}
\Code{1.3}{\}}
\Code{1.3}{while (!Empty(\&S)) Append(\&out, Pop(\&S));}
\Code{1.3}{return out;}
\Code{0.0}{\}}
\Code{0.0}{ShouldPop(top, incoming) \{}
\Code{1.3}{a = Priority(top);}
\Code{1.3}{b = Priority(incoming);}
\Hint{1.3}{优先级相同时，当前读到的运算符是左结合才弹栈。}\nopagebreak[4]
\Code{1.3}{return a>b || (a == b \&\& IsLeftAssoc(incoming));}
\Code{0.0}{\}}
\nopagebreak[4]
\Cost{\ensuremath{O\!\left(n\right)}}{\ensuremath{O\!\left(n\right)}}
\par\medskip

\Needspace{12\baselineskip}
\subsubsection{中缀表达式转换为前缀表达式}
\Code{0.0}{ToPrefix(tokens) \{}
\Hint{1.3}{后缀序列建出的运算符树，前序输出就是前缀式。}\nopagebreak[4]
\Code{1.3}{root = BuildInfixExpr(tokens);}
\Hint{1.3}{后缀中的操作数先成为叶，随后读到的运算符把两棵子树合成一棵。}\nopagebreak[4]
\Code{1.3}{PreOrder(root);}
\Code{1.3}{DestroyTree(root);}
\Code{0.0}{\}}
\nopagebreak[4]
\Cost{\ensuremath{O\!\left(n\right)}}{\ensuremath{O\!\left(n\right)}}
\par\medskip

\Needspace{19\baselineskip}
\subsubsection{后缀表达式求值}
\Code{0.0}{EvalPostfix(tokens) \{}
\Hint{1.3}{先弹出的是右操作数，减法和除法不能颠倒。}\nopagebreak[4]
\Code{1.3}{InitStack(\&S);}
\Code{1.3}{for (t in tokens) \{}
\Code{2.6}{if (IsValue(t)) Push(\&S, Value(t));}
\Code{2.6}{else \{}
\Code{3.9}{b = Pop(\&S);}
\Hint{3.9}{例如后缀 8 3 - 的第一次弹出是 3，运算顺序应为 8−3。}\nopagebreak[4]
\Code{3.9}{a = Pop(\&S);}
\Code{3.9}{Push(\&S, Apply(t, a, b));}
\Code{2.6}{\}}
\Code{1.3}{\}}
\Code{1.3}{return Pop(\&S);}
\Code{0.0}{\}}
\nopagebreak[4]
\Cost{\ensuremath{O\!\left(n\right)}}{\ensuremath{O\!\left(n\right)}}
\par\medskip

\Needspace{20\baselineskip}
\subsubsection{前缀表达式求值}
\Code{0.0}{EvalPrefix(tokens) \{}
\Hint{1.3}{从右向左扫描时，先弹出的值是左操作数。}\nopagebreak[4]
\Code{1.3}{InitStack(\&S);}
\Code{1.3}{for (i = len(tokens)-1; i >= 0; i--) \{}
\Code{2.6}{t = tokens[i];}
\Code{2.6}{if (IsValue(t)) Push(\&S, Value(t));}
\Code{2.6}{else \{}
\Code{3.9}{a = Pop(\&S);}
\Hint{3.9}{例如前缀 - 8 3 从右扫描，弹出的第一个数是 8。}\nopagebreak[4]
\Code{3.9}{b = Pop(\&S);}
\Code{3.9}{Push(\&S, Apply(t, a, b));}
\Code{2.6}{\}}
\Code{1.3}{\}}
\Code{1.3}{return Pop(\&S);}
\Code{0.0}{\}}
\nopagebreak[4]
\Cost{\ensuremath{O\!\left(n\right)}}{\ensuremath{O\!\left(n\right)}}
\par\medskip

\Needspace{31\baselineskip}
\subsubsection{中缀表达式直接求值（双栈法）}
\Code{0.0}{EvalInfix(tokens) \{}
\Hint{1.3}{Reduce 从值栈先取右值，再取左值，算完压回。}\nopagebreak[4]
\Code{1.3}{InitStack(\&V);}
\Code{1.3}{InitStack(\&S);}
\Code{1.3}{for (t in tokens) \{}
\Code{2.6}{if (IsValue(t)) Push(\&V, Value(t));}
\Code{2.6}{else if (t == '(') Push(\&S, t);}
\Code{2.6}{else if (t == ')') \{}
\Code{3.9}{while (Top(\&S) != '(') Reduce(\&S, \&V);}
\Code{3.9}{Pop(\&S);}
\Code{2.6}{\}}
\Code{2.6}{else \{}
\Code{3.9}{while (!Empty(\&S) \&\& Top(\&S) != '(' \&\& ShouldPop(Top(\&S), t)) Reduce(\&S, \&V);}
\Code{3.9}{Push(\&S, t);}
\Code{2.6}{\}}
\Code{1.3}{\}}
\Hint{1.3}{右括号要求算完括号内的全部运算，然后移除配对的左括号。}\nopagebreak[4]
\Code{1.3}{while (!Empty(\&S)) Reduce(\&S, \&V);}
\Code{1.3}{return Pop(\&V);}
\Code{0.0}{\}}
\Code{0.0}{Reduce(S, V) \{}
\Code{1.3}{op = Pop(S);}
\Code{1.3}{right = Pop(V);}
\Code{1.3}{left = Pop(V);}
\Code{1.3}{Push(V, Apply(op, left, right));}
\Code{0.0}{\}}
\nopagebreak[4]
\Cost{\ensuremath{O\!\left(n\right)}}{\ensuremath{O\!\left(n\right)}}
\par\medskip

\Needspace{21\baselineskip}
\subsubsection{利用栈模拟递归}
\Code{0.0}{SimPost(root) \{}
\Hint{1.3}{done 区分首次到达与两棵子树均已处理。}\nopagebreak[4]
\Code{1.3}{InitStack(\&S);}
\Code{1.3}{Push(\&S, (root, false));}
\Code{1.3}{while (!Empty(\&S)) \{}
\Code{2.6}{(p, done) = Pop(\&S);}
\Code{2.6}{if (!p) continue;}
\Code{2.6}{if (done) visit(p);}
\Hint{2.6}{栈后进先出，所以先压父结点的返回帧，再压右、左，接下来先处理左。}\nopagebreak[4]
\Code{2.6}{else \{}
\Code{3.9}{Push(\&S, (p, true));}
\Code{3.9}{Push(\&S, (p->rchild, false));}
\Code{3.9}{Push(\&S, (p->lchild, false));}
\Code{2.6}{\}}
\Code{1.3}{\}}
\Code{0.0}{\}}
\nopagebreak[4]
\Cost{\ensuremath{O\!\left(n\right)}}{\ensuremath{O\!\left(h\right)}}
\par\medskip

\Needspace{13\baselineskip}
\subsubsection{汉诺塔问题}
\Code{0.0}{Hanoi(n, from, via, to) \{}
\Hint{1.3}{把较小盘暂存到 via，才能移动最大盘。}\nopagebreak[4]
\Code{1.3}{if (n == 0) return;}
\Code{1.3}{Hanoi(n-1, from, to, via);}
\Hint{1.3}{第二次递归的起点已经变成 via，最大的盘不再移动。}\nopagebreak[4]
\Code{1.3}{EmitMove(from, to);}
\Code{1.3}{Hanoi(n-1, via, from, to);}
\Code{0.0}{\}}
\nopagebreak[4]
\Cost{\ensuremath{O\!\left(2^{n}\right)}}{\ensuremath{O\!\left(n\right)}}
\par\medskip

\Needspace{18\baselineskip}
\subsubsection{迷宫求解（回溯法）}
\Code{0.0}{MazeDFS(v) \{}
\Hint{1.3}{当前路径在尝试失败时撤销；已访问标记保留以免重复搜索。}\nopagebreak[4]
\Code{1.3}{if (!Open(v) || seen[v]) return false;}
\Code{1.3}{seen[v] = true;}
\Code{1.3}{Push(\&path, v);}
\Code{1.3}{if (v == target) return true;}
\Code{1.3}{for (w in Neighbors4(v)) \{}
\Hint{2.6}{若子递归找到出口，立即逐层返回，成功的路径就保留在 path 中。}\nopagebreak[4]
\Code{2.6}{if (MazeDFS(w)) return true;}
\Code{1.3}{\}}
\Code{1.3}{Pop(\&path);}
\Code{1.3}{return false;}
\Code{0.0}{\}}
\nopagebreak[4]
\Cost{\ensuremath{O\!\left(RC\right)}}{\ensuremath{O\!\left(RC\right)}}
\par\medskip

\Needspace{17\baselineskip}
\subsubsection{骑士巡游问题}
\Code{0.0}{Knight(v, k) \{}
\Hint{1.3}{step 为 0 表示可走，失败时撤销当前落点。}\nopagebreak[4]
\Code{1.3}{step[v] = k;}
\Code{1.3}{if (k == N) return true;}
\Code{1.3}{for (d = 0; d<8; d++) \{}
\Code{2.6}{w = Move(v, delta[d]);}
\Hint{2.6}{k 记录已经踩过的格数，而不是当前位置的棋盘坐标。}\nopagebreak[4]
\Code{2.6}{if (Inside(w) \&\& step[w] == 0 \&\& Knight(w, k+1)) return true;}
\Code{1.3}{\}}
\Code{1.3}{step[v] = 0;}
\Code{1.3}{return false;}
\Code{0.0}{\}}
\nopagebreak[4]
\Cost{\ensuremath{O\!\left(8^{N}\right)}}{\ensuremath{O\!\left(N\right)}}
\par\medskip

\Needspace{15\baselineskip}
\subsubsection{最小栈的实现}
\Code{0.0}{MinStack(S, op, x) \{}
\Hint{1.3}{每个栈项保存截至该层的最小值，出栈后自动恢复。}\nopagebreak[4]
\Code{1.3}{if (op == PUSH) \{}
\Code{2.6}{m = Empty(S)? x:min(x, Top(S).min);}
\Code{2.6}{Push(S, (x, m));}
\Code{1.3}{\}}
\Hint{1.3}{每次压入一对值，栈顶的 min 因而总与当前栈内容同步。}\nopagebreak[4]
\Code{1.3}{else if (op == MIN) return Top(S).min;}
\Code{1.3}{else return Pop(S).value;}
\Code{0.0}{\}}
\nopagebreak[4]
\Cost{\ensuremath{O\!\left(1\right)}/次}{\ensuremath{O\!\left(n\right)}（结构）}
\par\medskip

\Needspace{16\baselineskip}
\subsubsection{单调栈的实现与应用}
\Code{0.0}{NextGreater(A, n) \{}
\Hint{1.3}{被弹出的下标首次遇到右侧更大的数。}\nopagebreak[4]
\Code{1.3}{Fill(ans, -1);}
\Code{1.3}{InitStack(\&S);}
\Code{1.3}{for (i = 0; i<n; i++) \{}
\Code{2.6}{while (!Empty(\&S) \&\& A[Top(\&S)]<A[i]) ans[Pop(\&S)] = i;}
\Hint{2.6}{相等项不弹出，因为题目要求严格更大。}\nopagebreak[4]
\Code{2.6}{Push(\&S, i);}
\Code{1.3}{\}}
\Code{1.3}{return ans;}
\Code{0.0}{\}}
\nopagebreak[4]
\Cost{\ensuremath{O\!\left(n\right)}}{\ensuremath{O\!\left(n\right)}}
\par\medskip

\Needspace{16\baselineskip}
\subsubsection{寻找每个元素左侧或右侧第一个更大或更小的元素}
\Code{0.0}{Nearest(A, n, left, greater) \{}
\Hint{1.3}{寻找左邻从左往右扫，寻找右邻反过来。相等值须弹出。}\nopagebreak[4]
\Code{1.3}{InitStack(\&S);}
\Code{1.3}{for (i in ScanOrder(n, left)) \{}
\Code{2.6}{while (!Empty(\&S) \&\& (greater ? A[Top(\&S)] <= A[i] : A[Top(\&S)] >= A[i])) Pop(\&S);}
\Code{2.6}{ans[i] = Empty(\&S)?-1:Top(\&S);}
\Hint{2.6}{先弹掉不合格的候选，再读栈顶；剩下的是最近的合格位置。}\nopagebreak[4]
\Code{2.6}{Push(\&S, i);}
\Code{1.3}{\}}
\Code{1.3}{return ans;}
\Code{0.0}{\}}
\nopagebreak[4]
\Cost{\ensuremath{O\!\left(n\right)}}{\ensuremath{O\!\left(n\right)}}
\par\medskip

\Needspace{21\baselineskip}
\subsubsection{柱状图中的最大矩形面积}
\Code{0.0}{MaxRect(h, n) \{}
\Hint{1.3}{弹出柱的左右边界分别是新栈顶和当前下标，二者均不计入宽度。}\nopagebreak[4]
\Code{1.3}{InitStack(\&S);}
\Code{1.3}{Push(\&S, -1);}
\Code{1.3}{best = 0;}
\Code{1.3}{for (i = 0; i <= n; i++) \{}
\Code{2.6}{cur = (i == n ? -1:h[i]);}
\Code{2.6}{while (Top(\&S) != -1 \&\& h[Top(\&S)] >= cur) \{}
\Hint{3.9}{末尾虚拟高度为 −1，用来让还在栈里的柱也得到一次面积计算。}\nopagebreak[4]
\Code{3.9}{j = Pop(\&S);}
\Code{3.9}{best = max(best, h[j]*(i-Top(\&S)-1));}
\Code{2.6}{\}}
\Code{2.6}{Push(\&S, i);}
\Code{1.3}{\}}
\Code{1.3}{return best;}
\Code{0.0}{\}}
\nopagebreak[4]
\Cost{\ensuremath{O\!\left(n\right)}}{\ensuremath{O\!\left(n\right)}}
\par\medskip

\Needspace{21\baselineskip}
\subsubsection{最长有效括号子串}
\Code{0.0}{LongestValid(s) \{}
\Hint{1.3}{栈底存最近一次无法匹配的右括号位置，初值 −1。}\nopagebreak[4]
\Code{1.3}{InitStack(\&S);}
\Code{1.3}{Push(\&S, -1);}
\Code{1.3}{best = 0;}
\Code{1.3}{for (i = 0; i<len(s); i++) \{}
\Code{2.6}{if (s[i] == '(') Push(\&S, i);}
\Code{2.6}{else \{}
\Hint{3.9}{匹配右括号后若栈未空，用当前下标减栈顶得当前有效后缀长度。}\nopagebreak[4]
\Code{3.9}{Pop(\&S);}
\Code{3.9}{if (Empty(\&S)) Push(\&S, i);}
\Code{3.9}{else best = max(best, i-Top(\&S));}
\Code{2.6}{\}}
\Code{1.3}{\}}
\Code{1.3}{return best;}
\Code{0.0}{\}}
\nopagebreak[4]
\Cost{\ensuremath{O\!\left(n\right)}}{\ensuremath{O\!\left(n\right)}}
\par\medskip

\Part{队列}

\Group{顺序队列}

\Needspace{12\baselineskip}
\subsubsection{顺序队列的初始化}
\Code{0.0}{InitQueue(Q, A, C) \{}
\Hint{1.3}{半开区间 [front,rear) 存放有效元素。}\nopagebreak[4]
\Code{1.3}{Q->a = A;}
\Hint{1.3}{刚建立时两个下标都为 0，数组还没有可读的数据。}\nopagebreak[4]
\Code{1.3}{Q->cap = C;}
\Code{1.3}{Q->front = Q->rear = 0;}
\Code{0.0}{\}}
\nopagebreak[4]
\Cost{\ensuremath{O\!\left(1\right)}}{\ensuremath{O\!\left(1\right)}}
\par\medskip

\Needspace{17\baselineskip}
\subsubsection{入队与出队}
\Code{0.0}{LinearQueue(Q, op, x) \{}
\Hint{1.3}{rear 指向下一写入槽，front 指向队头。}\nopagebreak[4]
\Code{1.3}{if (op == IN) \{}
\Code{2.6}{require(Q->rear<Q->cap);}
\Code{2.6}{Q->a[Q->rear++] = x;}
\Code{1.3}{\}}
\Hint{1.3}{出队只推进 front，不搬动数组中剩下的元素。}\nopagebreak[4]
\Code{1.3}{else \{}
\Code{2.6}{require(Q->front<Q->rear);}
\Code{2.6}{return Q->a[Q->front++];}
\Code{1.3}{\}}
\Code{0.0}{\}}
\nopagebreak[4]
\Cost{\ensuremath{O\!\left(1\right)}/次}{\ensuremath{O\!\left(1\right)}}
\par\medskip

\Needspace{11\baselineskip}
\subsubsection{获取队头元素}
\Code{0.0}{QueueFront(Q) \{}
\Hint{1.3}{只读取队头，不推进读指针。}\nopagebreak[4]
\Code{1.3}{require(Q->front<Q->rear);}
\Hint{1.3}{front 与 rear 相等时没有首元，因此要先做非空检查。}\nopagebreak[4]
\Code{1.3}{return Q->a[Q->front];}
\Code{0.0}{\}}
\nopagebreak[4]
\Cost{\ensuremath{O\!\left(1\right)}}{\ensuremath{O\!\left(1\right)}}
\par\medskip

\Needspace{14\baselineskip}
\subsubsection{顺序队列的假溢出处理}
\Code{0.0}{CompactQueue(Q) \{}
\Hint{1.3}{队尾到达数组末端时，把仍有效的区间移到开头。}\nopagebreak[4]
\Code{1.3}{if (Q->rear != Q->cap || Q->front == 0) return;}
\Code{1.3}{n = Q->rear-Q->front;}
\Hint{1.3}{从低下标向高下标复制时，待复制的源位置尚未被覆盖。}\nopagebreak[4]
\Code{1.3}{for (i = 0; i<n; i++) Q->a[i] = Q->a[Q->front+i];}
\Code{1.3}{Q->front = 0;}
\Code{1.3}{Q->rear = n;}
\Code{0.0}{\}}
\nopagebreak[4]
\Cost{\ensuremath{O\!\left(n\right)}}{\ensuremath{O\!\left(1\right)}}
\par\medskip

\Group{循环队列}

\Needspace{13\baselineskip}
\subsubsection{循环队列的初始化}
\Code{0.0}{InitRingQ(Q, A, C) \{}
\Hint{1.3}{保留一个空槽，容量 C 的数组最多存 C−1 项。}\nopagebreak[4]
\Code{1.3}{require(C >= 2);}
\Code{1.3}{Q->a = A;}
\Hint{1.3}{这里 front 指队头，rear 指下一次写入位置，两者取模使用。}\nopagebreak[4]
\Code{1.3}{Q->cap = C;}
\Code{1.3}{Q->front = Q->rear = 0;}
\Code{0.0}{\}}
\nopagebreak[4]
\Cost{\ensuremath{O\!\left(1\right)}}{\ensuremath{O\!\left(1\right)}}
\par\medskip

\Needspace{22\baselineskip}
\subsubsection{循环队列的入队与出队}
\Code{0.0}{RingQ(Q, op, x) \{}
\Hint{1.3}{取模使越过末槽的下标回到 0。}\nopagebreak[4]
\Code{1.3}{C = Q->cap;}
\Code{1.3}{if (op == IN) \{}
\Code{2.6}{next = (Q->rear+1)\%C;}
\Code{2.6}{require(next != Q->front);}
\Code{2.6}{Q->a[Q->rear] = x;}
\Code{2.6}{Q->rear = next;}
\Code{1.3}{\}}
\Hint{1.3}{元素写入 rear 当前槽，然后 rear 才移动到下一槽。}\nopagebreak[4]
\Code{1.3}{else \{}
\Code{2.6}{require(Q->front != Q->rear);}
\Code{2.6}{x = Q->a[Q->front];}
\Code{2.6}{Q->front = (Q->front+1)\%C;}
\Code{2.6}{return x;}
\Code{1.3}{\}}
\Code{0.0}{\}}
\nopagebreak[4]
\Cost{\ensuremath{O\!\left(1\right)}/次}{\ensuremath{O\!\left(1\right)}}
\par\medskip

\Needspace{12\baselineskip}
\subsubsection{循环队列的判空与判满}
\Code{0.0}{RingState(Q) \{}
\Hint{1.3}{队尾再向前一步追上队头时为满。}\nopagebreak[4]
\Code{1.3}{empty = (Q->front == Q->rear);}
\Hint{1.3}{若不保留空槽，只靠 front 与 rear 相等就不能区分满和空。}\nopagebreak[4]
\Code{1.3}{full = ((Q->rear+1)\%Q->cap == Q->front);}
\Code{1.3}{return (empty, full);}
\Code{0.0}{\}}
\nopagebreak[4]
\Cost{\ensuremath{O\!\left(1\right)}}{\ensuremath{O\!\left(1\right)}}
\par\medskip

\Needspace{10\baselineskip}
\subsubsection{循环队列的长度计算}
\Code{0.0}{RingLength(Q) \{}
\Hint{1.3}{加容量后再取模，兼容队尾绕回起点。}\nopagebreak[4]
\Hint{1.3}{当 rear 小于 front，队中数据分布在数组末尾和开头两段。}\nopagebreak[4]
\Code{1.3}{return (Q->rear-Q->front+Q->cap)\%Q->cap;}
\Code{0.0}{\}}
\nopagebreak[4]
\Cost{\ensuremath{O\!\left(1\right)}}{\ensuremath{O\!\left(1\right)}}
\par\medskip

\Needspace{16\baselineskip}
\subsubsection{利用计数器或标志位区分队空与队满}
\Code{0.0}{FullCapacity(Q, useCount) \{}
\Hint{1.3}{计数或记录最近成功操作，均可使用全部 C 个槽位。}\nopagebreak[4]
\Code{1.3}{if (useCount) return (Q->size == 0, Q->size == Q->cap);}
\Code{1.3}{same = (Q->rear == Q->front);}
\Code{1.3}{return (same \&\& Q->last == OUT, same \&\& Q->last == IN);}
\Code{0.0}{\}}
\Hint{0.0}{初始 size=0、last=OUT。成功入队后 size++、last=IN；成功出队后反向更新。}\nopagebreak[4]
\Code{0.0}{AfterSuccess(Q, op) \{}
\Code{1.3}{Q->size += (op == IN ? 1:-1);}
\Code{1.3}{Q->last = op;}
\Code{0.0}{\}}
\nopagebreak[4]
\Cost{\ensuremath{O\!\left(1\right)}/次}{\ensuremath{O\!\left(1\right)}}
\par\medskip

\Group{链式队列}

\Needspace{17\baselineskip}
\subsubsection{链式队列的建立}
\Code{0.0}{BuildLinkedQ(input) \{}
\Hint{1.3}{空队列的尾指针指向哨兵。}\nopagebreak[4]
\Code{1.3}{Q.head = NewNode();}
\Code{1.3}{Q.tail = Q.head;}
\Code{1.3}{for (x in input) \{}
\Code{2.6}{p = NewNode(x);}
\Hint{2.6}{建队时 t 指最后一项，因此一次尾插无需遍历整条链。}\nopagebreak[4]
\Code{2.6}{Q.tail->next = p;}
\Code{2.6}{Q.tail = p;}
\Code{1.3}{\}}
\Code{1.3}{return Q;}
\Code{0.0}{\}}
\nopagebreak[4]
\Cost{\ensuremath{O\!\left(n\right)}}{\ensuremath{O\!\left(1\right)}}
\par\medskip

\Needspace{23\baselineskip}
\subsubsection{链式队列的入队与出队}
\Code{0.0}{LinkedQ(Q, op, x) \{}
\Hint{1.3}{删除最后一项后，尾指针必须回到哨兵。}\nopagebreak[4]
\Code{1.3}{if (op == IN) \{}
\Code{2.6}{p = NewNode(x);}
\Code{2.6}{Q->tail->next = p;}
\Code{2.6}{Q->tail = p;}
\Code{1.3}{\}}
\Code{1.3}{else \{}
\Code{2.6}{p = Q->head->next;}
\Code{2.6}{require(p);}
\Hint{2.6}{如果不重设尾指针，下一次入队将向刚释放的内存写数据。}\nopagebreak[4]
\Code{2.6}{Q->head->next = p->next;}
\Code{2.6}{if (Q->tail == p) Q->tail = Q->head;}
\Code{2.6}{x = p->data;}
\Code{2.6}{free(p);}
\Code{2.6}{return x;}
\Code{1.3}{\}}
\Code{0.0}{\}}
\nopagebreak[4]
\Cost{\ensuremath{O\!\left(1\right)}/次}{\ensuremath{O\!\left(1\right)}}
\par\medskip

\Needspace{12\baselineskip}
\subsubsection{链式队列的销毁}
\Code{0.0}{DestroyQ(Q) \{}
\Hint{1.3}{逐个释放数据结点，最后处理哨兵。}\nopagebreak[4]
\Code{1.3}{while (Q->head->next) LinkedQ(Q, OUT, 0);}
\Hint{1.3}{释放哨兵后两个外部指针都应清空。}\nopagebreak[4]
\Code{1.3}{free(Q->head);}
\Code{1.3}{Q->head = Q->tail = NULL;}
\Code{0.0}{\}}
\nopagebreak[4]
\Cost{\ensuremath{O\!\left(n\right)}}{\ensuremath{O\!\left(1\right)}}
\par\medskip

\Needspace{24\baselineskip}
\subsubsection{不带头结点的链式队列实现}
\Code{0.0}{BareQ(Q, op, x) \{}
\Hint{1.3}{没有哨兵时，首次入队和最后一次出队同时影响两端。}\nopagebreak[4]
\Code{1.3}{if (op == IN) \{}
\Code{2.6}{p = NewNode(x);}
\Code{2.6}{if (Q->tail) Q->tail->next = p;}
\Code{2.6}{else Q->front = p;}
\Code{2.6}{Q->tail = p;}
\Code{1.3}{\}}
\Code{1.3}{else \{}
\Code{2.6}{p = Q->front;}
\Hint{2.6}{插入空队时，新结点同时成为首元和尾元。}\nopagebreak[4]
\Code{2.6}{require(p);}
\Code{2.6}{Q->front = p->next;}
\Code{2.6}{if (!Q->front) Q->tail = NULL;}
\Code{2.6}{x = p->data;}
\Code{2.6}{free(p);}
\Code{2.6}{return x;}
\Code{1.3}{\}}
\Code{0.0}{\}}
\nopagebreak[4]
\Cost{\ensuremath{O\!\left(1\right)}/次}{\ensuremath{O\!\left(1\right)}}
\par\medskip

\Group{双端队列}

\Needspace{22\baselineskip}
\subsubsection{队头插入与删除}
\Code{0.0}{DequeFront(D, op, x) \{}
\Hint{1.3}{头插先退一个槽，头删先读再前进。}\nopagebreak[4]
\Code{1.3}{if (op == IN) \{}
\Code{2.6}{require(D.size<C);}
\Code{2.6}{D.f = (D.f-1+C)\%C;}
\Code{2.6}{D.a[D.f] = x;}
\Code{2.6}{D.size++;}
\Code{1.3}{\}}
\Code{1.3}{else \{}
\Hint{2.6}{减一前加 C，是为了从下标 0 合法绕到最后一格。}\nopagebreak[4]
\Code{2.6}{require(D.size>0);}
\Code{2.6}{x = D.a[D.f];}
\Code{2.6}{D.f = (D.f+1)\%C;}
\Code{2.6}{D.size--;}
\Code{2.6}{return x;}
\Code{1.3}{\}}
\Code{0.0}{\}}
\nopagebreak[4]
\Cost{\ensuremath{O\!\left(1\right)}/次}{\ensuremath{O\!\left(1\right)}}
\par\medskip

\Needspace{21\baselineskip}
\subsubsection{队尾插入与删除}
\Code{0.0}{DequeBack(D, op, x) \{}
\Hint{1.3}{r 指向队尾元素之后，尾删先退一步。}\nopagebreak[4]
\Code{1.3}{if (op == IN) \{}
\Code{2.6}{require(D.size<C);}
\Code{2.6}{D.a[D.r] = x;}
\Code{2.6}{D.r = (D.r+1)\%C;}
\Code{2.6}{D.size++;}
\Code{1.3}{\}}
\Hint{1.3}{尾删后 r 正好停在刚刚释放出来的写入槽位。}\nopagebreak[4]
\Code{1.3}{else \{}
\Code{2.6}{require(D.size>0);}
\Code{2.6}{D.r = (D.r-1+C)\%C;}
\Code{2.6}{D.size--;}
\Code{2.6}{return D.a[D.r];}
\Code{1.3}{\}}
\Code{0.0}{\}}
\nopagebreak[4]
\Cost{\ensuremath{O\!\left(1\right)}/次}{\ensuremath{O\!\left(1\right)}}
\par\medskip

\Needspace{25\baselineskip}
\subsubsection{判断受限双端队列的合法输出序列}
\Code{0.0}{ValidDeque(i, j, D) \{}
\Hint{1.3}{试探失败须还原端点；左右端权限按题设传入。}\nopagebreak[4]
\Code{1.3}{if (j == n) return true;}
\Code{1.3}{if (i<n) \{}
\Code{2.6}{for (e in inputEnds) \{}
\Code{3.9}{InsertEnd(D, e, in[i]);}
\Code{3.9}{if (ValidDeque(i+1, j, D)) return true;}
\Code{3.9}{RemoveEnd(D, e);}
\Code{2.6}{\}}
\Code{1.3}{\}}
\Code{1.3}{for (e in outputEnds) \{}
\Hint{2.6}{一次递归要么消费一个输入项，要么匹配一个目标项；失败操作必须对称撤销。}\nopagebreak[4]
\Code{2.6}{if (!Empty(D) \&\& PeekEnd(D, e) == out[j]) \{}
\Code{3.9}{x = RemoveEnd(D, e);}
\Code{3.9}{if (ValidDeque(i, j+1, D)) return true;}
\Code{3.9}{InsertEnd(D, e, x);}
\Code{2.6}{\}}
\Code{1.3}{\}}
\Code{1.3}{return false;}
\Code{0.0}{\}}
\nopagebreak[4]
\Cost{\ensuremath{O\!\left(4^{2n}\right)}（回溯上界）}{\ensuremath{O\!\left(n\right)}}
\par\medskip

\Needspace{16\baselineskip}
\subsubsection{单调队列的实现与应用}
\Code{0.0}{WindowMax(A, n, k) \{}
\Hint{1.3}{首端移除过期项，尾端移除以后不可能成为最大值的项。}\nopagebreak[4]
\Code{1.3}{require(1 <= k \&\& k <= n);}
\Code{1.3}{InitDeque(\&D);}
\Code{1.3}{for (i = 0; i<n; i++) \{}
\Code{2.6}{while (!Empty(\&D) \&\& Front(\&D) <= i-k) PopFront(\&D);}
\Code{2.6}{while (!Empty(\&D) \&\& A[Back(\&D)] <= A[i]) PopBack(\&D);}
\Hint{2.6}{队列存下标，才能同时判断该项是否过期以及对应数值的大小。}\nopagebreak[4]
\Code{2.6}{PushBack(\&D, i);}
\Code{2.6}{if (i >= k-1) Emit(A[Front(\&D)]);}
\Code{1.3}{\}}
\Code{0.0}{\}}
\nopagebreak[4]
\Cost{\ensuremath{O\!\left(n\right)}}{\ensuremath{O\!\left(k\right)}}
\par\medskip

\Group{队列的应用}

\Needspace{23\baselineskip}
\subsubsection{利用队列打印杨辉三角}
\Code{0.0}{Pascal(rows) \{}
\Hint{1.3}{只处理旧行长度个数，新写入的数留给下一轮。}\nopagebreak[4]
\Code{1.3}{InitQueue(\&Q);}
\Code{1.3}{EnQueue(\&Q, 1);}
\Code{1.3}{repeat (rows) \{}
\Code{2.6}{k = Size(\&Q);}
\Code{2.6}{prev = 0;}
\Code{2.6}{repeat (k) \{}
\Code{3.9}{x = DeQueue(\&Q);}
\Code{3.9}{Emit(x);}
\Hint{3.9}{prev 保存上一项，刚出队的 x 与它相加得到下一行的一项。}\nopagebreak[4]
\Code{3.9}{EnQueue(\&Q, prev+x);}
\Code{3.9}{prev = x;}
\Code{2.6}{\}}
\Code{2.6}{EnQueue(\&Q, prev);}
\Code{2.6}{NewLine();}
\Code{1.3}{\}}
\Code{0.0}{\}}
\nopagebreak[4]
\Cost{\ensuremath{O\!\left(r^{2}\right)}}{\ensuremath{O\!\left(r\right)}}
\par\medskip

\Needspace{31\baselineskip}
\subsubsection{利用队列进行车厢重排}
\Code{0.0}{Rail(in, k) \{}
\Hint{1.3}{优先放到尾号小于 x 且最大的缓冲轨，空轨留给无处可放的车。}\nopagebreak[4]
\Code{1.3}{need = 1;}
\Code{1.3}{for (x in in) \{}
\Code{2.6}{if (x == need) Emit(need++);}
\Code{2.6}{else \{}
\Code{3.9}{t = BestTailLessThan(x);}
\Code{3.9}{if (t == NONE) t = EmptyTrack();}
\Code{3.9}{if (t == NONE) return false;}
\Code{3.9}{EnQueue(\&track[t], x);}
\Code{2.6}{\}}
\Code{2.6}{while ((t = TrackWithFront(need)) != NONE) \{}
\Code{3.9}{DeQueue(\&track[t]);}
\Code{3.9}{Emit(need++);}
\Code{2.6}{\}}
\Code{1.3}{\}}
\Hint{1.3}{每次 need 增大后要反复检查缓冲轨，因为可能有一串车已经可以出轨。}\nopagebreak[4]
\Code{1.3}{return need == len(in)+1;}
\Code{0.0}{\}}
\Code{0.0}{BestTailLessThan(x) \{}
\Code{1.3}{best = NONE;}
\Code{1.3}{for (t = 0; t<k; t++) \{}
\Code{2.6}{if (!Empty(\&track[t]) \&\& Back(\&track[t])<x \&\& (best == NONE || Back(\&track[t])>Back(\&track[best]))) best = t;}
\Code{1.3}{\}}
\Code{1.3}{return best;}
\Code{0.0}{\}}
\nopagebreak[4]
\Cost{\ensuremath{O\!\left(n(k+1)\right)}}{\ensuremath{O\!\left(n+k\right)}}
\par\medskip

\Needspace{18\baselineskip}
\subsubsection{利用队列实现任务调度}
\Code{0.0}{FCFS(job, n) \{}
\Hint{1.3}{任务已按到达时刻排序。CPU 空闲时跳到下一任务到达时刻。}\nopagebreak[4]
\Code{1.3}{i = time = 0;}
\Code{1.3}{InitQueue(\&Q);}
\Code{1.3}{while (i<n || !Empty(\&Q)) \{}
\Code{2.6}{if (Empty(\&Q)) time = max(time, job[i].arrival);}
\Code{2.6}{while (i<n \&\& job[i].arrival <= time) EnQueue(\&Q, i++);}
\Hint{2.6}{结束时刻等于本次任务开始时刻加运行时长；等待中的任务保持先来先服务。}\nopagebreak[4]
\Code{2.6}{j = DeQueue(\&Q);}
\Code{2.6}{time += job[j].work;}
\Code{2.6}{finish[j] = time;}
\Code{1.3}{\}}
\Code{0.0}{\}}
\nopagebreak[4]
\Cost{\ensuremath{O\!\left(n\right)}}{\ensuremath{O\!\left(n\right)}}
\par\medskip

\Needspace{22\baselineskip}
\subsubsection{利用队列实现广度优先搜索}
\Code{0.0}{BFS(G, s) \{}
\Hint{1.3}{入队时标记，每个点至多进队一次。}\nopagebreak[4]
\Code{1.3}{InitQueue(\&Q);}
\Code{1.3}{seen[s] = true;}
\Code{1.3}{EnQueue(\&Q, s);}
\Code{1.3}{while (!Empty(\&Q)) \{}
\Code{2.6}{u = DeQueue(\&Q);}
\Code{2.6}{visit(u);}
\Code{2.6}{for (v in adj[u]) \{}
\Hint{3.9}{若出队后才首次标记，同一顶点可能被多个父顶点重复入队。}\nopagebreak[4]
\Code{3.9}{if (!seen[v]) \{}
\Code{5.2}{seen[v] = true;}
\Code{5.2}{EnQueue(\&Q, v);}
\Code{3.9}{\}}
\Code{2.6}{\}}
\Code{1.3}{\}}
\Code{0.0}{\}}
\nopagebreak[4]
\Cost{\ensuremath{O\!\left(V+E\right)}}{\ensuremath{O\!\left(V\right)}}
\par\medskip

\Needspace{8\baselineskip}
\subsubsection{利用队列求解迷宫最短路径}
\Code{0.0}{MazeShortest(start, end) \{}
\Hint{1.3}{pre 同时记录是否发现及最短路上的前一格。}\nopagebreak[4]
\Code{1.3}{Fill(pre, NONE);}
\Code{1.3}{pre[start] = start;}
\Code{1.3}{EnQueue(\&Q, start);}
\Code{1.3}{while (!Empty(\&Q)) \{}
\Code{2.6}{u = DeQueue(\&Q);}
\Code{2.6}{if (u == end) return RestoreGrid(pre, start, end);}
\Code{2.6}{for (v in Neighbors4(u)) \{}
\Code{3.9}{if (Open(v) \&\& pre[v] == NONE) \{}
\Code{5.2}{pre[v] = u;}
\Code{5.2}{EnQueue(\&Q, v);}
\Code{3.9}{\}}
\Code{2.6}{\}}
\Code{1.3}{\}}
\Hint{1.3}{从出口沿 pre 回溯，最后把逆序格子序列反转即可得到路线。}\nopagebreak[4]
\Code{1.3}{return FAIL;}
\Code{0.0}{\}}
\Code{0.0}{RestoreGrid(pre, s, t) \{}
\Hint{1.3}{调用前须已找到终点；返回含起点和终点的格子序列。}\nopagebreak[4]
\Code{1.3}{path = EmptyArray();}
\Code{1.3}{for (p = t; p != s; p = pre[p]) Append(\&path, p);}
\Code{1.3}{Append(\&path, s);}
\Code{1.3}{Reverse(path, len(path));}
\Code{1.3}{return path;}
\Code{0.0}{\}}
\nopagebreak[4]
\Cost{\ensuremath{O\!\left(RC\right)}}{\ensuremath{O\!\left(RC\right)}}
\par\medskip

\Needspace{23\baselineskip}
\subsubsection{循环队列实现时间片轮转调度}
\Code{0.0}{RoundRobin(job, n, q) \{}
\Hint{1.3}{job 按到达时间非降序排列，初始剩余运行时间非负。时间片结束时先接纳新到任务，再让当前任务回队。}\nopagebreak[4]
\Code{1.3}{require(q>0);}
\Code{1.3}{i = time = 0;}
\Code{1.3}{queueStorage = AllocElems(n+1);}
\Code{1.3}{InitRingQ(\&Q, queueStorage, n+1);}
\Code{1.3}{while (i<n || !Empty(\&Q)) \{}
\Code{2.6}{if (Empty(\&Q)) time = max(time, job[i].arrival);}
\Code{2.6}{while (i<n \&\& job[i].arrival <= time) EnQueue(\&Q, i++);}
\Code{2.6}{j = DeQueue(\&Q);}
\Code{2.6}{d = min(q, job[j].left);}
\Hint{2.6}{剩余时间小于一个时间片时只运行实际剩余量，完成任务不再回队。}\nopagebreak[4]
\Code{2.6}{time += d;}
\Code{2.6}{job[j].left -= d;}
\Code{2.6}{while (i<n \&\& job[i].arrival <= time) EnQueue(\&Q, i++);}
\Code{2.6}{if (job[j].left>0) EnQueue(\&Q, j);}
\Code{2.6}{else finish[j] = time;}
\Code{1.3}{\}}
\Code{0.0}{\}}
\nopagebreak[4]
\Cost{\ensuremath{O\!\left(n+s\right)}（$s$ 为调度轮数）}{\ensuremath{O\!\left(n\right)}}
\par\medskip

\Part{前四部分的一些综合算法}

\Group{数据结构的相互模拟}

\Needspace{17\baselineskip}
\subsubsection{利用两个栈实现队列}
\Code{0.0}{TwoStackQ(op, x) \{}
\Hint{1.3}{只有输出栈为空时才整体倒栈。}\nopagebreak[4]
\Code{1.3}{if (op == IN) Push(\&in, x);}
\Code{1.3}{else \{}
\Code{2.6}{if (Empty(\&out)) \{}
\Code{3.9}{while (!Empty(\&in)) Push(\&out, Pop(\&in));}
\Code{2.6}{\}}
\Hint{2.6}{两栈都空才是真的空队；倒栈结束后原来最早入队的项恰好在 out 顶部。}\nopagebreak[4]
\Code{2.6}{require(!Empty(\&out));}
\Code{2.6}{return Pop(\&out);}
\Code{1.3}{\}}
\Code{0.0}{\}}
\nopagebreak[4]
\Cost{均摊 \ensuremath{O\!\left(1\right)}，单次最坏 \ensuremath{O\!\left(n\right)}}{\ensuremath{O\!\left(n\right)}}
\par\medskip

\Needspace{15\baselineskip}
\subsubsection{利用两个队列实现栈}
\Code{0.0}{TwoQueueStack(op, x) \{}
\Hint{1.3}{新元素入空队，再把旧元素排到它后面。}\nopagebreak[4]
\Code{1.3}{if (op == PUSH) \{}
\Code{2.6}{EnQueue(\&b, x);}
\Code{2.6}{while (!Empty(\&a)) EnQueue(\&b, DeQueue(\&a));}
\Hint{2.6}{交换的是两个队列的角色，下一次出栈始终从主队队头取。}\nopagebreak[4]
\Code{2.6}{SwapQueues(\&a, \&b);}
\Code{1.3}{\}}
\Code{1.3}{else return DeQueue(\&a);}
\Code{0.0}{\}}
\nopagebreak[4]
\Cost{入 \ensuremath{O\!\left(n\right)}，出 \ensuremath{O\!\left(1\right)}}{\ensuremath{O\!\left(n\right)}}
\par\medskip

\Needspace{15\baselineskip}
\subsubsection{利用一个队列实现栈}
\Code{0.0}{OneQueueStack(op, x) \{}
\Hint{1.3}{轮转入队前已有的元素，使新元素成为队头。}\nopagebreak[4]
\Code{1.3}{if (op == PUSH) \{}
\Code{2.6}{k = Size(\&q);}
\Code{2.6}{EnQueue(\&q, x);}
\Hint{2.6}{先记录入队前长度，恰好转完旧元素，刚入队的 x 就留在最前面。}\nopagebreak[4]
\Code{2.6}{repeat (k) EnQueue(\&q, DeQueue(\&q));}
\Code{1.3}{\}}
\Code{1.3}{else return DeQueue(\&q);}
\Code{0.0}{\}}
\nopagebreak[4]
\Cost{入 \ensuremath{O\!\left(n\right)}，出 \ensuremath{O\!\left(1\right)}}{\ensuremath{O\!\left(n\right)}}
\par\medskip

\Needspace{12\baselineskip}
\subsubsection{利用栈实现队列逆置}
\Code{0.0}{ReverseQueue(Q) \{}
\Hint{1.3}{经过栈一次，取出次序恰好反转。}\nopagebreak[4]
\Code{1.3}{InitStack(\&S);}
\Hint{1.3}{第一轮把队首次序变成栈底到栈顶的次序，第二轮从顶到底取回。}\nopagebreak[4]
\Code{1.3}{while (!Empty(Q)) Push(\&S, DeQueue(Q));}
\Code{1.3}{while (!Empty(\&S)) EnQueue(Q, Pop(\&S));}
\Code{0.0}{\}}
\nopagebreak[4]
\Cost{\ensuremath{O\!\left(n\right)}}{\ensuremath{O\!\left(n\right)}}
\par\medskip

\Needspace{26\baselineskip}
\subsubsection{利用队列实现栈的排序}
\Code{0.0}{SortStack(S) \{}
\Hint{1.3}{每轮取出一个最大值先压回栈，最终栈顶最小。}\nopagebreak[4]
\Code{1.3}{while (!Empty(S)) EnQueue(\&Q, Pop(S));}
\Code{1.3}{while (!Empty(\&Q)) \{}
\Code{2.6}{k = Size(\&Q);}
\Code{2.6}{m = Front(\&Q);}
\Code{2.6}{repeat (k) \{}
\Code{3.9}{x = DeQueue(\&Q);}
\Code{3.9}{m = max(m, x);}
\Code{3.9}{EnQueue(\&Q, x);}
\Code{2.6}{\}}
\Hint{2.6}{第二次轮转仍按原长度计数，即使移除最大项，剩下的项也都只检查一次。}\nopagebreak[4]
\Code{2.6}{removed = false;}
\Code{2.6}{repeat (k) \{}
\Code{3.9}{x = DeQueue(\&Q);}
\Code{3.9}{if (!removed \&\& x == m) removed = true;}
\Code{3.9}{else EnQueue(\&Q, x);}
\Code{2.6}{\}}
\Code{2.6}{Push(S, m);}
\Code{1.3}{\}}
\Code{0.0}{\}}
\nopagebreak[4]
\Cost{\ensuremath{O\!\left(n^{2}\right)}}{\ensuremath{O\!\left(n\right)}}
\par\medskip

\Needspace{21\baselineskip}
\subsubsection{利用数组模拟链表}
\Code{0.0}{ArrayListEdit(p, x, insert) \{}
\Hint{1.3}{next 存下标，空闲槽另连一条链。}\nopagebreak[4]
\Code{1.3}{if (insert) \{}
\Code{2.6}{i = Slot(ALLOC, 0);}
\Code{2.6}{data[i] = x;}
\Code{2.6}{next[i] = next[p];}
\Code{2.6}{next[p] = i;}
\Code{1.3}{\}}
\Hint{1.3}{归还 i 之前要把前驱改接到 next[i]，因为归还操作会复用 next[i]。}\nopagebreak[4]
\Code{1.3}{else \{}
\Code{2.6}{i = next[p];}
\Code{2.6}{require(i != -1);}
\Code{2.6}{next[p] = next[i];}
\Code{2.6}{Slot(FREE, i);}
\Code{1.3}{\}}
\Code{0.0}{\}}
\nopagebreak[4]
\Cost{初始化 \ensuremath{O\!\left(C\right)}，已知前驱的增删 \ensuremath{O\!\left(1\right)}}{\ensuremath{O\!\left(C\right)}（结构）}
\par\medskip

\Needspace{15\baselineskip}
\subsubsection{利用链表实现栈和队列}
\Code{0.0}{LinkedAdapters(op, x) \{}
\Hint{1.3}{栈在同一端进出，队列在两端分别进出。}\nopagebreak[4]
\Code{1.3}{switch (op) \{}
\Code{2.6}{case PUSH: return LinkedStack(\&top, PUSH, x);}
\Code{2.6}{case POP: return LinkedStack(\&top, POP, 0);}
\Hint{2.6}{同为链式存储，栈只需保存 top，队列还需 tail 才能常数时间尾插。}\nopagebreak[4]
\Code{2.6}{case IN: return LinkedQ(\&q, IN, x);}
\Code{2.6}{case OUT: return LinkedQ(\&q, OUT, 0);}
\Code{1.3}{\}}
\Code{0.0}{\}}
\nopagebreak[4]
\Cost{\ensuremath{O\!\left(1\right)}/次}{\ensuremath{O\!\left(n\right)}（结构）}
\par\medskip

\Needspace{20\baselineskip}
\subsubsection{仅利用递归实现栈的逆置}
\Code{0.0}{ReverseStack(S) \{}
\Hint{1.3}{Bottom 取走栈底后会恢复上面的元素。}\nopagebreak[4]
\Code{1.3}{if (Empty(S)) return;}
\Code{1.3}{b = Bottom(S);}
\Code{1.3}{ReverseStack(S);}
\Code{1.3}{Push(S, b);}
\Code{0.0}{\}}
\Code{0.0}{Bottom(S) \{}
\Hint{1.3}{Bottom 返回时保存的上层元素逐个压回，只有最底项被取走。}\nopagebreak[4]
\Code{1.3}{x = Pop(S);}
\Code{1.3}{if (Empty(S)) return x;}
\Code{1.3}{b = Bottom(S);}
\Code{1.3}{Push(S, x);}
\Code{1.3}{return b;}
\Code{0.0}{\}}
\nopagebreak[4]
\Cost{\ensuremath{O\!\left(n^{2}\right)}}{\ensuremath{O\!\left(n\right)}}
\par\medskip

\Group{经典算法思想与问题}

\Needspace{18\baselineskip}
\subsubsection{双指针算法}
\Code{0.0}{PairSorted(A, n, t) \{}
\Hint{1.3}{和偏小只能增大左值，和偏大就减小右值。}\nopagebreak[4]
\Code{1.3}{l = 0;}
\Code{1.3}{r = n-1;}
\Code{1.3}{while (l<r) \{}
\Code{2.6}{s = A[l]+A[r];}
\Code{2.6}{if (s == t) return (l, r);}
\Hint{2.6}{若左右指针相遇，不能用同一位置上的数和自己配对。}\nopagebreak[4]
\Code{2.6}{if (s<t) l++;}
\Code{2.6}{else r--;}
\Code{1.3}{\}}
\Code{1.3}{return NONE;}
\Code{0.0}{\}}
\nopagebreak[4]
\Cost{\ensuremath{O\!\left(n\right)}}{\ensuremath{O\!\left(1\right)}}
\par\medskip

\Needspace{16\baselineskip}
\subsubsection{快慢指针算法}
\Code{0.0}{LeftMiddle(first) \{}
\Hint{1.3}{快指针还能走满两步才让慢指针前进。}\nopagebreak[4]
\Code{1.3}{if (!first) return NULL;}
\Code{1.3}{s = f = first;}
\Code{1.3}{while (f->next \&\& f->next->next) \{}
\Code{2.6}{s = s->next;}
\Hint{2.6}{本题与右中点版本只差快指针的循环条件，偶数长度时结果相差一位。}\nopagebreak[4]
\Code{2.6}{f = f->next->next;}
\Code{1.3}{\}}
\Code{1.3}{return s;}
\Code{0.0}{\}}
\nopagebreak[4]
\Cost{\ensuremath{O\!\left(n\right)}}{\ensuremath{O\!\left(1\right)}}
\par\medskip

\Needspace{19\baselineskip}
\subsubsection{滑动窗口算法}
\Code{0.0}{MinPositiveWindow(A, t) \{}
\Hint{1.3}{元素均为正，缩短左端会严格减小窗口和。}\nopagebreak[4]
\Code{1.3}{l = sum = 0;}
\Code{1.3}{best = INF;}
\Code{1.3}{for (r = 0; r<len(A); r++) \{}
\Code{2.6}{sum += A[r];}
\Code{2.6}{while (sum >= t) \{}
\Hint{3.9}{满足目标后仍要继续缩，直到再删一个元素会让总和不足。}\nopagebreak[4]
\Code{3.9}{best = min(best, r-l+1);}
\Code{3.9}{sum -= A[l++];}
\Code{2.6}{\}}
\Code{1.3}{\}}
\Code{1.3}{return best;}
\Code{0.0}{\}}
\nopagebreak[4]
\Cost{\ensuremath{O\!\left(n\right)}}{\ensuremath{O\!\left(1\right)}}
\par\medskip

\Needspace{12\baselineskip}
\subsubsection{原地算法}
\Code{0.0}{MoveZeros(A, n) \{}
\Hint{1.3}{先紧凑写入非零值，再补齐尾部。}\nopagebreak[4]
\Code{1.3}{w = 0;}
\Hint{1.3}{w 从不超过读指针，因此写入保留项不会覆盖尚未读取的数据。}\nopagebreak[4]
\Code{1.3}{for (i = 0; i<n; i++) if (A[i] != 0) A[w++] = A[i];}
\Code{1.3}{while (w<n) A[w++] = 0;}
\Code{0.0}{\}}
\nopagebreak[4]
\Cost{\ensuremath{O\!\left(n\right)}}{\ensuremath{O\!\left(1\right)}}
\par\medskip

\Needspace{11\baselineskip}
\subsubsection{递归算法}
\Code{0.0}{PrefixSumRec(A, n) \{}
\Hint{1.3}{每次减少一项，空前缀的和为 0。}\nopagebreak[4]
\Code{1.3}{if (n == 0) return 0;}
\Hint{1.3}{n=0 时直接返回，随后每个调用把刚才暂时搁置的一项加回来。}\nopagebreak[4]
\Code{1.3}{return PrefixSumRec(A, n-1)+A[n-1];}
\Code{0.0}{\}}
\nopagebreak[4]
\Cost{\ensuremath{O\!\left(n\right)}}{\ensuremath{O\!\left(n\right)}}
\par\medskip

\Needspace{20\baselineskip}
\subsubsection{分治算法}
\Code{0.0}{MergeSort(A, l, r, tmp) \{}
\Hint{1.3}{递归区间为左闭右开，tmp 在所有调用间共用。}\nopagebreak[4]
\Code{1.3}{if (r-l<2) return;}
\Code{1.3}{m = (l+r)/2;}
\Code{1.3}{MergeSort(A, l, m, tmp);}
\Code{1.3}{MergeSort(A, m, r, tmp);}
\Code{1.3}{i = l;}
\Code{1.3}{j = m;}
\Hint{1.3}{若一侧已经耗尽，直接从另一侧取值，不读越界位置。}\nopagebreak[4]
\Code{1.3}{for (k = l; k<r; k++) \{}
\Code{2.6}{if (j == r || (i<m \&\& A[i] <= A[j])) tmp[k] = A[i++];}
\Code{2.6}{else tmp[k] = A[j++];}
\Code{1.3}{\}}
\Code{1.3}{for (k = l; k<r; k++) A[k] = tmp[k];}
\Code{0.0}{\}}
\nopagebreak[4]
\Cost{\ensuremath{O\!\left(n \log (n+1)\right)}}{\ensuremath{O\!\left(n\right)}}
\par\medskip

\Needspace{17\baselineskip}
\subsubsection{回溯算法}
\Code{0.0}{Subsets(i) \{}
\Hint{1.3}{每层决定是否选 A[i]，离开选取分支时撤销。}\nopagebreak[4]
\Code{1.3}{if (i == n) \{}
\Code{2.6}{Emit(path);}
\Code{2.6}{return;}
\Code{1.3}{\}}
\Hint{1.3}{不选择分支先执行，返回后才把 A[i] 加入路径去试选择分支。}\nopagebreak[4]
\Code{1.3}{Subsets(i+1);}
\Code{1.3}{Push(\&path, A[i]);}
\Code{1.3}{Subsets(i+1);}
\Code{1.3}{Pop(\&path);}
\Code{0.0}{\}}
\nopagebreak[4]
\Cost{\ensuremath{O\!\left(n\,2^{n}\right)}}{\ensuremath{O\!\left(n\right)}}
\par\medskip

\Needspace{18\baselineskip}
\subsubsection{有序线性表的集合交集}
\Code{0.0}{Intersection(A, B) \{}
\Hint{1.3}{相等才输出，较小者无法再与另一侧当前值匹配。}\nopagebreak[4]
\Code{1.3}{i = j = 0;}
\Code{1.3}{while (i<m \&\& j<n) \{}
\Code{2.6}{if (A[i]<B[j]) i++;}
\Code{2.6}{else if (A[i]>B[j]) j++;}
\Code{2.6}{else \{}
\Hint{3.9}{相等时必须同时推进两个游标，否则同一项可能重复输出。}\nopagebreak[4]
\Code{3.9}{Emit(A[i++]);}
\Code{3.9}{j++;}
\Code{2.6}{\}}
\Code{1.3}{\}}
\Code{0.0}{\}}
\nopagebreak[4]
\Cost{\ensuremath{O\!\left(m+n\right)}}{\ensuremath{O\!\left(1\right)}}
\par\medskip

\Needspace{18\baselineskip}
\subsubsection{有序线性表的集合并集}
\Code{0.0}{UnionSorted(A, B) \{}
\Hint{1.3}{两输入均升序无重复，相等时仅输出一次。}\nopagebreak[4]
\Code{1.3}{i = j = 0;}
\Code{1.3}{while (i<m || j<n) \{}
\Code{2.6}{if (j == n || (i<m \&\& A[i]<B[j])) Emit(A[i++]);}
\Code{2.6}{else if (i == m || B[j]<A[i]) Emit(B[j++]);}
\Code{2.6}{else \{}
\Hint{3.9}{某一侧耗尽时，循环还需要把另一侧剩余项全部并入结果。}\nopagebreak[4]
\Code{3.9}{Emit(A[i++]);}
\Code{3.9}{j++;}
\Code{2.6}{\}}
\Code{1.3}{\}}
\Code{0.0}{\}}
\nopagebreak[4]
\Cost{\ensuremath{O\!\left(m+n\right)}}{\ensuremath{O\!\left(1\right)}}
\par\medskip

\Needspace{14\baselineskip}
\subsubsection{有序线性表的集合差集}
\Code{0.0}{Difference(A, B) \{}
\Hint{1.3}{B 的游标只前进，停在首个不小于当前 x 的值。}\nopagebreak[4]
\Code{1.3}{j = 0;}
\Code{1.3}{for (x in A) \{}
\Hint{2.6}{若 B 已经到末尾，A 的后续每项都应保留。}\nopagebreak[4]
\Code{2.6}{while (j<n \&\& B[j]<x) j++;}
\Code{2.6}{if (j == n || B[j] != x) Emit(x);}
\Code{1.3}{\}}
\Code{0.0}{\}}
\nopagebreak[4]
\Cost{\ensuremath{O\!\left(m+n\right)}}{\ensuremath{O\!\left(1\right)}}
\par\medskip

\Needspace{15\baselineskip}
\subsubsection{两数之和问题}
\Code{0.0}{TwoSum(A, t) \{}
\Hint{1.3}{先查补数再记录当前项，保证两个下标不同。}\nopagebreak[4]
\Code{1.3}{InitMap(\&M);}
\Code{1.3}{for (i = 0; i<len(A); i++) \{}
\Code{2.6}{if (Has(\&M, t-A[i])) return (Get(\&M, t-A[i]), i);}
\Hint{2.6}{表中保存的是下标，所以命中后可以直接返回两个位置。}\nopagebreak[4]
\Code{2.6}{Put(\&M, A[i], i);}
\Code{1.3}{\}}
\Code{1.3}{return NONE;}
\Code{0.0}{\}}
\nopagebreak[4]
\Cost{\ensuremath{O\!\left(n\right)}（期望）}{\ensuremath{O\!\left(n\right)}}
\par\medskip

\Needspace{27\baselineskip}
\subsubsection{三数之和问题}
\Code{0.0}{ThreeSum(A, n, t) \{}
\Hint{1.3}{固定首项后压缩双指针，相同三元组只输出一次。}\nopagebreak[4]
\Code{1.3}{HeapSort(A, n);}
\Code{1.3}{for (i = 0; i<n-2; i++) \{}
\Code{2.6}{if (i>0 \&\& A[i] == A[i-1]) continue;}
\Code{2.6}{l = i+1;}
\Code{2.6}{r = n-1;}
\Code{2.6}{while (l<r) \{}
\Code{3.9}{s = A[i]+A[l]+A[r];}
\Code{3.9}{if (s<t) l++;}
\Code{3.9}{else if (s>t) r--;}
\Code{3.9}{else \{}
\Hint{5.2}{记录当前 a、b 后跳过同值，避免双指针下一轮输出相同数值三元组。}\nopagebreak[4]
\Code{5.2}{Emit(A[i], A[l], A[r]);}
\Code{5.2}{a = A[l];}
\Code{5.2}{b = A[r];}
\Code{5.2}{while (l<r \&\& A[l] == a) l++;}
\Code{5.2}{while (l<r \&\& A[r] == b) r--;}
\Code{3.9}{\}}
\Code{2.6}{\}}
\Code{1.3}{\}}
\Code{0.0}{\}}
\nopagebreak[4]
\Cost{\ensuremath{O\!\left(n^{2}\right)}}{\ensuremath{O\!\left(1\right)}（原地堆排序）}
\par\medskip

\Needspace{17\baselineskip}
\subsubsection{多项式的链表表示}
\Code{0.0}{BuildPoly(terms) \{}
\Hint{1.3}{输入指数严格递增，零系数省略。零多项式的哨兵后继为 NULL。}\nopagebreak[4]
\Code{1.3}{H = NewTerm();}
\Code{1.3}{t = H;}
\Code{1.3}{for ((coef, exp) in terms) \{}
\Code{2.6}{if (coef == 0) continue;}
\Hint{2.6}{尾指针只用于常数时间接链，最后返回头结点才能访问完整多项式。}\nopagebreak[4]
\Code{2.6}{t->next = NewTerm(coef, exp);}
\Code{2.6}{t = t->next;}
\Code{1.3}{\}}
\Code{1.3}{return H;}
\Code{0.0}{\}}
\nopagebreak[4]
\Cost{\ensuremath{O\!\left(n\right)}}{\ensuremath{O\!\left(1\right)}}
\par\medskip

\Needspace{27\baselineskip}
\subsubsection{一元多项式的加法}
\Code{0.0}{AddPoly(a, b) \{}
\Hint{1.3}{输入 a、b 为首元，可为 NULL。若由 BuildPoly 得到哨兵 A、B，调用 AddPoly(A->next, B->next)。}\nopagebreak[4]
\Code{1.3}{C = NewTerm();}
\Code{1.3}{tail = C;}
\Code{1.3}{while (a || b) \{}
\Code{2.6}{if (!b || (a \&\& a->exp<b->exp)) \{}
\Code{3.9}{tail->next = NewTerm(a->coef, a->exp);}
\Code{3.9}{tail = tail->next;}
\Code{3.9}{a = a->next;}
\Code{2.6}{\}}
\Code{2.6}{else if (!a || b->exp<a->exp) \{}
\Code{3.9}{tail->next = NewTerm(b->coef, b->exp);}
\Code{3.9}{tail = tail->next;}
\Code{3.9}{b = b->next;}
\Code{2.6}{\}}
\Hint{2.6}{两条输入按指数递增，当前较小的指数在另一条链后面不可能再出现。}\nopagebreak[4]
\Code{2.6}{else \{}
\Code{3.9}{s = a->coef+b->coef;}
\Code{3.9}{if (s != 0) \{}
\Code{5.2}{tail->next = NewTerm(s, a->exp);}
\Code{5.2}{tail = tail->next;}
\Code{3.9}{\}}
\Code{3.9}{a = a->next;}
\Code{3.9}{b = b->next;}
\Code{2.6}{\}}
\Code{1.3}{\}}
\Hint{1.3}{返回结果的哨兵 C；若所有项抵消，C 的后继为 NULL。}\nopagebreak[4]
\Code{1.3}{return C;}
\Code{0.0}{\}}
\nopagebreak[4]
\Cost{\ensuremath{O\!\left(m+n\right)}}{\ensuremath{O\!\left(1\right)}}
\par\medskip

\Needspace{14\baselineskip}
\subsubsection{一元多项式的乘法}
\Code{0.0}{MulPoly(a, b) \{}
\Hint{1.3}{输入 a、b 为首元，可为 NULL。有序映射按指数累加，最后按指数升序输出。}\nopagebreak[4]
\Code{1.3}{InitOrderedMap(\&M, 0);}
\Code{1.3}{if (!a || !b) return;}
\Code{1.3}{for (p = a; p; p = p->next) \{}
\Hint{2.6}{同一指数可能来自多对乘积，须先累加再输出。}\nopagebreak[4]
\Code{2.6}{for (q = b; q; q = q->next) \{}
\Code{3.9}{M[p->exp+q->exp] += p->coef*q->coef;}
\Code{2.6}{\}}
\Code{1.3}{\}}
\Code{1.3}{for (e in SortedKeys(M)) if (M[e] != 0) Emit(M[e], e);}
\Code{0.0}{\}}
\nopagebreak[4]
\Cost{\ensuremath{O\!\left(1+mn \log (d+2)\right)}}{\ensuremath{O\!\left(d\right)}}
\par\medskip

\Needspace{20\baselineskip}
\subsubsection{区间合并问题}
\Code{0.0}{MergeIntervals(I, n) \{}
\Hint{1.3}{闭区间相接也可合并，结束时补输出最后一段。}\nopagebreak[4]
\Code{1.3}{if (n == 0) return;}
\Code{1.3}{HeapSortByStart(I, n);}
\Code{1.3}{cur = I[0];}
\Code{1.3}{for (i = 1; i<n; i++) \{}
\Code{2.6}{if (I[i].l <= cur.r) cur.r = max(cur.r, I[i].r);}
\Code{2.6}{else \{}
\Hint{3.9}{遇到不重叠区间时，之前的 cur 才确定不再延长，可以输出。}\nopagebreak[4]
\Code{3.9}{Emit(cur);}
\Code{3.9}{cur = I[i];}
\Code{2.6}{\}}
\Code{1.3}{\}}
\Code{1.3}{Emit(cur);}
\Code{0.0}{\}}
\nopagebreak[4]
\Cost{\ensuremath{O\!\left(n \log (n+1)\right)}}{\ensuremath{O\!\left(1\right)}（原地堆排序）}
\par\medskip

\Needspace{28\baselineskip}
\subsubsection{生成全排列与组合}
\Code{0.0}{Perm(i) \{}
\Hint{1.3}{排列通过交换固定当前位置；组合通过递增下标避免重复。}\nopagebreak[4]
\Code{1.3}{if (i == n) \{}
\Code{2.6}{Emit(A);}
\Code{2.6}{return;}
\Code{1.3}{\}}
\Code{1.3}{for (j = i; j<n; j++) \{}
\Code{2.6}{swap(A[i], A[j]);}
\Code{2.6}{Perm(i+1);}
\Code{2.6}{swap(A[i], A[j]);}
\Code{1.3}{\}}
\Code{0.0}{\}}
\Hint{0.0}{组合中 j 的上界为剩余项预留了足够位置，递归传 j+1 保证下标不重复。}\nopagebreak[4]
\Code{0.0}{Choose(start, d) \{}
\Code{1.3}{if (d == k) \{}
\Code{2.6}{Emit(path);}
\Code{2.6}{return;}
\Code{1.3}{\}}
\Code{1.3}{for (j = start; j <= n-k+d; j++) \{}
\Code{2.6}{path[d] = A[j];}
\Code{2.6}{Choose(j+1, d+1);}
\Code{1.3}{\}}
\Code{0.0}{\}}
\nopagebreak[4]
\Cost{排列 \ensuremath{O\!\left(n\,n!\right)}；组合 \ensuremath{O\!\left(k\,C(n,k)+1\right)}}{排列 \ensuremath{O\!\left(n\right)}；组合 \ensuremath{O\!\left(k+1\right)}}
\par\medskip

\Needspace{22\baselineskip}
\subsubsection{N皇后问题}
\Code{0.0}{Queens(r) \{}
\Hint{1.3}{列和两组斜线各有独立占用标记。}\nopagebreak[4]
\Code{1.3}{if (r == n) \{}
\Code{2.6}{Emit(pos);}
\Code{2.6}{return;}
\Code{1.3}{\}}
\Code{1.3}{for (c = 0; c<n; c++) \{}
\Code{2.6}{a = r-c+n-1;}
\Code{2.6}{b = r+c;}
\Hint{2.6}{对角线 r−c 可能为负，加 n−1 才能作为数组下标。}\nopagebreak[4]
\Code{2.6}{if (col[c] || diag[a] || anti[b]) continue;}
\Code{2.6}{col[c] = diag[a] = anti[b] = true;}
\Code{2.6}{pos[r] = c;}
\Code{2.6}{Queens(r+1);}
\Code{2.6}{col[c] = diag[a] = anti[b] = false;}
\Code{1.3}{\}}
\Code{0.0}{\}}
\nopagebreak[4]
\Cost{\ensuremath{O\!\left(n\,n!\right)}（搜索上界，另计输出）}{\ensuremath{O\!\left(n\right)}}
\par\medskip

\Part{串}

\Group{串的基本操作}

\Needspace{12\baselineskip}
\subsubsection{串的初始化与赋值}
\Code{0.0}{Assign(S, src, n) \{}
\Hint{1.3}{长度单独保存，字符逐个复制。}\nopagebreak[4]
\Code{1.3}{S->data = AllocChars(n);}
\Hint{1.3}{长度字段记录 n，因此这里的数据区不依赖末尾终止符。}\nopagebreak[4]
\Code{1.3}{for (i = 0; i<n; i++) S->data[i] = src[i];}
\Code{1.3}{S->len = n;}
\Code{0.0}{\}}
\nopagebreak[4]
\Cost{\ensuremath{O\!\left(n\right)}}{\ensuremath{O\!\left(1\right)}}
\par\medskip

\Needspace{12\baselineskip}
\subsubsection{串的复制}
\Code{0.0}{CopyString(S) \{}
\Hint{1.3}{目标具有独立存储，改动副本不会改动原串。}\nopagebreak[4]
\Code{1.3}{String T;}
\Hint{1.3}{Assign 通过目标地址写入新串的数据指针和长度。}\nopagebreak[4]
\Code{1.3}{Assign(\&T, S.data, S.len);}
\Code{1.3}{return T;}
\Code{0.0}{\}}
\nopagebreak[4]
\Cost{\ensuremath{O\!\left(n\right)}}{\ensuremath{O\!\left(1\right)}}
\par\medskip

\Needspace{13\baselineskip}
\subsubsection{串的比较}
\Code{0.0}{Compare(A, B) \{}
\Hint{1.3}{第一个不同字符决定顺序，公共前缀相同则短串靠前。}\nopagebreak[4]
\Code{1.3}{for (i = 0; i<min(m, n); i++) \{}
\Code{2.6}{if (A[i] != B[i]) return Sign(A[i]-B[i]);}
\Code{1.3}{\}}
\Hint{1.3}{首个不等字符出现后无需检查后面各位。}\nopagebreak[4]
\Code{1.3}{return Sign(m-n);}
\Code{0.0}{\}}
\nopagebreak[4]
\Cost{\ensuremath{O\!\left(\min (m,n)+1\right)}}{\ensuremath{O\!\left(1\right)}}
\par\medskip

\Needspace{13\baselineskip}
\subsubsection{串的连接}
\Code{0.0}{Concat(A, B) \{}
\Hint{1.3}{后串从结果下标 m 开始写。}\nopagebreak[4]
\Code{1.3}{T = AllocChars(m+n);}
\Code{1.3}{for (i = 0; i<m; i++) T[i] = A[i];}
\Hint{1.3}{第二个循环把 B 的第 0 项写到 T[m]，紧接 A 的最后一项。}\nopagebreak[4]
\Code{1.3}{for (j = 0; j<n; j++) T[m+j] = B[j];}
\Code{1.3}{return (T, m+n);}
\Code{0.0}{\}}
\nopagebreak[4]
\Cost{\ensuremath{O\!\left(m+n\right)}}{\ensuremath{O\!\left(1\right)}}
\par\medskip

\Needspace{12\baselineskip}
\subsubsection{串的长度计算}
\Code{0.0}{Length(char *s) \{}
\Hint{1.3}{终止符占存储空间，但不计入长度。}\nopagebreak[4]
\Code{1.3}{i = 0;}
\Hint{1.3}{长度为零的 C 串在首字节即遇到终止符。}\nopagebreak[4]
\Code{1.3}{while (s[i] != 0) i++;}
\Code{1.3}{return i;}
\Code{0.0}{\}}
\nopagebreak[4]
\Cost{\ensuremath{O\!\left(n\right)}（终止符串）}{\ensuremath{O\!\left(1\right)}}
\par\medskip

\Needspace{13\baselineskip}
\subsubsection{子串的提取}
\Code{0.0}{Slice(S, i, k) \{}
\Hint{1.3}{复制区间的末端 i+k 不包含在内。}\nopagebreak[4]
\Code{1.3}{require(i >= 0 \&\& k >= 0 \&\& i+k <= n);}
\Code{1.3}{T = AllocChars(k);}
\Hint{1.3}{目标下标从 0 开始，对应的源下标从 i 开始。}\nopagebreak[4]
\Code{1.3}{for (j = 0; j<k; j++) T[j] = S[i+j];}
\Code{1.3}{return (T, k);}
\Code{0.0}{\}}
\nopagebreak[4]
\Cost{\ensuremath{O\!\left(k\right)}}{\ensuremath{O\!\left(1\right)}}
\par\medskip

\Needspace{15\baselineskip}
\subsubsection{子串的插入与删除}
\Code{0.0}{Splice(S, i, k, T) \{}
\Hint{1.3}{以 T 替换 k 项：k 为零即插入，T 为空即删除。}\nopagebreak[4]
\Code{1.3}{require(0 <= i \&\& i+k <= n \&\& k >= 0);}
\Code{1.3}{m = len(T);}
\Code{1.3}{InitBuffer(\&out, n-k+m);}
\Code{1.3}{AppendRange(\&out, S, 0, i);}
\Hint{1.3}{三段的排列次序始终是原前缀、插入串、原后缀。}\nopagebreak[4]
\Code{1.3}{AppendRange(\&out, T, 0, m);}
\Code{1.3}{AppendRange(\&out, S, i+k, n);}
\Code{1.3}{return out;}
\Code{0.0}{\}}
\nopagebreak[4]
\Cost{\ensuremath{O\!\left(n-k+m+1\right)}}{\ensuremath{O\!\left(1\right)}（新串为输出）}
\par\medskip

\Needspace{21\baselineskip}
\subsubsection{子串的查找与替换}
\Code{0.0}{Replace(S, P, T) \{}
\Hint{1.3}{匹配后跨过整段模式，替换串不重新参与匹配。}\nopagebreak[4]
\Code{1.3}{require(m>0);}
\Code{1.3}{i = 0;}
\Code{1.3}{while (i<n) \{}
\Code{2.6}{j = 0;}
\Code{2.6}{while (j<m \&\& i+j<n \&\& S[i+j] == P[j]) j++;}
\Code{2.6}{if (j == m) \{}
\Hint{3.9}{末尾剩余长度小于模式长度时，比较循环自行停止，不会越界。}\nopagebreak[4]
\Code{3.9}{Append(\&out, T);}
\Code{3.9}{i += m;}
\Code{2.6}{\}}
\Code{2.6}{else AppendChar(\&out, S[i++]);}
\Code{1.3}{\}}
\Code{1.3}{return out;}
\Code{0.0}{\}}
\nopagebreak[4]
\Cost{\ensuremath{O\!\left(nm+L\right)}}{\ensuremath{O\!\left(1\right)}（输出可追加）}
\par\medskip

\Needspace{16\baselineskip}
\subsubsection{字符串逆置}
\Code{0.0}{ReverseChars(s, n) \{}
\Hint{1.3}{交换一对后，两端同时向内。}\nopagebreak[4]
\Code{1.3}{l = 0;}
\Code{1.3}{r = n-1;}
\Code{1.3}{while (l<r) \{}
\Code{2.6}{swap(s[l], s[r]);}
\Hint{2.6}{奇数长度时中间字符留在原位。}\nopagebreak[4]
\Code{2.6}{l++;}
\Code{2.6}{r--;}
\Code{1.3}{\}}
\Code{0.0}{\}}
\nopagebreak[4]
\Cost{\ensuremath{O\!\left(n\right)}}{\ensuremath{O\!\left(1\right)}}
\par\medskip

\Needspace{13\baselineskip}
\subsubsection{判断回文字符串}
\Code{0.0}{PalString(s, n) \{}
\Hint{1.3}{只要有一对对称位置不等就可提前结束。}\nopagebreak[4]
\Code{1.3}{for (i = 0; i<n/2; i++) \{}
\Code{2.6}{if (s[i] != s[n-1-i]) return false;}
\Code{1.3}{\}}
\Hint{1.3}{第 i 项的对称位置是 n−1−i，只需检查前半段。}\nopagebreak[4]
\Code{1.3}{return true;}
\Code{0.0}{\}}
\nopagebreak[4]
\Cost{\ensuremath{O\!\left(n\right)}}{\ensuremath{O\!\left(1\right)}}
\par\medskip

\Group{字符串模式匹配}

\Needspace{15\baselineskip}
\subsubsection{BF暴力匹配算法}
\Code{0.0}{Brute(T, n, P, m) \{}
\Hint{1.3}{每次失配只换文本起点，从模式首字符重试。}\nopagebreak[4]
\Code{1.3}{for (i = 0; i <= n-m; i++) \{}
\Code{2.6}{j = 0;}
\Code{2.6}{while (j<m \&\& T[i+j] == P[j]) j++;}
\Hint{2.6}{j 达到 m 说明所有字符都已匹配，返回文本起点 i。}\nopagebreak[4]
\Code{2.6}{if (j == m) return i;}
\Code{1.3}{\}}
\Code{1.3}{return -1;}
\Code{0.0}{\}}
\nopagebreak[4]
\Cost{\ensuremath{O\!\left(nm+1\right)}}{\ensuremath{O\!\left(1\right)}}
\par\medskip

\Needspace{20\baselineskip}
\subsubsection{KMP算法}
\Code{0.0}{KMP(T, n, P, m) \{}
\Hint{1.3}{失配时只回退 j，已经扫描的文本无需重读。}\nopagebreak[4]
\Code{1.3}{if (m == 0) return 0;}
\Code{1.3}{BuildNext(P, m, next);}
\Code{1.3}{i = j = 0;}
\Code{1.3}{while (i<n \&\& j<m) \{}
\Hint{2.6}{j 为 −1 时不要读取 P[j]，短路条件先处理这个哨兵值。}\nopagebreak[4]
\Code{2.6}{if (j == -1 || T[i] == P[j]) \{}
\Code{3.9}{i++;}
\Code{3.9}{j++;}
\Code{2.6}{\}}
\Hint{2.6}{保留文本下标 i，只换模式串的比较位置。}\nopagebreak[4]
\Code{2.6}{else j = next[j];}
\Code{1.3}{\}}
\Code{1.3}{return j == m ? i-m : -1;}
\Code{0.0}{\}}
\nopagebreak[4]
\Cost{\ensuremath{O\!\left(n+m\right)}}{\ensuremath{O\!\left(m\right)}}
\par\medskip

\Needspace{20\baselineskip}
\subsubsection{KMP算法的next数组构建}
\Code{0.0}{BuildNext(P, m, next) \{}
\Hint{1.3}{j 表示当前可复用前缀长度；−1 让空前缀失配也能统一推进。}\nopagebreak[4]
\Code{1.3}{next[0] = -1;}
\Code{1.3}{i = 0;}
\Code{1.3}{j = -1;}
\Code{1.3}{while (i<m) \{}
\Code{2.6}{if (j == -1 || P[i] == P[j]) \{}
\Code{3.9}{i++;}
\Code{3.9}{j++;}
\Hint{3.9}{next[i] 描述 P 的前 i 个字符。刚匹配到下一字符后，先增加 i、j 再写表。}\nopagebreak[4]
\Hint{3.9}{前 i 个字符的最长真边界长度已确定，写入这一项。}\nopagebreak[4]
\Code{3.9}{next[i] = j;}
\Code{2.6}{\}}
\Code{2.6}{else j = next[j];}
\Code{1.3}{\}}
\Code{0.0}{\}}
\nopagebreak[4]
\Cost{\ensuremath{O\!\left(m\right)}}{\ensuremath{O\!\left(1\right)}}
\par\medskip

\Needspace{16\baselineskip}
\subsubsection{KMP算法的nextval数组构建}
\Code{0.0}{BuildNextVal(P, m, val) \{}
\Hint{1.3}{回退后若仍将比较同一字符，可继续跳过该位置。}\nopagebreak[4]
\Code{1.3}{if (m == 0) return;}
\Code{1.3}{BuildNext(P, m, next);}
\Code{1.3}{val[0] = -1;}
\Code{1.3}{for (i = 1; i<m; i++) \{}
\Hint{2.6}{普通 next 已指示失配后还能保留多少字符，val 是继续压缩后的失配跳转表。}\nopagebreak[4]
\Code{2.6}{j = next[i];}
\Code{2.6}{val[i] = (j >= 0 \&\& P[i] == P[j]) ? val[j]:j;}
\Code{1.3}{\}}
\Code{0.0}{\}}
\nopagebreak[4]
\Cost{\ensuremath{O\!\left(m\right)}}{\ensuremath{O\!\left(m\right)}}
\par\medskip

\Needspace{19\baselineskip}
\subsubsection{KR算法（Rabin-Karp算法）}
\Code{0.0}{RabinKarp(T, n, P, m) \{}
\Hint{1.3}{哈希相同后逐字核对，排除碰撞。}\nopagebreak[4]
\Code{1.3}{if (m == 0) return 0;}
\Code{1.3}{if (m>n) return -1;}
\Code{1.3}{hp = Hash(P, m);}
\Code{1.3}{h = Hash(T, m);}
\Code{1.3}{power = PowMod(base, m-1, q);}
\Hint{1.3}{最后一个窗口后不再更新哈希，以免读取 T[n]。}\nopagebreak[4]
\Code{1.3}{for (i = 0; i <= n-m; i++) \{}
\Code{2.6}{if (h == hp \&\& EqualRange(T+i, P, m)) return i;}
\Code{2.6}{if (i<n-m) h = Mod((h-T[i]*power)*base+T[i+m], q);}
\Code{1.3}{\}}
\Code{1.3}{return -1;}
\Code{0.0}{\}}
\nopagebreak[4]
\Cost{最坏 \ensuremath{O\!\left(nm\right)}，初始化 \ensuremath{O\!\left(m\right)}}{\ensuremath{O\!\left(1\right)}}
\par\medskip

\Needspace{21\baselineskip}
\subsubsection{BM算法（Boyer-Moore算法）}
\Code{0.0}{BoyerMoore(T, n, P, m) \{}
\Hint{1.3}{两个启发规则各给出安全位移，取较大者。}\nopagebreak[4]
\Code{1.3}{if (m == 0) return 0;}
\Code{1.3}{BuildLast(P, m, last);}
\Code{1.3}{BuildGood(P, m, suf, pre);}
\Code{1.3}{s = 0;}
\Code{1.3}{while (s <= n-m) \{}
\Code{2.6}{j = m-1;}
\Hint{2.6}{从模式末尾向左比，j 变为 −1 表示整串都相等。}\nopagebreak[4]
\Code{2.6}{while (j >= 0 \&\& T[s+j] == P[j]) j--;}
\Code{2.6}{if (j<0) return s;}
\Code{2.6}{bad = j-last[T[s+j]];}
\Code{2.6}{s += max(1, max(bad, GoodShift(j, m, suf, pre)));}
\Code{1.3}{\}}
\Code{1.3}{return -1;}
\Code{0.0}{\}}
\nopagebreak[4]
\Cost{\ensuremath{O\!\left(n m+m^{2}+\sigma \right)}}{\ensuremath{O\!\left(m+\sigma \right)}}
\par\medskip

\Needspace{14\baselineskip}
\subsubsection{BM算法的坏字符规则}
\Code{0.0}{BuildLast(P, m, last) \{}
\Hint{1.3}{同一字符多次出现时最后一次写入的是最右位置。}\nopagebreak[4]
\Code{1.3}{Fill(last, -1);}
\Code{1.3}{for (i = 0; i<m; i++) last[P[i]] = i;}
\Code{0.0}{\}}
\Hint{0.0}{从未出现的字符记录为 −1，坏字符位移可以跨过它。}\nopagebreak[4]
\Code{0.0}{BadShift(j, c, last) \{}
\Code{1.3}{return max(1, j-last[c]);}
\Code{0.0}{\}}
\nopagebreak[4]
\Cost{预处理 \ensuremath{O\!\left(m+\sigma \right)}，查询 \ensuremath{O\!\left(1\right)}}{\ensuremath{O\!\left(\sigma \right)}}
\par\medskip

\Needspace{28\baselineskip}
\subsubsection{BM算法的好后缀规则}
\Code{0.0}{BuildGood(P, m, suf, pre) \{}
\Hint{1.3}{suf[k] 存长度 k 的后缀在前部的起点；pre[k] 标出它也是前缀的情形。}\nopagebreak[4]
\Code{1.3}{Fill(suf, -1);}
\Code{1.3}{Fill(pre, false);}
\Code{1.3}{for (i = 0; i<m-1; i++) \{}
\Code{2.6}{j = i;}
\Code{2.6}{k = 0;}
\Code{2.6}{while (j >= 0 \&\& P[j] == P[m-1-k]) \{}
\Code{3.9}{j--;}
\Code{3.9}{k++;}
\Code{3.9}{suf[k] = j+1;}
\Code{2.6}{\}}
\Hint{2.6}{若好后缀未在前部完整出现，转找能和其尾部重合的最长模式前缀。}\nopagebreak[4]
\Code{2.6}{if (j<0) pre[k] = true;}
\Code{1.3}{\}}
\Code{0.0}{\}}
\Code{0.0}{GoodShift(j, m, suf, pre) \{}
\Code{1.3}{k = m-1-j;}
\Code{1.3}{if (k == 0) return 1;}
\Code{1.3}{if (suf[k] >= 0) return j+1-suf[k];}
\Code{1.3}{for (t = k-1; t>0; t--) if (pre[t]) return m-t;}
\Code{1.3}{return m;}
\Code{0.0}{\}}
\nopagebreak[4]
\Cost{预处理 \ensuremath{O\!\left(m^{2}\right)}，单次最坏 \ensuremath{O\!\left(m\right)}}{\ensuremath{O\!\left(m\right)}}
\par\medskip

\Needspace{20\baselineskip}
\subsubsection{Z算法（扩展KMP的基础）}
\Code{0.0}{BuildZ(s, n, Z) \{}
\Hint{1.3}{已知窗口 [l,r) 内可借用已有匹配长度，右边界外再逐字扩展。}\nopagebreak[4]
\Code{1.3}{if (!n) return;}
\Code{1.3}{Z[0] = n;}
\Code{1.3}{l = r = 0;}
\Code{1.3}{for (i = 1; i<n; i++) \{}
\Code{2.6}{Z[i] = i<r ? min(r-i, Z[i-l]):0;}
\Code{2.6}{while (i+Z[i]<n \&\& s[Z[i]] == s[i+Z[i]]) Z[i]++;}
\Hint{2.6}{r 是不包含的右端点；i 恰好等于 r 时没有可复用的部分。}\nopagebreak[4]
\Code{2.6}{if (i+Z[i]>r) \{}
\Code{3.9}{l = i;}
\Code{3.9}{r = i+Z[i];}
\Code{2.6}{\}}
\Code{1.3}{\}}
\Code{0.0}{\}}
\nopagebreak[4]
\Cost{\ensuremath{O\!\left(n\right)}}{\ensuremath{O\!\left(1\right)}}
\par\medskip

\Group{字符串综合算法}

\Needspace{12\baselineskip}
\subsubsection{求两个字符串的最长公共前缀}
\Code{0.0}{CommonPrefix(A, B) \{}
\Hint{1.3}{停下的位置恰是前缀长度。}\nopagebreak[4]
\Code{1.3}{i = 0;}
\Hint{1.3}{任何一个串先耗尽时，已经比较相等的部分就是答案。}\nopagebreak[4]
\Code{1.3}{while (i<min(m, n) \&\& A[i] == B[i]) i++;}
\Code{1.3}{return Slice(A, 0, i);}
\Code{0.0}{\}}
\nopagebreak[4]
\Cost{\ensuremath{O\!\left(\min (m,n)+1\right)}}{\ensuremath{O\!\left(1\right)}}
\par\medskip

\Needspace{12\baselineskip}
\subsubsection{求字符串的最长真前缀与真后缀的公共部分}
\Code{0.0}{Border(S, n) \{}
\Hint{1.3}{next[n] 给出整串的最长真边界长度。}\nopagebreak[4]
\Code{1.3}{if (n == 0) return EmptyString();}
\Hint{1.3}{由于要求真前后缀，不能把整个串本身当答案。}\nopagebreak[4]
\Code{1.3}{BuildNext(S, n, next);}
\Code{1.3}{return Slice(S, 0, next[n]);}
\Code{0.0}{\}}
\nopagebreak[4]
\Cost{\ensuremath{O\!\left(n\right)}}{\ensuremath{O\!\left(n\right)}}
\par\medskip

\Needspace{13\baselineskip}
\subsubsection{求字符串的最小循环节（KMP）}
\Code{0.0}{Period(S, n) \{}
\Hint{1.3}{候选周期还须整除串长才是完整重复。}\nopagebreak[4]
\Code{1.3}{if (!n) return EmptyString();}
\Code{1.3}{BuildNext(S, n, next);}
\Hint{1.3}{若不能整除，只能把整串当作一次重复的循环节。}\nopagebreak[4]
\Code{1.3}{p = n-next[n];}
\Code{1.3}{return Slice(S, 0, n\%p == 0 ? p:n);}
\Code{0.0}{\}}
\nopagebreak[4]
\Cost{\ensuremath{O\!\left(n\right)}}{\ensuremath{O\!\left(n\right)}}
\par\medskip

\Needspace{22\baselineskip}
\subsubsection{求字符串的最长公共子串}
\Code{0.0}{CommonSubstring(A, B) \{}
\Hint{1.3}{逆序更新使 D[j−1] 仍属于上一行，失配时连续长度清零。}\nopagebreak[4]
\Code{1.3}{if (len(A)<len(B)) return CommonSubstring(B, A);}
\Code{1.3}{Fill(D, 0);}
\Code{1.3}{best = end = 0;}
\Code{1.3}{for (i = 1; i <= len(A); i++) \{}
\Code{2.6}{for (j = len(B); j >= 1; j--) \{}
\Code{3.9}{D[j] = A[i-1] == B[j-1]?D[j-1]+1:0;}
\Code{3.9}{if (D[j]>best) \{}
\Hint{5.2}{end 指最长匹配段在 A 中的下一位置，倒退 best 项找到段首。}\nopagebreak[4]
\Code{5.2}{best = D[j];}
\Code{5.2}{end = i;}
\Code{3.9}{\}}
\Code{2.6}{\}}
\Code{1.3}{\}}
\Code{1.3}{return Slice(A, end-best, best);}
\Code{0.0}{\}}
\nopagebreak[4]
\Cost{\ensuremath{O\!\left(mn\right)}}{\ensuremath{O\!\left(\min (m,n)\right)}}
\par\medskip

\Needspace{25\baselineskip}
\subsubsection{求字符串的最长公共子序列}
\Code{0.0}{LCS(A, m, B, n) \{}
\Hint{1.3}{表格存长度，回溯只在匹配时收集字符。}\nopagebreak[4]
\Code{1.3}{Zero(D);}
\Code{1.3}{for (i = 1; i <= m; i++) \{}
\Code{2.6}{for (j = 1; j <= n; j++) D[i][j] = A[i-1] == B[j-1]?1+D[i-1][j-1]:max(D[i-1][j], D[i][j-1]);}
\Code{1.3}{\}}
\Code{1.3}{i = m;}
\Code{1.3}{j = n;}
\Code{1.3}{while (i \&\& j) \{}
\Code{2.6}{if (A[i-1] == B[j-1]) \{}
\Code{3.9}{AppendChar(\&out, A[--i]);}
\Hint{3.9}{两串当前字符不同就向较大邻格回走，分数相同任选一边都能恢复一条答案。}\nopagebreak[4]
\Code{3.9}{j--;}
\Code{2.6}{\}}
\Code{2.6}{else if (D[i-1][j] >= D[i][j-1]) i--;}
\Code{2.6}{else j--;}
\Code{1.3}{\}}
\Code{1.3}{ReverseChars(out, len(out));}
\Code{1.3}{return out;}
\Code{0.0}{\}}
\nopagebreak[4]
\Cost{\ensuremath{O\!\left(mn\right)}}{\ensuremath{O\!\left(mn\right)}}
\par\medskip

\Needspace{21\baselineskip}
\subsubsection{求字符串的最长重复子串}
\Code{0.0}{LongestRepeat(S, n) \{}
\Hint{1.3}{仅比较不同终点，两个重复片段允许重叠。}\nopagebreak[4]
\Code{1.3}{Zero(D);}
\Code{1.3}{best = end = 0;}
\Code{1.3}{for (i = 1; i <= n; i++) \{}
\Code{2.6}{for (j = n; j>i; j--) \{}
\Code{3.9}{D[j] = S[i-1] == S[j-1]?D[j-1]+1:0;}
\Code{3.9}{if (D[j]>best) \{}
\Hint{5.2}{j 从右往左更新，避免读到本轮刚写的状态。}\nopagebreak[4]
\Code{5.2}{best = D[j];}
\Code{5.2}{end = j;}
\Code{3.9}{\}}
\Code{2.6}{\}}
\Code{1.3}{\}}
\Code{1.3}{return Slice(S, end-best, best);}
\Code{0.0}{\}}
\nopagebreak[4]
\Cost{\ensuremath{O\!\left(n^{2}\right)}}{\ensuremath{O\!\left(n\right)}}
\par\medskip

\Needspace{25\baselineskip}
\subsubsection{求字符串的最长回文子串（中心扩展法）}
\Code{0.0}{CenterPal(S, n) \{}
\Hint{1.3}{每个位置同时尝试单字符中心和字符间中心。}\nopagebreak[4]
\Code{1.3}{best = start = 0;}
\Code{1.3}{for (i = 0; i<n; i++) \{}
\Code{2.6}{for (gap = 0; gap <= 1; gap++) \{}
\Code{3.9}{l = i;}
\Code{3.9}{r = i+gap;}
\Code{3.9}{while (l >= 0 \&\& r<n \&\& S[l] == S[r]) \{}
\Code{5.2}{if (r-l+1>best) \{}
\Code{6.5}{start = l;}
\Code{6.5}{best = r-l+1;}
\Code{5.2}{\}}
\Hint{5.2}{gap=0 从一个字符开始，gap=1 从相邻两个字符开始。}\nopagebreak[4]
\Code{5.2}{l--;}
\Code{5.2}{r++;}
\Code{3.9}{\}}
\Code{2.6}{\}}
\Code{1.3}{\}}
\Code{1.3}{return Slice(S, start, best);}
\Code{0.0}{\}}
\nopagebreak[4]
\Cost{\ensuremath{O\!\left(n^{2}\right)}}{\ensuremath{O\!\left(1\right)}}
\par\medskip

\Needspace{13\baselineskip}
\subsubsection{求字符串的最长回文子串（Manacher算法）}
\Code{0.0}{Manacher(S, n) \{}
\Hint{1.3}{插入专用分隔符后，奇偶长度用同一套半径表示。}\nopagebreak[4]
\Code{1.3}{T = InsertSeparators(S, n);}
\Code{1.3}{Radius(T, len(T), R);}
\Hint{1.3}{变换串的最大半径恰好等于原串中该回文的字符数。}\nopagebreak[4]
\Code{1.3}{c = ArgMax(R);}
\Code{1.3}{return Slice(S, (c-R[c])/2, R[c]);}
\Code{0.0}{\}}
\nopagebreak[4]
\Cost{\ensuremath{O\!\left(n\right)}}{\ensuremath{O\!\left(n\right)}}
\par\medskip

\Needspace{19\baselineskip}
\subsubsection{Manacher算法的回文半径数组构建}
\Code{0.0}{Radius(T, n, R) \{}
\Hint{1.3}{right 是当前最靠右的回文端点，镜像值不能直接越过它。}\nopagebreak[4]
\Code{1.3}{c = right = 0;}
\Code{1.3}{for (i = 0; i<n; i++) \{}
\Code{2.6}{r = (i<right ? min(right-i, R[2*c-i]):0);}
\Code{2.6}{while (i-r-1 >= 0 \&\& i+r+1<n \&\& T[i-r-1] == T[i+r+1]) r++;}
\Code{2.6}{R[i] = r;}
\Hint{2.6}{只有新回文突破 right 时，才有必要替换当前中心 c。}\nopagebreak[4]
\Code{2.6}{if (i+r>right) \{}
\Code{3.9}{c = i;}
\Code{3.9}{right = i+r;}
\Code{2.6}{\}}
\Code{1.3}{\}}
\Code{0.0}{\}}
\nopagebreak[4]
\Cost{\ensuremath{O\!\left(n\right)}}{\ensuremath{O\!\left(1\right)}（已给变换串）}
\par\medskip

\Needspace{20\baselineskip}
\subsubsection{求无重复字符的最长子串}
\Code{0.0}{UniqueWindow(S, n) \{}
\Hint{1.3}{左边界不能因读到窗口外的旧字符而倒退。}\nopagebreak[4]
\Code{1.3}{InitMap(\&last, -1);}
\Code{1.3}{l = start = best = 0;}
\Code{1.3}{for (r = 0; r<n; r++) \{}
\Code{2.6}{l = max(l, last[S[r]]+1);}
\Code{2.6}{last[S[r]] = r;}
\Code{2.6}{if (r-l+1>best) \{}
\Hint{3.9}{last 存字符最近下标，初始 −1，第一次见到时不会移动左边界。}\nopagebreak[4]
\Code{3.9}{start = l;}
\Code{3.9}{best = r-l+1;}
\Code{2.6}{\}}
\Code{1.3}{\}}
\Code{1.3}{return Slice(S, start, best);}
\Code{0.0}{\}}
\nopagebreak[4]
\Cost{\ensuremath{O\!\left(n\right)}（期望）}{\ensuremath{O\!\left(\min (n,\sigma )\right)}}
\par\medskip

\Needspace{25\baselineskip}
\subsubsection{求最小覆盖子串}
\Code{0.0}{MinCover(S, T) \{}
\Hint{1.3}{miss 计尚缺的字符总数，重复字符按出现次数计。}\nopagebreak[4]
\Code{1.3}{if (!len(T)) return EmptyString();}
\Code{1.3}{need = Counts(T);}
\Code{1.3}{miss = len(T);}
\Code{1.3}{l = 0;}
\Code{1.3}{best = INF;}
\Code{1.3}{for (r = 0; r<len(S); r++) \{}
\Code{2.6}{if (need[S[r]]-->0) miss--;}
\Code{2.6}{while (miss == 0) \{}
\Code{3.9}{if (r-l+1<best) \{}
\Hint{5.2}{左端字符还回需求计数时，若其计数刚变为正，窗口便不再覆盖目标。}\nopagebreak[4]
\Code{5.2}{start = l;}
\Code{5.2}{best = r-l+1;}
\Code{3.9}{\}}
\Code{3.9}{if (++need[S[l++]]>0) miss++;}
\Code{2.6}{\}}
\Code{1.3}{\}}
\Code{1.3}{return best == INF ? EmptyString():Slice(S, start, best);}
\Code{0.0}{\}}
\nopagebreak[4]
\Cost{\ensuremath{O\!\left(m+n\right)}（期望）}{\ensuremath{O\!\left(\sigma \right)}}
\par\medskip

\Needspace{12\baselineskip}
\subsubsection{判断字符串是否为循环移位所得}
\Code{0.0}{IsRotation(A, B) \{}
\Hint{1.3}{A 的每次循环移位都出现在两份 A 的拼接串中。}\nopagebreak[4]
\Code{1.3}{if (len(A) != len(B)) return false;}
\Hint{1.3}{先比较长度，避免误把 A+A 中一个更短的子串判为完整旋转。}\nopagebreak[4]
\Code{1.3}{T = Concat(A, A);}
\Code{1.3}{return KMP(T, len(T), B, len(B)) >= 0;}
\Code{0.0}{\}}
\nopagebreak[4]
\Cost{\ensuremath{O\!\left(n\right)}}{\ensuremath{O\!\left(n\right)}}
\par\medskip

\Needspace{22\baselineskip}
\subsubsection{利用滚动哈希实现字符串匹配}
\Code{0.0}{RollingMatch(T, n, P, m) \{}
\Hint{1.3}{先扣掉旧首字符占据的最高位，再乘基数并加入新末位。}\nopagebreak[4]
\Code{1.3}{if (!m) return 0;}
\Code{1.3}{if (m>n) return -1;}
\Code{1.3}{power = PowMod(base, m-1, q);}
\Code{1.3}{h = hp = 0;}
\Code{1.3}{for (j = 0; j<m; j++) \{}
\Code{2.6}{h = Mod(h*base+T[j], q);}
\Code{2.6}{hp = Mod(hp*base+P[j], q);}
\Code{1.3}{\}}
\Hint{1.3}{Mod 表示非负模，扣除旧首字符之后即使中间结果为负也能正确归一化。}\nopagebreak[4]
\Code{1.3}{for (i = 0; i <= n-m; i++) \{}
\Code{2.6}{if (h == hp \&\& EqualRange(T+i, P, m)) return i;}
\Code{2.6}{if (i<n-m) h = Mod((h-T[i]*power)*base+T[i+m], q);}
\Code{1.3}{\}}
\Code{1.3}{return -1;}
\Code{0.0}{\}}
\nopagebreak[4]
\Cost{最坏 \ensuremath{O\!\left(nm\right)}，滚动更新 \ensuremath{O\!\left(1\right)}/次}{\ensuremath{O\!\left(1\right)}}
\par\medskip

\Needspace{21\baselineskip}
\subsubsection{字符串的编辑距离}
\Code{0.0}{EditDistance(A, B) \{}
\Hint{1.3}{diag 保存左上角旧值，up 保存覆盖前的正上方。}\nopagebreak[4]
\Code{1.3}{if (len(A)<len(B)) return EditDistance(B, A);}
\Code{1.3}{for (j = 0; j <= len(B); j++) D[j] = j;}
\Code{1.3}{for (i = 1; i <= len(A); i++) \{}
\Code{2.6}{diag = D[0];}
\Code{2.6}{D[0] = i;}
\Code{2.6}{for (j = 1; j <= len(B); j++) \{}
\Hint{3.9}{删除、插入和替换分别读上方、左方和左上方的状态。}\nopagebreak[4]
\Code{3.9}{up = D[j];}
\Code{3.9}{D[j] = min(1+min(up, D[j-1]), diag+(A[i-1] != B[j-1]));}
\Code{3.9}{diag = up;}
\Code{2.6}{\}}
\Code{1.3}{\}}
\Code{1.3}{return D[len(B)];}
\Code{0.0}{\}}
\nopagebreak[4]
\Cost{\ensuremath{O\!\left(mn\right)}}{\ensuremath{O\!\left(\min (m,n)\right)}}
\par\medskip

\Needspace{22\baselineskip}
\subsubsection{字符串的压缩与解压缩}
\Code{0.0}{EncodeRuns(S, n) \{}
\Hint{1.3}{每段输出字符与次数，解码时按次数重复写出。}\nopagebreak[4]
\Code{1.3}{i = 0;}
\Code{1.3}{while (i<n) \{}
\Code{2.6}{j = i+1;}
\Code{2.6}{while (j<n \&\& S[j] == S[i]) j++;}
\Code{2.6}{Emit(S[i], j-i);}
\Code{2.6}{i = j;}
\Code{1.3}{\}}
\Code{0.0}{\}}
\Hint{0.0}{一段输出后 i 直接移动到 j，下一段从第一个不同字符开始。}\nopagebreak[4]
\Code{0.0}{DecodeRuns(pairs) \{}
\Code{1.3}{for ((c, k) in pairs) \{}
\Code{2.6}{require(k>0);}
\Code{2.6}{repeat (k) Emit(c);}
\Code{1.3}{\}}
\Code{0.0}{\}}
\nopagebreak[4]
\Cost{压 \ensuremath{O\!\left(n\right)}，解 \ensuremath{O\!\left(L\right)}}{\ensuremath{O\!\left(1\right)}}
\par\medskip

\Part{数组}

\Group{数组的基本算法}

\Needspace{18\baselineskip}
\subsubsection{一维数组与多维数组的下标转换}
\Code{0.0}{Flatten(idx, size, d) \{}
\Hint{1.3}{行主序最后一维变化最快，逆变换从末维开始取余。}\nopagebreak[4]
\Code{1.3}{k = 0;}
\Code{1.3}{for (i = 0; i<d; i++) k = k*size[i]+idx[i];}
\Code{1.3}{return k;}
\Code{0.0}{\}}
\Code{0.0}{Unflatten(k, size, d, idx) \{}
\Hint{1.3}{反推时余数是当前维坐标，整数商交给前一维继续分解。}\nopagebreak[4]
\Code{1.3}{for (i = d-1; i >= 0; i--) \{}
\Code{2.6}{idx[i] = k\%size[i];}
\Code{2.6}{k /= size[i];}
\Code{1.3}{\}}
\Code{0.0}{\}}
\nopagebreak[4]
\Cost{\ensuremath{O\!\left(d\right)}}{\ensuremath{O\!\left(1\right)}（输出坐标除外）}
\par\medskip

\Needspace{12\baselineskip}
\subsubsection{多维数组的存储地址计算}
\Code{0.0}{Address(base, idx, low, size, d, bytes) \{}
\Hint{1.3}{先扣除每一维下界，再把元素偏移换算成字节。}\nopagebreak[4]
\Code{1.3}{k = 0;}
\Hint{1.3}{base 是首元素的字节地址；这里按行主序处理，最后一维步长为一。}\nopagebreak[4]
\Code{1.3}{for (i = 0; i<d; i++) k = k*size[i]+idx[i]-low[i];}
\Code{1.3}{return base+bytes*k;}
\Code{0.0}{\}}
\nopagebreak[4]
\Cost{\ensuremath{O\!\left(d\right)}}{\ensuremath{O\!\left(1\right)}}
\par\medskip

\Needspace{12\baselineskip}
\subsubsection{二维数组的遍历}
\Code{0.0}{VisitMatrix(A, R, C) \{}
\Hint{1.3}{固定行号后依次访问各列。}\nopagebreak[4]
\Code{1.3}{for (i = 0; i<R; i++) \{}
\Hint{2.6}{内层列号先耗尽，再把外层行号加一。}\nopagebreak[4]
\Code{2.6}{for (j = 0; j<C; j++) visit(A[i][j]);}
\Code{1.3}{\}}
\Code{0.0}{\}}
\nopagebreak[4]
\Cost{\ensuremath{O\!\left(RC\right)}}{\ensuremath{O\!\left(1\right)}}
\par\medskip

\Needspace{14\baselineskip}
\subsubsection{二维数组的转置}
\Code{0.0}{Transpose(A, R, C) \{}
\Hint{1.3}{新矩阵行数为 C，列数为 R。}\nopagebreak[4]
\Code{1.3}{B = NewMatrix(C, R);}
\Code{1.3}{for (i = 0; i<R; i++) \{}
\Hint{2.6}{原来第 i 行第 j 列写入新矩阵第 j 行第 i 列。}\nopagebreak[4]
\Code{2.6}{for (j = 0; j<C; j++) B[j][i] = A[i][j];}
\Code{1.3}{\}}
\Code{1.3}{return B;}
\Code{0.0}{\}}
\nopagebreak[4]
\Cost{\ensuremath{O\!\left(RC\right)}}{\ensuremath{O\!\left(1\right)}}
\par\medskip

\Needspace{13\baselineskip}
\subsubsection{矩阵的旋转}
\Code{0.0}{Rotate90(A, n) \{}
\Hint{1.3}{沿主对角线交换一次，再逆置每一行得到顺时针旋转。}\nopagebreak[4]
\Code{1.3}{for (i = 0; i<n; i++) \{}
\Code{2.6}{for (j = i+1; j<n; j++) swap(A[i][j], A[j][i]);}
\Code{1.3}{\}}
\Hint{1.3}{两重循环只交换严格上三角的项，否则同一对会被交换两次。}\nopagebreak[4]
\Code{1.3}{for (i = 0; i<n; i++) Reverse(A[i], n);}
\Code{0.0}{\}}
\nopagebreak[4]
\Cost{\ensuremath{O\!\left(n^{2}\right)}}{\ensuremath{O\!\left(1\right)}}
\par\medskip

\Needspace{26\baselineskip}
\subsubsection{矩阵的螺旋遍历}
\Code{0.0}{Spiral(A, R, C) \{}
\Hint{1.3}{下边和左边开始前重验边界，避免单行或单列重复输出。}\nopagebreak[4]
\Code{1.3}{u = l = 0;}
\Code{1.3}{d = R-1;}
\Code{1.3}{r = C-1;}
\Code{1.3}{while (u <= d \&\& l <= r) \{}
\Code{2.6}{for (j = l; j <= r; j++) Emit(A[u][j]);}
\Code{2.6}{u++;}
\Code{2.6}{for (i = u; i <= d; i++) Emit(A[i][r]);}
\Code{2.6}{r--;}
\Code{2.6}{if (u <= d) \{}
\Hint{3.9}{每扫完一边就收缩相应边界，剩余小矩形交给下一轮。}\nopagebreak[4]
\Code{3.9}{for (j = r; j >= l; j--) Emit(A[d][j]);}
\Code{3.9}{d--;}
\Code{2.6}{\}}
\Code{2.6}{if (l <= r) \{}
\Code{3.9}{for (i = d; i >= u; i--) Emit(A[i][l]);}
\Code{3.9}{l++;}
\Code{2.6}{\}}
\Code{1.3}{\}}
\Code{0.0}{\}}
\nopagebreak[4]
\Cost{\ensuremath{O\!\left(RC\right)}}{\ensuremath{O\!\left(1\right)}}
\par\medskip

\Needspace{17\baselineskip}
\subsubsection{有序二维矩阵的查找}
\Code{0.0}{MatrixFind(A, R, C, x) \{}
\Hint{1.3}{从右上角出发，每次可以排除一行或一列。}\nopagebreak[4]
\Code{1.3}{i = 0;}
\Code{1.3}{j = C-1;}
\Code{1.3}{while (i<R \&\& j >= 0) \{}
\Code{2.6}{if (A[i][j] == x) return (i, j);}
\Hint{2.6}{若当前值太大，其所在列在当前行以下的值也不会更小。}\nopagebreak[4]
\Code{2.6}{if (A[i][j]>x) j--;}
\Code{2.6}{else i++;}
\Code{1.3}{\}}
\Code{1.3}{return NONE;}
\Code{0.0}{\}}
\nopagebreak[4]
\Cost{\ensuremath{O\!\left(R+C\right)}}{\ensuremath{O\!\left(1\right)}}
\par\medskip

\Needspace{21\baselineskip}
\subsubsection{矩阵的加法与乘法}
\Code{0.0}{MatrixAdd(A, B, R, C) \{}
\Hint{1.3}{加法逐格对应，乘法每个结果格对应一行与一列的内积。}\nopagebreak[4]
\Code{1.3}{for (i = 0; i<R; i++) for (j = 0; j<C; j++) Z[i][j] = A[i][j]+B[i][j];}
\Code{1.3}{return Z;}
\Code{0.0}{\}}
\Code{0.0}{MatrixMul(A, B, r, k, t) \{}
\Code{1.3}{for (i = 0; i<r; i++) \{}
\Code{2.6}{for (j = 0; j<t; j++) \{}
\Hint{3.9}{内层累加器 s 每换一个结果格就要清零，不能沿用上一格的和。}\nopagebreak[4]
\Code{3.9}{s = 0;}
\Code{3.9}{for (h = 0; h<k; h++) s += A[i][h]*B[h][j];}
\Code{3.9}{Z[i][j] = s;}
\Code{2.6}{\}}
\Code{1.3}{\}}
\Code{1.3}{return Z;}
\Code{0.0}{\}}
\nopagebreak[4]
\Cost{加 \ensuremath{O\!\left(RC\right)}；乘 \ensuremath{O\!\left(rkt\right)}}{\ensuremath{O\!\left(1\right)}}
\par\medskip

\Needspace{14\baselineskip}
\subsubsection{一维前缀和}
\Code{0.0}{Prefix1(A, n, P) \{}
\Hint{1.3}{P[r] 包含下标 r 之前的元素。}\nopagebreak[4]
\Code{1.3}{P[0] = 0;}
\Code{1.3}{for (i = 0; i<n; i++) P[i+1] = P[i]+A[i];}
\Code{0.0}{\}}
\Hint{0.0}{P[0] 为 0，所以查询从下标 0 开始的区间也无需特殊判断。}\nopagebreak[4]
\Code{0.0}{RangeSum(P, l, r) \{}
\Code{1.3}{return P[r]-P[l];}
\Code{0.0}{\}}
\nopagebreak[4]
\Cost{建 \ensuremath{O\!\left(n\right)}，查 \ensuremath{O\!\left(1\right)}}{\ensuremath{O\!\left(n\right)}（索引）}
\par\medskip

\Needspace{16\baselineskip}
\subsubsection{二维前缀和}
\Code{0.0}{Prefix2(A, R, C, P) \{}
\Hint{1.3}{多出的零行和零列让首行、首列也使用相同公式。}\nopagebreak[4]
\Code{1.3}{Zero(P);}
\Code{1.3}{for (i = 0; i<R; i++) \{}
\Code{2.6}{for (j = 0; j<C; j++) P[i+1][j+1] = A[i][j]+P[i][j+1]+P[i+1][j]-P[i][j];}
\Code{1.3}{\}}
\Code{0.0}{\}}
\Hint{0.0}{左上区域被两个前缀各减了一次，查询时须加还一次。}\nopagebreak[4]
\Code{0.0}{RectSum(P, a, b, c, d) \{}
\Code{1.3}{return P[b][d]-P[a][d]-P[b][c]+P[a][c];}
\Code{0.0}{\}}
\nopagebreak[4]
\Cost{建 \ensuremath{O\!\left(RC\right)}，查 \ensuremath{O\!\left(1\right)}}{\ensuremath{O\!\left(RC\right)}（索引）}
\par\medskip

\Needspace{20\baselineskip}
\subsubsection{一维差分数组}
\Code{0.0}{Diff1(A, n, updates) \{}
\Hint{1.3}{左端加值、右端抵消，使前缀累加时改动限于区间内。}\nopagebreak[4]
\Code{1.3}{D[n] = 0;}
\Code{1.3}{for (i = 0; i<n; i++) D[i] = A[i]-(i?A[i-1]:0);}
\Code{1.3}{for ((l, r, v) in updates) \{}
\Code{2.6}{D[l] += v;}
\Code{2.6}{D[r] -= v;}
\Code{1.3}{\}}
\Hint{1.3}{恢复时的运行和 s 包含从首位到当前位置的全部边界变化。}\nopagebreak[4]
\Code{1.3}{s = 0;}
\Code{1.3}{for (i = 0; i<n; i++) \{}
\Code{2.6}{s += D[i];}
\Code{2.6}{A[i] = s;}
\Code{1.3}{\}}
\Code{0.0}{\}}
\nopagebreak[4]
\Cost{建与还原 \ensuremath{O\!\left(n\right)}，每次更新 \ensuremath{O\!\left(1\right)}}{\ensuremath{O\!\left(n\right)}}
\par\medskip

\Needspace{18\baselineskip}
\subsubsection{二维差分数组}
\Code{0.0}{Diff2(A, R, C, updates) \{}
\Hint{1.3}{D 预留末行末列；Get0 在越界时返回 0。}\nopagebreak[4]
\Code{1.3}{Zero(D);}
\Code{1.3}{for (i = 0; i<R; i++) for (j = 0; j<C; j++) D[i][j] = A[i][j]-Get0(A, i-1, j)-Get0(A, i, j-1)+Get0(A, i-1, j-1);}
\Code{1.3}{for ((a, b, c, d, v) in updates) \{}
\Code{2.6}{D[a][c] += v;}
\Code{2.6}{D[b][c] -= v;}
\Hint{2.6}{四角符号依次为加、减、减、加，使矩形外多余的变化恰好抵消。}\nopagebreak[4]
\Code{2.6}{D[a][d] -= v;}
\Code{2.6}{D[b][d] += v;}
\Code{1.3}{\}}
\Code{1.3}{for (i = 0; i<R; i++) for (j = 0; j<C; j++) A[i][j] = D[i][j]+Get0(A, i-1, j)+Get0(A, i, j-1)-Get0(A, i-1, j-1);}
\Code{0.0}{\}}
\nopagebreak[4]
\Cost{建与还原 \ensuremath{O\!\left(RC\right)}，每次更新 \ensuremath{O\!\left(1\right)}}{\ensuremath{O\!\left(RC\right)}}
\par\medskip

\Group{特殊矩阵的压缩存储}

\Needspace{12\baselineskip}
\subsubsection{对称矩阵的压缩存储与下标转换}
\Code{0.0}{SymGet(B, i, j) \{}
\Hint{1.3}{只存下三角，先把上三角坐标镜像过来。}\nopagebreak[4]
\Code{1.3}{if (i<j) swap(i, j);}
\Hint{1.3}{下标公式的前半部分给出本行起点，最后再加列号。}\nopagebreak[4]
\Code{1.3}{base = i*(i+1)/2;}
\Code{1.3}{return B[base+j];}
\Code{0.0}{\}}
\nopagebreak[4]
\Cost{访问 \ensuremath{O\!\left(1\right)}，建 \ensuremath{O\!\left(n^{2}\right)}}{\ensuremath{O\!\left(1\right)}}
\par\medskip

\Needspace{12\baselineskip}
\subsubsection{上三角矩阵的压缩存储与下标转换}
\Code{0.0}{UpperGet(B, n, i, j) \{}
\Hint{1.3}{前 i 行的长度依次为 n、n−1 等。}\nopagebreak[4]
\Code{1.3}{if (i>j) return 0;}
\Hint{1.3}{第 i 行从列 i 开始，因此行内偏移是 j−i。}\nopagebreak[4]
\Code{1.3}{base = i*n-i*(i-1)/2;}
\Code{1.3}{return B[base+j-i];}
\Code{0.0}{\}}
\nopagebreak[4]
\Cost{访问 \ensuremath{O\!\left(1\right)}，建 \ensuremath{O\!\left(n^{2}\right)}}{\ensuremath{O\!\left(1\right)}}
\par\medskip

\Needspace{11\baselineskip}
\subsubsection{下三角矩阵的压缩存储与下标转换}
\Code{0.0}{LowerGet(B, i, j) \{}
\Hint{1.3}{第 i 行前面共存 1 到 i 的和。}\nopagebreak[4]
\Code{1.3}{if (j>i) return 0;}
\Hint{1.3}{上三角区域的零不会占用槽位，访问时直接返回 0。}\nopagebreak[4]
\Code{1.3}{return B[i*(i+1)/2+j];}
\Code{0.0}{\}}
\nopagebreak[4]
\Cost{访问 \ensuremath{O\!\left(1\right)}，建 \ensuremath{O\!\left(n^{2}\right)}}{\ensuremath{O\!\left(1\right)}}
\par\medskip

\Needspace{13\baselineskip}
\subsubsection{对角矩阵的压缩存储与下标转换}
\Code{0.0}{Diagonal(A, n, B) \{}
\Hint{1.3}{主对角线第 i 项的压缩下标仍为 i。}\nopagebreak[4]
\Code{1.3}{for (i = 0; i<n; i++) B[i] = A[i][i];}
\Code{0.0}{\}}
\Hint{0.0}{非对角位置统一返回零，不能用列号去读取不存在的数据。}\nopagebreak[4]
\Code{0.0}{DiagGet(B, i, j) \{}
\Code{1.3}{return i == j ? B[i]:0;}
\Code{0.0}{\}}
\nopagebreak[4]
\Cost{访问 \ensuremath{O\!\left(1\right)}，建 \ensuremath{O\!\left(n\right)}}{\ensuremath{O\!\left(1\right)}}
\par\medskip

\Needspace{17\baselineskip}
\subsubsection{三对角矩阵的压缩存储与下标转换}
\Code{0.0}{TriDiag(A, n, B) \{}
\Hint{1.3}{按行只存主对角线及紧邻的两条对角线。}\nopagebreak[4]
\Code{1.3}{k = 0;}
\Code{1.3}{for (i = 0; i<n; i++) \{}
\Code{2.6}{for (j = max(0, i-1); j <= min(n-1, i+1); j++) B[k++] = A[i][j];}
\Code{1.3}{\}}
\Code{0.0}{\}}
\Hint{0.0}{首行只有两项，最后一行也少一项。公式已经计入首行的缺项。}\nopagebreak[4]
\Code{0.0}{BandGet(B, i, j) \{}
\Code{1.3}{if (abs(i-j)>1) return 0;}
\Code{1.3}{return B[2*i+j];}
\Code{0.0}{\}}
\nopagebreak[4]
\Cost{访问 \ensuremath{O\!\left(1\right)}，建 \ensuremath{O\!\left(n\right)}}{\ensuremath{O\!\left(1\right)}}
\par\medskip

\Group{稀疏矩阵的三元组表示}

\Needspace{16\baselineskip}
\subsubsection{稀疏矩阵与三元组表的相互转换}
\Code{0.0}{ToTriples(A, R, C) \{}
\Hint{1.3}{零元由缺省值表达，恢复时必须先清零。}\nopagebreak[4]
\Code{1.3}{for (i = 0; i<R; i++) for (j = 0; j<C; j++) if (A[i][j] != 0) Append(\&T, (i, j, A[i][j]));}
\Code{1.3}{return T;}
\Code{0.0}{\}}
\Code{0.0}{ToDense(T, R, C) \{}
\Hint{1.3}{每条三元组保存坐标和值；逆转换时覆盖到对应槽位即可。}\nopagebreak[4]
\Code{1.3}{B = ZeroMatrix(R, C);}
\Code{1.3}{for ((i, j, v) in T) B[i][j] = v;}
\Code{1.3}{return B;}
\Code{0.0}{\}}
\nopagebreak[4]
\Cost{\ensuremath{O\!\left(RC+t\right)}}{\ensuremath{O\!\left(1\right)}}
\par\medskip

\Needspace{13\baselineskip}
\subsubsection{稀疏矩阵的普通转置}
\Code{0.0}{SlowTranspose(T, C) \{}
\Hint{1.3}{原列号就是新行号，逐列筛出对应项。}\nopagebreak[4]
\Code{1.3}{for (j = 0; j<C; j++) \{}
\Code{2.6}{for (t in T) if (t.col == j) Append(\&B, (j, t.row, t.val));}
\Code{1.3}{\}}
\Hint{1.3}{结果按原列号递增输出，所以自然是新矩阵的行主序。}\nopagebreak[4]
\Code{1.3}{return B;}
\Code{0.0}{\}}
\nopagebreak[4]
\Cost{\ensuremath{O\!\left(C(t+1)\right)}}{\ensuremath{O\!\left(1\right)}}
\par\medskip

\Needspace{21\baselineskip}
\subsubsection{稀疏矩阵的快速转置}
\Code{0.0}{FastTranspose(T, C) \{}
\Hint{1.3}{先求每个新行的起始槽，pos 在写入后前进一步。}\nopagebreak[4]
\Code{1.3}{Zero(count);}
\Code{1.3}{for (t in T) count[t.col]++;}
\Code{1.3}{sum = 0;}
\Code{1.3}{for (j = 0; j<C; j++) \{}
\Code{2.6}{pos[j] = sum;}
\Code{2.6}{sum += count[j];}
\Code{1.3}{\}}
\Hint{1.3}{count 保存每个新行的长度；前缀累加得到每个新行的第一空槽。}\nopagebreak[4]
\Code{1.3}{for (t in T) \{}
\Code{2.6}{k = pos[t.col]++;}
\Code{2.6}{B[k] = (t.col, t.row, t.val);}
\Code{1.3}{\}}
\Code{1.3}{return B;}
\Code{0.0}{\}}
\nopagebreak[4]
\Cost{\ensuremath{O\!\left(C+t\right)}}{\ensuremath{O\!\left(C\right)}}
\par\medskip

\Needspace{21\baselineskip}
\subsubsection{稀疏矩阵的加法}
\Code{0.0}{SparseAdd(A, B) \{}
\Hint{1.3}{按行列二元坐标归并，同位置抵消后不存零项。}\nopagebreak[4]
\Code{1.3}{i = j = 0;}
\Code{1.3}{while (i<m || j<n) \{}
\Code{2.6}{if (j == n || (i<m \&\& Key(A[i])<Key(B[j]))) Append(\&Z, A[i++]);}
\Code{2.6}{else if (i == m || Key(B[j])<Key(A[i])) Append(\&Z, B[j++]);}
\Code{2.6}{else \{}
\Code{3.9}{v = A[i].val+B[j].val;}
\Hint{3.9}{遇同坐标时两条输入都要推进，即使相加为零也已经处理完这两项。}\nopagebreak[4]
\Code{3.9}{if (v) Append(\&Z, (A[i].row, A[i].col, v));}
\Code{3.9}{i++;}
\Code{3.9}{j++;}
\Code{2.6}{\}}
\Code{1.3}{\}}
\Code{1.3}{return Z;}
\Code{0.0}{\}}
\nopagebreak[4]
\Cost{\ensuremath{O\!\left(t_{1}+t_{2}\right)}}{\ensuremath{O\!\left(1\right)}}
\par\medskip

\Needspace{17\baselineskip}
\subsubsection{稀疏矩阵的乘法}
\Code{0.0}{SparseMul(A, B) \{}
\Hint{1.3}{A 的列号直接定位 B 的对应行，用 acc 合并同一输出列的乘积。}\nopagebreak[4]
\Code{1.3}{IndexRows(B);}
\Code{1.3}{for (i = 0; i<R; i++) \{}
\Code{2.6}{Zero(acc);}
\Code{2.6}{for (a in Row(A, i)) \{}
\Hint{3.9}{每换一个结果行都要清 acc，避免把上一行的积加进来。}\nopagebreak[4]
\Code{3.9}{for (b in Row(B, a.col)) acc[b.col] += a.val*b.val;}
\Code{2.6}{\}}
\Code{2.6}{for (j = 0; j<C; j++) if (acc[j]) Emit(i, j, acc[j]);}
\Code{1.3}{\}}
\Code{0.0}{\}}
\nopagebreak[4]
\Cost{\ensuremath{O\!\left(R\,C+t_{1}+t_{2}+k+P\right)}}{\ensuremath{O\!\left(k+C\right)}}
\par\medskip

\Group{稀疏矩阵的十字链表表示}

\Needspace{16\baselineskip}
\subsubsection{十字链表的建立}
\Code{0.0}{BuildCross(triples) \{}
\Hint{1.3}{按行主序输入；同一结点分别接入行尾与列尾。哨兵作为初始尾。}\nopagebreak[4]
\Code{1.3}{InitRowColHeads();}
\Code{1.3}{for ((i, j, v) in triples) \{}
\Code{2.6}{p = NewCross(i, j, v);}
\Code{2.6}{rt[i]->right = p;}
\Hint{2.6}{rt 和 ct 只移动尾指针，行头与列头仍保留，查找时从相应哨兵开始。}\nopagebreak[4]
\Code{2.6}{ct[j]->down = p;}
\Code{2.6}{rt[i] = ct[j] = p;}
\Code{1.3}{\}}
\Code{0.0}{\}}
\nopagebreak[4]
\Cost{\ensuremath{O\!\left(R+C+t\right)}（行主序输入）}{\ensuremath{O\!\left(R+C\right)}}
\par\medskip

\Needspace{15\baselineskip}
\subsubsection{十字链表的遍历}
\Code{0.0}{VisitCross(M, byRow) \{}
\Hint{1.3}{每次只沿一种链接，可保证每个非零项恰访问一次。}\nopagebreak[4]
\Code{1.3}{if (byRow) \{}
\Code{2.6}{for (i = 0; i<R; i++) for (p = M.row[i]->right; p; p = p->right) visit(p);}
\Code{1.3}{\}}
\Hint{1.3}{横向访问用 right，纵向访问用 down，两种后继不能混用。}\nopagebreak[4]
\Code{1.3}{else \{}
\Code{2.6}{for (j = 0; j<C; j++) for (p = M.col[j]->down; p; p = p->down) visit(p);}
\Code{1.3}{\}}
\Code{0.0}{\}}
\nopagebreak[4]
\Cost{按行 \ensuremath{O\!\left(R+t\right)}，按列 \ensuremath{O\!\left(C+t\right)}}{\ensuremath{O\!\left(1\right)}}
\par\medskip

\Needspace{12\baselineskip}
\subsubsection{非零元素的查找}
\Code{0.0}{CrossFind(M, i, j) \{}
\Hint{1.3}{行链列号递增，超过 j 即可停止。}\nopagebreak[4]
\Code{1.3}{p = M.row[i]->right;}
\Hint{1.3}{只走到首个列号不小于 j 的项，再检查是否恰好等于 j。}\nopagebreak[4]
\Code{1.3}{while (p \&\& p->col<j) p = p->right;}
\Code{1.3}{return p \&\& p->col == j ? p:NULL;}
\Code{0.0}{\}}
\nopagebreak[4]
\Cost{\ensuremath{O\!\left(r+1\right)}（r 为行非零数）}{\ensuremath{O\!\left(1\right)}}
\par\medskip

\Needspace{28\baselineskip}
\subsubsection{非零元素的插入与删除}
\Code{0.0}{CrossSet(M, i, j, v) \{}
\Hint{1.3}{a、b 分别是行链与列链的前驱。删除时两条链都要断开。}\nopagebreak[4]
\Code{1.3}{a = M.row[i];}
\Code{1.3}{b = M.col[j];}
\Code{1.3}{while (a->right \&\& a->right->col<j) a = a->right;}
\Code{1.3}{while (b->down \&\& b->down->row<i) b = b->down;}
\Code{1.3}{p = a->right;}
\Code{1.3}{if (p \&\& p->col == j) \{}
\Code{2.6}{if (v) p->val = v;}
\Code{2.6}{else \{}
\Code{3.9}{a->right = p->right;}
\Code{3.9}{b->down = p->down;}
\Hint{3.9}{替换已有非零值只改 val，不会改变行链和列链的顺序。}\nopagebreak[4]
\Code{3.9}{free(p);}
\Code{2.6}{\}}
\Code{1.3}{\}}
\Code{1.3}{else if (v) \{}
\Code{2.6}{p = NewCross(i, j, v);}
\Code{2.6}{p->right = a->right;}
\Code{2.6}{p->down = b->down;}
\Code{2.6}{a->right = b->down = p;}
\Code{1.3}{\}}
\Code{0.0}{\}}
\nopagebreak[4]
\Cost{\ensuremath{O\!\left(r+c+1\right)}（行、列非零数）}{\ensuremath{O\!\left(1\right)}}
\par\medskip

\Needspace{8\baselineskip}
\subsubsection{基于十字链表的稀疏矩阵加法}
\Code{0.0}{CrossAdd(A, B) \{}
\Hint{1.3}{行内同步前进，AppendCross 尾接当前行与对应列，二者共享结点。}\nopagebreak[4]
\Code{1.3}{InitCross(\&Z);}
\Code{1.3}{for (i = 0; i<R; i++) \{}
\Code{2.6}{a = A.row[i]->right;}
\Code{2.6}{b = B.row[i]->right;}
\Code{2.6}{while (a || b) \{}
\Code{3.9}{if (!b || (a \&\& a->col<b->col)) \{}
\Code{5.2}{AppendCross(\&Z, i, a->col, a->val);}
\Code{5.2}{a = a->right;}
\Code{3.9}{\}}
\Code{3.9}{else if (!a || b->col<a->col) \{}
\Code{5.2}{AppendCross(\&Z, i, b->col, b->val);}
\Code{5.2}{b = b->right;}
\Code{3.9}{\}}
\Code{3.9}{else \{}
\Code{5.2}{v = a->val+b->val;}
\Code{5.2}{if (v) AppendCross(\&Z, i, a->col, v);}
\Hint{5.2}{两侧当前列号相等时共同前进；若其中一行先耗尽，则复制另一行余项。}\nopagebreak[4]
\Code{5.2}{a = a->right;}
\Code{5.2}{b = b->right;}
\Code{3.9}{\}}
\Code{2.6}{\}}
\Code{1.3}{\}}
\Code{1.3}{return Z;}
\Code{0.0}{\}}
\Code{0.0}{AppendCross(Z, i, j, v) \{}
\Code{1.3}{p = NewCross(i, j, v);}
\Code{1.3}{Z.rt[i]->right = p;}
\Code{1.3}{Z.ct[j]->down = p;}
\Code{1.3}{Z.rt[i] = Z.ct[j] = p;}
\Code{0.0}{\}}
\nopagebreak[4]
\Cost{\ensuremath{O\!\left(R+C+t_{1}+t_{2}\right)}}{\ensuremath{O\!\left(C+1\right)}}
\par\medskip

\Needspace{22\baselineskip}
\subsubsection{基于十字链表的稀疏矩阵转置}
\Code{0.0}{CrossTranspose(M) \{}
\Hint{1.3}{交换域之前保存原行后继；哨兵本身的首边也要换向。}\nopagebreak[4]
\Code{1.3}{for (i = 0; i<R; i++) \{}
\Code{2.6}{p = M.row[i]->right;}
\Code{2.6}{while (p) \{}
\Code{3.9}{next = p->right;}
\Code{3.9}{swap(p->row, p->col);}
\Code{3.9}{swap(p->right, p->down);}
\Code{3.9}{p = next;}
\Code{2.6}{\}}
\Hint{2.6}{行头数组和列头数组也要互换，否则入口指针与已经换向的链域不匹配。}\nopagebreak[4]
\Code{2.6}{M.row[i]->down = M.row[i]->right;}
\Code{1.3}{\}}
\Code{1.3}{for (j = 0; j<C; j++) M.col[j]->right = M.col[j]->down;}
\Code{1.3}{swap(M.row, M.col);}
\Code{1.3}{swap(R, C);}
\Code{0.0}{\}}
\nopagebreak[4]
\Cost{\ensuremath{O\!\left(R+C+t\right)}}{\ensuremath{O\!\left(1\right)}}
\par\medskip

\Part{广义表}

\Group{广义表的基本操作}

\Needspace{13\baselineskip}
\subsubsection{广义表的头尾链表表示}
\Code{0.0}{BuildHT(e) \{}
\Hint{1.3}{Cons 的 hp 保存当前元素，tp 指向后面的同层表。}\nopagebreak[4]
\Code{1.3}{if (IsAtom(e)) return NewAtom(Value(e));}
\Code{1.3}{tail = NIL;}
\Hint{1.3}{从最后一项往前构造，头插后的结果才与原来的书写次序一致。}\nopagebreak[4]
\Code{1.3}{for (i = Length(e)-1; i >= 0; i--) tail = Cons(BuildHT(e[i]), tail);}
\Code{1.3}{return tail;}
\Code{0.0}{\}}
\nopagebreak[4]
\Cost{建 \ensuremath{O\!\left(N\right)}}{\ensuremath{O\!\left(N\right)}（结构）}
\par\medskip

\Needspace{19\baselineskip}
\subsubsection{广义表的扩展线性链表表示}
\Code{0.0}{BuildExt(e) \{}
\Hint{1.3}{child 连下层首元素，next 连同层兄弟。}\nopagebreak[4]
\Code{1.3}{p = NewTaggedNode(Tag(e));}
\Code{1.3}{if (IsAtom(e)) p->value = Value(e);}
\Code{1.3}{else \{}
\Code{2.6}{slot = \&p->child;}
\Code{2.6}{for (x in e) \{}
\Hint{3.9}{slot 在每次接入后移动到新结点的 next 域，为下一兄弟留好入口。}\nopagebreak[4]
\Code{3.9}{*slot = BuildExt(x);}
\Code{3.9}{slot = \&(*slot)->next;}
\Code{2.6}{\}}
\Code{1.3}{\}}
\Code{1.3}{return p;}
\Code{0.0}{\}}
\nopagebreak[4]
\Cost{建 \ensuremath{O\!\left(N\right)}}{\ensuremath{O\!\left(N\right)}（结构）}
\par\medskip

\Needspace{22\baselineskip}
\subsubsection{广义表的创建}
\Code{0.0}{ParseList(stream) \{}
\Hint{1.3}{每次子调用消费一个完整元素，逗号只分隔本层元素。}\nopagebreak[4]
\Code{1.3}{t = ReadToken(stream);}
\Code{1.3}{if (IsAtomToken(t)) return NewAtom(t);}
\Code{1.3}{require(t == '(');}
\Code{1.3}{InitList(\&L);}
\Code{1.3}{if (!IsRightParen(PeekToken(stream))) \{}
\Code{2.6}{Append(\&L, ParseList(stream));}
\Code{2.6}{while (PeekToken(stream) == ', ') \{}
\Hint{3.9}{读到左括号后紧接右括号，说明这是空表，一个子项也不用读。}\nopagebreak[4]
\Code{3.9}{ReadToken(stream);}
\Code{3.9}{Append(\&L, ParseList(stream));}
\Code{2.6}{\}}
\Code{1.3}{\}}
\Code{1.3}{require(ReadToken(stream) == ')');}
\Code{1.3}{return L;}
\Code{0.0}{\}}
\nopagebreak[4]
\Cost{\ensuremath{O\!\left(L\right)}}{\ensuremath{O\!\left(h\right)}}
\par\medskip

\Needspace{11\baselineskip}
\subsubsection{求广义表的表头（GetHead）}
\Code{0.0}{GetHead(L) \{}
\Hint{1.3}{表头是第一个元素本身，可为原子或子表。}\nopagebreak[4]
\Code{1.3}{require(L != NIL \&\& IsCons(L));}
\Hint{1.3}{若首元素本身是子表，函数直接返回这个子表的引用。}\nopagebreak[4]
\Code{1.3}{return L->hp;}
\Code{0.0}{\}}
\nopagebreak[4]
\Cost{\ensuremath{O\!\left(1\right)}}{\ensuremath{O\!\left(1\right)}}
\par\medskip

\Needspace{11\baselineskip}
\subsubsection{求广义表的表尾（GetTail）}
\Code{0.0}{GetTail(L) \{}
\Hint{1.3}{摘去首元素后仍为表，单元素表的尾是空表。}\nopagebreak[4]
\Code{1.3}{require(L != NIL \&\& IsCons(L));}
\Hint{1.3}{返回值是余下所有元素组成的表，而非原表的第二个元素。}\nopagebreak[4]
\Code{1.3}{return L->tp;}
\Code{0.0}{\}}
\nopagebreak[4]
\Cost{\ensuremath{O\!\left(1\right)}}{\ensuremath{O\!\left(1\right)}}
\par\medskip

\Needspace{12\baselineskip}
\subsubsection{求广义表的长度}
\Code{0.0}{ListLength(L) \{}
\Hint{1.3}{只沿同层 tp 计数，子表整体算一个元素。}\nopagebreak[4]
\Code{1.3}{k = 0;}
\Hint{1.3}{每跨过一个 Cons 单元计数一次，无需深入 hp。}\nopagebreak[4]
\Code{1.3}{for (p = L; p != NIL; p = p->tp) k++;}
\Code{1.3}{return k;}
\Code{0.0}{\}}
\nopagebreak[4]
\Cost{\ensuremath{O\!\left(k\right)}}{\ensuremath{O\!\left(1\right)}}
\par\medskip

\Needspace{13\baselineskip}
\subsubsection{求广义表的深度}
\Code{0.0}{ListDepth(e) \{}
\Hint{1.3}{原子深度 0，空表深度 1。}\nopagebreak[4]
\Code{1.3}{if (IsAtom(e)) return 0;}
\Code{1.3}{m = 0;}
\Hint{1.3}{求所有孩子深度的最大值，然后加上包住它们的这一层。}\nopagebreak[4]
\Code{1.3}{for (p = e; p != NIL; p = p->tp) m = max(m, ListDepth(p->hp));}
\Code{1.3}{return 1+m;}
\Code{0.0}{\}}
\nopagebreak[4]
\Cost{\ensuremath{O\!\left(N\right)}}{\ensuremath{O\!\left(h\right)}}
\par\medskip

\Needspace{13\baselineskip}
\subsubsection{统计广义表中的原子个数}
\Code{0.0}{CountAtoms(e) \{}
\Hint{1.3}{遇原子加一，空表循环零次。}\nopagebreak[4]
\Code{1.3}{if (IsAtom(e)) return 1;}
\Code{1.3}{n = 0;}
\Hint{1.3}{只遇到原子时返回 1；括号层本身不增加计数。}\nopagebreak[4]
\Code{1.3}{for (p = e; p != NIL; p = p->tp) n += CountAtoms(p->hp);}
\Code{1.3}{return n;}
\Code{0.0}{\}}
\nopagebreak[4]
\Cost{\ensuremath{O\!\left(N\right)}}{\ensuremath{O\!\left(h\right)}}
\par\medskip

\Group{广义表的递归算法}

\Needspace{14\baselineskip}
\subsubsection{广义表的递归遍历}
\Code{0.0}{WalkList(e) \{}
\Hint{1.3}{遇子表先深入，返回后再访问同层后续元素。}\nopagebreak[4]
\Code{1.3}{if (IsAtom(e)) \{}
\Code{2.6}{visit(e->value);}
\Hint{2.6}{遇到空表时没有任何表单元，因此本次调用不会产生输出。}\nopagebreak[4]
\Code{2.6}{return;}
\Code{1.3}{\}}
\Code{1.3}{for (p = e; p != NIL; p = p->tp) WalkList(p->hp);}
\Code{0.0}{\}}
\nopagebreak[4]
\Cost{\ensuremath{O\!\left(N\right)}}{\ensuremath{O\!\left(h\right)}}
\par\medskip

\Needspace{17\baselineskip}
\subsubsection{广义表的复制}
\Code{0.0}{CopyGL(e) \{}
\Hint{1.3}{复制每个表单元，尾链接引用指向新分配的槽。}\nopagebreak[4]
\Code{1.3}{if (IsAtom(e)) return NewAtom(e->value);}
\Code{1.3}{head = NIL;}
\Code{1.3}{slot = \&head;}
\Code{1.3}{for (p = e; p != NIL; p = p->tp) \{}
\Hint{2.6}{slot 保存可写指针域的位置，使空表与非空表的追加操作保持一致。}\nopagebreak[4]
\Code{2.6}{*slot = Cons(CopyGL(p->hp), NIL);}
\Code{2.6}{slot = \&(*slot)->tp;}
\Code{1.3}{\}}
\Code{1.3}{return head;}
\Code{0.0}{\}}
\nopagebreak[4]
\Cost{\ensuremath{O\!\left(N\right)}}{\ensuremath{O\!\left(h\right)}}
\par\medskip

\Needspace{19\baselineskip}
\subsubsection{广义表的销毁}
\Code{0.0}{FreeGL(e) \{}
\Hint{1.3}{只适用于无共享、无环结构；先释放内容后释放表单元。}\nopagebreak[4]
\Code{1.3}{if (IsAtom(e)) \{}
\Code{2.6}{free(e);}
\Code{2.6}{return;}
\Code{1.3}{\}}
\Code{1.3}{while (e != NIL) \{}
\Hint{2.6}{当前单元一旦释放就不能再访问其 tp，因此必须预存 next。}\nopagebreak[4]
\Code{2.6}{next = e->tp;}
\Code{2.6}{FreeGL(e->hp);}
\Code{2.6}{free(e);}
\Code{2.6}{e = next;}
\Code{1.3}{\}}
\Code{0.0}{\}}
\nopagebreak[4]
\Cost{\ensuremath{O\!\left(N\right)}}{\ensuremath{O\!\left(h\right)}}
\par\medskip

\Needspace{19\baselineskip}
\subsubsection{广义表的逆置}
\Code{0.0}{ReverseGL(e) \{}
\Hint{1.3}{递归翻转每个子表，同时改向本层的 tp。}\nopagebreak[4]
\Code{1.3}{if (e == NIL || IsAtom(e)) return e;}
\Code{1.3}{prev = NIL;}
\Code{1.3}{while (e != NIL) \{}
\Code{2.6}{next = e->tp;}
\Code{2.6}{e->hp = ReverseGL(e->hp);}
\Hint{2.6}{子表内容与同层链接分别处理，不能只反转最外层。}\nopagebreak[4]
\Code{2.6}{e->tp = prev;}
\Code{2.6}{prev = e;}
\Code{2.6}{e = next;}
\Code{1.3}{\}}
\Code{1.3}{return prev;}
\Code{0.0}{\}}
\nopagebreak[4]
\Cost{\ensuremath{O\!\left(N\right)}}{\ensuremath{O\!\left(h\right)}}
\par\medskip

\Needspace{14\baselineskip}
\subsubsection{广义表的扁平化}
\Code{0.0}{FlattenGL(e, out) \{}
\Hint{1.3}{沿原次序收集叶端原子，各层共用同一个输出表。}\nopagebreak[4]
\Code{1.3}{if (IsAtom(e)) \{}
\Code{2.6}{Append(out, e->value);}
\Hint{2.6}{共同使用一个输出容器，子调用返回后无需再拼接临时表。}\nopagebreak[4]
\Code{2.6}{return;}
\Code{1.3}{\}}
\Code{1.3}{for (p = e; p != NIL; p = p->tp) FlattenGL(p->hp, out);}
\Code{0.0}{\}}
\nopagebreak[4]
\Cost{\ensuremath{O\!\left(N\right)}}{\ensuremath{O\!\left(h\right)}}
\par\medskip

\Needspace{15\baselineskip}
\subsubsection{查找广义表中的指定原子}
\Code{0.0}{FindAtom(e, x) \{}
\Hint{1.3}{子调用返回非空时立即把匹配地址传回上层。}\nopagebreak[4]
\Code{1.3}{if (IsAtom(e)) return e->value == x ? e:NULL;}
\Code{1.3}{for (p = e; p != NIL; p = p->tp) \{}
\Code{2.6}{hit = FindAtom(p->hp, x);}
\Hint{2.6}{两个分支都找不到时才返回空，调用方由此继续检查下一项。}\nopagebreak[4]
\Code{2.6}{if (hit) return hit;}
\Code{1.3}{\}}
\Code{1.3}{return NULL;}
\Code{0.0}{\}}
\nopagebreak[4]
\Cost{\ensuremath{O\!\left(N\right)}}{\ensuremath{O\!\left(h\right)}}
\par\medskip

\Needspace{22\baselineskip}
\subsubsection{删除广义表中的指定原子}
\Code{0.0}{DeleteAtoms(Node **L, x) \{}
\Hint{1.3}{slot 指向进入当前表单元的链接，删除后继续检查这个位置。}\nopagebreak[4]
\Code{1.3}{slot = \&(*L);}
\Code{1.3}{while (*slot != NIL) \{}
\Code{2.6}{p = *slot;}
\Code{2.6}{if (IsAtom(p->hp) \&\& p->hp->value == x) \{}
\Code{3.9}{*slot = p->tp;}
\Code{3.9}{free(p->hp);}
\Code{3.9}{free(p);}
\Code{2.6}{\}}
\Hint{2.6}{只在保留当前项时移动 slot，避免漏删连续的目标原子。}\nopagebreak[4]
\Code{2.6}{else \{}
\Code{3.9}{if (!IsAtom(p->hp)) DeleteAtoms(\&p->hp, x);}
\Code{3.9}{slot = \&p->tp;}
\Code{2.6}{\}}
\Code{1.3}{\}}
\Code{0.0}{\}}
\nopagebreak[4]
\Cost{\ensuremath{O\!\left(N\right)}}{\ensuremath{O\!\left(h\right)}}
\par\medskip

\Needspace{16\baselineskip}
\subsubsection{判断两个广义表是否相等}
\Code{0.0}{EqualGL(a, b) \{}
\Hint{1.3}{对应元素都相等还不够，两层链也必须同时耗尽。}\nopagebreak[4]
\Code{1.3}{if (IsAtom(a) || IsAtom(b)) return IsAtom(a) \&\& IsAtom(b) \&\& a->value == b->value;}
\Code{1.3}{while (a != NIL \&\& b != NIL) \{}
\Code{2.6}{if (!EqualGL(a->hp, b->hp)) return false;}
\Code{2.6}{a = a->tp;}
\Hint{2.6}{元素顺序也要相同，比较是逐对进行的。}\nopagebreak[4]
\Code{2.6}{b = b->tp;}
\Code{1.3}{\}}
\Code{1.3}{return a == NIL \&\& b == NIL;}
\Code{0.0}{\}}
\nopagebreak[4]
\Cost{\ensuremath{O\!\left(N_{1}+N_{2}\right)}}{\ensuremath{O\!\left(h\right)}}
\par\medskip

\Needspace{24\baselineskip}
\subsubsection{广义表与字符串的相互转换}
\Code{0.0}{WriteGL(e, out) \{}
\Hint{1.3}{逗号只写在相邻元素之间；读回使用 ParseList。}\nopagebreak[4]
\Code{1.3}{if (IsAtom(e)) \{}
\Code{2.6}{WriteToken(out, e->value);}
\Code{2.6}{return;}
\Code{1.3}{\}}
\Code{1.3}{WriteToken(out, '(');}
\Code{1.3}{for (p = e; p != NIL; p = p->tp) \{}
\Code{2.6}{if (p != e) WriteToken(out, ', ');}
\Code{2.6}{WriteGL(p->hp, out);}
\Code{1.3}{\}}
\Hint{1.3}{原子需使用可区分于逗号和括号的词元编码，顶层解析后应到输入末尾。}\nopagebreak[4]
\Code{1.3}{WriteToken(out, ')');}
\Code{0.0}{\}}
\Code{0.0}{ReadGL(in) \{}
\Code{1.3}{e = ParseList(in);}
\Code{1.3}{require(AtEnd(in));}
\Code{1.3}{return e;}
\Code{0.0}{\}}
\nopagebreak[4]
\Cost{\ensuremath{O\!\left(L+N\right)}}{\ensuremath{O\!\left(h\right)}}
\par\medskip

\Needspace{22\baselineskip}
\subsubsection{广义表与普通树的相互转换}
\Code{0.0}{TreeToGL(t) \{}
\Hint{1.3}{编码首项存根值，后续每一项各代表一棵完整子树。}\nopagebreak[4]
\Code{1.3}{InitList(\&L);}
\Code{1.3}{Append(\&L, NewAtom(t->value));}
\Code{1.3}{for (c = t->child; c; c = c->sibling) Append(\&L, TreeToGL(c));}
\Code{1.3}{return L;}
\Code{0.0}{\}}
\Code{0.0}{GLToTree(L) \{}
\Code{1.3}{t = NewTree(GetHead(L)->value);}
\Hint{1.3}{恢复时从表尾枚举孩子，每次把新孩子接到上一孩子的兄弟指针上。}\nopagebreak[4]
\Code{1.3}{slot = \&t->child;}
\Code{1.3}{for (p = GetTail(L); p != NIL; p = p->tp) \{}
\Code{2.6}{*slot = GLToTree(p->hp);}
\Code{2.6}{slot = \&(*slot)->sibling;}
\Code{1.3}{\}}
\Code{1.3}{return t;}
\Code{0.0}{\}}
\nopagebreak[4]
\Cost{\ensuremath{O\!\left(N\right)}}{\ensuremath{O\!\left(h\right)}}
\par\medskip

\Part{树}

\Group{二叉树的遍历}

\Needspace{27\baselineskip}
\subsubsection{前序遍历（递归、非递归）}
\Code{0.0}{void PreOrder(BiTNode *T) \{}
\Hint{1.3}{前序在第一次到达结点时访问。栈中保留的是尚待转向右子树的祖先。}\nopagebreak[4]
\Code{1.3}{if (!T) return;}
\Code{1.3}{visit(T);}
\Code{1.3}{PreOrder(T->lchild);}
\Code{1.3}{PreOrder(T->rchild);}
\Code{0.0}{\}}
\Code{0.0}{void PreOrderIter(BiTNode *T) \{}
\Code{1.3}{Stack S;}
\Code{1.3}{InitStack(\&S);}
\Code{1.3}{BiTNode *p = T;}
\Code{1.3}{while (p || !Empty(\&S)) \{}
\Code{2.6}{while (p) \{}
\Code{3.9}{visit(p);}
\Hint{3.9}{把当前结点的地址留下，左侧走完后才能找回它的右孩子。}\nopagebreak[4]
\Code{3.9}{Push(\&S, p);}
\Code{3.9}{p = p->lchild;}
\Code{2.6}{\}}
\Hint{2.6}{从栈中取回 p 时它已经访问过，下一步直接进入它的右孩子。}\nopagebreak[4]
\Code{2.6}{Pop(\&S, \&p);}
\Code{2.6}{p = p->rchild;}
\Code{1.3}{\}}
\Code{0.0}{\}}
\nopagebreak[4]
\Cost{\ensuremath{O\!\left(n\right)}}{\ensuremath{O\!\left(h\right)}}
\par\medskip

\Needspace{26\baselineskip}
\subsubsection{中序遍历（递归、非递归）}
\Code{0.0}{void InOrder(BiTNode *T) \{}
\Hint{1.3}{先压左链。左子树结束后出栈，此时才访问根。}\nopagebreak[4]
\Code{1.3}{if (!T) return;}
\Code{1.3}{InOrder(T->lchild);}
\Code{1.3}{visit(T);}
\Code{1.3}{InOrder(T->rchild);}
\Code{0.0}{\}}
\Code{0.0}{void InOrderIter(BiTNode *T) \{}
\Code{1.3}{InitStack(\&S);}
\Code{1.3}{p = T;}
\Code{1.3}{while (p || !Empty(\&S)) \{}
\Code{2.6}{while (p) \{}
\Code{3.9}{Push(\&S, p);}
\Code{3.9}{p = p->lchild;}
\Code{2.6}{\}}
\Hint{2.6}{与前序迭代的区别在 visit 的位置。中序在弹栈后访问，此时左子树已经完成。}\nopagebreak[4]
\Code{2.6}{p = Pop(\&S);}
\Hint{2.6}{左子树已经处理完，这里访问根。}\nopagebreak[4]
\Code{2.6}{visit(p);}
\Code{2.6}{p = p->rchild;}
\Code{1.3}{\}}
\Code{0.0}{\}}
\nopagebreak[4]
\Cost{\ensuremath{O\!\left(n\right)}}{\ensuremath{O\!\left(h\right)}}
\par\medskip

\Needspace{8\baselineskip}
\subsubsection{后序遍历（递归、非递归）}
\Code{0.0}{void PostOrder(BiTNode *T) \{}
\Hint{1.3}{后序要等右子树也返回。last 记录上一个完成访问的结点。}\nopagebreak[4]
\Code{1.3}{if (!T) return;}
\Code{1.3}{PostOrder(T->lchild);}
\Code{1.3}{PostOrder(T->rchild);}
\Code{1.3}{visit(T);}
\Code{0.0}{\}}
\Code{0.0}{void PostOrderIter(BiTNode *T) \{}
\Code{1.3}{InitStack(\&S);}
\Code{1.3}{p = T;}
\Code{1.3}{last = NULL;}
\Code{1.3}{while (p || !Empty(\&S)) \{}
\Code{2.6}{if (p) \{}
\Code{3.9}{Push(\&S, p);}
\Code{3.9}{p = p->lchild;}
\Code{2.6}{\}}
\Code{2.6}{else \{}
\Code{3.9}{q = Top(\&S);}
\Code{3.9}{if (q->rchild \&\& last != q->rchild) p = q->rchild;}
\Code{3.9}{else \{}
\Code{5.2}{visit(q);}
\Hint{5.2}{仅当右孩子不存在或恰是 last 时，栈顶的两个分支才都已结束。}\nopagebreak[4]
\Hint{5.2}{该结点已处理完，回到父亲时可据 last 判断右子树是否结束。}\nopagebreak[4]
\Code{5.2}{last = Pop(\&S);}
\Code{3.9}{\}}
\Code{2.6}{\}}
\Code{1.3}{\}}
\Code{0.0}{\}}
\nopagebreak[4]
\Cost{\ensuremath{O\!\left(n\right)}}{\ensuremath{O\!\left(h\right)}}
\par\medskip

\Needspace{18\baselineskip}
\subsubsection{层序遍历（队列实现）}
\Code{0.0}{void LevelOrder(BiTNode *T) \{}
\Hint{1.3}{取出本层结点时，按左、右顺序把它的孩子放到队尾。}\nopagebreak[4]
\Code{1.3}{if (!T) return;}
\Code{1.3}{InitQueue(\&Q);}
\Code{1.3}{EnQueue(\&Q, T);}
\Code{1.3}{while (!Empty(\&Q)) \{}
\Hint{2.6}{取出队列中最早进入的结点，该轮访问它。}\nopagebreak[4]
\Code{2.6}{BiTNode *p = DeQueue(\&Q);}
\Code{2.6}{visit(p);}
\Hint{2.6}{这里仅把非空孩子入队，所以每次出队都可以直接 visit。}\nopagebreak[4]
\Code{2.6}{if (p->lchild) EnQueue(\&Q, p->lchild);}
\Code{2.6}{if (p->rchild) EnQueue(\&Q, p->rchild);}
\Code{1.3}{\}}
\Code{0.0}{\}}
\nopagebreak[4]
\Cost{\ensuremath{O\!\left(n\right)}}{\ensuremath{O\!\left(w\right)}}
\par\medskip

\Needspace{26\baselineskip}
\subsubsection{二叉树的之字形层序遍历}
\Code{0.0}{Zigzag(T) \{}
\Hint{1.3}{入队顺序保持从左到右，只改变本层结果的写入下标。}\nopagebreak[4]
\Code{1.3}{if (!T) return;}
\Code{1.3}{EnQueue(\&Q, T);}
\Code{1.3}{forward = true;}
\Code{1.3}{while (!Empty(\&Q)) \{}
\Code{2.6}{k = Size(\&Q);}
\Code{2.6}{for (i = 0; i<k; i++) \{}
\Code{3.9}{p = DeQueue(\&Q);}
\Code{3.9}{row[forward?i:k-1-i] = p->data;}
\Code{3.9}{EnqueueChildren(\&Q, p);}
\Code{2.6}{\}}
\Hint{2.6}{本层长度 k 在出队前固定，新入队的孩子不会混入本轮。}\nopagebreak[4]
\Code{2.6}{Emit(row, k);}
\Code{2.6}{forward = !forward;}
\Code{1.3}{\}}
\Code{0.0}{\}}
\Code{0.0}{EnqueueChildren(Q, p) \{}
\Code{1.3}{if (p->lchild) EnQueue(Q, p->lchild);}
\Code{1.3}{if (p->rchild) EnQueue(Q, p->rchild);}
\Code{0.0}{\}}
\nopagebreak[4]
\Cost{\ensuremath{O\!\left(n\right)}}{\ensuremath{O\!\left(w\right)}}
\par\medskip

\Needspace{29\baselineskip}
\subsubsection{Morris遍历（O(1)辅助空间遍历）}
\Code{0.0}{Morris(T) \{}
\Hint{1.3}{第一次遇到前驱时搭临时回边，第二次回来时拆边并访问根。}\nopagebreak[4]
\Code{1.3}{p = T;}
\Code{1.3}{while (p) \{}
\Code{2.6}{if (!p->lchild) \{}
\Code{3.9}{visit(p);}
\Code{3.9}{p = p->rchild;}
\Code{2.6}{\}}
\Code{2.6}{else \{}
\Code{3.9}{q = p->lchild;}
\Code{3.9}{while (q->rchild \&\& q->rchild != p) q = q->rchild;}
\Code{3.9}{if (!q->rchild) \{}
\Code{5.2}{q->rchild = p;}
\Hint{5.2}{q 原本为空的右域才可用来搭桥；离开时恢复为空，原树结构得以保留。}\nopagebreak[4]
\Code{5.2}{p = p->lchild;}
\Code{3.9}{\}}
\Code{3.9}{else \{}
\Code{5.2}{q->rchild = NULL;}
\Code{5.2}{visit(p);}
\Code{5.2}{p = p->rchild;}
\Code{3.9}{\}}
\Code{2.6}{\}}
\Code{1.3}{\}}
\Code{0.0}{\}}
\nopagebreak[4]
\Cost{\ensuremath{O\!\left(n\right)}}{\ensuremath{O\!\left(1\right)}}
\par\medskip

\Needspace{19\baselineskip}
\subsubsection{线索二叉树的遍历}
\Code{0.0}{ThreadWalk(T) \{}
\Hint{1.3}{tag 为 1 表示线索，不能沿左线索继续下降。}\nopagebreak[4]
\Code{1.3}{p = T;}
\Code{1.3}{while (p \&\& p->ltag == 0 \&\& p->lchild) p = p->lchild;}
\Code{1.3}{while (p) \{}
\Code{2.6}{visit(p);}
\Code{2.6}{if (p->rtag == 1) p = p->rchild;}
\Hint{2.6}{右域若是线索就直达后继，若是真孩子则后继在右子树的最左端。}\nopagebreak[4]
\Code{2.6}{else \{}
\Code{3.9}{p = p->rchild;}
\Code{3.9}{while (p \&\& p->ltag == 0 \&\& p->lchild) p = p->lchild;}
\Code{2.6}{\}}
\Code{1.3}{\}}
\Code{0.0}{\}}
\nopagebreak[4]
\Cost{\ensuremath{O\!\left(n\right)}}{\ensuremath{O\!\left(1\right)}}
\par\medskip

\Group{二叉树的构建与转换}

\Needspace{16\baselineskip}
\subsubsection{根据带空结点标记的前序序列建树}
\Code{0.0}{BiTNode *ReadTree(Stream *in) \{}
\Hint{1.3}{一次调用消费一棵完整子树，专用符号 \# 立即结束当前分支。}\nopagebreak[4]
\Code{1.3}{t = ReadToken(in);}
\Code{1.3}{require(t != EOF);}
\Code{1.3}{if (t == HASH) return NULL;}
\Code{1.3}{p = NewTree(t);}
\Hint{1.3}{顶层读完一棵树还应检查没有剩余词元，可发现多余输入。}\nopagebreak[4]
\Code{1.3}{p->lchild = ReadTree(in);}
\Code{1.3}{p->rchild = ReadTree(in);}
\Code{1.3}{return p;}
\Code{0.0}{\}}
\nopagebreak[4]
\Cost{\ensuremath{O\!\left(n\right)}}{\ensuremath{O\!\left(h\right)}}
\par\medskip

\Needspace{17\baselineskip}
\subsubsection{根据前序、中序序列恢复二叉树}
\Code{0.0}{BuildPre(l, r) \{}
\Hint{1.3}{中序区间为左闭右开，i 是共享前序游标；互异值的位置已建立索引。}\nopagebreak[4]
\Code{1.3}{if (l == r) return NULL;}
\Code{1.3}{x = pre[i++];}
\Code{1.3}{k = pos[x];}
\Code{1.3}{require(l <= k \&\& k<r);}
\Hint{1.3}{先建左子树会推进共享游标 i，接下来读取的位置自然是右子树的根。}\nopagebreak[4]
\Code{1.3}{p = NewTree(x);}
\Code{1.3}{p->lchild = BuildPre(l, k);}
\Code{1.3}{p->rchild = BuildPre(k+1, r);}
\Code{1.3}{return p;}
\Code{0.0}{\}}
\nopagebreak[4]
\Cost{\ensuremath{O\!\left(n\right)}（期望）}{\ensuremath{O\!\left(n\right)}}
\par\medskip

\Needspace{17\baselineskip}
\subsubsection{根据后序、中序序列恢复二叉树}
\Code{0.0}{BuildPost(l, r) \{}
\Hint{1.3}{从后向前读后序，先得到根，再得到右子树。共享游标初值 n−1。}\nopagebreak[4]
\Code{1.3}{if (l == r) return NULL;}
\Code{1.3}{x = post[i--];}
\Code{1.3}{k = pos[x];}
\Code{1.3}{require(l <= k \&\& k<r);}
\Hint{1.3}{右子树先读不代表它挂到左边，递归结果仍分别接入对应孩子域。}\nopagebreak[4]
\Code{1.3}{p = NewTree(x);}
\Code{1.3}{p->rchild = BuildPost(k+1, r);}
\Code{1.3}{p->lchild = BuildPost(l, k);}
\Code{1.3}{return p;}
\Code{0.0}{\}}
\nopagebreak[4]
\Cost{\ensuremath{O\!\left(n\right)}（期望）}{\ensuremath{O\!\left(n\right)}}
\par\medskip

\Needspace{16\baselineskip}
\subsubsection{根据层序、中序序列恢复二叉树}
\Code{0.0}{BuildLevel(l, r) \{}
\Hint{1.3}{rank 是互异值在层序中的名次，区间内最早出现者为根。}\nopagebreak[4]
\Code{1.3}{if (l == r) return NULL;}
\Code{1.3}{k = l;}
\Code{1.3}{for (i = l+1; i<r; i++) if (rank[in[i]]<rank[in[k]]) k = i;}
\Code{1.3}{p = NewTree(in[k]);}
\Hint{1.3}{层序输入中排名最小的那个结点，一定先于同一子树内的后代出现。}\nopagebreak[4]
\Code{1.3}{p->lchild = BuildLevel(l, k);}
\Code{1.3}{p->rchild = BuildLevel(k+1, r);}
\Code{1.3}{return p;}
\Code{0.0}{\}}
\nopagebreak[4]
\Cost{\ensuremath{O\!\left(n^{2}\right)}}{\ensuremath{O\!\left(n\right)}}
\par\medskip

\Needspace{25\baselineskip}
\subsubsection{二叉树的序列化与反序列化}
\Code{0.0}{void Serialize(BiTNode *T, Stream *out) \{}
\Hint{1.3}{空孩子也要输出一个 \#，读回时才能判断缺的是左孩子还是右孩子。}\nopagebreak[4]
\Code{1.3}{if (T == NULL) \{}
\Code{2.6}{WriteToken(out, "\#");}
\Code{2.6}{return;}
\Code{1.3}{\}}
\Code{1.3}{WriteToken(out, T->data);}
\Code{1.3}{Serialize(T->lchild, out);}
\Code{1.3}{Serialize(T->rchild, out);}
\Code{0.0}{\}}
\Code{0.0}{BiTNode *Deserialize(Stream *in) \{}
\Code{1.3}{t = ReadToken(in);}
\Code{1.3}{require(t != EOF);}
\Hint{1.3}{两个空孩子需要两个标记；省去其中任一标记都会丢失形状信息。}\nopagebreak[4]
\Code{1.3}{if (t == HASH) return NULL;}
\Code{1.3}{p = NewTree(t);}
\Code{1.3}{p->lchild = Deserialize(in);}
\Code{1.3}{p->rchild = Deserialize(in);}
\Code{1.3}{return p;}
\Code{0.0}{\}}
\nopagebreak[4]
\Cost{\ensuremath{O\!\left(n\right)}}{\ensuremath{O\!\left(h\right)}}
\par\medskip

\Needspace{21\baselineskip}
\subsubsection{树、森林与二叉树的相互转换}
\Code{0.0}{ForestToBinary(p) \{}
\Hint{1.3}{p 是森林的第一个根；左指针接首孩子，右指针接下一个兄弟。}\nopagebreak[4]
\Code{1.3}{if (!p) return NULL;}
\Code{1.3}{q = NewTree(p->data);}
\Code{1.3}{q->lchild = ForestToBinary(p->child);}
\Code{1.3}{q->rchild = ForestToBinary(p->sibling);}
\Code{1.3}{return q;}
\Code{0.0}{\}}
\Hint{0.0}{二叉树的右边在这里代表同层兄弟，转回森林时不能当作孩子。}\nopagebreak[4]
\Code{0.0}{BinaryToForest(q) \{}
\Code{1.3}{if (!q) return NULL;}
\Code{1.3}{p = NewGeneralNode(q->data);}
\Code{1.3}{p->child = BinaryToForest(q->lchild);}
\Code{1.3}{p->sibling = BinaryToForest(q->rchild);}
\Code{1.3}{return p;}
\Code{0.0}{\}}
\nopagebreak[4]
\Cost{\ensuremath{O\!\left(n\right)}}{\ensuremath{O\!\left(h\right)}（递归深度，最坏 n）}
\par\medskip

\Needspace{27\baselineskip}
\subsubsection{二叉树的中序线索化}
\Code{0.0}{ThreadIn(p, Node **prev) \{}
\Hint{1.3}{prev 始终指向刚访问的中序前驱，末结点的右线索须由顶层补上。}\nopagebreak[4]
\Code{1.3}{if (!p) return;}
\Hint{1.3}{递归前保存原右孩子，避免递归返回后沿新线索误入已处理结点。}\nopagebreak[4]
\Code{1.3}{r = p->rchild;}
\Code{1.3}{ThreadIn(p->lchild, prev);}
\Code{1.3}{if (!p->lchild) \{}
\Code{2.6}{p->ltag = 1;}
\Code{2.6}{p->lchild = (*prev);}
\Code{1.3}{\}}
\Code{1.3}{if ((*prev) \&\& !(*prev)->rchild) \{}
\Code{2.6}{(*prev)->rtag = 1;}
\Code{2.6}{(*prev)->rchild = p;}
\Code{1.3}{\}}
\Code{1.3}{(*prev) = p;}
\Code{1.3}{ThreadIn(r, prev);}
\Code{0.0}{\}}
\Code{0.0}{MakeThreads(root) \{}
\Hint{1.3}{输入为尚未线索化的二叉树，所有结点的 ltag 和 rtag 初始均为 0。}\nopagebreak[4]
\Code{1.3}{Node *prev = NULL;}
\Code{1.3}{ThreadIn(root, \&prev);}
\Code{1.3}{if (prev) prev->rtag = 1;}
\Code{0.0}{\}}
\nopagebreak[4]
\Cost{\ensuremath{O\!\left(n\right)}}{\ensuremath{O\!\left(h\right)}}
\par\medskip

\Needspace{20\baselineskip}
\subsubsection{中缀表达式转换为表达式树}
\Code{0.0}{BuildInfixExpr(tokens) \{}
\Code{1.3}{return BuildExpr(ToPostfix(tokens));}
\Code{0.0}{\}}
\Code{0.0}{BuildExpr(post) \{}
\Hint{1.3}{输入为合法的二元运算后缀序列。先弹出的子树作为右孩子。}\nopagebreak[4]
\Code{1.3}{InitStack(\&S);}
\Code{1.3}{for (t in post) \{}
\Code{2.6}{p = NewTree(t);}
\Code{2.6}{if (IsOperator(t)) \{}
\Code{3.9}{p->rchild = Pop(\&S);}
\Code{3.9}{p->lchild = Pop(\&S);}
\Code{2.6}{\}}
\Hint{2.6}{运算符接好两棵子树后，把新建的父结点压栈。}\nopagebreak[4]
\Code{2.6}{Push(\&S, p);}
\Code{1.3}{\}}
\Code{1.3}{require(Size(\&S) == 1);}
\Code{1.3}{return Pop(\&S);}
\Code{0.0}{\}}
\nopagebreak[4]
\Cost{\ensuremath{O\!\left(n\right)}}{\ensuremath{O\!\left(n\right)}}
\par\medskip

\Needspace{13\baselineskip}
\subsubsection{表达式树的求值}
\Code{0.0}{EvalTree(p) \{}
\Hint{1.3}{叶子提供数值，内部结点负责合并左右两个递归结果。}\nopagebreak[4]
\Code{1.3}{if (IsLeaf(p)) return Value(p);}
\Code{1.3}{a = EvalTree(p->lchild);}
\Hint{1.3}{先区分数值叶与运算符，再递归，叶子不应继续取孩子。}\nopagebreak[4]
\Code{1.3}{b = EvalTree(p->rchild);}
\Code{1.3}{return Apply(p->data, a, b);}
\Code{0.0}{\}}
\nopagebreak[4]
\Cost{\ensuremath{O\!\left(n\right)}}{\ensuremath{O\!\left(h\right)}}
\par\medskip

\Needspace{15\baselineskip}
\subsubsection{二叉树展开为单链表}
\Code{0.0}{FlattenTree(p, Node **prev) \{}
\Hint{1.3}{按右、左、根逆序处理，prev 指向已展开的后缀。初值为空。}\nopagebreak[4]
\Code{1.3}{if (!p) return;}
\Code{1.3}{FlattenTree(p->rchild, prev);}
\Code{1.3}{FlattenTree(p->lchild, prev);}
\Hint{1.3}{右、左递归完成后再改连当前结点，原来的子树入口此时已处理完。}\nopagebreak[4]
\Code{1.3}{p->rchild = (*prev);}
\Code{1.3}{p->lchild = NULL;}
\Code{1.3}{(*prev) = p;}
\Code{0.0}{\}}
\nopagebreak[4]
\Cost{\ensuremath{O\!\left(n\right)}}{\ensuremath{O\!\left(h\right)}}
\par\medskip

\Group{二叉树的基本递归算法}

\Needspace{13\baselineskip}
\subsubsection{求二叉树的深度}
\Code{0.0}{Height(p) \{}
\Hint{1.3}{叶结点左右高度均为 0，自身高度为 1。}\nopagebreak[4]
\Code{1.3}{if (!p) return 0;}
\Code{1.3}{l = Height(p->lchild);}
\Hint{1.3}{这里只返回最长分支，左右高度不相加。}\nopagebreak[4]
\Code{1.3}{r = Height(p->rchild);}
\Code{1.3}{return 1+max(l, r);}
\Code{0.0}{\}}
\nopagebreak[4]
\Cost{\ensuremath{O\!\left(n\right)}}{\ensuremath{O\!\left(h\right)}}
\par\medskip

\Needspace{13\baselineskip}
\subsubsection{求二叉树的最小深度}
\Code{0.0}{MinDepth(p) \{}
\Hint{1.3}{只有一边有孩子时必须向那一边走，空指针不代表叶。}\nopagebreak[4]
\Code{1.3}{if (!p) return 0;}
\Code{1.3}{if (!p->lchild) return 1+MinDepth(p->rchild);}
\Hint{1.3}{只有两侧孩子都存在时才能直接比较两侧深度的较小值。}\nopagebreak[4]
\Code{1.3}{if (!p->rchild) return 1+MinDepth(p->lchild);}
\Code{1.3}{return 1+min(MinDepth(p->lchild), MinDepth(p->rchild));}
\Code{0.0}{\}}
\nopagebreak[4]
\Cost{\ensuremath{O\!\left(n\right)}}{\ensuremath{O\!\left(h\right)}}
\par\medskip

\Needspace{11\baselineskip}
\subsubsection{统计结点总数}
\Code{0.0}{TreeSize(p) \{}
\Hint{1.3}{两棵子树计数后，加当前根结点一项。}\nopagebreak[4]
\Code{1.3}{if (!p) return 0;}
\Hint{1.3}{空子树贡献零，因此不用区分当前根有一个孩子还是两个孩子。}\nopagebreak[4]
\Code{1.3}{return 1+TreeSize(p->lchild)+TreeSize(p->rchild);}
\Code{0.0}{\}}
\nopagebreak[4]
\Cost{\ensuremath{O\!\left(n\right)}}{\ensuremath{O\!\left(h\right)}}
\par\medskip

\Needspace{12\baselineskip}
\subsubsection{统计叶结点数}
\Code{0.0}{LeafCount(p) \{}
\Hint{1.3}{左右均空才是叶；空树本身没有叶。}\nopagebreak[4]
\Code{1.3}{if (!p) return 0;}
\Hint{1.3}{有一个孩子的结点也不算叶，它的叶数由非空子树提供。}\nopagebreak[4]
\Code{1.3}{if (!p->lchild \&\& !p->rchild) return 1;}
\Code{1.3}{return LeafCount(p->lchild)+LeafCount(p->rchild);}
\Code{0.0}{\}}
\nopagebreak[4]
\Cost{\ensuremath{O\!\left(n\right)}}{\ensuremath{O\!\left(h\right)}}
\par\medskip

\Needspace{14\baselineskip}
\subsubsection{统计度为0、1、2的结点数}
\Code{0.0}{DegreeCount(p, count) \{}
\Hint{1.3}{两个布尔条件相加得到当前结点的度，计数数组初始为零。}\nopagebreak[4]
\Code{1.3}{if (!p) return;}
\Code{1.3}{d = (p->lchild != NULL)+(p->rchild != NULL);}
\Hint{1.3}{同一结点只更新其中一个桶，遍历后三个桶相加等于结点总数。}\nopagebreak[4]
\Code{1.3}{count[d]++;}
\Code{1.3}{DegreeCount(p->lchild, count);}
\Code{1.3}{DegreeCount(p->rchild, count);}
\Code{0.0}{\}}
\nopagebreak[4]
\Cost{\ensuremath{O\!\left(n\right)}}{\ensuremath{O\!\left(h\right)}}
\par\medskip

\Needspace{12\baselineskip}
\subsubsection{统计第k层的结点数}
\Code{0.0}{CountLevel(p, k) \{}
\Hint{1.3}{从第 1 层开始，每进入孩子，剩余目标层数减一。}\nopagebreak[4]
\Code{1.3}{if (!p || k<1) return 0;}
\Hint{1.3}{只剩一层时当前结点自身就是应统计的一项，无需再向下找。}\nopagebreak[4]
\Code{1.3}{if (k == 1) return 1;}
\Code{1.3}{return CountLevel(p->lchild, k-1)+CountLevel(p->rchild, k-1);}
\Code{0.0}{\}}
\nopagebreak[4]
\Cost{\ensuremath{O\!\left(n\right)}}{\ensuremath{O\!\left(\min (k,h)\right)}}
\par\medskip

\Needspace{13\baselineskip}
\subsubsection{查找指定结点}
\Code{0.0}{TreeFind(p, x) \{}
\Hint{1.3}{左边成功时直接返回匹配地址。}\nopagebreak[4]
\Code{1.3}{if (!p || p->data == x) return p;}
\Code{1.3}{q = TreeFind(p->lchild, x);}
\Hint{1.3}{这里按值匹配，值重复时得到前序次序下第一个匹配结点。}\nopagebreak[4]
\Code{1.3}{if (q) return q;}
\Code{1.3}{return TreeFind(p->rchild, x);}
\Code{0.0}{\}}
\nopagebreak[4]
\Cost{\ensuremath{O\!\left(n\right)}}{\ensuremath{O\!\left(h\right)}}
\par\medskip

\Needspace{14\baselineskip}
\subsubsection{求指定结点所在的层次}
\Code{0.0}{FindLevel(p, target, d) \{}
\Hint{1.3}{根传 1；返回 0 代表该分支没有目标。}\nopagebreak[4]
\Code{1.3}{if (!p) return 0;}
\Code{1.3}{if (p == target) return d;}
\Hint{1.3}{按结点地址寻找，避免树内存在相同值时返回错误的层次。}\nopagebreak[4]
\Code{1.3}{k = FindLevel(p->lchild, target, d+1);}
\Code{1.3}{if (k) return k;}
\Code{1.3}{return FindLevel(p->rchild, target, d+1);}
\Code{0.0}{\}}
\nopagebreak[4]
\Cost{\ensuremath{O\!\left(n\right)}}{\ensuremath{O\!\left(h\right)}}
\par\medskip

\Needspace{13\baselineskip}
\subsubsection{交换所有结点的左右子树（镜像翻转）}
\Code{0.0}{Mirror(p) \{}
\Hint{1.3}{交换的是左右孩子指针，结点数据无需改动。}\nopagebreak[4]
\Code{1.3}{if (!p) return;}
\Code{1.3}{swap(p->lchild, p->rchild);}
\Hint{1.3}{即使左右只有一边非空，也要交换两域，使独生孩子换到另一侧。}\nopagebreak[4]
\Code{1.3}{Mirror(p->lchild);}
\Code{1.3}{Mirror(p->rchild);}
\Code{0.0}{\}}
\nopagebreak[4]
\Cost{\ensuremath{O\!\left(n\right)}}{\ensuremath{O\!\left(h\right)}}
\par\medskip

\Needspace{11\baselineskip}
\subsubsection{判断两棵二叉树是否相同}
\Code{0.0}{Same(a, b) \{}
\Hint{1.3}{两个空位置相同，单边为空意味着形状不同。}\nopagebreak[4]
\Code{1.3}{if (!a || !b) return a == b;}
\Hint{1.3}{当前两个值相等还不够，两侧各自的孩子也必须逐一匹配。}\nopagebreak[4]
\Code{1.3}{return a->data == b->data \&\& Same(a->lchild, b->lchild) \&\& Same(a->rchild, b->rchild);}
\Code{0.0}{\}}
\nopagebreak[4]
\Cost{\ensuremath{O\!\left(n_{1}+n_{2}\right)}}{\ensuremath{O\!\left(\min (h_{1},h_{2})+1\right)}}
\par\medskip

\Needspace{25\baselineskip}
\subsubsection{判断两棵二叉树是否镜像对称（递归、非递归）}
\Code{0.0}{MirrorEq(a, b) \{}
\Hint{1.3}{交叉配对比较对应位置，队列版本保存待比较的结点对。}\nopagebreak[4]
\Code{1.3}{if (!a || !b) return a == b;}
\Code{1.3}{return a->data == b->data \&\& MirrorEq(a->lchild, b->rchild) \&\& MirrorEq(a->rchild, b->lchild);}
\Code{0.0}{\}}
\Code{0.0}{MirrorIter(a, b) \{}
\Code{1.3}{EnQueue(\&Q, (a, b));}
\Code{1.3}{while (!Empty(\&Q)) \{}
\Code{2.6}{(u, v) = DeQueue(\&Q);}
\Code{2.6}{if (!u || !v) \{}
\Code{3.9}{if (u != v) return false;}
\Hint{3.9}{空指针对应真实结点就立即失败，两个都是空才跳过这对。}\nopagebreak[4]
\Code{3.9}{continue;}
\Code{2.6}{\}}
\Code{2.6}{if (u->data != v->data) return false;}
\Code{2.6}{EnQueue(\&Q, (u->lchild, v->rchild));}
\Code{2.6}{EnQueue(\&Q, (u->rchild, v->lchild));}
\Code{1.3}{\}}
\Code{1.3}{return true;}
\Code{0.0}{\}}
\nopagebreak[4]
\Cost{\ensuremath{O\!\left(n_{1}+n_{2}\right)}}{递归 \ensuremath{O\!\left(\min (h_{1},h_{2})+1\right)}；迭代 \ensuremath{O\!\left(w\right)}}
\par\medskip

\Needspace{11\baselineskip}
\subsubsection{判断一棵二叉树是否轴对称}
\Code{0.0}{Symmetric(T) \{}
\Hint{1.3}{只比较根两侧的镜像位置。}\nopagebreak[4]
\Code{1.3}{if (!T) return true;}
\Hint{1.3}{空树以及只有一个根的树都满足轴对称。}\nopagebreak[4]
\Code{1.3}{return MirrorEq(T->lchild, T->rchild);}
\Code{0.0}{\}}
\nopagebreak[4]
\Cost{\ensuremath{O\!\left(n\right)}}{\ensuremath{O\!\left(h\right)}}
\par\medskip

\Needspace{20\baselineskip}
\subsubsection{复制、销毁二叉树}
\Code{0.0}{CloneTree(p) \{}
\Hint{1.3}{复制需要分配新根；销毁须在左右孩子处理后释放父结点。}\nopagebreak[4]
\Code{1.3}{if (!p) return NULL;}
\Code{1.3}{q = NewTree(p->data);}
\Code{1.3}{q->lchild = CloneTree(p->lchild);}
\Code{1.3}{q->rchild = CloneTree(p->rchild);}
\Code{1.3}{return q;}
\Code{0.0}{\}}
\Hint{0.0}{叶结点的复制也必须分配新地址；销毁后调用方要清空原根指针。}\nopagebreak[4]
\Code{0.0}{DestroyTree(p) \{}
\Code{1.3}{if (!p) return;}
\Code{1.3}{DestroyTree(p->lchild);}
\Code{1.3}{DestroyTree(p->rchild);}
\Code{1.3}{free(p);}
\Code{0.0}{\}}
\nopagebreak[4]
\Cost{\ensuremath{O\!\left(n\right)}}{\ensuremath{O\!\left(h\right)}}
\par\medskip

\Needspace{13\baselineskip}
\subsubsection{求孩子—兄弟链表表示的树的深度}
\Code{0.0}{GeneralHeight(p) \{}
\Hint{1.3}{沿兄弟链枚举同层孩子；只有进入孩子才增加深度。}\nopagebreak[4]
\Code{1.3}{if (!p) return 0;}
\Code{1.3}{h = 0;}
\Hint{1.3}{沿 sibling 前进是在枚举孩子，不能给树高额外加一。}\nopagebreak[4]
\Code{1.3}{for (c = p->child; c; c = c->sibling) h = max(h, GeneralHeight(c));}
\Code{1.3}{return h+1;}
\Code{0.0}{\}}
\nopagebreak[4]
\Cost{\ensuremath{O\!\left(n\right)}}{\ensuremath{O\!\left(h\right)}}
\par\medskip

\Needspace{16\baselineskip}
\subsubsection{求孩子—兄弟链表表示的树的最大度数}
\Code{0.0}{GeneralDegree(p) \{}
\Hint{1.3}{局部 d 统计直接孩子数，best 同时考虑各子树。}\nopagebreak[4]
\Code{1.3}{if (!p) return 0;}
\Code{1.3}{d = best = 0;}
\Code{1.3}{for (c = p->child; c; c = c->sibling) \{}
\Code{2.6}{d++;}
\Hint{2.6}{遍历完全部孩子后才能知道此结点的度，应与子树最大度再比较一次。}\nopagebreak[4]
\Code{2.6}{best = max(best, GeneralDegree(c));}
\Code{1.3}{\}}
\Code{1.3}{return max(d, best);}
\Code{0.0}{\}}
\nopagebreak[4]
\Cost{\ensuremath{O\!\left(n\right)}}{\ensuremath{O\!\left(h\right)}}
\par\medskip

\Needspace{14\baselineskip}
\subsubsection{统计孩子—兄弟链表表示的树的叶结点数}
\Code{0.0}{GeneralLeaves(p) \{}
\Hint{1.3}{有兄弟不影响叶子身份，只看自己有无孩子。}\nopagebreak[4]
\Code{1.3}{if (!p) return 0;}
\Code{1.3}{if (!p->child) return 1;}
\Hint{1.3}{空树返回零；非空且没有孩子才贡献一片叶。}\nopagebreak[4]
\Code{1.3}{n = 0;}
\Code{1.3}{for (c = p->child; c; c = c->sibling) n += GeneralLeaves(c);}
\Code{1.3}{return n;}
\Code{0.0}{\}}
\nopagebreak[4]
\Cost{\ensuremath{O\!\left(n\right)}}{\ensuremath{O\!\left(h\right)}}
\par\medskip

\Group{二叉树的综合算法}

\Needspace{16\baselineskip}
\subsubsection{求指定结点的所有祖先}
\Code{0.0}{Ancestors(p, t) \{}
\Hint{1.3}{目标所在分支返回 true 时才输出父结点，顺序为由近到远。}\nopagebreak[4]
\Code{1.3}{if (!p) return false;}
\Code{1.3}{if (p == t) return true;}
\Code{1.3}{if (Ancestors(p->lchild, t) || Ancestors(p->rchild, t)) \{}
\Code{2.6}{Emit(p);}
\Hint{2.6}{目标自身只报告命中，不应把它打印成自己的祖先。}\nopagebreak[4]
\Code{2.6}{return true;}
\Code{1.3}{\}}
\Code{1.3}{return false;}
\Code{0.0}{\}}
\nopagebreak[4]
\Cost{\ensuremath{O\!\left(n\right)}}{\ensuremath{O\!\left(h\right)}}
\par\medskip

\Needspace{14\baselineskip}
\subsubsection{求两个结点的最近公共祖先（LCA）}
\Code{0.0}{LCA(p, a, b) \{}
\Hint{1.3}{假定两点均在树内。左右均有命中时当前根把两者分开。}\nopagebreak[4]
\Code{1.3}{if (!p || p == a || p == b) return p;}
\Code{1.3}{l = LCA(p->lchild, a, b);}
\Hint{1.3}{若其中一个目标就是当前根，返回它允许祖先与后代这一情形得到正确答案。}\nopagebreak[4]
\Code{1.3}{r = LCA(p->rchild, a, b);}
\Code{1.3}{if (l \&\& r) return p;}
\Code{1.3}{return l?l:r;}
\Code{0.0}{\}}
\nopagebreak[4]
\Cost{\ensuremath{O\!\left(n\right)}}{\ensuremath{O\!\left(h\right)}}
\par\medskip

\Needspace{14\baselineskip}
\subsubsection{求根结点到指定结点的路径}
\Code{0.0}{RootPath(p, t, path) \{}
\Hint{1.3}{成功时不弹栈，整条根到目标的路径便保留下来。}\nopagebreak[4]
\Code{1.3}{if (!p) return false;}
\Code{1.3}{Push(path, p);}
\Hint{1.3}{若两个子树都失败，弹出本结点，避免在别的路径中留下错误前缀。}\nopagebreak[4]
\Code{1.3}{if (p == t || RootPath(p->lchild, t, path) || RootPath(p->rchild, t, path)) return true;}
\Code{1.3}{Pop(path);}
\Code{1.3}{return false;}
\Code{0.0}{\}}
\nopagebreak[4]
\Cost{\ensuremath{O\!\left(n\right)}}{\ensuremath{O\!\left(h\right)}}
\par\medskip

\Needspace{17\baselineskip}
\subsubsection{输出根结点到所有叶结点的路径}
\Code{0.0}{LeafPaths(p, path) \{}
\Hint{1.3}{退出前弹出当前结点，下一条分支才能复用 path。}\nopagebreak[4]
\Code{1.3}{if (!p) return;}
\Code{1.3}{Push(path, p);}
\Code{1.3}{if (IsLeaf(p)) Emit(path);}
\Code{1.3}{else \{}
\Hint{2.6}{Emit 需要当场输出或复制路径，不能长期保存随后还会变动的栈地址。}\nopagebreak[4]
\Code{2.6}{LeafPaths(p->lchild, path);}
\Code{2.6}{LeafPaths(p->rchild, path);}
\Code{1.3}{\}}
\Code{1.3}{Pop(path);}
\Code{0.0}{\}}
\nopagebreak[4]
\Cost{\ensuremath{O\!\left(n+L\right)}}{\ensuremath{O\!\left(h\right)}}
\par\medskip

\Needspace{14\baselineskip}
\subsubsection{求两个指定结点之间的距离}
\Code{0.0}{TreeDistance(root, a, b) \{}
\Hint{1.3}{从最近公共祖先分别向下计边，公共前缀恰扣除两次。}\nopagebreak[4]
\Code{1.3}{c = LCA(root, a, b);}
\Code{1.3}{da = FindLevel(root, a, 1);}
\Hint{1.3}{两点相同则三者层次相同，代入得到距离零。}\nopagebreak[4]
\Code{1.3}{db = FindLevel(root, b, 1);}
\Code{1.3}{dc = FindLevel(root, c, 1);}
\Code{1.3}{return da+db-2*dc;}
\Code{0.0}{\}}
\nopagebreak[4]
\Cost{\ensuremath{O\!\left(n\right)}}{\ensuremath{O\!\left(h\right)}}
\par\medskip

\Needspace{21\baselineskip}
\subsubsection{判断是否为完全二叉树}
\Code{0.0}{Complete(root) \{}
\Hint{1.3}{层序首次遇到缺位后，不应再遇到真实结点。}\nopagebreak[4]
\Code{1.3}{if (!root) return true;}
\Code{1.3}{gap = false;}
\Code{1.3}{EnQueue(\&Q, root);}
\Code{1.3}{while (!Empty(\&Q)) \{}
\Code{2.6}{p = DeQueue(\&Q);}
\Code{2.6}{for (c in [p->lchild, p->rchild]) \{}
\Hint{3.9}{上一层出现空孩子后，队列中已经排入的点仍需检查它们的孩子。}\nopagebreak[4]
\Code{3.9}{if (!c) gap = true;}
\Code{3.9}{else if (gap) return false;}
\Code{3.9}{else EnQueue(\&Q, c);}
\Code{2.6}{\}}
\Code{1.3}{\}}
\Code{1.3}{return true;}
\Code{0.0}{\}}
\nopagebreak[4]
\Cost{\ensuremath{O\!\left(n\right)}}{\ensuremath{O\!\left(w\right)}}
\par\medskip

\Needspace{13\baselineskip}
\subsubsection{判断是否为满二叉树}
\Code{0.0}{Perfect(p) \{}
\Hint{1.3}{满二叉树要求左右均满且等高，函数同时返回高度。}\nopagebreak[4]
\Code{1.3}{if (!p) return (true, 0);}
\Code{1.3}{(a, l) = Perfect(p->lchild);}
\Hint{1.3}{左右子树即使各自都满，高度不同时整棵树仍不是这里定义的满二叉树。}\nopagebreak[4]
\Code{1.3}{(b, r) = Perfect(p->rchild);}
\Code{1.3}{return (a \&\& b \&\& l == r, 1+max(l, r));}
\Code{0.0}{\}}
\nopagebreak[4]
\Cost{\ensuremath{O\!\left(n\right)}}{\ensuremath{O\!\left(h\right)}}
\par\medskip

\Needspace{14\baselineskip}
\subsubsection{判断是否为平衡二叉树}
\Code{0.0}{BalancedHeight(p) \{}
\Hint{1.3}{−1 向上报告失衡，正常时返回高度；顶层检验结果是否非负。}\nopagebreak[4]
\Code{1.3}{if (!p) return 0;}
\Code{1.3}{l = BalancedHeight(p->lchild);}
\Hint{1.3}{递归返回负值表示失衡信号，不可拿它继续当真实高度计算。}\nopagebreak[4]
\Code{1.3}{r = BalancedHeight(p->rchild);}
\Code{1.3}{if (l<0 || r<0 || abs(l-r)>1) return -1;}
\Code{1.3}{return 1+max(l, r);}
\Code{0.0}{\}}
\nopagebreak[4]
\Cost{\ensuremath{O\!\left(n\right)}}{\ensuremath{O\!\left(h\right)}}
\par\medskip

\Needspace{21\baselineskip}
\subsubsection{求二叉树的最大宽度}
\Code{0.0}{Width(root) \{}
\Hint{1.3}{每轮开始的队列大小就是本层实际结点数。}\nopagebreak[4]
\Code{1.3}{if (!root) return 0;}
\Code{1.3}{EnQueue(\&Q, root);}
\Code{1.3}{best = 0;}
\Code{1.3}{while (!Empty(\&Q)) \{}
\Code{2.6}{k = Size(\&Q);}
\Code{2.6}{best = max(k, best);}
\Hint{2.6}{k 轮出队结束时，队列恰好只剩下一层的孩子。}\nopagebreak[4]
\Code{2.6}{repeat (k) \{}
\Code{3.9}{p = DeQueue(\&Q);}
\Code{3.9}{EnqueueChildren(\&Q, p);}
\Code{2.6}{\}}
\Code{1.3}{\}}
\Code{1.3}{return best;}
\Code{0.0}{\}}
\nopagebreak[4]
\Cost{\ensuremath{O\!\left(n\right)}}{\ensuremath{O\!\left(w\right)}}
\par\medskip

\Needspace{14\baselineskip}
\subsubsection{求二叉树的直径（任意两结点间最长路径）}
\Code{0.0}{DiameterHeight(p, int *best) \{}
\Hint{1.3}{左右高度之和是经过当前根的最长路径边数；best 初始 0。}\nopagebreak[4]
\Code{1.3}{if (!p) return 0;}
\Code{1.3}{l = DiameterHeight(p->lchild, best);}
\Hint{1.3}{递归返回的是单支高度，best 才是整棵树目前最长的两点间路径。}\nopagebreak[4]
\Code{1.3}{r = DiameterHeight(p->rchild, best);}
\Code{1.3}{(*best) = max((*best), l+r);}
\Code{1.3}{return 1+max(l, r);}
\Code{0.0}{\}}
\nopagebreak[4]
\Cost{\ensuremath{O\!\left(n\right)}}{\ensuremath{O\!\left(h\right)}}
\par\medskip

\Needspace{13\baselineskip}
\subsubsection{判断一棵树是否为另一棵树的子树}
\Code{0.0}{Subtree(a, b) \{}
\Hint{1.3}{候选根的整棵子树须和值、形状完全匹配。}\nopagebreak[4]
\Code{1.3}{if (!b) return true;}
\Code{1.3}{if (!a) return false;}
\Hint{1.3}{先检查根是否匹配；失败只排除这个根，仍要检查它的左右后代。}\nopagebreak[4]
\Code{1.3}{if (Same(a, b)) return true;}
\Code{1.3}{return Subtree(a->lchild, b) || Subtree(a->rchild, b);}
\Code{0.0}{\}}
\nopagebreak[4]
\Cost{\ensuremath{O\!\left(nm+n\right)}}{\ensuremath{O\!\left(h_{1}+h_{2}\right)}}
\par\medskip

\Needspace{14\baselineskip}
\subsubsection{求二叉树的最大路径和}
\Code{0.0}{Gain(p, int *best) \{}
\Hint{1.3}{路径可在当前根连接两侧，递交给父亲时却只能延伸一侧。best 初值负无穷。}\nopagebreak[4]
\Code{1.3}{if (!p) return 0;}
\Code{1.3}{l = max(0, Gain(p->lchild, best));}
\Hint{1.3}{负贡献直接舍弃，但当前根必须计入，因此全负树也会返回最大的结点值。}\nopagebreak[4]
\Code{1.3}{r = max(0, Gain(p->rchild, best));}
\Code{1.3}{(*best) = max((*best), p->data+l+r);}
\Code{1.3}{return p->data+max(l, r);}
\Code{0.0}{\}}
\nopagebreak[4]
\Cost{\ensuremath{O\!\left(n\right)}}{\ensuremath{O\!\left(h\right)}}
\par\medskip

\Needspace{16\baselineskip}
\subsubsection{查找路径和等于指定值的所有根到叶路径}
\Code{0.0}{TargetPaths(p, left, path) \{}
\Hint{1.3}{只在叶结点检验剩余和，避免把内部结点误当终点。}\nopagebreak[4]
\Code{1.3}{if (!p) return;}
\Code{1.3}{Push(path, p);}
\Code{1.3}{left -= p->data;}
\Code{1.3}{if (IsLeaf(p) \&\& left == 0) Emit(path);}
\Hint{1.3}{left 是该分支自己的局部剩余量，路径栈才需要显式撤销。}\nopagebreak[4]
\Code{1.3}{TargetPaths(p->lchild, left, path);}
\Code{1.3}{TargetPaths(p->rchild, left, path);}
\Code{1.3}{Pop(path);}
\Code{0.0}{\}}
\nopagebreak[4]
\Cost{\ensuremath{O\!\left(n+L\right)}}{\ensuremath{O\!\left(h\right)}}
\par\medskip

\Group{二叉搜索树（BST）}

\Needspace{12\baselineskip}
\subsubsection{二叉搜索树的构建}
\Code{0.0}{BuildBST(keys) \{}
\Hint{1.3}{树根变量的地址传给插入函数，空树的第一次插入也能改根。}\nopagebreak[4]
\Code{1.3}{root = NULL;}
\Hint{1.3}{输入顺序会决定形状，单调输入可使这棵普通搜索树退化为长链。}\nopagebreak[4]
\Code{1.3}{for (x in keys) BSTInsert(\&root, x);}
\Code{1.3}{return root;}
\Code{0.0}{\}}
\nopagebreak[4]
\Cost{最坏 \ensuremath{O\!\left(n^{2}\right)}}{\ensuremath{O\!\left(1\right)}}
\par\medskip

\Needspace{12\baselineskip}
\subsubsection{结点查找}
\Code{0.0}{BSTFind(root, x) \{}
\Hint{1.3}{每次比较后只沿一侧孩子继续查找。}\nopagebreak[4]
\Code{1.3}{p = root;}
\Hint{1.3}{走到空位置说明相应子树中不存在 x，此时直接返回空指针。}\nopagebreak[4]
\Code{1.3}{while (p \&\& p->key != x) p = (x<p->key ? p->lchild:p->rchild);}
\Code{1.3}{return p;}
\Code{0.0}{\}}
\nopagebreak[4]
\Cost{\ensuremath{O\!\left(h\right)}}{\ensuremath{O\!\left(1\right)}}
\par\medskip

\Needspace{16\baselineskip}
\subsubsection{结点插入}
\Code{0.0}{BSTInsert(Node **link, x) \{}
\Hint{1.3}{link 指向可写的根变量或孩子指针域，找到空槽直接填入。}\nopagebreak[4]
\Code{1.3}{while (*link) \{}
\Code{2.6}{p = *link;}
\Code{2.6}{if (p->key == x) return false;}
\Code{2.6}{link = x<p->key ? \&p->lchild:\&p->rchild;}
\Code{1.3}{\}}
\Hint{1.3}{重复键在本实现中跳过，因此验证 BST 时使用严格大小关系。}\nopagebreak[4]
\Code{1.3}{*link = NewTree(x);}
\Code{1.3}{return true;}
\Code{0.0}{\}}
\nopagebreak[4]
\Cost{\ensuremath{O\!\left(h+1\right)}}{\ensuremath{O\!\left(1\right)}}
\par\medskip

\Needspace{25\baselineskip}
\subsubsection{结点删除（三种情况）}
\Code{0.0}{BSTDelete(Node **link, x) \{}
\Hint{1.3}{有两个孩子时先复制后继，再删后继原结点；后继没有左孩子。}\nopagebreak[4]
\Code{1.3}{while (*link \&\& (*link)->key != x) \{}
\Code{2.6}{p = *link;}
\Code{2.6}{link = x<p->key ? \&p->lchild:\&p->rchild;}
\Code{1.3}{\}}
\Code{1.3}{if (!*link) return false;}
\Code{1.3}{p = *link;}
\Code{1.3}{if (p->lchild \&\& p->rchild) \{}
\Code{2.6}{slot = \&p->rchild;}
\Code{2.6}{while ((*slot)->lchild) slot = \&(*slot)->lchild;}
\Hint{2.6}{只复制业务键值和数据，不要覆盖左右孩子指针。}\nopagebreak[4]
\Code{2.6}{CopyPayload(p, *slot);}
\Code{2.6}{link = slot;}
\Code{2.6}{p = *link;}
\Code{1.3}{\}}
\Code{1.3}{*link = p->lchild ? p->lchild:p->rchild;}
\Code{1.3}{free(p);}
\Code{1.3}{return true;}
\Code{0.0}{\}}
\nopagebreak[4]
\Cost{\ensuremath{O\!\left(h\right)}}{\ensuremath{O\!\left(1\right)}}
\par\medskip

\Needspace{14\baselineskip}
\subsubsection{查找最大值、最小值}
\Code{0.0}{Extreme(root, wantMin) \{}
\Hint{1.3}{最小值在最左结点，最大值在最右结点。}\nopagebreak[4]
\Code{1.3}{p = root;}
\Code{1.3}{if (!p) return NULL;}
\Hint{1.3}{选 wantMin=true 时停止在没有左孩子的位置，该结点仍可有右孩子。}\nopagebreak[4]
\Code{1.3}{if (wantMin) while (p->lchild) p = p->lchild;}
\Code{1.3}{else while (p->rchild) p = p->rchild;}
\Code{1.3}{return p;}
\Code{0.0}{\}}
\nopagebreak[4]
\Cost{\ensuremath{O\!\left(h\right)}}{\ensuremath{O\!\left(1\right)}}
\par\medskip

\Needspace{25\baselineskip}
\subsubsection{查找中序前驱与后继}
\Code{0.0}{Neighbors(root, x) \{}
\Hint{1.3}{向右走前的祖先是前驱候选。命中后，左子树取最大键，右子树取最小键。}\nopagebreak[4]
\Code{1.3}{p = root;}
\Code{1.3}{pred = succ = NULL;}
\Code{1.3}{while (p \&\& p->key != x) \{}
\Code{2.6}{if (x<p->key) \{}
\Code{3.9}{succ = p;}
\Code{3.9}{p = p->lchild;}
\Code{2.6}{\}}
\Code{2.6}{else \{}
\Code{3.9}{pred = p;}
\Hint{3.9}{缺少前驱或后继时返回对应的 NULL，不能用根值代替。}\nopagebreak[4]
\Code{3.9}{p = p->rchild;}
\Code{2.6}{\}}
\Code{1.3}{\}}
\Code{1.3}{require(p);}
\Code{1.3}{if (p->lchild) pred = Extreme(p->lchild, false);}
\Code{1.3}{if (p->rchild) succ = Extreme(p->rchild, true);}
\Code{1.3}{return (pred, succ);}
\Code{0.0}{\}}
\nopagebreak[4]
\Cost{\ensuremath{O\!\left(h\right)}}{\ensuremath{O\!\left(1\right)}}
\par\medskip

\Needspace{12\baselineskip}
\subsubsection{判断一棵树是否为二叉搜索树}
\Code{0.0}{ValidBST(p, lo, hi) \{}
\Hint{1.3}{用祖先传下来的严格界，同时限制整棵子树。初始界为正负无穷。}\nopagebreak[4]
\Code{1.3}{if (!p) return true;}
\Hint{1.3}{进入左子树要收紧上界，进入右子树要抬高下界。}\nopagebreak[4]
\Code{1.3}{if (p->key <= lo || p->key >= hi) return false;}
\Code{1.3}{return ValidBST(p->lchild, lo, p->key) \&\& ValidBST(p->rchild, p->key, hi);}
\Code{0.0}{\}}
\nopagebreak[4]
\Cost{\ensuremath{O\!\left(n\right)}}{\ensuremath{O\!\left(h\right)}}
\par\medskip

\Needspace{22\baselineskip}
\subsubsection{查找第k小的元素}
\Code{0.0}{KthKey(root, k) \{}
\Hint{1.3}{中序每弹出一个结点才算访问一个键。}\nopagebreak[4]
\Code{1.3}{require(k>0);}
\Code{1.3}{p = root;}
\Code{1.3}{InitStack(\&S);}
\Code{1.3}{while (p || !Empty(\&S)) \{}
\Code{2.6}{while (p) \{}
\Code{3.9}{Push(\&S, p);}
\Code{3.9}{p = p->lchild;}
\Code{2.6}{\}}
\Hint{2.6}{k 减到零的这一刻，已经恰好按升序访问了 k 项。}\nopagebreak[4]
\Code{2.6}{p = Pop(\&S);}
\Code{2.6}{if (--k == 0) return p->key;}
\Code{2.6}{p = p->rchild;}
\Code{1.3}{\}}
\Code{1.3}{return FAIL;}
\Code{0.0}{\}}
\nopagebreak[4]
\Cost{\ensuremath{O\!\left(h+k\right)}}{\ensuremath{O\!\left(h\right)}}
\par\medskip

\Needspace{17\baselineskip}
\subsubsection{二叉搜索树的最近公共祖先}
\Code{0.0}{BSTAncestor(root, a, b) \{}
\Hint{1.3}{假设 a 和 b 均存在。两键分居当前键两侧或其一等于当前键时停止。}\nopagebreak[4]
\Code{1.3}{require(a <= b);}
\Code{1.3}{p = root;}
\Code{1.3}{while (p) \{}
\Code{2.6}{if (p->key<a) p = p->rchild;}
\Hint{2.6}{当前键小于 a 时两个目标都在右边，反之都在左边。}\nopagebreak[4]
\Code{2.6}{else if (p->key>b) p = p->lchild;}
\Code{2.6}{else return p;}
\Code{1.3}{\}}
\Code{1.3}{return NULL;}
\Code{0.0}{\}}
\nopagebreak[4]
\Cost{\ensuremath{O\!\left(h\right)}}{\ensuremath{O\!\left(1\right)}}
\par\medskip

\Needspace{15\baselineskip}
\subsubsection{将有序数组转换为平衡二叉搜索树}
\Code{0.0}{ArrayBST(A, l, r) \{}
\Hint{1.3}{在键值严格递增的区间 [l,r) 取中点，左右段规模至多相差一。}\nopagebreak[4]
\Code{1.3}{if (l >= r) return NULL;}
\Code{1.3}{m = (l+r)/2;}
\Code{1.3}{p = NewTree(A[m]);}
\Hint{1.3}{把中点从两个子区间中排除，递归规模才会严格缩小。}\nopagebreak[4]
\Code{1.3}{p->lchild = ArrayBST(A, l, m);}
\Code{1.3}{p->rchild = ArrayBST(A, m+1, r);}
\Code{1.3}{return p;}
\Code{0.0}{\}}
\nopagebreak[4]
\Cost{\ensuremath{O\!\left(n\right)}}{\ensuremath{O\!\left(\log (n+1)\right)}}
\par\medskip

\Needspace{17\baselineskip}
\subsubsection{将二叉搜索树转换为有序双向链表}
\Code{0.0}{LinkBST(p, Node **prev, Node **head) \{}
\Hint{1.3}{中序访问时 prev 就是前驱，left 与 right 被重用作双链。}\nopagebreak[4]
\Code{1.3}{if (!p) return;}
\Code{1.3}{r = p->rchild;}
\Code{1.3}{LinkBST(p->lchild, prev, head);}
\Code{1.3}{p->lchild = (*prev);}
\Hint{1.3}{调用前将 head、prev 清空；第一次访问的最小结点成为双链表头。}\nopagebreak[4]
\Code{1.3}{if ((*prev)) (*prev)->rchild = p;}
\Code{1.3}{else (*head) = p;}
\Code{1.3}{(*prev) = p;}
\Code{1.3}{LinkBST(r, prev, head);}
\Code{0.0}{\}}
\nopagebreak[4]
\Cost{\ensuremath{O\!\left(n\right)}}{\ensuremath{O\!\left(h\right)}}
\par\medskip

\Group{平衡二叉树（AVL）}

\Needspace{13\baselineskip}
\subsubsection{计算平衡因子}
\Code{0.0}{BF(p) \{}
\Hint{1.3}{孩子的高度已经维护，空孩子高度取 0。}\nopagebreak[4]
\Code{1.3}{return H(p->lchild)-H(p->rchild);}
\Code{0.0}{\}}
\Hint{0.0}{BF 正值表示左高，负值表示右高，旋转时要使用同一套符号约定。}\nopagebreak[4]
\Code{0.0}{Update(p) \{}
\Code{1.3}{p->h = 1+max(H(p->lchild), H(p->rchild));}
\Code{0.0}{\}}
\nopagebreak[4]
\Cost{\ensuremath{O\!\left(1\right)}（维护高度）}{\ensuremath{O\!\left(1\right)}}
\par\medskip

\Needspace{15\baselineskip}
\subsubsection{LL型旋转}
\Code{0.0}{RotateRightAVL(y) \{}
\Hint{1.3}{x 的右子树先交给 y，再把 y 挂到 x 右侧；高度从下向上更新。}\nopagebreak[4]
\Code{1.3}{x = y->lchild;}
\Code{1.3}{y->lchild = x->rchild;}
\Code{1.3}{x->rchild = y;}
\Hint{1.3}{先更新旧根 y，再更新新根 x，因为后者的高度依赖前者。}\nopagebreak[4]
\Code{1.3}{Update(y);}
\Code{1.3}{Update(x);}
\Code{1.3}{return x;}
\Code{0.0}{\}}
\nopagebreak[4]
\Cost{\ensuremath{O\!\left(1\right)}}{\ensuremath{O\!\left(1\right)}}
\par\medskip

\Needspace{15\baselineskip}
\subsubsection{RR型旋转}
\Code{0.0}{RotateLeftAVL(x) \{}
\Hint{1.3}{y 的左子树移交给 x，y 成为新根。}\nopagebreak[4]
\Code{1.3}{y = x->rchild;}
\Code{1.3}{x->rchild = y->lchild;}
\Code{1.3}{y->lchild = x;}
\Hint{1.3}{旋转返回新根 y，上层须用其替换旧的子树入口。}\nopagebreak[4]
\Code{1.3}{Update(x);}
\Code{1.3}{Update(y);}
\Code{1.3}{return y;}
\Code{0.0}{\}}
\nopagebreak[4]
\Cost{\ensuremath{O\!\left(1\right)}}{\ensuremath{O\!\left(1\right)}}
\par\medskip

\Needspace{11\baselineskip}
\subsubsection{LR型旋转}
\Code{0.0}{RotateLR(p) \{}
\Hint{1.3}{先把左孩子的右侧重心转到外侧，再对根右旋。}\nopagebreak[4]
\Code{1.3}{p->lchild = RotateLeftAVL(p->lchild);}
\Hint{1.3}{第一步的返回值要接回 p 的左域，第二步才在正确结构上旋转。}\nopagebreak[4]
\Code{1.3}{return RotateRightAVL(p);}
\Code{0.0}{\}}
\nopagebreak[4]
\Cost{\ensuremath{O\!\left(1\right)}}{\ensuremath{O\!\left(1\right)}}
\par\medskip

\Needspace{11\baselineskip}
\subsubsection{RL型旋转}
\Code{0.0}{RotateRL(p) \{}
\Hint{1.3}{先对右孩子右旋，再对当前根左旋。}\nopagebreak[4]
\Code{1.3}{p->rchild = RotateRightAVL(p->rchild);}
\Hint{1.3}{第一步完成后，原来的折线已经变成可以左旋修复的直线。}\nopagebreak[4]
\Code{1.3}{return RotateLeftAVL(p);}
\Code{0.0}{\}}
\nopagebreak[4]
\Cost{\ensuremath{O\!\left(1\right)}}{\ensuremath{O\!\left(1\right)}}
\par\medskip

\Needspace{20\baselineskip}
\subsubsection{AVL树的插入与平衡调整}
\Code{0.0}{AVLInsert(p, x) \{}
\Hint{1.3}{每个递归调用都返回新的局部根，父结点必须接回它。}\nopagebreak[4]
\Code{1.3}{if (!p) return NewLeaf(x, 1);}
\Code{1.3}{if (x<p->key) p->lchild = AVLInsert(p->lchild, x);}
\Code{1.3}{else if (x>p->key) p->rchild = AVLInsert(p->rchild, x);}
\Code{1.3}{else return p;}
\Code{1.3}{return Balance(p);}
\Code{0.0}{\}}
\Hint{0.0}{重复键直接返回；真正新增叶子后才需要沿调用栈重新计算高度。}\nopagebreak[4]
\Code{0.0}{Balance(p) \{}
\Code{1.3}{Update(p);}
\Code{1.3}{if (BF(p)>1) return BF(p->lchild) >= 0 ? RotateRightAVL(p):RotateLR(p);}
\Code{1.3}{if (BF(p)<-1) return BF(p->rchild) <= 0 ? RotateLeftAVL(p):RotateRL(p);}
\Code{1.3}{return p;}
\Code{0.0}{\}}
\nopagebreak[4]
\Cost{\ensuremath{O\!\left(\log (n+1)\right)}}{\ensuremath{O\!\left(\log (n+1)\right)}}
\par\medskip

\Needspace{23\baselineskip}
\subsubsection{AVL树的删除与平衡调整}
\Code{0.0}{AVLDelete(p, x) \{}
\Hint{1.3}{删除可能使多层同时失衡，沿整个返回路径调用 Balance。}\nopagebreak[4]
\Code{1.3}{if (!p) return NULL;}
\Code{1.3}{if (x<p->key) p->lchild = AVLDelete(p->lchild, x);}
\Code{1.3}{else if (x>p->key) p->rchild = AVLDelete(p->rchild, x);}
\Code{1.3}{else if (p->lchild \&\& p->rchild) \{}
\Code{2.6}{q = Extreme(p->rchild, true);}
\Code{2.6}{CopyPayload(p, q);}
\Code{2.6}{p->rchild = AVLDelete(p->rchild, q->key);}
\Code{1.3}{\}}
\Hint{1.3}{有两个孩子时，直接释放当前根会丢失子树，应转而删除右子树中的后继。}\nopagebreak[4]
\Code{1.3}{else \{}
\Code{2.6}{q = p->lchild?p->lchild:p->rchild;}
\Code{2.6}{free(p);}
\Code{2.6}{return q;}
\Code{1.3}{\}}
\Code{1.3}{return Balance(p);}
\Code{0.0}{\}}
\nopagebreak[4]
\Cost{\ensuremath{O\!\left(\log (n+1)\right)}}{\ensuremath{O\!\left(\log (n+1)\right)}}
\par\medskip

\Group{二叉堆}

\Needspace{15\baselineskip}
\subsubsection{上浮调整}
\Code{0.0}{SiftUp(A, i) \{}
\Hint{1.3}{大顶堆中只可能违反与父结点的关系。}\nopagebreak[4]
\Code{1.3}{while (i>0) \{}
\Code{2.6}{p = (i-1)/2;}
\Code{2.6}{if (A[p] >= A[i]) break;}
\Hint{2.6}{根的下标为零，走到根就不再计算父下标。}\nopagebreak[4]
\Code{2.6}{swap(A[p], A[i]);}
\Code{2.6}{i = p;}
\Code{1.3}{\}}
\Code{0.0}{\}}
\nopagebreak[4]
\Cost{\ensuremath{O\!\left(\log (n+1)\right)}}{\ensuremath{O\!\left(1\right)}}
\par\medskip

\Needspace{16\baselineskip}
\subsubsection{下沉调整}
\Code{0.0}{SiftDown(A, n, i) \{}
\Hint{1.3}{每轮先选较大孩子，保证交换后父结点压住两侧。}\nopagebreak[4]
\Code{1.3}{while (2*i+1<n) \{}
\Code{2.6}{c = 2*i+1;}
\Code{2.6}{if (c+1<n \&\& A[c+1]>A[c]) c++;}
\Code{2.6}{if (A[i] >= A[c]) break;}
\Hint{2.6}{只有左孩子而没有右孩子时，c 保持在左孩子的位置。}\nopagebreak[4]
\Code{2.6}{swap(A[i], A[c]);}
\Code{2.6}{i = c;}
\Code{1.3}{\}}
\Code{0.0}{\}}
\nopagebreak[4]
\Cost{\ensuremath{O\!\left(\log (n+1)\right)}}{\ensuremath{O\!\left(1\right)}}
\par\medskip

\Needspace{10\baselineskip}
\subsubsection{逐个插入建堆}
\Code{0.0}{BuildByInsert(A, n) \{}
\Hint{1.3}{把 A 的前缀当成堆，边界每次扩大一项。}\nopagebreak[4]
\Hint{1.3}{前 i 项已经构成堆，插入新项只可能破坏其向上的祖先链。}\nopagebreak[4]
\Code{1.3}{for (i = 1; i<n; i++) SiftUp(A, i);}
\Code{0.0}{\}}
\nopagebreak[4]
\Cost{\ensuremath{O\!\left(n \log (n+1)\right)}}{\ensuremath{O\!\left(1\right)}}
\par\medskip

\Needspace{10\baselineskip}
\subsubsection{自底向上建堆}
\Code{0.0}{BuildHeap(A, n) \{}
\Hint{1.3}{叶子本身已是堆，从最后一个有孩子的结点开始。}\nopagebreak[4]
\Hint{1.3}{倒序循环的下标用有符号整数，减到 −1 时应结束。}\nopagebreak[4]
\Code{1.3}{for (int i = n/2-1; i >= 0; i--) SiftDown(A, n, i);}
\Code{0.0}{\}}
\nopagebreak[4]
\Cost{\ensuremath{O\!\left(n\right)}}{\ensuremath{O\!\left(1\right)}}
\par\medskip

\Needspace{13\baselineskip}
\subsubsection{插入新元素}
\Code{0.0}{HeapPush(A, int *n, x) \{}
\Hint{1.3}{新值写在第一空槽，再向上恢复堆序。}\nopagebreak[4]
\Code{1.3}{require((*n)<capacity);}
\Code{1.3}{A[(*n)] = x;}
\Hint{1.3}{插入需要可写槽位；扩容另计，表中给的是容量已足够的复杂度。}\nopagebreak[4]
\Code{1.3}{(*n)++;}
\Code{1.3}{SiftUp(A, (*n)-1);}
\Code{0.0}{\}}
\nopagebreak[4]
\Cost{\ensuremath{O\!\left(\log (n+1)\right)}}{\ensuremath{O\!\left(1\right)}}
\par\medskip

\Needspace{17\baselineskip}
\subsubsection{删除堆顶元素}
\Code{0.0}{HeapPop(A, int *n) \{}
\Hint{1.3}{尾元素补到根，旧尾随长度缩短而退出堆。}\nopagebreak[4]
\Code{1.3}{require((*n)>0);}
\Code{1.3}{x = A[0];}
\Code{1.3}{(*n)--;}
\Code{1.3}{if ((*n)>0) \{}
\Hint{2.6}{删除单元素堆后长度为零，此时不应再调用下沉。}\nopagebreak[4]
\Code{2.6}{A[0] = A[(*n)];}
\Code{2.6}{SiftDown(A, (*n), 0);}
\Code{1.3}{\}}
\Code{1.3}{return x;}
\Code{0.0}{\}}
\nopagebreak[4]
\Cost{\ensuremath{O\!\left(\log (n+1)\right)}}{\ensuremath{O\!\left(1\right)}}
\par\medskip

\Needspace{15\baselineskip}
\subsubsection{删除任意指定元素}
\Code{0.0}{HeapErase(A, int *n, i) \{}
\Hint{1.3}{替代值可能破坏父边或子边，只需选一个方向。}\nopagebreak[4]
\Code{1.3}{require(0 <= i \&\& i<(*n));}
\Code{1.3}{(*n)--;}
\Code{1.3}{if (i == (*n)) return;}
\Hint{1.3}{删的是末项时缩短长度就已完成，不能继续访问原下标。}\nopagebreak[4]
\Code{1.3}{A[i] = A[(*n)];}
\Code{1.3}{if (i>0 \&\& A[i]>A[(i-1)/2]) SiftUp(A, i);}
\Code{1.3}{else SiftDown(A, (*n), i);}
\Code{0.0}{\}}
\nopagebreak[4]
\Cost{\ensuremath{O\!\left(\log (n+1)\right)}（下标已知）}{\ensuremath{O\!\left(1\right)}}
\par\medskip

\Needspace{14\baselineskip}
\subsubsection{修改指定元素的权重}
\Code{0.0}{HeapChange(A, n, i, x) \{}
\Hint{1.3}{权重增大向上检查，减小向下检查。}\nopagebreak[4]
\Code{1.3}{require(0 <= i \&\& i<n);}
\Code{1.3}{old = A[i];}
\Hint{1.3}{其他元素之间的父子关系未改，只需修正与当前元素有关的一条路径。}\nopagebreak[4]
\Code{1.3}{A[i] = x;}
\Code{1.3}{if (x>old) SiftUp(A, i);}
\Code{1.3}{else if (x<old) SiftDown(A, n, i);}
\Code{0.0}{\}}
\nopagebreak[4]
\Cost{\ensuremath{O\!\left(\log (n+1)\right)}（下标已知）}{\ensuremath{O\!\left(1\right)}}
\par\medskip

\Needspace{14\baselineskip}
\subsubsection{判断一个数组是否满足堆序性质}
\Code{0.0}{IsMaxHeap(A, n) \{}
\Hint{1.3}{每条父子边只从孩子端检查一次。}\nopagebreak[4]
\Code{1.3}{for (i = 1; i<n; i++) \{}
\Code{2.6}{p = (i-1)/2;}
\Hint{2.6}{数组本身给出完全树形，因此这里只需逐边检查大小关系。}\nopagebreak[4]
\Code{2.6}{if (A[p]<A[i]) return false;}
\Code{1.3}{\}}
\Code{1.3}{return true;}
\Code{0.0}{\}}
\nopagebreak[4]
\Cost{\ensuremath{O\!\left(n\right)}}{\ensuremath{O\!\left(1\right)}}
\par\medskip

\Needspace{14\baselineskip}
\subsubsection{堆排序}
\Code{0.0}{HeapSort(A, n) \{}
\Hint{1.3}{末端是已排好的后缀，下沉只作用于剩余前缀。}\nopagebreak[4]
\Code{1.3}{BuildHeap(A, n);}
\Code{1.3}{for (end = n-1; end>0; end--) \{}
\Hint{2.6}{交换后最大项已经在最终位置，下一次下沉不得把它重新纳入堆。}\nopagebreak[4]
\Code{2.6}{swap(A[0], A[end]);}
\Code{2.6}{SiftDown(A, end, 0);}
\Code{1.3}{\}}
\Code{0.0}{\}}
\nopagebreak[4]
\Cost{\ensuremath{O\!\left(n \log (n+1)\right)}}{\ensuremath{O\!\left(1\right)}}
\par\medskip

\Needspace{19\baselineskip}
\subsubsection{利用堆解决Top-K问题}
\Code{0.0}{TopK(A, n, k) \{}
\Hint{1.3}{小顶堆的根是当前入选下限，只有更大的新值才替换它。}\nopagebreak[4]
\Code{1.3}{require(k >= 1 \&\& k <= n);}
\Code{1.3}{InitMinHeap(\&H);}
\Code{1.3}{for (i = 0; i<n; i++) \{}
\Code{2.6}{if (Size(\&H)<k) MinPush(\&H, A[i]);}
\Code{2.6}{else if (A[i]>Min(\&H)) \{}
\Hint{3.9}{最终结果是堆内的 k 项，如果还要求有序展示，需要额外排序。}\nopagebreak[4]
\Code{3.9}{ExtractMin(\&H);}
\Code{3.9}{MinPush(\&H, A[i]);}
\Code{2.6}{\}}
\Code{1.3}{\}}
\Code{1.3}{return H;}
\Code{0.0}{\}}
\nopagebreak[4]
\Cost{\ensuremath{O\!\left(n \log (k+1)\right)}}{\ensuremath{O\!\left(k\right)}}
\par\medskip

\Group{哈夫曼树}

\Needspace{8\baselineskip}
\subsubsection{哈夫曼树的静态三叉链表构建}
\Code{0.0}{HTNode *CreateHuffmanTree(int w[], int n) \{}
\Hint{1.3}{下标 0 到 n−1 为叶。parent 为 −1 表示仍可参加下一次合并。}\nopagebreak[4]
\Code{1.3}{if (n <= 0) return NULL;}
\Code{1.3}{int m = 2*n-1;}
\Code{1.3}{HTNode *HT = malloc(m*sizeof(*HT));}
\Code{1.3}{if (HT == NULL) return NULL;}
\Code{1.3}{for (int i = 0; i<m; i++) \{}
\Code{2.6}{HT[i].weight = i<n?w[i]:0;}
\Code{2.6}{HT[i].parent = HT[i].lchild = HT[i].rchild = -1;}
\Code{1.3}{\}}
\Code{1.3}{for (int i = n; i<m; i++) \{}
\Code{2.6}{int s1 = -1, s2 = -1;}
\Code{2.6}{for (int j = 0; j<i; j++) \{}
\Code{3.9}{if (HT[j].parent != -1) continue;}
\Code{3.9}{if (s1 == -1 || HT[j].weight<HT[s1].weight) \{}
\Hint{5.2}{发现新的最小值时，原第一名仍是第二小候选。}\nopagebreak[4]
\Code{5.2}{s2 = s1;}
\Code{5.2}{s1 = j;}
\Code{3.9}{\}}
\Code{3.9}{else if (s2 == -1 || HT[j].weight<HT[s2].weight) s2 = j;}
\Code{2.6}{\}}
\Hint{2.6}{s1 更新为新最小时，原 s1 要退到 s2。两个孩子写入 parent 后，后续轮次会跳过它们。}\nopagebreak[4]
\Code{2.6}{HT[i].weight = HT[s1].weight+HT[s2].weight;}
\Code{2.6}{HT[i].lchild = s1;}
\Code{2.6}{HT[i].rchild = s2;}
\Code{2.6}{HT[s1].parent = HT[s2].parent = i;}
\Code{1.3}{\}}
\Code{1.3}{return HT;}
\Code{0.0}{\}}
\nopagebreak[4]
\Cost{\ensuremath{O\!\left(n^{2}\right)}}{\ensuremath{O\!\left(1\right)}}
\par\medskip

\Needspace{19\baselineskip}
\subsubsection{基于最小堆的哈夫曼树构建}
\Code{0.0}{HuffmanHeap(w, n) \{}
\Hint{1.3}{堆中每个元素都是一棵尚未合并的子树的根。}\nopagebreak[4]
\Code{1.3}{if (!n) return NULL;}
\Code{1.3}{InitMinHeap(\&H);}
\Code{1.3}{for (i = 0; i<n; i++) MinPush(\&H, NewWeightedLeaf(i, w[i]));}
\Code{1.3}{while (Size(\&H)>1) \{}
\Code{2.6}{a = ExtractMin(\&H);}
\Hint{2.6}{每次取两棵树、放回一棵，堆大小减少一。最终剩下的结点就是全树根。}\nopagebreak[4]
\Code{2.6}{b = ExtractMin(\&H);}
\Code{2.6}{p = NewParent(a, b, a->weight+b->weight);}
\Code{2.6}{MinPush(\&H, p);}
\Code{1.3}{\}}
\Code{1.3}{return ExtractMin(\&H);}
\Code{0.0}{\}}
\nopagebreak[4]
\Cost{\ensuremath{O\!\left(n \log (n+1)\right)}}{\ensuremath{O\!\left(n\right)}}
\par\medskip

\Needspace{12\baselineskip}
\subsubsection{计算WPL}
\Code{0.0}{WPL(p, depth) \{}
\Hint{1.3}{只有叶参与加权求和，根深度传 0。}\nopagebreak[4]
\Code{1.3}{if (!p) return 0;}
\Hint{1.3}{递归进入孩子时深度加一，同一叶子的权重只在该叶被计算一次。}\nopagebreak[4]
\Code{1.3}{if (IsLeaf(p)) return p->weight*depth;}
\Code{1.3}{return WPL(p->lchild, depth+1)+WPL(p->rchild, depth+1);}
\Code{0.0}{\}}
\nopagebreak[4]
\Cost{\ensuremath{O\!\left(n\right)}}{\ensuremath{O\!\left(h\right)}}
\par\medskip

\Needspace{19\baselineskip}
\subsubsection{生成哈夫曼编码}
\Code{0.0}{WriteCodes(p, bits, d) \{}
\Hint{1.3}{路径共用数组，左子树回来后直接把本位改成 1。}\nopagebreak[4]
\Code{1.3}{if (!p) return;}
\Code{1.3}{if (IsLeaf(p)) \{}
\Code{2.6}{if (d == 0) EmitCode(p, "0");}
\Code{2.6}{else EmitCode(p, bits, d);}
\Code{2.6}{return;}
\Code{1.3}{\}}
\Hint{1.3}{d 表示已填入的位数；走到叶后只输出前 d 位，不输出数组未用区域。}\nopagebreak[4]
\Code{1.3}{bits[d] = 0;}
\Code{1.3}{WriteCodes(p->lchild, bits, d+1);}
\Code{1.3}{bits[d] = 1;}
\Code{1.3}{WriteCodes(p->rchild, bits, d+1);}
\Code{0.0}{\}}
\nopagebreak[4]
\Cost{\ensuremath{O\!\left(n+L\right)}}{\ensuremath{O\!\left(h\right)}}
\par\medskip

\Needspace{25\baselineskip}
\subsubsection{哈夫曼编码的解码}
\Code{0.0}{DecodeBits(root, bits) \{}
\Hint{1.3}{每次到叶立即输出并回到根；最后停在内部结点表示码流截断。}\nopagebreak[4]
\Code{1.3}{p = root;}
\Code{1.3}{for (b in bits) \{}
\Code{2.6}{require(b == 0 || b == 1);}
\Code{2.6}{if (IsLeaf(root)) \{}
\Code{3.9}{require(b == 0);}
\Code{3.9}{Emit(root->symbol);}
\Code{3.9}{continue;}
\Code{2.6}{\}}
\Code{2.6}{p = b?p->rchild:p->lchild;}
\Hint{2.6}{对于单符号树，本手册约定其码为一位 0，需单独处理。}\nopagebreak[4]
\Code{2.6}{require(p);}
\Code{2.6}{if (IsLeaf(p)) \{}
\Code{3.9}{Emit(p->symbol);}
\Code{3.9}{p = root;}
\Code{2.6}{\}}
\Code{1.3}{\}}
\Code{1.3}{require(p == root);}
\Code{0.0}{\}}
\nopagebreak[4]
\Cost{\ensuremath{O\!\left(b\right)}}{\ensuremath{O\!\left(1\right)}}
\par\medskip

\Needspace{22\baselineskip}
\subsubsection{判断一组编码是否为合法前缀码}
\Code{0.0}{PrefixFree(codes) \{}
\Hint{1.3}{途中遇词尾表示旧码是前缀，走完仍有孩子表示新码是前缀。}\nopagebreak[4]
\Code{1.3}{root = NewTrie();}
\Code{1.3}{for (s in codes) \{}
\Code{2.6}{if (len(s) == 0) return false;}
\Code{2.6}{p = root;}
\Code{2.6}{for (b in s) \{}
\Code{3.9}{if (p->end || !IsBit(b)) return false;}
\Code{3.9}{if (!p->child[b]) p->child[b] = NewTrie();}
\Hint{3.9}{终点已有词尾标记表示重复码，终点已有孩子表示它是旧码的前缀。}\nopagebreak[4]
\Code{3.9}{p = p->child[b];}
\Code{2.6}{\}}
\Code{2.6}{if (p->end || HasChild(p)) return false;}
\Code{2.6}{p->end = true;}
\Code{1.3}{\}}
\Code{1.3}{return true;}
\Code{0.0}{\}}
\nopagebreak[4]
\Cost{\ensuremath{O\!\left(L\right)}}{\ensuremath{O\!\left(L\right)}}
\par\medskip

\Needspace{14\baselineskip}
\subsubsection{判断一组编码是否为最优前缀码}
\Code{0.0}{OptimalCode(w, codes, n) \{}
\Hint{1.3}{先验符号一一对应与前缀性，再与哈夫曼最小加权码长比较。}\nopagebreak[4]
\Code{1.3}{if (!ExactSymbols(codes, w) || !PrefixFree(codes)) return false;}
\Code{1.3}{if (n == 1) return len(codes[0]) == 1;}
\Hint{1.3}{不同哈夫曼树可能同样最优，因此比较加权总码长，不逐个比较标准码字。}\nopagebreak[4]
\Code{1.3}{total = 0;}
\Code{1.3}{for (i = 0; i<n; i++) total += w[i]*len(codes[i]);}
\Code{1.3}{return total == WPL(HuffmanHeap(w, n), 0);}
\Code{0.0}{\}}
\nopagebreak[4]
\Cost{\ensuremath{O\!\left(L+n \log (n+1)\right)}}{\ensuremath{O\!\left(L+n\right)}}
\par\medskip

\Group{前缀树（Trie）}

\Needspace{14\baselineskip}
\subsubsection{Trie的静态数组实现}
\Code{0.0}{InitTrie(C, sigma) \{}
\Hint{1.3}{编号 0 为空链接，根使用编号 1。}\nopagebreak[4]
\Code{1.3}{Zero(next);}
\Code{1.3}{Zero(pass);}
\Hint{1.3}{每个新分配的结点也要把子边和计数清零，容量不足时报告失败。}\nopagebreak[4]
\Code{1.3}{Zero(end);}
\Code{1.3}{used = 1;}
\Code{1.3}{root = 1;}
\Code{0.0}{\}}
\nopagebreak[4]
\Cost{初始化 \ensuremath{O\!\left(C\sigma \right)}}{\ensuremath{O\!\left(C\sigma \right)}（结构）}
\par\medskip

\Needspace{17\baselineskip}
\subsubsection{Trie的插入}
\Code{0.0}{TrieInsert(word) \{}
\Hint{1.3}{pass 在经过的每个结点累加，end 只在词尾累加。}\nopagebreak[4]
\Code{1.3}{p = root;}
\Code{1.3}{pass[p]++;}
\Code{1.3}{for (c in word) \{}
\Code{2.6}{if (!next[p][c]) next[p][c] = AllocZeroTrieSlot();}
\Hint{2.6}{支持重复插入，因此同一词的每次出现都要给词尾计数加一。}\nopagebreak[4]
\Code{2.6}{p = next[p][c];}
\Code{2.6}{pass[p]++;}
\Code{1.3}{\}}
\Code{1.3}{end[p]++;}
\Code{0.0}{\}}
\nopagebreak[4]
\Cost{\ensuremath{O\!\left(\ell \right)}}{\ensuremath{O\!\left(1\right)}（预分配）}
\par\medskip

\Needspace{15\baselineskip}
\subsubsection{Trie的查找}
\Code{0.0}{TrieFind(word) \{}
\Hint{1.3}{路径存在仍可能只是前缀，须检查词尾计数。}\nopagebreak[4]
\Code{1.3}{p = root;}
\Code{1.3}{for (c in word) \{}
\Code{2.6}{p = next[p][c];}
\Hint{2.6}{删除过一个词后该路径可能仍为其他词保留；end 能区分这种情况。}\nopagebreak[4]
\Code{2.6}{if (!p) return false;}
\Code{1.3}{\}}
\Code{1.3}{return end[p]>0;}
\Code{0.0}{\}}
\nopagebreak[4]
\Cost{\ensuremath{O\!\left(\ell \right)}}{\ensuremath{O\!\left(1\right)}}
\par\medskip

\Needspace{22\baselineskip}
\subsubsection{Trie的删除}
\Code{0.0}{TrieDelete(word) \{}
\Hint{1.3}{先确认整词存在；逆向清除 pass 降到零的分支。}\nopagebreak[4]
\Code{1.3}{if (!TrieFind(word)) return false;}
\Code{1.3}{path[0] = root;}
\Code{1.3}{for (i = 0; i<len(word); i++) path[i+1] = next[path[i]][word[i]];}
\Code{1.3}{end[path[len(word)]]--;}
\Code{1.3}{for (i = len(word); i >= 0; i--) \{}
\Code{2.6}{p = path[i];}
\Code{2.6}{pass[p]--;}
\Hint{2.6}{先走完整查询再更新计数，避免删除不存在的词时部分改坏路径。}\nopagebreak[4]
\Code{2.6}{if (i>0 \&\& pass[p] == 0) \{}
\Code{3.9}{next[path[i-1]][word[i-1]] = 0;}
\Code{3.9}{RecycleSlot(p);}
\Code{2.6}{\}}
\Code{1.3}{\}}
\Code{1.3}{return true;}
\Code{0.0}{\}}
\nopagebreak[4]
\Cost{\ensuremath{O\!\left(\ell \right)}}{\ensuremath{O\!\left(\ell \right)}}
\par\medskip

\Needspace{15\baselineskip}
\subsubsection{统计具有指定前缀的字符串数量}
\Code{0.0}{PrefixCount(prefix) \{}
\Hint{1.3}{走完前缀时 pass 包含所有经过该点的词，包括重复词。}\nopagebreak[4]
\Code{1.3}{p = root;}
\Code{1.3}{for (c in prefix) \{}
\Code{2.6}{p = next[p][c];}
\Hint{2.6}{空前缀不会移动 p，此时返回根上的总词数。}\nopagebreak[4]
\Code{2.6}{if (!p) return 0;}
\Code{1.3}{\}}
\Code{1.3}{return pass[p];}
\Code{0.0}{\}}
\nopagebreak[4]
\Cost{\ensuremath{O\!\left(\ell \right)}}{\ensuremath{O\!\left(1\right)}}
\par\medskip

\Needspace{21\baselineskip}
\subsubsection{基于二进制Trie的最大异或值查询}
\Code{0.0}{MaxXor(x, W) \{}
\Hint{1.3}{高位的贡献大于全部更低位之和，优先取相反位。}\nopagebreak[4]
\Code{1.3}{p = root;}
\Code{1.3}{ans = 0;}
\Code{1.3}{for (i = W-1; i >= 0; i--) \{}
\Code{2.6}{b = (x>>i)\&1;}
\Code{2.6}{q = next[p][1-b];}
\Code{2.6}{if (q \&\& pass[q]>0) \{}
\Hint{3.9}{若相反分支不存在，仍须沿相同位进入下一层，当前答案位保持零。}\nopagebreak[4]
\Code{3.9}{ans|=(1U<<i);}
\Code{3.9}{p = q;}
\Code{2.6}{\}}
\Code{2.6}{else p = next[p][b];}
\Code{1.3}{\}}
\Code{1.3}{return ans;}
\Code{0.0}{\}}
\nopagebreak[4]
\Cost{\ensuremath{O\!\left(W\right)}}{\ensuremath{O\!\left(1\right)}}
\par\medskip

\Needspace{26\baselineskip}
\subsubsection{后缀Trie的构建与子串查找}
\Code{0.0}{SuffixTrie(s, n) \{}
\Hint{1.3}{子串恰为某个后缀的前缀，查询不用检查词尾。}\nopagebreak[4]
\Code{1.3}{for (i = 0; i<n; i++) \{}
\Code{2.6}{p = root;}
\Code{2.6}{for (j = i; j<n; j++) \{}
\Code{3.9}{c = s[j];}
\Code{3.9}{if (!next[p][c]) next[p][c] = AllocZeroTrieSlot();}
\Code{3.9}{p = next[p][c];}
\Code{2.6}{\}}
\Code{2.6}{end[p] = 1;}
\Code{1.3}{\}}
\Code{0.0}{\}}
\Hint{0.0}{每个后缀都从根重新出发，相同前缀自动共享已有边。}\nopagebreak[4]
\Code{0.0}{Contains(P) \{}
\Code{1.3}{p = root;}
\Code{1.3}{for (c in P) \{}
\Code{2.6}{p = next[p][c];}
\Code{2.6}{if (!p) return false;}
\Code{1.3}{\}}
\Code{1.3}{return true;}
\Code{0.0}{\}}
\nopagebreak[4]
\Cost{建 \ensuremath{O\!\left(n^{2}\right)}，查 \ensuremath{O\!\left(\ell \right)}}{建 \ensuremath{O\!\left(n^{2}\right)}（结点）；查 \ensuremath{O\!\left(1\right)}}
\par\medskip

\Needspace{8\baselineskip}
\subsubsection{后缀树}
\Code{0.0}{SuffixTree(s) \{}
\Hint{1.3}{加唯一结束符后没有后缀是另一后缀的前缀。边标签存左闭右开区间。}\nopagebreak[4]
\Code{1.3}{AppendTerminator(\&s);}
\Code{1.3}{n = len(s);}
\Code{1.3}{for (i = 0; i<n; i++) \{}
\Code{2.6}{p = root;}
\Code{2.6}{j = i;}
\Code{2.6}{while (j<n) \{}
\Code{3.9}{e = Edge(p, s[j]);}
\Code{3.9}{if (!e) \{}
\Code{5.2}{AddEdge(p, j, n, NewLeaf(i));}
\Code{5.2}{break;}
\Code{3.9}{\}}
\Code{3.9}{k = 0;}
\Code{3.9}{while (e.lo+k<e.hi \&\& s[e.lo+k] == s[j+k]) k++;}
\Code{3.9}{j += k;}
\Code{3.9}{if (e.lo+k == e.hi) \{}
\Code{5.2}{p = e.child;}
\Code{5.2}{continue;}
\Code{3.9}{\}}
\Code{3.9}{mid = SplitEdge(p, e, k);}
\Code{3.9}{AddEdge(mid, j, n, NewLeaf(i));}
\Hint{3.9}{SplitEdge 在原边内创建分叉，旧边的剩余标签和新后缀余段分别成为它的两条出边。}\nopagebreak[4]
\Code{3.9}{break;}
\Code{2.6}{\}}
\Code{1.3}{\}}
\Code{1.3}{return root;}
\Code{0.0}{\}}
\Code{0.0}{SuffixContains(P) \{}
\Code{1.3}{p = root;}
\Code{1.3}{i = 0;}
\Code{1.3}{while (i<len(P)) \{}
\Code{2.6}{e = Edge(p, P[i]);}
\Code{2.6}{if (!e) return false;}
\Code{2.6}{for (k = e.lo; k<e.hi \&\& i<len(P); k++, i++) if (s[k] != P[i]) return false;}
\Code{2.6}{p = e.child;}
\Code{1.3}{\}}
\Code{1.3}{return true;}
\Code{0.0}{\}}
\nopagebreak[4]
\Cost{建 \ensuremath{O\!\left(n^{2}\right)}，查 \ensuremath{O\!\left(\ell \right)}（固定字母表）}{建 \ensuremath{O\!\left(n\right)}；查 \ensuremath{O\!\left(1\right)}}
\par\medskip

\Needspace{8\baselineskip}
\subsubsection{后缀自动机}
\Code{0.0}{ExtendSAM(c) \{}
\Hint{1.3}{对每个字符调用一次。clone 复制转移与后缀链接，只缩短最大长度。}\nopagebreak[4]
\Code{1.3}{cur = NewState();}
\Code{1.3}{len[cur] = len[last]+1;}
\Code{1.3}{p = last;}
\Code{1.3}{while (p != NONE \&\& !HasEdge(p, c)) \{}
\Code{2.6}{next[p][c] = cur;}
\Code{2.6}{p = link[p];}
\Code{1.3}{\}}
\Code{1.3}{if (p == NONE) link[cur] = root;}
\Code{1.3}{else \{}
\Code{2.6}{q = next[p][c];}
\Code{2.6}{if (len[p]+1 == len[q]) link[cur] = q;}
\Code{2.6}{else \{}
\Code{3.9}{clone = CopyState(q);}
\Code{3.9}{len[clone] = len[p]+1;}
\Code{3.9}{while (p != NONE \&\& next[p][c] == q) \{}
\Code{5.2}{next[p][c] = clone;}
\Hint{5.2}{沿 link 上行重定向时只改原先指向 q 的边，其他转移不属于这次拆分。}\nopagebreak[4]
\Code{5.2}{p = link[p];}
\Code{3.9}{\}}
\Code{3.9}{link[q] = link[cur] = clone;}
\Code{2.6}{\}}
\Code{1.3}{\}}
\Code{1.3}{last = cur;}
\Code{0.0}{\}}
\Code{0.0}{BuildSAM(s) \{}
\Code{1.3}{root = NewState();}
\Code{1.3}{len[root] = 0;}
\Code{1.3}{link[root] = NONE;}
\Code{1.3}{last = root;}
\Code{1.3}{for (c in s) ExtendSAM(c);}
\Code{0.0}{\}}
\nopagebreak[4]
\Cost{建 \ensuremath{O\!\left(n\right)}，查 \ensuremath{O\!\left(\ell \right)}（固定字母表）}{\ensuremath{O\!\left(n\right)}}
\par\medskip

\Part{图}

\Group{图的存储结构}

\Needspace{14\baselineskip}
\subsubsection{邻接矩阵的建立}
\Code{0.0}{BuildMatrix(G, edges) \{}
\Hint{1.3}{无边使用专用标记，避免与合法零权边混淆。}\nopagebreak[4]
\Code{1.3}{Fill(G.mat, NOEDGE);}
\Code{1.3}{for ((u, v, w) in edges) \{}
\Hint{2.6}{有向图只填起点行终点列；无向图还要补上对称位置。}\nopagebreak[4]
\Code{2.6}{G.mat[u][v] = w;}
\Code{2.6}{if (!G.directed) G.mat[v][u] = w;}
\Code{1.3}{\}}
\Code{0.0}{\}}
\nopagebreak[4]
\Cost{\ensuremath{O\!\left(V^{2}+E\right)}}{\ensuremath{O\!\left(1\right)}}
\par\medskip

\Needspace{14\baselineskip}
\subsubsection{邻接表的建立}
\Code{0.0}{BuildAdj(G, edges) \{}
\Hint{1.3}{无向边的两条邻接记录共享 id，可辨别平行边。}\nopagebreak[4]
\Code{1.3}{Fill(G.first, NULL);}
\Code{1.3}{for ((u, v, w, id) in edges) \{}
\Hint{2.6}{头插只改变该顶点的首边，不要求输入边按端点排好。}\nopagebreak[4]
\Code{2.6}{PrependArc(\&G.first[u], v, w, id);}
\Code{2.6}{if (!G.directed) PrependArc(\&G.first[v], u, w, id);}
\Code{1.3}{\}}
\Code{0.0}{\}}
\nopagebreak[4]
\Cost{\ensuremath{O\!\left(V+E\right)}}{\ensuremath{O\!\left(1\right)}}
\par\medskip

\Needspace{16\baselineskip}
\subsubsection{邻接矩阵与邻接表的相互转换}
\Code{0.0}{ConvertGraph(G, toList) \{}
\Hint{1.3}{无向矩阵须对称，只处理上三角及对角线；同一条边的两条记录共享编号。}\nopagebreak[4]
\Code{1.3}{if (toList) \{}
\Code{2.6}{Fill(G.first, NULL); id = 0;}
\Code{2.6}{for (u in vertices) for (v in vertices) \{}
\Code{3.9}{w = G.mat[u][v];}
\Code{3.9}{if (w == NOEDGE || (!G.directed \&\& v<u)) continue;}
\Code{3.9}{PrependArc(\&G.first[u], v, w, id);}
\Code{3.9}{if (!G.directed) PrependArc(\&G.first[v], u, w, id);}
\Code{3.9}{id++;}
\Code{2.6}{\}}
\Code{1.3}{\} else \{}
\Hint{2.6}{矩阵默认表示简单图；若允许平行边，事先约定聚合权值或改用多值单元。}\nopagebreak[4]
\Code{2.6}{Fill(G.mat, NOEDGE);}
\Code{2.6}{for (u in vertices) for (e in G.first[u]) G.mat[u][e.to] = e.weight;}
\Code{1.3}{\}}
\Code{0.0}{\}}
\nopagebreak[4]
\Cost{\ensuremath{O\!\left(V^{2}+E\right)}}{\ensuremath{O\!\left(1\right)}}
\par\medskip

\Needspace{14\baselineskip}
\subsubsection{有向图的十字链表表示}
\Code{0.0}{InsertCrossArc(G, u, v) \{}
\Hint{1.3}{一条弧同时进入起点的出弧链与终点的入弧链。}\nopagebreak[4]
\Code{1.3}{e = NewArc(u, v);}
\Code{1.3}{e->nextOut = G.firstOut[u];}
\Hint{1.3}{先保存旧的两条表头，再改表头，否则会把原有弧链丢掉。}\nopagebreak[4]
\Code{1.3}{e->nextIn = G.firstIn[v];}
\Code{1.3}{G.firstOut[u] = e;}
\Code{1.3}{G.firstIn[v] = e;}
\Code{0.0}{\}}
\nopagebreak[4]
\Cost{单次插边 \ensuremath{O\!\left(1\right)}，初始化表头并逐边建图 \ensuremath{O\!\left(V+E\right)}}{\ensuremath{O\!\left(1\right)}}
\par\medskip

\Needspace{13\baselineskip}
\subsubsection{无向图的邻接多重表表示}
\Code{0.0}{InsertMultiEdge(G, a, b) \{}
\Hint{1.3}{每个端点都有自己的后继指针；此版本处理无自环图。}\nopagebreak[4]
\Code{1.3}{e = NewEdge(a, b);}
\Code{1.3}{e->aNext = G.first[a];}
\Hint{1.3}{遍历某顶点的邻边时，根据它位于 a 端还是 b 端选择相应后继域。}\nopagebreak[4]
\Code{1.3}{e->bNext = G.first[b];}
\Code{1.3}{G.first[a] = G.first[b] = e;}
\Code{0.0}{\}}
\nopagebreak[4]
\Cost{单次插边 \ensuremath{O\!\left(1\right)}，初始化表头并逐边建图 \ensuremath{O\!\left(V+E\right)}}{\ensuremath{O\!\left(1\right)}}
\par\medskip

\Needspace{27\baselineskip}
\subsubsection{顶点的插入与删除}
\Code{0.0}{VertexEdit(G, v, insert) \{}
\Hint{1.3}{编号保持稳定。删除顶点时必须清除所有指向它的入边。}\nopagebreak[4]
\Code{1.3}{if (insert) \{}
\Code{2.6}{G.used[v] = true;}
\Code{2.6}{G.first[v] = NULL;}
\Code{2.6}{return;}
\Code{1.3}{\}}
\Code{1.3}{FreeChain(\&G.first[v]);}
\Code{1.3}{for (u in vertices) \{}
\Code{2.6}{slot = \&G.first[u];}
\Code{2.6}{while (*slot) \{}
\Code{3.9}{e = *slot;}
\Hint{3.9}{从邻接表删除中间项后，slot 所在的入边域已指向下一待检查项。}\nopagebreak[4]
\Code{3.9}{if (e->to == v) \{}
\Code{5.2}{*slot = e->next;}
\Code{5.2}{free(e);}
\Code{3.9}{\}}
\Code{3.9}{else slot = \&e->next;}
\Code{2.6}{\}}
\Code{1.3}{\}}
\Code{1.3}{G.used[v] = false;}
\Code{0.0}{\}}
\nopagebreak[4]
\Cost{插 \ensuremath{O\!\left(1\right)}（预留槽），删 \ensuremath{O\!\left(V+E\right)}}{\ensuremath{O\!\left(1\right)}}
\par\medskip

\Needspace{25\baselineskip}
\subsubsection{边的插入与删除}
\Code{0.0}{EdgeEdit(G, u, v, w, id, insert) \{}
\Hint{1.3}{无向图同时更新两端；删除按编号匹配。}\nopagebreak[4]
\Code{1.3}{if (insert) \{}
\Code{2.6}{PrependArc(\&G.first[u], v, w, id);}
\Code{2.6}{if (!G.directed) PrependArc(\&G.first[v], u, w, id);}
\Code{1.3}{\}}
\Code{1.3}{else \{}
\Code{2.6}{UnlinkId(\&G.first[u], id);}
\Code{2.6}{if (!G.directed) UnlinkId(\&G.first[v], id);}
\Code{1.3}{\}}
\Code{0.0}{\}}
\Hint{0.0}{若按终点而不是编号删除，平行边情形下可能删掉另一条边。}\nopagebreak[4]
\Code{0.0}{UnlinkId(slot, id) \{}
\Code{1.3}{while (*slot \&\& (*slot)->id != id) slot = \&(*slot)->next;}
\Code{1.3}{if (*slot) \{}
\Code{2.6}{p = *slot;}
\Code{2.6}{*slot = p->next;}
\Code{2.6}{free(p);}
\Code{1.3}{\}}
\Code{0.0}{\}}
\nopagebreak[4]
\Cost{插 \ensuremath{O\!\left(1\right)}，删 \ensuremath{O\!\left(\deg (u)+\deg (v)+1\right)}}{\ensuremath{O\!\left(1\right)}}
\par\medskip

\Needspace{22\baselineskip}
\subsubsection{计算顶点的度、入度与出度}
\Code{0.0}{Degrees(G) \{}
\Hint{1.3}{无向图按物理边遍历一次，自环两端相同，度数自然加二。}\nopagebreak[4]
\Code{1.3}{Zero(in);}
\Code{1.3}{Zero(out);}
\Code{1.3}{Zero(deg);}
\Code{1.3}{for ((u, v) in edges) \{}
\Code{2.6}{if (G.directed) \{}
\Code{3.9}{out[u]++;}
\Code{3.9}{in[v]++;}
\Code{2.6}{\}}
\Hint{2.6}{in 与 out 分开累计，有向顶点的总度是二者之和。}\nopagebreak[4]
\Code{2.6}{else \{}
\Code{3.9}{deg[u]++;}
\Code{3.9}{deg[v]++;}
\Code{2.6}{\}}
\Code{1.3}{\}}
\Code{0.0}{\}}
\nopagebreak[4]
\Cost{\ensuremath{O\!\left(V+E\right)}}{\ensuremath{O\!\left(V\right)}}
\par\medskip

\Needspace{14\baselineskip}
\subsubsection{有向图的转置}
\Code{0.0}{TransposeGraph(G) \{}
\Hint{1.3}{保留所有顶点，包括无边的孤立点。}\nopagebreak[4]
\Code{1.3}{T = EmptyGraph(V);}
\Code{1.3}{for (u in vertices) \{}
\Hint{2.6}{边权和边数保持，只有起点与终点对调。}\nopagebreak[4]
\Code{2.6}{for (e in adj[u]) AddArc(\&T, e.to, u, e.weight);}
\Code{1.3}{\}}
\Code{1.3}{return T;}
\Code{0.0}{\}}
\nopagebreak[4]
\Cost{\ensuremath{O\!\left(V+E\right)}}{\ensuremath{O\!\left(1\right)}}
\par\medskip

\Group{图的遍历}

\Needspace{15\baselineskip}
\subsubsection{深度优先搜索（DFS，递归实现）}
\Code{0.0}{DFS(u) \{}
\Hint{1.3}{进入结点立即标记，再向未访问的邻居深入。}\nopagebreak[4]
\Code{1.3}{seen[u] = true;}
\Code{1.3}{visit(u);}
\Code{1.3}{for (e = first[u]; e; e = e->next) \{}
\Hint{2.6}{子调用完成后继续检查 u 的后续邻边，访问标记在各次递归间共享。}\nopagebreak[4]
\Code{2.6}{v = e->to;}
\Code{2.6}{if (!seen[v]) DFS(v);}
\Code{1.3}{\}}
\Code{0.0}{\}}
\nopagebreak[4]
\Cost{\ensuremath{O\!\left(V+E\right)}}{\ensuremath{O\!\left(V\right)}}
\par\medskip

\Needspace{27\baselineskip}
\subsubsection{深度优先搜索（DFS，非递归实现）}
\Code{0.0}{IterDFS(s) \{}
\Hint{1.3}{栈帧保存下一条待查边，回到父结点可从原位置继续。}\nopagebreak[4]
\Code{1.3}{seen[s] = true;}
\Code{1.3}{visit(s);}
\Code{1.3}{Push(\&S, (s, first[s]));}
\Code{1.3}{while (!Empty(\&S)) \{}
\Code{2.6}{f = TopPtr(\&S);}
\Code{2.6}{e = f->edge;}
\Code{2.6}{if (!e) \{}
\Code{3.9}{Pop(\&S);}
\Code{3.9}{continue;}
\Code{2.6}{\}}
\Hint{2.6}{先推进当前帧的游标再压入孩子帧，弹回后不会重查同一条边。}\nopagebreak[4]
\Code{2.6}{f->edge = e->next;}
\Code{2.6}{v = e->to;}
\Code{2.6}{if (!seen[v]) \{}
\Code{3.9}{seen[v] = true;}
\Code{3.9}{visit(v);}
\Code{3.9}{Push(\&S, (v, first[v]));}
\Code{2.6}{\}}
\Code{1.3}{\}}
\Code{0.0}{\}}
\nopagebreak[4]
\Cost{\ensuremath{O\!\left(V+E\right)}}{\ensuremath{O\!\left(V\right)}}
\par\medskip

\Needspace{22\baselineskip}
\subsubsection{广度优先搜索（BFS，队列实现）}
\Code{0.0}{GraphBFS(s) \{}
\Hint{1.3}{标记放在入队处，防止多个邻点把同一结点反复送入队列。}\nopagebreak[4]
\Code{1.3}{seen[s] = true;}
\Code{1.3}{EnQueue(\&Q, s);}
\Code{1.3}{while (!Empty(\&Q)) \{}
\Code{2.6}{u = DeQueue(\&Q);}
\Code{2.6}{visit(u);}
\Code{2.6}{for (e = first[u]; e; e = e->next) \{}
\Code{3.9}{v = e->to;}
\Hint{3.9}{队头取出的点负责把它的未访邻点加入队尾，次序与层数一致。}\nopagebreak[4]
\Code{3.9}{if (!seen[v]) \{}
\Code{5.2}{seen[v] = true;}
\Code{5.2}{EnQueue(\&Q, v);}
\Code{3.9}{\}}
\Code{2.6}{\}}
\Code{1.3}{\}}
\Code{0.0}{\}}
\nopagebreak[4]
\Cost{\ensuremath{O\!\left(V+E\right)}}{\ensuremath{O\!\left(V\right)}}
\par\medskip

\Needspace{26\baselineskip}
\subsubsection{基于邻接矩阵的DFS与BFS}
\Code{0.0}{MatrixDFS(u) \{}
\Hint{1.3}{矩阵需要检查整行，BFS 把首次发现的点加入队尾。}\nopagebreak[4]
\Code{1.3}{seen[u] = true;}
\Code{1.3}{visit(u);}
\Code{1.3}{for (v = 0; v<V; v++) if (M[u][v] != NOEDGE \&\& !seen[v]) MatrixDFS(v);}
\Code{0.0}{\}}
\Code{0.0}{MatrixBFS(s) \{}
\Code{1.3}{seen[s] = true;}
\Code{1.3}{EnQueue(\&Q, s);}
\Code{1.3}{while (!Empty(\&Q)) \{}
\Code{2.6}{u = DeQueue(\&Q);}
\Hint{2.6}{矩阵中顶点对即使没有边也要检查，因此扫描一行始终需要 V 项。}\nopagebreak[4]
\Code{2.6}{visit(u);}
\Code{2.6}{for (v = 0; v<V; v++) \{}
\Code{3.9}{if (M[u][v] != NOEDGE \&\& !seen[v]) \{}
\Code{5.2}{seen[v] = true;}
\Code{5.2}{EnQueue(\&Q, v);}
\Code{3.9}{\}}
\Code{2.6}{\}}
\Code{1.3}{\}}
\Code{0.0}{\}}
\nopagebreak[4]
\Cost{\ensuremath{O\!\left(V^{2}\right)}}{\ensuremath{O\!\left(V\right)}}
\par\medskip

\Needspace{12\baselineskip}
\subsubsection{基于邻接表的DFS与BFS}
\Code{0.0}{ListTraversal(G, s, depthFirst) \{}
\Hint{1.3}{first[u] 指向邻接表首边；两个搜索函数都只沿实际存在的边扫描。}\nopagebreak[4]
\Code{1.3}{Fill(seen, false);}
\Hint{1.3}{这里给出普通入口，具体 DFS 与 BFS 循环见本节前面的两个邻接表实现。}\nopagebreak[4]
\Code{1.3}{if (depthFirst) DFS(s);}
\Code{1.3}{else GraphBFS(s);}
\Code{0.0}{\}}
\nopagebreak[4]
\Cost{\ensuremath{O\!\left(V+E\right)}}{\ensuremath{O\!\left(V\right)}}
\par\medskip

\Needspace{13\baselineskip}
\subsubsection{非连通图的遍历}
\Code{0.0}{TraverseAll(G) \{}
\Hint{1.3}{上一轮的访问标记保留，下一轮只从尚未覆盖的顶点出发。}\nopagebreak[4]
\Code{1.3}{Fill(seen, false);}
\Code{1.3}{for (s in vertices) \{}
\Hint{2.6}{若只从 0 号点开始，其他不与它相连的分量就不会被访问。}\nopagebreak[4]
\Code{2.6}{if (!seen[s]) DFS(s);}
\Code{1.3}{\}}
\Code{0.0}{\}}
\nopagebreak[4]
\Cost{\ensuremath{O\!\left(V+E\right)}}{\ensuremath{O\!\left(V\right)}}
\par\medskip

\Needspace{16\baselineskip}
\subsubsection{深度优先生成树与生成森林}
\Code{0.0}{DFSForest(G) \{}
\Hint{1.3}{只有发现新顶点的那条边记录为树边。}\nopagebreak[4]
\Code{1.3}{Fill(seen, false);}
\Code{1.3}{for (s in vertices) if (!seen[s]) GrowDFS(s, NONE);}
\Code{0.0}{\}}
\Code{0.0}{GrowDFS(u, p) \{}
\Hint{1.3}{每个新根的 parent 取 NONE，识别森林时可据此数树。}\nopagebreak[4]
\Code{1.3}{seen[u] = true;}
\Code{1.3}{parent[u] = p;}
\Code{1.3}{for (v in adj[u]) if (!seen[v]) GrowDFS(v, u);}
\Code{0.0}{\}}
\nopagebreak[4]
\Cost{\ensuremath{O\!\left(V+E\right)}}{\ensuremath{O\!\left(V\right)}}
\par\medskip

\Needspace{24\baselineskip}
\subsubsection{广度优先生成树与生成森林}
\Code{0.0}{BFSForest(G) \{}
\Hint{1.3}{每个新的队列起点都是一棵生成树的根。}\nopagebreak[4]
\Code{1.3}{Fill(seen, false);}
\Code{1.3}{for (s in vertices) \{}
\Code{2.6}{if (seen[s]) continue;}
\Code{2.6}{seen[s] = true;}
\Code{2.6}{parent[s] = NONE;}
\Code{2.6}{EnQueue(\&Q, s);}
\Code{2.6}{while (!Empty(\&Q)) \{}
\Code{3.9}{u = DeQueue(\&Q);}
\Hint{3.9}{遇到已访邻点不再写 parent，否则会把最初发现它的树边覆盖。}\nopagebreak[4]
\Code{3.9}{for (v in adj[u]) if (!seen[v]) \{}
\Code{5.2}{seen[v] = true;}
\Code{5.2}{parent[v] = u;}
\Code{5.2}{EnQueue(\&Q, v);}
\Code{3.9}{\}}
\Code{2.6}{\}}
\Code{1.3}{\}}
\Code{0.0}{\}}
\nopagebreak[4]
\Cost{\ensuremath{O\!\left(V+E\right)}}{\ensuremath{O\!\left(V\right)}}
\par\medskip

\Group{图的基本与综合算法}

\Needspace{13\baselineskip}
\subsubsection{判断两个顶点之间是否存在路径}
\Code{0.0}{Reach(s, t) \{}
\Hint{1.3}{首次命中即可返回，seen 防止环上反复递归。}\nopagebreak[4]
\Code{1.3}{if (s == t) return true;}
\Code{1.3}{seen[s] = true;}
\Hint{1.3}{目标与起点相同应立即成功，对应一条零长路径。}\nopagebreak[4]
\Code{1.3}{for (v in adj[s]) if (!seen[v] \&\& Reach(v, t)) return true;}
\Code{1.3}{return false;}
\Code{0.0}{\}}
\nopagebreak[4]
\Cost{\ensuremath{O\!\left(V+E\right)}}{\ensuremath{O\!\left(V\right)}}
\par\medskip

\Needspace{14\baselineskip}
\subsubsection{判断无向图是否连通}
\Code{0.0}{Connected(G) \{}
\Hint{1.3}{从任一顶点出发，无向连通图的所有顶点都会被标记。}\nopagebreak[4]
\Code{1.3}{require(V>0);}
\Code{1.3}{Fill(seen, false);}
\Hint{1.3}{非空无向图若有多个分量，任取一个根的一次搜索总会漏掉其他分量。}\nopagebreak[4]
\Code{1.3}{DFS(0);}
\Code{1.3}{for (u in vertices) if (!seen[u]) return false;}
\Code{1.3}{return true;}
\Code{0.0}{\}}
\nopagebreak[4]
\Cost{\ensuremath{O\!\left(V+E\right)}}{\ensuremath{O\!\left(V\right)}}
\par\medskip

\Needspace{17\baselineskip}
\subsubsection{求无向图的连通分量}
\Code{0.0}{Components(G) \{}
\Hint{1.3}{用分量编号本身代替 visited 标记。}\nopagebreak[4]
\Code{1.3}{Fill(id, NONE);}
\Code{1.3}{k = 0;}
\Code{1.3}{for (s in vertices) if (id[s] == NONE) Paint(s, k++);}
\Code{1.3}{return k;}
\Code{0.0}{\}}
\Hint{0.0}{每次顶层 Paint 结束才使用新编号，与它连通的所有点共用本轮编号。}\nopagebreak[4]
\Code{0.0}{Paint(u, c) \{}
\Code{1.3}{id[u] = c;}
\Code{1.3}{for (v in adj[u]) if (id[v] == NONE) Paint(v, c);}
\Code{0.0}{\}}
\nopagebreak[4]
\Cost{\ensuremath{O\!\left(V+E\right)}}{\ensuremath{O\!\left(V\right)}}
\par\medskip

\Needspace{28\baselineskip}
\subsubsection{求有向图的强连通分量（Tarjan算法）}
\Code{0.0}{Tarjan(u) \{}
\Hint{1.3}{dfn 为首次发现时间；只有仍在活动栈中的非树边终点可以降低 low。从每个未发现点调用。}\nopagebreak[4]
\Code{1.3}{dfn[u] = low[u] = ++clock;}
\Code{1.3}{Push(\&S, u);}
\Code{1.3}{on[u] = true;}
\Code{1.3}{for (v in adj[u]) \{}
\Code{2.6}{if (!dfn[v]) \{}
\Code{3.9}{Tarjan(v);}
\Code{3.9}{low[u] = min(low[u], low[v]);}
\Code{2.6}{\}}
\Code{2.6}{else if (on[v]) low[u] = min(low[u], dfn[v]);}
\Code{1.3}{\}}
\Hint{1.3}{low[u] 与 dfn[u] 相等时，弹栈直到 u，刚弹出的一组顶点恰构成一个强连通分量。}\nopagebreak[4]
\Hint{1.3}{没有更早的活动祖先可达，u 是一个分量的弹栈边界。}\nopagebreak[4]
\Code{1.3}{if (low[u] == dfn[u]) \{}
\Code{2.6}{do \{}
\Code{3.9}{x = Pop(\&S);}
\Code{3.9}{on[x] = false;}
\Code{3.9}{id[x] = component;}
\Code{2.6}{\}}
\Code{2.6}{while (x != u);}
\Code{2.6}{component++;}
\Code{1.3}{\}}
\Code{0.0}{\}}
\nopagebreak[4]
\Cost{\ensuremath{O\!\left(V+E\right)}}{\ensuremath{O\!\left(V\right)}}
\par\medskip

\Needspace{22\baselineskip}
\subsubsection{求有向图的强连通分量（Kosaraju算法）}
\Code{0.0}{Kosaraju(G) \{}
\Hint{1.3}{第一遍按退出时刻压栈；第二遍在转置图按相反顺序搜索。}\nopagebreak[4]
\Code{1.3}{for (u in vertices) if (!seen[u]) FinishDFS(G, u);}
\Code{1.3}{T = TransposeGraph(G);}
\Code{1.3}{Fill(id, NONE);}
\Code{1.3}{k = 0;}
\Code{1.3}{while (!Empty(\&S)) \{}
\Code{2.6}{u = Pop(\&S);}
\Code{2.6}{if (id[u] == NONE) PaintOn(T, u, k++);}
\Code{1.3}{\}}
\Code{0.0}{\}}
\Hint{0.0}{逆序取点用的是第一遍完成时刻，而不是被第一次发现的时刻。}\nopagebreak[4]
\Code{0.0}{FinishDFS(G, u) \{}
\Code{1.3}{seen[u] = true;}
\Code{1.3}{for (v in adj[u]) if (!seen[v]) FinishDFS(G, v);}
\Code{1.3}{Push(\&S, u);}
\Code{0.0}{\}}
\nopagebreak[4]
\Cost{\ensuremath{O\!\left(V+E\right)}}{\ensuremath{O\!\left(V+E\right)}}
\par\medskip

\Needspace{16\baselineskip}
\subsubsection{判断有向图是否存在环}
\Code{0.0}{HasDirectedCycle(G) \{}
\Hint{1.3}{所有顶点初始为白色，从每个仍为白色的顶点启动搜索。}\nopagebreak[4]
\Code{1.3}{Fill(color, WHITE);}
\Code{1.3}{for (u in vertices of G) \{}
\Code{2.6}{if (color[u] == WHITE \&\& DirectedCycle(G, u, color)) return true;}
\Code{1.3}{\}}
\Code{1.3}{return false;}
\Code{0.0}{\}}
\Code{0.0}{DirectedCycle(G, u, color) \{}
\Hint{1.3}{灰色点位于当前递归路径，指向灰点的边闭合一个有向环。}\nopagebreak[4]
\Code{1.3}{color[u] = GRAY;}
\Code{1.3}{for (e in G.first[u]) \{}
\Code{2.6}{v = e.to;}
\Code{2.6}{if (color[v] == GRAY) return true;}
\Code{2.6}{if (color[v] == WHITE \&\& DirectedCycle(G, v, color)) return true;}
\Code{1.3}{\}}
\Hint{1.3}{退出 u 时将其染黑，后来再遇黑色点不代表路径绕回当前递归栈。}\nopagebreak[4]
\Code{1.3}{color[u] = BLACK;}
\Code{1.3}{return false;}
\Code{0.0}{\}}
\nopagebreak[4]
\Cost{\ensuremath{O\!\left(V+E\right)}}{\ensuremath{O\!\left(V\right)}}
\par\medskip

\Needspace{16\baselineskip}
\subsubsection{判断无向图是否存在环}
\Code{0.0}{UndirectedCycle(u, parentEdge) \{}
\Hint{1.3}{只跳过同一条父边，另一条平行边仍须检查。对各未访根调用。}\nopagebreak[4]
\Code{1.3}{seen[u] = true;}
\Code{1.3}{for (e in adj[u]) \{}
\Code{2.6}{if (e.id == parentEdge) continue;}
\Code{2.6}{v = e.to;}
\Hint{2.6}{无向边反向回到父亲很正常，只有参与 DFS 的同一条父边需要跳过。}\nopagebreak[4]
\Code{2.6}{if (seen[v] || UndirectedCycle(v, e.id)) return true;}
\Code{1.3}{\}}
\Code{1.3}{return false;}
\Code{0.0}{\}}
\nopagebreak[4]
\Cost{\ensuremath{O\!\left(V+E\right)}}{\ensuremath{O\!\left(V\right)}}
\par\medskip

\Needspace{11\baselineskip}
\subsubsection{判断无向图是否为一棵树}
\Code{0.0}{IsTree(G) \{}
\Hint{1.3}{无向图同时满足连通且边数为顶点数减一即可。}\nopagebreak[4]
\Code{1.3}{if (V == 0 || E != V-1) return false;}
\Hint{1.3}{只数边不足以判断树，图可能同时有一个环和一个孤立点。}\nopagebreak[4]
\Code{1.3}{return Connected(G);}
\Code{0.0}{\}}
\nopagebreak[4]
\Cost{\ensuremath{O\!\left(V+E\right)}}{\ensuremath{O\!\left(V\right)}}
\par\medskip

\Needspace{26\baselineskip}
\subsubsection{判断一个图是否为二分图}
\Code{0.0}{Bipartite(G) \{}
\Hint{1.3}{新邻点染相反颜色，已染邻点用于检验冲突。}\nopagebreak[4]
\Code{1.3}{Fill(color, -1);}
\Code{1.3}{for (s in vertices) \{}
\Code{2.6}{if (color[s] != -1) continue;}
\Code{2.6}{color[s] = 0;}
\Code{2.6}{EnQueue(\&Q, s);}
\Code{2.6}{while (!Empty(\&Q)) \{}
\Code{3.9}{u = DeQueue(\&Q);}
\Code{3.9}{for (v in adj[u]) \{}
\Code{5.2}{if (color[v] == color[u]) return false;}
\Hint{5.2}{无色分量重新从颜色 0 开始，不要求整张图本来连通。}\nopagebreak[4]
\Code{5.2}{if (color[v]<0) \{}
\Code{6.5}{color[v] = 1-color[u];}
\Code{6.5}{EnQueue(\&Q, v);}
\Code{5.2}{\}}
\Code{3.9}{\}}
\Code{2.6}{\}}
\Code{1.3}{\}}
\Code{1.3}{return true;}
\Code{0.0}{\}}
\nopagebreak[4]
\Cost{\ensuremath{O\!\left(V+E\right)}}{\ensuremath{O\!\left(V\right)}}
\par\medskip

\Needspace{15\baselineskip}
\subsubsection{求两个顶点之间的所有简单路径}
\Code{0.0}{Paths(u, t) \{}
\Hint{1.3}{used 只表示本条路径占用，回溯后其他路径还可以经过 u。}\nopagebreak[4]
\Code{1.3}{used[u] = true;}
\Code{1.3}{Push(\&path, u);}
\Code{1.3}{if (u == t) Emit(path);}
\Hint{1.3}{到达 t 时输出 path 的当前副本，若继续深入可能错误地重复经过终点。}\nopagebreak[4]
\Code{1.3}{else for (v in adj[u]) if (!used[v]) Paths(v, t);}
\Code{1.3}{Pop(\&path);}
\Code{1.3}{used[u] = false;}
\Code{0.0}{\}}
\nopagebreak[4]
\Cost{\ensuremath{O\!\left(n!\,n\right)}（简单图上界）}{\ensuremath{O\!\left(V\right)}}
\par\medskip

\Needspace{19\baselineskip}
\subsubsection{求两个顶点之间长度为k的简单路径}
\Code{0.0}{PathsK(u, t, k) \{}
\Hint{1.3}{k 表示还需走几条边。到达目标过早也不能重复进入目标。}\nopagebreak[4]
\Code{1.3}{used[u] = true;}
\Code{1.3}{Push(\&path, u);}
\Code{1.3}{if (k == 0) \{}
\Code{2.6}{if (u == t) Emit(path);}
\Code{1.3}{\}}
\Hint{1.3}{第一次调用从 s 出发，k 是目标边数；一条 k 边简单路径包含 k+1 个不同顶点。}\nopagebreak[4]
\Code{1.3}{else if (u != t) \{}
\Code{2.6}{for (v in adj[u]) if (!used[v]) PathsK(v, t, k-1);}
\Code{1.3}{\}}
\Code{1.3}{Pop(\&path);}
\Code{1.3}{used[u] = false;}
\Code{0.0}{\}}
\nopagebreak[4]
\Cost{\ensuremath{O\!\left((\Delta +1)^{k+1}+L\right)}}{\ensuremath{O\!\left(V+k\right)}}
\par\medskip

\Needspace{17\baselineskip}
\subsubsection{求两个顶点之间长度为k的路径数量}
\Code{0.0}{WalkCount(s, t, k) \{}
\Hint{1.3}{每轮从上一层转移，next 必须清零；此处允许重复经过顶点。}\nopagebreak[4]
\Code{1.3}{Zero(dp);}
\Code{1.3}{dp[s] = 1;}
\Code{1.3}{repeat (k) \{}
\Code{2.6}{Zero(next);}
\Hint{2.6}{每条平行弧都会贡献一次转移。如果题目要求简单路径，应使用前面的回溯版。}\nopagebreak[4]
\Code{2.6}{for ((u, v) in arcs) next[v] += dp[u];}
\Code{2.6}{SwapBuffers(\&dp, \&next);}
\Code{1.3}{\}}
\Code{1.3}{return dp[t];}
\Code{0.0}{\}}
\nopagebreak[4]
\Cost{\ensuremath{O\!\left(k(V+E)+V\right)}}{\ensuremath{O\!\left(V\right)}}
\par\medskip

\Needspace{16\baselineskip}
\subsubsection{利用邻接矩阵的幂计算定长游走数量}
\Code{0.0}{WalkPower(A, k, s, t) \{}
\Hint{1.3}{邻接矩阵的乘法按中间顶点汇总方案，单位矩阵对应零步。}\nopagebreak[4]
\Code{1.3}{R = Identity(V);}
\Code{1.3}{while (k>0) \{}
\Code{2.6}{if (k\%2) R = MatrixMul(R, A);}
\Code{2.6}{A = MatrixMul(A, A);}
\Hint{2.6}{乘法按整数算术进行，不要把加法替换为布尔或，否则只能判断有无路线。}\nopagebreak[4]
\Code{2.6}{k /= 2;}
\Code{1.3}{\}}
\Code{1.3}{return R[s][t];}
\Code{0.0}{\}}
\nopagebreak[4]
\Cost{\ensuremath{O\!\left(V^{3} \log (k+2)\right)}}{\ensuremath{O\!\left(V^{2}\right)}}
\par\medskip

\Needspace{20\baselineskip}
\subsubsection{求图中顶点的可达集合}
\Code{0.0}{ReachSet(s) \{}
\Hint{1.3}{零步可达使起点本身也在结果中。}\nopagebreak[4]
\Code{1.3}{Fill(seen, false);}
\Code{1.3}{seen[s] = true;}
\Code{1.3}{Push(\&S, s);}
\Code{1.3}{while (!Empty(\&S)) \{}
\Code{2.6}{u = Pop(\&S);}
\Code{2.6}{for (v in adj[u]) if (!seen[v]) \{}
\Hint{3.9}{若只要布尔访问结果，则无需存前驱或完整路径。}\nopagebreak[4]
\Code{3.9}{seen[v] = true;}
\Code{3.9}{Push(\&S, v);}
\Code{2.6}{\}}
\Code{1.3}{\}}
\Code{1.3}{return seen;}
\Code{0.0}{\}}
\nopagebreak[4]
\Cost{\ensuremath{O\!\left(V+E\right)}}{\ensuremath{O\!\left(V\right)}}
\par\medskip

\Needspace{16\baselineskip}
\subsubsection{求图的传递闭包（Warshall算法）}
\Code{0.0}{Warshall(R, V) \{}
\Hint{1.3}{k 必须是最外层，每一轮允许多使用一个中间点。}\nopagebreak[4]
\Code{1.3}{for (i = 0; i<V; i++) R[i][i] = true;}
\Code{1.3}{for (k = 0; k<V; k++) \{}
\Code{2.6}{for (i = 0; i<V; i++) \{}
\Code{3.9}{if (!R[i][k]) continue;}
\Hint{3.9}{令 R[i][i] 为真表示把零长路也算作可达；若题意排除零长路可省略这一步。}\nopagebreak[4]
\Code{3.9}{for (j = 0; j<V; j++) R[i][j] = R[i][j] || R[k][j];}
\Code{2.6}{\}}
\Code{1.3}{\}}
\Code{0.0}{\}}
\nopagebreak[4]
\Cost{\ensuremath{O\!\left(V^{3}\right)}}{\ensuremath{O\!\left(1\right)}（原地更新结果）}
\par\medskip

\Needspace{24\baselineskip}
\subsubsection{求割点与桥（Tarjan算法）}
\Code{0.0}{CutDFS(u, parentEdge) \{}
\Hint{1.3}{桥用严格大于；非根割点允许等于。根单独按 DFS 孩子数判断。}\nopagebreak[4]
\Code{1.3}{dfn[u] = low[u] = ++clock;}
\Code{1.3}{kids = 0;}
\Code{1.3}{for (e in adj[u]) \{}
\Code{2.6}{if (e.id == parentEdge) continue;}
\Code{2.6}{v = e.to;}
\Code{2.6}{if (!dfn[v]) \{}
\Code{3.9}{kids++;}
\Code{3.9}{CutDFS(v, e.id);}
\Hint{3.9}{low[v] 表示 v 的 DFS 子树能回到的最早发现点，检查应放在对子树递归完成之后。}\nopagebreak[4]
\Code{3.9}{low[u] = min(low[u], low[v]);}
\Code{3.9}{if (low[v]>dfn[u]) EmitBridge(e.id);}
\Code{3.9}{if (parentEdge != NONE \&\& low[v] >= dfn[u]) cut[u] = true;}
\Code{2.6}{\}}
\Code{2.6}{else low[u] = min(low[u], dfn[v]);}
\Code{1.3}{\}}
\Code{1.3}{if (parentEdge == NONE \&\& kids >= 2) cut[u] = true;}
\Code{0.0}{\}}
\nopagebreak[4]
\Cost{\ensuremath{O\!\left(V+E\right)}}{\ensuremath{O\!\left(V\right)}}
\par\medskip

\Group{并查集}

\Needspace{14\baselineskip}
\subsubsection{并查集的初始化}
\Code{0.0}{InitDSU(n, byRank) \{}
\Code{1.3}{useRank = byRank;}
\Hint{1.3}{初始每个点既是根也是单元素集合。}\nopagebreak[4]
\Code{1.3}{for (i = 0; i<n; i++) \{}
\Code{2.6}{parent[i] = i;}
\Hint{2.6}{parent 的初值必须是自己的编号，不能把所有点的 parent 都设为 0。}\nopagebreak[4]
\Code{2.6}{size[i] = 1;}
\Code{2.6}{rank[i] = 0;}
\Code{1.3}{\}}
\Code{0.0}{\}}
\nopagebreak[4]
\Cost{\ensuremath{O\!\left(n\right)}}{\ensuremath{O\!\left(n\right)}（结构）}
\par\medskip

\Needspace{11\baselineskip}
\subsubsection{查找集合的代表元素（Find）}
\Code{0.0}{Find(x) \{}
\Hint{1.3}{只有根结点的父指针指向自己。}\nopagebreak[4]
\Code{1.3}{while (parent[x] != x) x = parent[x];}
\Hint{1.3}{查找只沿父边向上，不需要遍历集合里的全部成员。}\nopagebreak[4]
\Code{1.3}{return x;}
\Code{0.0}{\}}
\nopagebreak[4]
\Cost{\ensuremath{O\!\left(h\right)}}{\ensuremath{O\!\left(1\right)}}
\par\medskip

\Needspace{15\baselineskip}
\subsubsection{合并两个集合（Union）}
\Code{0.0}{Union(x, y) \{}
\Hint{1.3}{代表相同表示已经连通，不再改变树。}\nopagebreak[4]
\Code{1.3}{a = Find(x);}
\Code{1.3}{b = Find(y);}
\Code{1.3}{if (a == b) return false;}
\Hint{1.3}{基本 Union 直接挂接。使用优化版时，从初始化起通过 WeightedUnion 合并。}\nopagebreak[4]
\Code{1.3}{parent[b] = a;}
\Code{1.3}{size[a] += size[b];}
\Code{1.3}{return true;}
\Code{0.0}{\}}
\nopagebreak[4]
\Cost{\ensuremath{O\!\left(h\right)}}{\ensuremath{O\!\left(1\right)}}
\par\medskip

\Needspace{16\baselineskip}
\subsubsection{路径压缩}
\Code{0.0}{Compress(x) \{}
\Hint{1.3}{先定位根，再沿原链回走一遍；改父指针前保存旧值。}\nopagebreak[4]
\Code{1.3}{r = Find(x);}
\Code{1.3}{while (x != r) \{}
\Code{2.6}{next = parent[x];}
\Code{2.6}{parent[x] = r;}
\Hint{2.6}{第二遍结束后，路径上的点均直接指向根，方便以后再查。}\nopagebreak[4]
\Code{2.6}{x = next;}
\Code{1.3}{\}}
\Code{1.3}{return r;}
\Code{0.0}{\}}
\nopagebreak[4]
\Cost{单次 \ensuremath{O\!\left(h\right)}，配合按秩合并均摊 \ensuremath{O\!\left(\alpha (n)\right)}}{\ensuremath{O\!\left(1\right)}}
\par\medskip

\Needspace{23\baselineskip}
\subsubsection{按秩合并与按大小合并}
\Code{0.0}{WeightedUnion(x, y) \{}
\Hint{1.3}{初始化时固定合并策略，并全程使用本入口。按秩只在同秩时增秩；按大小累加结点数。}\nopagebreak[4]
\Code{1.3}{a = Compress(x);}
\Code{1.3}{b = Compress(y);}
\Code{1.3}{if (a == b) return false;}
\Code{1.3}{if (useRank) \{}
\Code{2.6}{if (rank[a]<rank[b]) swap(a, b);}
\Code{2.6}{parent[b] = a;}
\Code{2.6}{if (rank[a] == rank[b]) rank[a]++;}
\Code{1.3}{\}}
\Hint{1.3}{如果本来已经同根，就不能再增加大小或秩，否则会重复统计。}\nopagebreak[4]
\Code{1.3}{else \{}
\Code{2.6}{if (size[a]<size[b]) swap(a, b);}
\Code{2.6}{parent[b] = a;}
\Code{2.6}{size[a] += size[b];}
\Code{1.3}{\}}
\Code{1.3}{return true;}
\Code{0.0}{\}}
\nopagebreak[4]
\Cost{与压缩结合均摊 \ensuremath{O\!\left(\alpha (n)\right)}}{\ensuremath{O\!\left(1\right)}}
\par\medskip

\Needspace{13\baselineskip}
\subsubsection{利用并查集判断图的连通性}
\Code{0.0}{DSUConnected(G) \{}
\Hint{1.3}{只有一次成功合并才减少分量数。}\nopagebreak[4]
\Code{1.3}{InitDSU(V, false);}
\Code{1.3}{groups = V;}
\Hint{1.3}{若只问两点是否连通，处理完全部边后直接比较两点代表即可。}\nopagebreak[4]
\Code{1.3}{for ((u, v) in edges) if (WeightedUnion(u, v)) groups--;}
\Code{1.3}{return groups == 1;}
\Code{0.0}{\}}
\nopagebreak[4]
\Cost{\ensuremath{O\!\left(V+E\,\alpha (V)\right)}}{\ensuremath{O\!\left(V\right)}}
\par\medskip

\Needspace{14\baselineskip}
\subsubsection{利用并查集判断无向图是否存在环}
\Code{0.0}{DSUCycle(G) \{}
\Hint{1.3}{若一条物理边两端已经连通，这条边会闭合一条环。}\nopagebreak[4]
\Code{1.3}{InitDSU(V, false);}
\Code{1.3}{for ((u, v) in edges) \{}
\Hint{2.6}{每条无向物理边只处理一次，否则第二次处理其反向记录就会误报环。}\nopagebreak[4]
\Code{2.6}{if (!WeightedUnion(u, v)) return true;}
\Code{1.3}{\}}
\Code{1.3}{return false;}
\Code{0.0}{\}}
\nopagebreak[4]
\Cost{\ensuremath{O\!\left(V+E\,\alpha (V)\right)}}{\ensuremath{O\!\left(V\right)}}
\par\medskip

\Group{最小生成树}

\Needspace{31\baselineskip}
\subsubsection{Prim算法}
\Code{0.0}{Prim(G, s) \{}
\Hint{1.3}{key[v] 记录从已选顶点集合连接到 v 的最轻边；parent 记录这条边的另一端。}\nopagebreak[4]
\Code{1.3}{Fill(key, INF);}
\Code{1.3}{Fill(used, false);}
\Code{1.3}{key[s] = 0;}
\Code{1.3}{parent[s] = NONE;}
\Code{1.3}{repeat (V) \{}
\Code{2.6}{u = MinUnused(key, used);}
\Code{2.6}{if (key[u] == INF) return DISCONNECTED;}
\Code{2.6}{used[u] = true;}
\Code{2.6}{if (parent[u] != NONE) Emit(parent[u], u);}
\Code{2.6}{for (e in adj[u]) \{}
\Code{3.9}{v = e.to;}
\Code{3.9}{if (!used[v] \&\& e.w<key[v]) \{}
\Hint{5.2}{key 初值正无穷表示树内尚无边接到它，选到这样的点意味着图不连通。}\nopagebreak[4]
\Code{5.2}{key[v] = e.w;}
\Code{5.2}{parent[v] = u;}
\Code{3.9}{\}}
\Code{2.6}{\}}
\Code{1.3}{\}}
\Code{0.0}{\}}
\Code{0.0}{MinUnused(key, used) \{}
\Code{1.3}{best = NONE;}
\Code{1.3}{for (v in vertices) if (!used[v] \&\& (best == NONE || key[v]<key[best])) best = v;}
\Code{1.3}{return best;}
\Code{0.0}{\}}
\nopagebreak[4]
\Cost{\ensuremath{O\!\left(V^{2}+E\right)}}{\ensuremath{O\!\left(V\right)}}
\par\medskip

\Needspace{25\baselineskip}
\subsubsection{Prim算法的邻接矩阵实现}
\Code{0.0}{PrimMatrix(M, s) \{}
\Hint{1.3}{线性选最轻候选点，再扫描该点的整行来更新候选边。}\nopagebreak[4]
\Code{1.3}{Fill(key, INF);}
\Code{1.3}{Fill(used, false);}
\Code{1.3}{key[s] = 0;}
\Code{1.3}{parent[s] = NONE;}
\Code{1.3}{repeat (V) \{}
\Code{2.6}{u = MinUnused(key, used);}
\Code{2.6}{if (key[u] == INF) return DISCONNECTED;}
\Code{2.6}{used[u] = true;}
\Code{2.6}{for (v = 0; v<V; v++) \{}
\Hint{3.9}{parent 最终记录的是整棵树的边，可从每个非根顶点读到它的树内父亲。}\nopagebreak[4]
\Code{3.9}{if (!used[v] \&\& M[u][v] != NOEDGE \&\& M[u][v]<key[v]) \{}
\Code{5.2}{key[v] = M[u][v];}
\Code{5.2}{parent[v] = u;}
\Code{3.9}{\}}
\Code{2.6}{\}}
\Code{1.3}{\}}
\Code{1.3}{return parent;}
\Code{0.0}{\}}
\nopagebreak[4]
\Cost{\ensuremath{O\!\left(V^{2}\right)}}{\ensuremath{O\!\left(V\right)}}
\par\medskip

\Needspace{26\baselineskip}
\subsubsection{Prim算法的优先队列优化}
\Code{0.0}{PrimHeap(G, s) \{}
\Hint{1.3}{堆中每个未选顶点只有一项，降低键值时调整这一项的位置。}\nopagebreak[4]
\Code{1.3}{Fill(key, INF);}
\Code{1.3}{key[s] = 0;}
\Code{1.3}{parent[s] = NONE;}
\Code{1.3}{BuildIndexedHeap(\&H, key);}
\Code{1.3}{while (!Empty(\&H)) \{}
\Code{2.6}{u = ExtractMin(\&H);}
\Code{2.6}{if (key[u] == INF) return DISCONNECTED;}
\Code{2.6}{for (e in adj[u]) \{}
\Code{3.9}{v = e.to;}
\Hint{3.9}{取出顶点后不会重新入堆，更新时只考虑还保留在堆里的端点。}\nopagebreak[4]
\Code{3.9}{if (InHeap(\&H, v) \&\& e.w<key[v]) \{}
\Code{5.2}{key[v] = e.w;}
\Code{5.2}{parent[v] = u;}
\Code{5.2}{DecreaseKey(\&H, v, e.w);}
\Code{3.9}{\}}
\Code{2.6}{\}}
\Code{1.3}{\}}
\Code{1.3}{return parent;}
\Code{0.0}{\}}
\nopagebreak[4]
\Cost{\ensuremath{O\!\left((V+E) \log (V+1)\right)}}{\ensuremath{O\!\left(V\right)}}
\par\medskip

\Needspace{22\baselineskip}
\subsubsection{Kruskal算法}
\Code{0.0}{KruskalLabels(G) \{}
\Hint{1.3}{标签相同说明已有路径，再加边会形成环；每次成功选边后统一旧标签。}\nopagebreak[4]
\Code{1.3}{SortEdgesByWeight(edges);}
\Code{1.3}{for (u in vertices) id[u] = u;}
\Code{1.3}{k = 0;}
\Code{1.3}{for (e in edges) \{}
\Code{2.6}{a = id[e.u];}
\Code{2.6}{b = id[e.v];}
\Code{2.6}{if (a == b) continue;}
\Hint{2.6}{保存旧编号 b 再扫描标签，否则扫描中改值可能改变判断参照。}\nopagebreak[4]
\Code{2.6}{Emit(e);}
\Code{2.6}{k++;}
\Code{2.6}{for (v in vertices) if (id[v] == b) id[v] = a;}
\Code{2.6}{if (k == V-1) break;}
\Code{1.3}{\}}
\Code{1.3}{return k == V-1;}
\Code{0.0}{\}}
\nopagebreak[4]
\Cost{\ensuremath{O\!\left(E \log (E+1)+VE+V\right)}}{\ensuremath{O\!\left(V\right)}}
\par\medskip

\Needspace{17\baselineskip}
\subsubsection{基于并查集的Kruskal算法}
\Code{0.0}{KruskalDSU(G) \{}
\Hint{1.3}{并查集合并成功才把边放入树，合并失败的边两端已经连通。}\nopagebreak[4]
\Code{1.3}{SortEdgesByWeight(edges);}
\Code{1.3}{InitDSU(V, false);}
\Code{1.3}{tree = EmptyList();}
\Code{1.3}{for (e in edges) \{}
\Hint{2.6}{有负权边不影响此算法，边只按权重排序，是否选入取决于有无成环。}\nopagebreak[4]
\Code{2.6}{if (WeightedUnion(e.u, e.v)) Append(\&tree, e);}
\Code{2.6}{if (Size(\&tree) == V-1) break;}
\Code{1.3}{\}}
\Code{1.3}{return Size(\&tree) == V-1 ? tree:FAIL;}
\Code{0.0}{\}}
\nopagebreak[4]
\Cost{\ensuremath{O\!\left(E \log (E+1)+V+E\,\alpha (V)\right)}}{\ensuremath{O\!\left(V\right)}}
\par\medskip

\Needspace{23\baselineskip}
\subsubsection{最小生成树的唯一性判断}
\Code{0.0}{UniqueMST(G) \{}
\Hint{1.3}{等权的一批边要先一起检查。收缩低权边后若出现环，就能替换一条等权树边。}\nopagebreak[4]
\Code{1.3}{require(Connected(G));}
\Code{1.3}{SortEdgesByWeight(edges);}
\Code{1.3}{InitDSU(V, false);}
\Code{1.3}{for (batch in EqualWeightGroups(edges)) \{}
\Code{2.6}{InitSparseDSU(\&tmp);}
\Code{2.6}{for (e in batch) \{}
\Code{3.9}{a = Compress(e.u);}
\Code{3.9}{b = Compress(e.v);}
\Hint{3.9}{临时集合必须每个权值组重新建立；主集合要等该组检查结束再统一合并。}\nopagebreak[4]
\Code{3.9}{if (a == b) continue;}
\Code{3.9}{if (!UnionIn(\&tmp, a, b)) return false;}
\Code{2.6}{\}}
\Code{2.6}{for (e in batch) WeightedUnion(e.u, e.v);}
\Code{1.3}{\}}
\Code{1.3}{return true;}
\Code{0.0}{\}}
\nopagebreak[4]
\Cost{\ensuremath{O\!\left(E \log (E+1)+V+E\,\alpha (V)\right)}}{\ensuremath{O\!\left(V+E\right)}}
\par\medskip

\Needspace{14\baselineskip}
\subsubsection{最小生成树的总权重计算}
\Code{0.0}{MSTWeight(G) \{}
\Hint{1.3}{只有真正进入生成树的边才累计权重。}\nopagebreak[4]
\Code{1.3}{T = KruskalDSU(G);}
\Code{1.3}{if (T == FAIL) return DISCONNECTED;}
\Hint{1.3}{若选不到 V−1 条边，得到的是森林，不能当作整个图的生成树总重。}\nopagebreak[4]
\Code{1.3}{sum = 0;}
\Code{1.3}{for (e in T) sum += e.w;}
\Code{1.3}{return sum;}
\Code{0.0}{\}}
\nopagebreak[4]
\Cost{\ensuremath{O\!\left(E \log (E+1)+V+E\,\alpha (V)\right)}}{\ensuremath{O\!\left(V\right)}}
\par\medskip

\Group{最短路径}

\Needspace{22\baselineskip}
\subsubsection{无权图的最短路径（BFS）}
\Code{0.0}{Unweighted(G, s) \{}
\Hint{1.3}{首次到达的层数就是最少边数；前驱用来恢复具体走法。}\nopagebreak[4]
\Code{1.3}{Fill(d, INF);}
\Code{1.3}{d[s] = 0;}
\Code{1.3}{prev[s] = NONE;}
\Code{1.3}{EnQueue(\&Q, s);}
\Code{1.3}{while (!Empty(\&Q)) \{}
\Code{2.6}{u = DeQueue(\&Q);}
\Code{2.6}{for (v in adj[u]) if (d[v] == INF) \{}
\Hint{3.9}{对于无权图一次出边增加一步，第一次设置 d[v] 后不需要反复松弛。}\nopagebreak[4]
\Code{3.9}{d[v] = d[u]+1;}
\Code{3.9}{prev[v] = u;}
\Code{3.9}{EnQueue(\&Q, v);}
\Code{2.6}{\}}
\Code{1.3}{\}}
\Code{1.3}{return d;}
\Code{0.0}{\}}
\nopagebreak[4]
\Cost{\ensuremath{O\!\left(V+E\right)}}{\ensuremath{O\!\left(V\right)}}
\par\medskip

\Needspace{26\baselineskip}
\subsubsection{单源最短路径（Dijkstra算法）}
\Code{0.0}{Dijkstra(G, s) \{}
\Hint{1.3}{仅用于非负权边。确定 u 后尝试通过 u 缩短邻点的暂定距离。}\nopagebreak[4]
\Code{1.3}{require(AllWeightsNonnegative(G));}
\Code{1.3}{Fill(d, INF);}
\Code{1.3}{Fill(done, false);}
\Code{1.3}{d[s] = 0;}
\Code{1.3}{repeat (V) \{}
\Code{2.6}{u = MinUnused(d, done);}
\Code{2.6}{if (d[u] == INF) break;}
\Code{2.6}{done[u] = true;}
\Code{2.6}{for (e in adj[u]) \{}
\Hint{3.9}{d[u]+w 是经过当前 u 的候选总距离，不是这条边本身的权重。}\nopagebreak[4]
\Code{3.9}{v = e.to;}
\Code{3.9}{if (d[u]+e.w<d[v]) \{}
\Code{5.2}{d[v] = d[u]+e.w;}
\Code{5.2}{prev[v] = u;}
\Code{3.9}{\}}
\Code{2.6}{\}}
\Code{1.3}{\}}
\Code{1.3}{return d;}
\Code{0.0}{\}}
\nopagebreak[4]
\Cost{\ensuremath{O\!\left(V^{2}+E\right)}}{\ensuremath{O\!\left(V\right)}}
\par\medskip

\Needspace{24\baselineskip}
\subsubsection{Dijkstra算法的邻接矩阵实现}
\Code{0.0}{DijkstraMatrix(M, s) \{}
\Hint{1.3}{无边标记和零权边需要区分。不可达点的距离一直保持正无穷。}\nopagebreak[4]
\Code{1.3}{Fill(d, INF);}
\Code{1.3}{Fill(done, false);}
\Code{1.3}{d[s] = 0;}
\Code{1.3}{repeat (V) \{}
\Code{2.6}{u = MinUnused(d, done);}
\Code{2.6}{if (d[u] == INF) break;}
\Code{2.6}{done[u] = true;}
\Code{2.6}{for (v = 0; v<V; v++) \{}
\Hint{3.9}{标为 done 的是已经确定最短距离的点；遇到暂定距离为无穷的点应提前退出。}\nopagebreak[4]
\Code{3.9}{if (M[u][v] != NOEDGE \&\& d[u]+M[u][v]<d[v]) \{}
\Code{5.2}{d[v] = d[u]+M[u][v];}
\Code{5.2}{prev[v] = u;}
\Code{3.9}{\}}
\Code{2.6}{\}}
\Code{1.3}{\}}
\Code{1.3}{return d;}
\Code{0.0}{\}}
\nopagebreak[4]
\Cost{\ensuremath{O\!\left(V^{2}\right)}}{\ensuremath{O\!\left(V\right)}}
\par\medskip

\Needspace{27\baselineskip}
\subsubsection{Dijkstra算法的优先队列优化}
\Code{0.0}{DijkstraHeap(G, s) \{}
\Hint{1.3}{索引堆中每点至多一项；距离变小时同步更新堆键。边权须非负。}\nopagebreak[4]
\Code{1.3}{require(AllWeightsNonnegative(G));}
\Code{1.3}{Fill(d, INF);}
\Code{1.3}{d[s] = 0;}
\Code{1.3}{BuildIndexedHeap(\&H, d);}
\Code{1.3}{while (!Empty(\&H)) \{}
\Code{2.6}{u = ExtractMin(\&H);}
\Code{2.6}{if (d[u] == INF) break;}
\Code{2.6}{for (e in adj[u]) \{}
\Code{3.9}{v = e.to;}
\Code{3.9}{cand = d[u]+e.w;}
\Hint{3.9}{这里用能按顶点定位的堆，DecreaseKey 直接调整原记录。}\nopagebreak[4]
\Code{3.9}{if (cand<d[v]) \{}
\Code{5.2}{prev[v] = u;}
\Code{5.2}{d[v] = cand;}
\Code{5.2}{DecreaseKey(\&H, v, cand);}
\Code{3.9}{\}}
\Code{2.6}{\}}
\Code{1.3}{\}}
\Code{1.3}{return d;}
\Code{0.0}{\}}
\nopagebreak[4]
\Cost{\ensuremath{O\!\left((V+E) \log (V+1)\right)}}{\ensuremath{O\!\left(V\right)}}
\par\medskip

\Needspace{24\baselineskip}
\subsubsection{单源最短路径（Bellman-Ford算法）}
\Code{0.0}{BellmanFord(G, s) \{}
\Hint{1.3}{无负环时最短简单路至多经过 V−1 条边。再能改进就存在源可达负环。}\nopagebreak[4]
\Code{1.3}{Fill(d, INF);}
\Code{1.3}{d[s] = 0;}
\Code{1.3}{repeat (V-1) \{}
\Code{2.6}{changed = false;}
\Code{2.6}{for (e in arcs) \{}
\Code{3.9}{if (d[e.u] != INF \&\& d[e.u]+e.w<d[e.v]) \{}
\Code{5.2}{d[e.v] = d[e.u]+e.w;}
\Code{5.2}{prev[e.v] = e.u;}
\Hint{5.2}{每轮只从可达点松弛，否则把无穷当普通整数相加可能产生错误距离。}\nopagebreak[4]
\Code{5.2}{changed = true;}
\Code{3.9}{\}}
\Code{2.6}{\}}
\Code{2.6}{if (!changed) break;}
\Code{1.3}{\}}
\Code{1.3}{for (e in arcs) if (d[e.u] != INF \&\& d[e.u]+e.w<d[e.v]) return NEGATIVE\_CYCLE;}
\Code{1.3}{return d;}
\Code{0.0}{\}}
\nopagebreak[4]
\Cost{\ensuremath{O\!\left(VE+V\right)}}{\ensuremath{O\!\left(V\right)}}
\par\medskip

\Needspace{31\baselineskip}
\subsubsection{Bellman-Ford算法的队列优化（SPFA）}
\Code{0.0}{SPFA(G, s) \{}
\Hint{1.3}{队列只放距离有改进的点；inQ 避免重复入队，cnt 记录当前路径边数。}\nopagebreak[4]
\Code{1.3}{Fill(d, INF);}
\Code{1.3}{Fill(inQ, false);}
\Code{1.3}{d[s] = 0;}
\Code{1.3}{cnt[s] = 0;}
\Code{1.3}{EnQueue(\&Q, s);}
\Code{1.3}{inQ[s] = true;}
\Code{1.3}{while (!Empty(\&Q)) \{}
\Code{2.6}{u = DeQueue(\&Q);}
\Code{2.6}{inQ[u] = false;}
\Code{2.6}{for (e in adj[u]) \{}
\Code{3.9}{v = e.to;}
\Code{3.9}{if (d[u]+e.w >= d[v]) continue;}
\Hint{3.9}{同一个点出队后若再次被改善仍可重新入队，因此出队时要清 inQ。}\nopagebreak[4]
\Code{3.9}{d[v] = d[u]+e.w;}
\Code{3.9}{cnt[v] = cnt[u]+1;}
\Code{3.9}{if (cnt[v] >= V) return NEGATIVE\_CYCLE;}
\Code{3.9}{if (!inQ[v]) \{}
\Code{5.2}{EnQueue(\&Q, v);}
\Code{5.2}{inQ[v] = true;}
\Code{3.9}{\}}
\Code{2.6}{\}}
\Code{1.3}{\}}
\Code{1.3}{return d;}
\Code{0.0}{\}}
\nopagebreak[4]
\Cost{\ensuremath{O\!\left(VE+V\right)}}{\ensuremath{O\!\left(V\right)}}
\par\medskip

\Needspace{18\baselineskip}
\subsubsection{多源最短路径（Floyd算法）}
\Code{0.0}{Floyd(G) \{}
\Hint{1.3}{第 k 轮尝试让路径经过 k；若对角线为负，该点可走负环。}\nopagebreak[4]
\Code{1.3}{Fill(D, INF);}
\Code{1.3}{for (i in vertices) D[i][i] = 0;}
\Code{1.3}{for (e in arcs) D[e.u][e.v] = min(D[e.u][e.v], e.w);}
\Code{1.3}{for (k in vertices) \{}
\Code{2.6}{for (i in vertices) for (j in vertices) \{}
\Hint{3.9}{若 D[k][k]<0，能到 k 且能从 k 到目标的点对都没有有限最短距离。}\nopagebreak[4]
\Code{3.9}{if (D[i][k] != INF \&\& D[k][j] != INF) D[i][j] = min(D[i][j], D[i][k]+D[k][j]);}
\Code{2.6}{\}}
\Code{1.3}{\}}
\Code{1.3}{return D;}
\Code{0.0}{\}}
\nopagebreak[4]
\Cost{\ensuremath{O\!\left(V^{3}\right)}}{\ensuremath{O\!\left(V^{2}\right)}（距离表）}
\par\medskip

\Needspace{21\baselineskip}
\subsubsection{有向无环图的最短路径}
\Code{0.0}{DAGShortest(G, s) \{}
\Hint{1.3}{拓扑序保证在展开 u 前，它所有可能的前驱已经处理完。}\nopagebreak[4]
\Code{1.3}{order = Topo(G);}
\Code{1.3}{require(order != FAIL);}
\Code{1.3}{Fill(d, INF);}
\Code{1.3}{d[s] = 0;}
\Code{1.3}{for (u in order) \{}
\Code{2.6}{if (d[u] == INF) continue;}
\Hint{2.6}{不可达点不展开，负边也不会把这种点误当成一条新路径的起点。}\nopagebreak[4]
\Code{2.6}{for (e in adj[u]) if (d[u]+e.w<d[e.to]) \{}
\Code{3.9}{d[e.to] = d[u]+e.w;}
\Code{3.9}{prev[e.to] = u;}
\Code{2.6}{\}}
\Code{1.3}{\}}
\Code{1.3}{return d;}
\Code{0.0}{\}}
\nopagebreak[4]
\Cost{\ensuremath{O\!\left(V+E\right)}}{\ensuremath{O\!\left(V\right)}}
\par\medskip

\Needspace{17\baselineskip}
\subsubsection{最短路径的记录与恢复}
\Code{0.0}{Restore(prev, s, t) \{}
\Hint{1.3}{每次严格改善距离时记录前驱；从终点逆走后用栈还原正向次序。}\nopagebreak[4]
\Code{1.3}{if (d[t] == INF || affected[t]) return FAIL;}
\Code{1.3}{p = t;}
\Code{1.3}{while (p != s) \{}
\Code{2.6}{Push(\&S, p);}
\Hint{2.6}{s 到自身的路径只含 s，恢复时循环零次即可。}\nopagebreak[4]
\Code{2.6}{p = prev[p];}
\Code{1.3}{\}}
\Code{1.3}{Emit(s);}
\Code{1.3}{while (!Empty(\&S)) Emit(Pop(\&S));}
\Code{0.0}{\}}
\nopagebreak[4]
\Cost{恢复 \ensuremath{O\!\left(V\right)}，记录 \ensuremath{O\!\left(1\right)}/次松弛}{\ensuremath{O\!\left(V\right)}}
\par\medskip

\Needspace{29\baselineskip}
\subsubsection{最短路径的方案数统计}
\Code{0.0}{ShortestCount(G, s) \{}
\Hint{1.3}{严格正权保证最短路前驱先于后继出堆。更短时覆盖计数，等长时追加。}\nopagebreak[4]
\Code{1.3}{Fill(d, INF);}
\Code{1.3}{Zero(ways);}
\Code{1.3}{d[s] = 0;}
\Code{1.3}{ways[s] = 1;}
\Code{1.3}{BuildIndexedHeap(\&H, d);}
\Code{1.3}{while (!Empty(\&H)) \{}
\Code{2.6}{u = ExtractMin(\&H);}
\Code{2.6}{if (d[u] == INF) break;}
\Code{2.6}{for (e in adj[u]) \{}
\Code{3.9}{v = e.to;}
\Code{3.9}{t = d[u]+e.w;}
\Hint{3.9}{当候选比旧距离更短时旧方案全部失效；只有相等时才把两种来源的计数相加。}\nopagebreak[4]
\Code{3.9}{if (t<d[v]) \{}
\Code{5.2}{d[v] = t;}
\Code{5.2}{ways[v] = ways[u];}
\Code{5.2}{DecreaseKey(\&H, v, t);}
\Code{3.9}{\}}
\Code{3.9}{else if (t == d[v]) ways[v] += ways[u];}
\Code{2.6}{\}}
\Code{1.3}{\}}
\Code{1.3}{return ways;}
\Code{0.0}{\}}
\nopagebreak[4]
\Cost{\ensuremath{O\!\left((V+E) \log (V+1)\right)}}{\ensuremath{O\!\left(V\right)}}
\par\medskip

\Needspace{26\baselineskip}
\subsubsection{负权边的处理}
\Code{0.0}{NegativeEdges(G, s) \{}
\Hint{1.3}{从仍能改善的边的终点向前传播，标记所有距离没有有限下界的点。}\nopagebreak[4]
\Code{1.3}{Fill(bad, false);}
\Code{1.3}{if (G.directed \&\& !HasDirectedCycle(G)) \{}
\Code{2.6}{return (DAGShortest(G, s), bad);}
\Code{1.3}{\}}
\Code{1.3}{d = BFRelaxRounds(G, s, V-1);}
\Code{1.3}{for (e in arcs) \{}
\Code{2.6}{if (d[e.u] != INF \&\& d[e.u]+e.w<d[e.v] \&\& !bad[e.v]) \{}
\Code{3.9}{bad[e.v] = true;}
\Code{3.9}{EnQueue(\&Q, e.v);}
\Code{2.6}{\}}
\Code{1.3}{\}}
\Hint{1.3}{bad 为真表示距离没有有限下界，此时不读取 d 中对应值。其余顶点的 d 为最短距离，不可达时为 INF。}\nopagebreak[4]
\Code{1.3}{while (!Empty(\&Q)) \{}
\Code{2.6}{u = DeQueue(\&Q);}
\Code{2.6}{for (v in adj[u]) if (!bad[v]) \{}
\Code{3.9}{bad[v] = true;}
\Code{3.9}{EnQueue(\&Q, v);}
\Code{2.6}{\}}
\Code{1.3}{\}}
\Code{1.3}{return (d, bad);}
\Code{0.0}{\}}
\nopagebreak[4]
\Cost{DAG 为 \ensuremath{O\!\left(V+E\right)}，一般图 \ensuremath{O\!\left(VE+V\right)}}{\ensuremath{O\!\left(V\right)}}
\par\medskip

\Needspace{21\baselineskip}
\subsubsection{负权环的判断}
\Code{0.0}{HasNegCycle(G) \{}
\Hint{1.3}{全图距离初始为 0，相当于从超级源可到所有点。}\nopagebreak[4]
\Code{1.3}{Zero(d);}
\Code{1.3}{for (round = 1; round <= V; round++) \{}
\Code{2.6}{changed = false;}
\Code{2.6}{for (e in arcs) \{}
\Code{3.9}{if (d[e.u]+e.w<d[e.v]) \{}
\Code{5.2}{d[e.v] = d[e.u]+e.w;}
\Hint{5.2}{若只关心从一个指定源可达的负环，其他点初始应设为无穷并跳过。}\nopagebreak[4]
\Code{5.2}{changed = true;}
\Code{3.9}{\}}
\Code{2.6}{\}}
\Code{2.6}{if (!changed) return false;}
\Code{1.3}{\}}
\Code{1.3}{return V>0;}
\Code{0.0}{\}}
\nopagebreak[4]
\Cost{\ensuremath{O\!\left(VE+V\right)}}{\ensuremath{O\!\left(V\right)}}
\par\medskip

\Needspace{24\baselineskip}
\subsubsection{边权仅为0或1的最短路径（0-1 BFS）}
\Code{0.0}{ZeroOneBFS(G, s) \{}
\Hint{1.3}{走 0 权边不增加层数，放队首；走 1 权边放队尾。}\nopagebreak[4]
\Code{1.3}{Fill(d, INF);}
\Code{1.3}{d[s] = 0;}
\Code{1.3}{PushFront(\&Q, (s, 0));}
\Code{1.3}{while (!Empty(\&Q)) \{}
\Code{2.6}{(u, old) = PopFront(\&Q);}
\Code{2.6}{if (old != d[u]) continue;}
\Code{2.6}{for (e in adj[u]) \{}
\Code{3.9}{t = old+e.w;}
\Hint{3.9}{弹出时的 old 保存入队那刻的距离；若后来已有更短方案就跳过这条旧记录。}\nopagebreak[4]
\Code{3.9}{if (t >= d[e.to]) continue;}
\Code{3.9}{d[e.to] = t;}
\Code{3.9}{if (e.w == 0) PushFront(\&Q, (e.to, t));}
\Code{3.9}{else PushBack(\&Q, (e.to, t));}
\Code{2.6}{\}}
\Code{1.3}{\}}
\Code{1.3}{return d;}
\Code{0.0}{\}}
\nopagebreak[4]
\Cost{\ensuremath{O\!\left(V+E\right)}}{\ensuremath{O\!\left(V+E\right)}}
\par\medskip

\Group{拓扑排序}

\Needspace{16\baselineskip}
\subsubsection{AOV网的建立}
\Code{0.0}{BuildAOV(tasks, deps) \{}
\Hint{1.3}{活动是顶点，先修关系 a→b 给 b 的入度加一。}\nopagebreak[4]
\Code{1.3}{G = EmptyGraph(len(tasks));}
\Code{1.3}{Zero(indeg);}
\Code{1.3}{for ((a, b) in deps) \{}
\Code{2.6}{AddArc(\&G, a, b);}
\Hint{2.6}{若有向依赖中出现环，就无法安排符合全部先后约束的顺序。}\nopagebreak[4]
\Code{2.6}{indeg[b]++;}
\Code{1.3}{\}}
\Code{1.3}{return (G, indeg);}
\Code{0.0}{\}}
\nopagebreak[4]
\Cost{\ensuremath{O\!\left(V+E\right)}}{\ensuremath{O\!\left(V\right)}（入度）}
\par\medskip

\Needspace{18\baselineskip}
\subsubsection{基于入度与队列的拓扑排序（Kahn算法）}
\Code{0.0}{Topo(G) \{}
\Hint{1.3}{入度为零表示所有前置活动都已取出，入度副本避免破坏图。}\nopagebreak[4]
\Code{1.3}{degree = CopyIndegree(G);}
\Code{1.3}{for (u in vertices) if (degree[u] == 0) EnQueue(\&Q, u);}
\Code{1.3}{InitList(\&order);}
\Code{1.3}{while (!Empty(\&Q)) \{}
\Code{2.6}{u = DeQueue(\&Q);}
\Hint{2.6}{出队顶点的每条出边只处理一次，因此每个后继入度正好减少其已完成前驱数。}\nopagebreak[4]
\Code{2.6}{Append(\&order, u);}
\Code{2.6}{for (v in adj[u]) if (--degree[v] == 0) EnQueue(\&Q, v);}
\Code{1.3}{\}}
\Code{1.3}{return len(order) == V ? order:FAIL;}
\Code{0.0}{\}}
\nopagebreak[4]
\Cost{\ensuremath{O\!\left(V+E\right)}}{\ensuremath{O\!\left(V\right)}}
\par\medskip

\Needspace{22\baselineskip}
\subsubsection{基于DFS的拓扑排序}
\Code{0.0}{TopoDFS(G) \{}
\Hint{1.3}{回溯退出时记录顶点，逆序即拓扑序；遇灰点报告环。}\nopagebreak[4]
\Code{1.3}{Fill(color, WHITE);}
\Code{1.3}{for (u in vertices) if (color[u] == WHITE \&\& !Finish(u)) return FAIL;}
\Code{1.3}{return Reverse(order);}
\Code{0.0}{\}}
\Code{0.0}{Finish(u) \{}
\Code{1.3}{color[u] = GRAY;}
\Code{1.3}{for (v in adj[u]) \{}
\Hint{2.6}{灰色表示尚未退出的调用；指向灰点构成回边，不能输出拓扑结果。}\nopagebreak[4]
\Code{2.6}{if (color[v] == GRAY) return false;}
\Code{2.6}{if (color[v] == WHITE \&\& !Finish(v)) return false;}
\Code{1.3}{\}}
\Code{1.3}{color[u] = BLACK;}
\Code{1.3}{Append(\&order, u);}
\Code{1.3}{return true;}
\Code{0.0}{\}}
\nopagebreak[4]
\Cost{\ensuremath{O\!\left(V+E\right)}}{\ensuremath{O\!\left(V\right)}}
\par\medskip

\Needspace{18\baselineskip}
\subsubsection{判断拓扑序列是否合法}
\Code{0.0}{ValidTopo(G, seq) \{}
\Hint{1.3}{先排除重复或缺失的顶点，再检查每条弧的先后关系。}\nopagebreak[4]
\Code{1.3}{if (len(seq) != V) return false;}
\Code{1.3}{Fill(pos, -1);}
\Code{1.3}{for (i = 0; i<V; i++) \{}
\Code{2.6}{u = seq[i];}
\Code{2.6}{if (!ValidVertex(u) || pos[u] != -1) return false;}
\Hint{2.6}{只检查相邻两个序列元素不足够，所有边的起点都必须排在终点前。}\nopagebreak[4]
\Code{2.6}{pos[u] = i;}
\Code{1.3}{\}}
\Code{1.3}{for (e in arcs) if (pos[e.u] >= pos[e.v]) return false;}
\Code{1.3}{return true;}
\Code{0.0}{\}}
\nopagebreak[4]
\Cost{\ensuremath{O\!\left(V+E\right)}}{\ensuremath{O\!\left(V\right)}}
\par\medskip

\Needspace{21\baselineskip}
\subsubsection{判断拓扑序列是否唯一}
\Code{0.0}{UniqueTopo(G) \{}
\Hint{1.3}{任何一步同时有两个零入度点，都可以交换优先次序。}\nopagebreak[4]
\Code{1.3}{degree = CopyIndegree(G);}
\Code{1.3}{Q = ZeroDegreeQueue(degree);}
\Code{1.3}{unique = true;}
\Code{1.3}{k = 0;}
\Code{1.3}{while (!Empty(\&Q)) \{}
\Code{2.6}{if (Size(\&Q)>1) unique = false;}
\Hint{2.6}{先验是否成功取满 V 个点，再讨论唯一性；有环时根本没有拓扑序。}\nopagebreak[4]
\Code{2.6}{u = DeQueue(\&Q);}
\Code{2.6}{k++;}
\Code{2.6}{for (v in adj[u]) if (--degree[v] == 0) EnQueue(\&Q, v);}
\Code{1.3}{\}}
\Code{1.3}{if (k<V) return NO\_ORDER;}
\Code{1.3}{return unique;}
\Code{0.0}{\}}
\nopagebreak[4]
\Cost{\ensuremath{O\!\left(V+E\right)}}{\ensuremath{O\!\left(V\right)}}
\par\medskip

\Needspace{18\baselineskip}
\subsubsection{利用拓扑排序判断有向图是否存在环}
\Code{0.0}{TopoCycle(G) \{}
\Hint{1.3}{Kahn 无法取出全部顶点时，剩余子图含环。}\nopagebreak[4]
\Code{1.3}{degree = CopyIndegree(G);}
\Code{1.3}{Q = ZeroDegreeQueue(degree);}
\Code{1.3}{count = 0;}
\Code{1.3}{while (!Empty(\&Q)) \{}
\Code{2.6}{u = DeQueue(\&Q);}
\Hint{2.6}{环上的顶点各保留至少一条未删除入边，所以算法会停在取完全部点之前。}\nopagebreak[4]
\Code{2.6}{count++;}
\Code{2.6}{for (v in adj[u]) if (--degree[v] == 0) EnQueue(\&Q, v);}
\Code{1.3}{\}}
\Code{1.3}{return count<V;}
\Code{0.0}{\}}
\nopagebreak[4]
\Cost{\ensuremath{O\!\left(V+E\right)}}{\ensuremath{O\!\left(V\right)}}
\par\medskip

\Needspace{22\baselineskip}
\subsubsection{求所有合法的拓扑序列}
\Code{0.0}{AllTopo(k) \{}
\Hint{1.3}{试探一个零入度点后，退出该分支时必须恢复所有后继的入度。}\nopagebreak[4]
\Code{1.3}{if (k == V) \{}
\Code{2.6}{Emit(path, V);}
\Code{2.6}{return;}
\Code{1.3}{\}}
\Code{1.3}{for (u in vertices) \{}
\Code{2.6}{if (used[u] || degree[u] != 0) continue;}
\Code{2.6}{used[u] = true;}
\Hint{2.6}{每次递归对应路径里已经排好的 k 个点，同一层尝试的是下一个位置的替代选择。}\nopagebreak[4]
\Code{2.6}{path[k] = u;}
\Code{2.6}{for (v in adj[u]) degree[v]--;}
\Code{2.6}{AllTopo(k+1);}
\Code{2.6}{for (v in adj[u]) degree[v]++;}
\Code{2.6}{used[u] = false;}
\Code{1.3}{\}}
\Code{0.0}{\}}
\nopagebreak[4]
\Cost{\ensuremath{O\!\left((V+E+1)\,V!\right)}}{\ensuremath{O\!\left(V\right)}}
\par\medskip

\Needspace{21\baselineskip}
\subsubsection{基于拓扑排序求有向无环图的最长路径}
\Code{0.0}{DAGLongest(G, s) \{}
\Hint{1.3}{不可达用负无穷，松弛时取较大值。}\nopagebreak[4]
\Code{1.3}{order = Topo(G);}
\Code{1.3}{require(order != FAIL);}
\Code{1.3}{Fill(d, -INF);}
\Code{1.3}{d[s] = 0;}
\Code{1.3}{for (u in order) \{}
\Code{2.6}{if (d[u] == -INF) continue;}
\Hint{2.6}{若想允许路径从任意顶点开始并允许零长路径，可以把每点初值都设为零。}\nopagebreak[4]
\Code{2.6}{for (e in adj[u]) if (d[u]+e.w>d[e.to]) \{}
\Code{3.9}{d[e.to] = d[u]+e.w;}
\Code{3.9}{prev[e.to] = u;}
\Code{2.6}{\}}
\Code{1.3}{\}}
\Code{1.3}{return d;}
\Code{0.0}{\}}
\nopagebreak[4]
\Cost{\ensuremath{O\!\left(V+E\right)}}{\ensuremath{O\!\left(V\right)}}
\par\medskip

\Group{关键路径}

\Needspace{22\baselineskip}
\subsubsection{AOE网的建立}
\Code{0.0}{BuildAOE(events, activities) \{}
\Hint{1.3}{活动是带工期的弧，超级源汇用零工期边连接多个起止事件。}\nopagebreak[4]
\Code{1.3}{G = EmptyGraph(len(events));}
\Code{1.3}{for ((u, v, w) in activities) \{}
\Code{2.6}{require(w >= 0);}
\Code{2.6}{AddArc(\&G, u, v, w);}
\Code{1.3}{\}}
\Code{1.3}{sources = ZeroInVertices(G);}
\Code{1.3}{sinks = ZeroOutVertices(G);}
\Hint{1.3}{添加超级源汇前先保存原有源点和汇点，避免将新增顶点自身也连入。}\nopagebreak[4]
\Code{1.3}{s = NewVertex(\&G);}
\Code{1.3}{t = NewVertex(\&G);}
\Code{1.3}{for (u in sources) AddArc(\&G, s, u, 0);}
\Code{1.3}{for (u in sinks) AddArc(\&G, u, t, 0);}
\Code{1.3}{require(Topo(G) != FAIL);}
\Code{1.3}{return (G, s, t);}
\Code{0.0}{\}}
\nopagebreak[4]
\Cost{\ensuremath{O\!\left(V+E\right)}}{\ensuremath{O\!\left(V\right)}}
\par\medskip

\Needspace{15\baselineskip}
\subsubsection{求事件的最早发生时间}
\Code{0.0}{Earliest(G, s) \{}
\Hint{1.3}{一个事件必须等待所有前置活动完成，因此取最大到达时间。}\nopagebreak[4]
\Code{1.3}{Fill(ve, -INF);}
\Code{1.3}{ve[s] = 0;}
\Code{1.3}{for (u in Topo(G)) \{}
\Hint{2.6}{按拓扑序计算确保更新 ve[u] 的所有活动已经先于 u 处理。}\nopagebreak[4]
\Code{2.6}{for (e in adj[u]) ve[e.to] = max(ve[e.to], ve[u]+e.w);}
\Code{1.3}{\}}
\Code{1.3}{return ve;}
\Code{0.0}{\}}
\nopagebreak[4]
\Cost{\ensuremath{O\!\left(V+E\right)}}{\ensuremath{O\!\left(V\right)}}
\par\medskip

\Needspace{15\baselineskip}
\subsubsection{求事件的最迟发生时间}
\Code{0.0}{Latest(G, t, ve) \{}
\Hint{1.3}{从截止工期倒推，取所有后继给出的最紧时限。}\nopagebreak[4]
\Code{1.3}{Fill(vl, INF);}
\Code{1.3}{vl[t] = ve[t];}
\Code{1.3}{for (u in Reverse(Topo(G))) \{}
\Hint{2.6}{逆拓扑保证算 u 时，所有后继 v 的最迟时刻 vl[v] 已经确定。}\nopagebreak[4]
\Code{2.6}{for (e in adj[u]) vl[u] = min(vl[u], vl[e.to]-e.w);}
\Code{1.3}{\}}
\Code{1.3}{return vl;}
\Code{0.0}{\}}
\nopagebreak[4]
\Cost{\ensuremath{O\!\left(V+E\right)}}{\ensuremath{O\!\left(V\right)}}
\par\medskip

\Needspace{13\baselineskip}
\subsubsection{求活动的最早开始时间}
\Code{0.0}{EarlyStart(arcs, ve) \{}
\Hint{1.3}{弧的起点事件发生后活动即可开始。}\nopagebreak[4]
\Code{1.3}{for (e in arcs) \{}
\Code{2.6}{start = ve[e.u];}
\Hint{2.6}{同一起点事件发出的活动，共用这个最早开始时刻。}\nopagebreak[4]
\Code{2.6}{Emit(e.id, start);}
\Code{1.3}{\}}
\Code{0.0}{\}}
\nopagebreak[4]
\Cost{\ensuremath{O\!\left(E\right)}（已有 ve）}{\ensuremath{O\!\left(1\right)}}
\par\medskip

\Needspace{13\baselineskip}
\subsubsection{求活动的最迟开始时间}
\Code{0.0}{LateStart(arcs, vl) \{}
\Hint{1.3}{从终点最迟时间减去活动本身工期。}\nopagebreak[4]
\Code{1.3}{for (e in arcs) \{}
\Code{2.6}{start = vl[e.v]-e.w;}
\Hint{2.6}{两个活动即使终点相同，也可能因为工期不同而有不同的最迟开始时刻。}\nopagebreak[4]
\Code{2.6}{Emit(e.id, start);}
\Code{1.3}{\}}
\Code{0.0}{\}}
\nopagebreak[4]
\Cost{\ensuremath{O\!\left(E\right)}（已有 vl）}{\ensuremath{O\!\left(1\right)}}
\par\medskip

\Needspace{14\baselineskip}
\subsubsection{求活动的时间余量}
\Code{0.0}{Slack(arcs, ve, vl) \{}
\Hint{1.3}{最迟开工减最早开工，是不拖延总工期所允许的延迟。}\nopagebreak[4]
\Code{1.3}{for (e in arcs) \{}
\Code{2.6}{early = ve[e.u];}
\Hint{2.6}{这项时间余量是总时差，计算中最早时间用起点、最迟时间用终点。}\nopagebreak[4]
\Code{2.6}{late = vl[e.v]-e.w;}
\Code{2.6}{Emit(e.id, late-early);}
\Code{1.3}{\}}
\Code{0.0}{\}}
\nopagebreak[4]
\Cost{\ensuremath{O\!\left(E\right)}（已有 ve,vl）}{\ensuremath{O\!\left(1\right)}}
\par\medskip

\Needspace{13\baselineskip}
\subsubsection{确定关键活动}
\Code{0.0}{CriticalEdges(arcs, ve, vl) \{}
\Hint{1.3}{零余量表示该活动无法延后。}\nopagebreak[4]
\Code{1.3}{for (e in arcs) \{}
\Code{2.6}{late = vl[e.v]-e.w;}
\Hint{2.6}{关键弧的起点和终点都在至少一条关键路径上。}\nopagebreak[4]
\Code{2.6}{if (late == ve[e.u]) Emit(e.id);}
\Code{1.3}{\}}
\Code{0.0}{\}}
\nopagebreak[4]
\Cost{\ensuremath{O\!\left(E\right)}（已有 ve,vl）}{\ensuremath{O\!\left(1\right)}}
\par\medskip

\Needspace{21\baselineskip}
\subsubsection{求关键路径}
\Code{0.0}{CriticalPath(G, s, t) \{}
\Hint{1.3}{统一源汇后，从源出发只沿零余量弧选择一支，直至汇点。}\nopagebreak[4]
\Code{1.3}{ve = Earliest(G, s);}
\Code{1.3}{vl = Latest(G, t, ve);}
\Code{1.3}{u = s;}
\Code{1.3}{while (u != t) \{}
\Code{2.6}{for (e in adj[u]) \{}
\Code{3.9}{if (ve[u] == vl[e.to]-e.w) \{}
\Hint{5.2}{同一位置可能有多个可选关键活动，选择其中一条只给出一条关键路径。}\nopagebreak[4]
\Code{5.2}{Emit(e.id);}
\Code{5.2}{u = e.to;}
\Code{5.2}{break;}
\Code{3.9}{\}}
\Code{2.6}{\}}
\Code{1.3}{\}}
\Code{0.0}{\}}
\nopagebreak[4]
\Cost{\ensuremath{O\!\left(V+E\right)}}{\ensuremath{O\!\left(V\right)}}
\par\medskip

\Needspace{11\baselineskip}
\subsubsection{求工程的最短完成时间}
\Code{0.0}{ProjectTime(G, s, t) \{}
\Hint{1.3}{工期由源到汇的最长工期链决定，假设互不依赖的活动可以同时进行。}\nopagebreak[4]
\Code{1.3}{ve = Earliest(G, s);}
\Hint{1.3}{超级汇的最早发生时间已考虑所有结束事件，因此是整个工程的工期。}\nopagebreak[4]
\Code{1.3}{return ve[t];}
\Code{0.0}{\}}
\nopagebreak[4]
\Cost{\ensuremath{O\!\left(V+E\right)}}{\ensuremath{O\!\left(V\right)}}
\par\medskip

\Needspace{24\baselineskip}
\subsubsection{多条关键路径的查找}
\Code{0.0}{AllCritical(G, s, t) \{}
\Hint{1.3}{先建只保留关键弧的邻接表，再回溯全部源汇路径。}\nopagebreak[4]
\Code{1.3}{ve = Earliest(G, s);}
\Code{1.3}{vl = Latest(G, t, ve);}
\Code{1.3}{for (e in arcs) if (ve[e.u] == vl[e.v]-e.w) Append(\&critical[e.u], e);}
\Code{1.3}{Walk(s);}
\Code{0.0}{\}}
\Code{0.0}{Walk(u) \{}
\Code{1.3}{if (u == t) \{}
\Code{2.6}{Emit(path);}
\Hint{2.6}{Emit 应复制或立即打印 path，回溯会复用同一段路径存储。}\nopagebreak[4]
\Code{2.6}{return;}
\Code{1.3}{\}}
\Code{1.3}{for (e in critical[u]) \{}
\Code{2.6}{Push(\&path, e.id);}
\Code{2.6}{Walk(e.to);}
\Code{2.6}{Pop(\&path);}
\Code{1.3}{\}}
\Code{0.0}{\}}
\nopagebreak[4]
\Cost{\ensuremath{O\!\left(V+E+L\right)}（建关键子图后）}{\ensuremath{O\!\left(V+E\right)}}
\par\medskip

\Group{图的经典综合问题}

\Needspace{24\baselineskip}
\subsubsection{欧拉路径与欧拉回路的判断}
\Code{0.0}{EulerType(G) \{}
\Hint{1.3}{先检验非孤立点弱连通。开放有向通路的首尾分别多一条出边和入边。}\nopagebreak[4]
\Code{1.3}{if (!WeakConnectedNonisolates(G)) return NONE;}
\Code{1.3}{if (!G.directed) \{}
\Code{2.6}{odd = 0;}
\Code{2.6}{for (u in vertices) odd += degree[u]\%2;}
\Code{2.6}{return odd == 0?CIRCUIT:(odd == 2?OPEN:NONE);}
\Code{1.3}{\}}
\Code{1.3}{plus = minus = 0;}
\Code{1.3}{for (u in vertices) \{}
\Hint{2.6}{有边的图至少选一个非孤立点开始检验连通；孤立点不影响走遍边。}\nopagebreak[4]
\Code{2.6}{d = out[u]-in[u];}
\Code{2.6}{if (d == 1) plus++;}
\Code{2.6}{else if (d == -1) minus++;}
\Code{2.6}{else if (d != 0) return NONE;}
\Code{1.3}{\}}
\Code{1.3}{if (plus == 0 \&\& minus == 0) return CIRCUIT;}
\Code{1.3}{return plus == 1 \&\& minus == 1 ? OPEN:NONE;}
\Code{0.0}{\}}
\nopagebreak[4]
\Cost{\ensuremath{O\!\left(V+E\right)}}{\ensuremath{O\!\left(V+E\right)}}
\par\medskip

\Needspace{23\baselineskip}
\subsubsection{欧拉路径与欧拉回路的构造}
\Code{0.0}{EulerWalk(G, s) \{}
\Hint{1.3}{起点按度数条件选。顶点已无未用边才弹出，因此最后要逆序输出。}\nopagebreak[4]
\Code{1.3}{require(EulerType(G) != NONE);}
\Code{1.3}{InitEdgeCursors(G);}
\Code{1.3}{Push(\&S, s);}
\Code{1.3}{while (!Empty(\&S)) \{}
\Code{2.6}{u = Top(\&S);}
\Code{2.6}{e = NextUnusedEdge(u);}
\Code{2.6}{if (e == NONE) Append(\&out, Pop(\&S));}
\Code{2.6}{else \{}
\Hint{3.9}{边使用后以编号标记，无向图的两个邻接记录要共享同一标志。}\nopagebreak[4]
\Code{3.9}{used[e.id] = true;}
\Code{3.9}{Push(\&S, OtherEnd(e, u));}
\Code{2.6}{\}}
\Code{1.3}{\}}
\Code{1.3}{require(len(out) == E+1);}
\Code{1.3}{return Reverse(out);}
\Code{0.0}{\}}
\nopagebreak[4]
\Cost{\ensuremath{O\!\left(V+E\right)}}{\ensuremath{O\!\left(V+E\right)}}
\par\medskip

\Needspace{29\baselineskip}
\subsubsection{哈密顿路径与哈密顿回路（回溯法）}
\Code{0.0}{Hamilton(u, k, cycle) \{}
\Hint{1.3}{path[0] 是起点。回路模式在覆盖全图后还要验证回到起点的边。}\nopagebreak[4]
\Code{1.3}{if (k == V) return !cycle || Edge(u, path[0]);}
\Code{1.3}{for (v in vertices) \{}
\Code{2.6}{if (used[v] || !Edge(u, v)) continue;}
\Code{2.6}{used[v] = true;}
\Code{2.6}{path[k] = v;}
\Code{2.6}{if (Hamilton(v, k+1, cycle)) return true;}
\Code{2.6}{used[v] = false;}
\Code{1.3}{\}}
\Code{1.3}{return false;}
\Code{0.0}{\}}
\Code{0.0}{HamiltonSearch(cycle) \{}
\Hint{1.3}{首次调用设 path[0]=s、used[s]=true、k=1。找开放路径时尝试各起点，找回路可固定起点。}\nopagebreak[4]
\Code{1.3}{Fill(used, false);}
\Code{1.3}{for (s in vertices) \{}
\Code{2.6}{path[0] = s;}
\Code{2.6}{used[s] = true;}
\Code{2.6}{if (Hamilton(s, 1, cycle)) return true;}
\Code{2.6}{used[s] = false;}
\Code{2.6}{if (cycle) break;}
\Code{1.3}{\}}
\Code{1.3}{return false;}
\Code{0.0}{\}}
\nopagebreak[4]
\Cost{\ensuremath{O\!\left(V\,V!\right)}（邻接矩阵）}{\ensuremath{O\!\left(V\right)}}
\par\medskip

\Needspace{19\baselineskip}
\subsubsection{图的着色问题}
\Code{0.0}{ColorGraph(i, k) \{}
\Hint{1.3}{只尝试邻点尚未使用的颜色，失败后清除本点颜色。自环直接判无解。}\nopagebreak[4]
\Code{1.3}{if (i == V) return true;}
\Code{1.3}{for (c = 1; c <= k; c++) \{}
\Code{2.6}{ok = true;}
\Code{2.6}{for (v in adj[i]) if (v == i || color[v] == c) ok = false;}
\Code{2.6}{if (!ok) continue;}
\Hint{2.6}{某一色不合适只跳过该色；所有颜色均试过才判这个分支失败。}\nopagebreak[4]
\Code{2.6}{color[i] = c;}
\Code{2.6}{if (ColorGraph(i+1, k)) return true;}
\Code{2.6}{color[i] = 0;}
\Code{1.3}{\}}
\Code{1.3}{return false;}
\Code{0.0}{\}}
\nopagebreak[4]
\Cost{\ensuremath{O\!\left(k^V\,(V+E+1)\right)}}{\ensuremath{O\!\left(V\right)}}
\par\medskip

\Needspace{21\baselineskip}
\subsubsection{迷宫路径搜索}
\Code{0.0}{GridPath(s, t) \{}
\Hint{1.3}{沿深度优先找到一条路径，pre 保存每格首次发现时的来路。}\nopagebreak[4]
\Code{1.3}{Fill(pre, NONE);}
\Code{1.3}{pre[s] = s;}
\Code{1.3}{Push(\&S, s);}
\Code{1.3}{while (!Empty(\&S)) \{}
\Code{2.6}{u = Pop(\&S);}
\Code{2.6}{if (u == t) return RestoreGrid(pre, s, t);}
\Hint{2.6}{pre 同时充当访问标记，记录的始终是首次抵达该格的来路。}\nopagebreak[4]
\Code{2.6}{for (v in Neighbors4(u)) if (Open(v) \&\& pre[v] == NONE) \{}
\Code{3.9}{pre[v] = u;}
\Code{3.9}{Push(\&S, v);}
\Code{2.6}{\}}
\Code{1.3}{\}}
\Code{1.3}{return FAIL;}
\Code{0.0}{\}}
\nopagebreak[4]
\Cost{\ensuremath{O\!\left(RC\right)}}{\ensuremath{O\!\left(RC\right)}}
\par\medskip

\Needspace{21\baselineskip}
\subsubsection{迷宫最短路径搜索}
\Code{0.0}{GridShortest(s, t) \{}
\Hint{1.3}{换成先进先出队列，先到的一定来自较短层。每步代价为 1。}\nopagebreak[4]
\Code{1.3}{Fill(pre, NONE);}
\Code{1.3}{pre[s] = s;}
\Code{1.3}{EnQueue(\&Q, s);}
\Code{1.3}{while (!Empty(\&Q)) \{}
\Code{2.6}{u = DeQueue(\&Q);}
\Code{2.6}{if (u == t) return RestoreGrid(pre, s, t);}
\Hint{2.6}{只在上下左右移动且代价相等时，按边数最少才等同于代价最小。}\nopagebreak[4]
\Code{2.6}{for (v in Neighbors4(u)) if (Open(v) \&\& pre[v] == NONE) \{}
\Code{3.9}{pre[v] = u;}
\Code{3.9}{EnQueue(\&Q, v);}
\Code{2.6}{\}}
\Code{1.3}{\}}
\Code{1.3}{return FAIL;}
\Code{0.0}{\}}
\nopagebreak[4]
\Cost{\ensuremath{O\!\left(RC\right)}}{\ensuremath{O\!\left(RC\right)}}
\par\medskip

\Needspace{25\baselineskip}
\subsubsection{网格图的连通区域统计（Flood Fill）}
\Code{0.0}{FloodFill(grid) \{}
\Hint{1.3}{一次扩展标记整个连通块，因此只在发现新起点时给计数加一。}\nopagebreak[4]
\Code{1.3}{Zero(seen);}
\Code{1.3}{count = 0;}
\Code{1.3}{for (s in cells) \{}
\Code{2.6}{if (!Target(s) || seen[s]) continue;}
\Code{2.6}{count++;}
\Code{2.6}{seen[s] = true;}
\Code{2.6}{Push(\&S, s);}
\Code{2.6}{while (!Empty(\&S)) \{}
\Code{3.9}{u = Pop(\&S);}
\Hint{3.9}{同一连通块后来再次被扫描到时已标记，不会把一块区域重复计数。}\nopagebreak[4]
\Code{3.9}{for (v in Neighbors4(u)) if (Inside(v) \&\& Target(v) \&\& !seen[v]) \{}
\Code{5.2}{seen[v] = true;}
\Code{5.2}{Push(\&S, v);}
\Code{3.9}{\}}
\Code{2.6}{\}}
\Code{1.3}{\}}
\Code{1.3}{return count;}
\Code{0.0}{\}}
\nopagebreak[4]
\Cost{\ensuremath{O\!\left(RC\right)}}{\ensuremath{O\!\left(RC\right)}}
\par\medskip

\Needspace{22\baselineskip}
\subsubsection{多源广度优先搜索}
\Code{0.0}{MultiBFS(G, sources) \{}
\Hint{1.3}{把全部源点同时放在第 0 层，避免对每个源单独搜索。}\nopagebreak[4]
\Code{1.3}{Fill(d, INF);}
\Code{1.3}{for (s in sources) if (d[s] == INF) \{}
\Code{2.6}{d[s] = 0;}
\Code{2.6}{EnQueue(\&Q, s);}
\Code{1.3}{\}}
\Code{1.3}{while (!Empty(\&Q)) \{}
\Code{2.6}{u = DeQueue(\&Q);}
\Hint{2.6}{一个点若与多个源等距，距离仍然唯一；需要具体最近源时再额外记录来源。}\nopagebreak[4]
\Code{2.6}{for (v in adj[u]) if (d[v] == INF) \{}
\Code{3.9}{d[v] = d[u]+1;}
\Code{3.9}{EnQueue(\&Q, v);}
\Code{2.6}{\}}
\Code{1.3}{\}}
\Code{1.3}{return d;}
\Code{0.0}{\}}
\nopagebreak[4]
\Cost{\ensuremath{O\!\left(V+E\right)}}{\ensuremath{O\!\left(V\right)}}
\par\medskip

\Needspace{25\baselineskip}
\subsubsection{双向广度优先搜索}
\Code{0.0}{BiBFS(G, s, t) \{}
\Hint{1.3}{两端每轮扩展完整一层；反向搜索使用反向边，取相遇候选的最短值。}\nopagebreak[4]
\Code{1.3}{if (s == t) return 0;}
\Code{1.3}{F = StartLayer(s, G);}
\Code{1.3}{B = StartLayer(t, TransposeGraph(G));}
\Code{1.3}{best = INF;}
\Code{1.3}{while (!Empty(F) \&\& !Empty(B)) \{}
\Code{2.6}{if (Depth(F)+Depth(B) >= best) break;}
\Code{2.6}{(a, b) = SmallerFrontFirst(\&F, \&B);}
\Code{2.6}{for (u in CurrentLayer(a)) \{}
\Code{3.9}{for (v in Neighbors(a, u)) \{}
\Hint{5.2}{AdvanceLayer 只在本轮整层完成后交换前沿；反向图不能继续使用正向邻居。}\nopagebreak[4]
\Code{5.2}{if (!Visited(a, v)) DiscoverNext(a, v, Dist(a, u)+1);}
\Code{5.2}{if (Visited(b, v)) best = min(best, Dist(a, v)+Dist(b, v));}
\Code{3.9}{\}}
\Code{2.6}{\}}
\Code{2.6}{AdvanceLayer(a);}
\Code{1.3}{\}}
\Code{1.3}{return best;}
\Code{0.0}{\}}
\nopagebreak[4]
\Cost{\ensuremath{O\!\left(V+E\right)}}{\ensuremath{O\!\left(V+E\right)}}
\par\medskip

\Needspace{26\baselineskip}
\subsubsection{启发式路径搜索（A算法）}
\Code{0.0}{AStar(G, s, t, h) \{}
\Hint{1.3}{堆键是已走代价 g 加尚需代价的估计 h。非负边权且 h 一致、h(t)=0。}\nopagebreak[4]
\Code{1.3}{Fill(g, INF);}
\Code{1.3}{g[s] = 0;}
\Code{1.3}{Insert(\&H, s, h(s));}
\Code{1.3}{while (!Empty(\&H)) \{}
\Code{2.6}{u = ExtractMin(\&H);}
\Code{2.6}{if (u == t) return g[u];}
\Code{2.6}{for (e in adj[u]) \{}
\Code{3.9}{v = e.to;}
\Code{3.9}{cand = g[u]+e.w;}
\Hint{3.9}{仅当目标从堆中取出时才返回，第一次发现目标时其距离仍可能被改进。}\nopagebreak[4]
\Code{3.9}{if (cand<g[v]) \{}
\Code{5.2}{g[v] = cand;}
\Code{5.2}{prev[v] = u;}
\Code{5.2}{InsertOrDecrease(\&H, v, cand+h(v));}
\Code{3.9}{\}}
\Code{2.6}{\}}
\Code{1.3}{\}}
\Code{1.3}{return INF;}
\Code{0.0}{\}}
\nopagebreak[4]
\Cost{\ensuremath{O\!\left((V+E) \log (V+1)\right)}（一致启发式）}{\ensuremath{O\!\left(V\right)}}
\par\medskip

\Needspace{28\baselineskip}
\subsubsection{二分图的最大匹配（匈牙利算法）}
\Code{0.0}{BipartiteMatch(G) \{}
\Hint{1.3}{右点已有配偶时，尝试为旧配偶另找一个位置。轮次标记避免同一次增广重复尝试。}\nopagebreak[4]
\Code{1.3}{Fill(match, NONE);}
\Code{1.3}{Zero(mark);}
\Code{1.3}{round = count = 0;}
\Code{1.3}{for (u in left) \{}
\Code{2.6}{round++;}
\Code{2.6}{if (Aug(u)) count++;}
\Code{1.3}{\}}
\Code{1.3}{return count;}
\Code{0.0}{\}}
\Code{0.0}{Aug(u) \{}
\Hint{1.3}{递归重配成功时，旧左点已经换到另一个右点，因此当前右点可以交给 u。}\nopagebreak[4]
\Code{1.3}{for (v in adj[u]) \{}
\Code{2.6}{if (mark[v] == round) continue;}
\Code{2.6}{mark[v] = round;}
\Code{2.6}{if (match[v] == NONE || Aug(match[v])) \{}
\Code{3.9}{match[v] = u;}
\Code{3.9}{return true;}
\Code{2.6}{\}}
\Code{1.3}{\}}
\Code{1.3}{return false;}
\Code{0.0}{\}}
\nopagebreak[4]
\Cost{\ensuremath{O\!\left(V\,E+V\right)}}{\ensuremath{O\!\left(V\right)}}
\par\medskip

\Needspace{8\baselineskip}
\subsubsection{网络最大流（Ford-Fulkerson算法）}
\Code{0.0}{FordFulkerson(G, s, t) \{}
\Hint{1.3}{容量为非负整数，源点与汇点不同。前驱记录残量边，以便沿配对边更新反向残量。}\nopagebreak[4]
\Code{1.3}{R = BuildResidualPairs(G);}
\Code{1.3}{flow = 0;}
\Code{1.3}{while (ResidualDFS(R, s, t, pre)) \{}
\Code{2.6}{d = INF;}
\Code{2.6}{for (v = t; v != s; v = pre[v].from) d = min(d, pre[v].cap);}
\Code{2.6}{for (v = t; v != s; v = pre[v].from) \{}
\Code{3.9}{e = pre[v];}
\Code{3.9}{e.cap -= d;}
\Code{3.9}{e.rev.cap += d;}
\Code{2.6}{\}}
\Code{2.6}{flow += d;}
\Code{1.3}{\}}
\Code{1.3}{return flow;}
\Code{0.0}{\}}
\Code{0.0}{ResidualDFS(R, s, t, pre) \{}
\Code{1.3}{Fill(seen, false);}
\Hint{1.3}{反向边的残量增加表示后续允许撤销以前的流量，不能只更新正向边。}\nopagebreak[4]
\Code{1.3}{Push(\&S, s);}
\Code{1.3}{seen[s] = true;}
\Code{1.3}{while (!Empty(\&S)) \{}
\Code{2.6}{u = Pop(\&S);}
\Code{2.6}{if (u == t) return true;}
\Code{2.6}{for (e in R[u]) if (e.cap>0 \&\& !seen[e.to]) \{}
\Code{3.9}{seen[e.to] = true;}
\Code{3.9}{pre[e.to] = e;}
\Code{3.9}{Push(\&S, e.to);}
\Code{2.6}{\}}
\Code{1.3}{\}}
\Code{1.3}{return false;}
\Code{0.0}{\}}
\nopagebreak[4]
\Cost{\ensuremath{O\!\left((V+E)(F+1)\right)}（$F$ 为最大流值，含最后一次搜索）}{\ensuremath{O\!\left(V+E\right)}}
\par\medskip

\Needspace{30\baselineskip}
\subsubsection{网络最大流（Edmonds-Karp算法）}
\Code{0.0}{EdmondsKarp(G, s, t) \{}
\Hint{1.3}{按残量边做 BFS 选择边数最少的增广路；每轮重建前驱。}\nopagebreak[4]
\Code{1.3}{R = BuildResidualPairs(G);}
\Code{1.3}{flow = 0;}
\Code{1.3}{while (true) \{}
\Code{2.6}{Fill(pre, NONE);}
\Code{2.6}{seen = ZeroArray(V);}
\Code{2.6}{seen[s] = true;}
\Code{2.6}{EnQueue(\&Q, s);}
\Code{2.6}{while (!Empty(\&Q)) \{}
\Code{3.9}{u = DeQueue(\&Q);}
\Code{3.9}{for (e in R[u]) if (e.cap>0 \&\& !seen[e.to]) \{}
\Code{5.2}{seen[e.to] = true;}
\Code{5.2}{pre[e.to] = e;}
\Hint{5.2}{Bottleneck 沿 pre 取最小残量，AugmentPairs 对每条路径边扣 d，并给配对反向边加 d。}\nopagebreak[4]
\Code{5.2}{EnQueue(\&Q, e.to);}
\Code{3.9}{\}}
\Code{2.6}{\}}
\Code{2.6}{if (!seen[t]) break;}
\Code{2.6}{d = Bottleneck(pre, s, t);}
\Code{2.6}{AugmentPairs(pre, s, t, d);}
\Code{2.6}{flow += d;}
\Code{1.3}{\}}
\Code{1.3}{return flow;}
\Code{0.0}{\}}
\nopagebreak[4]
\Cost{\ensuremath{O\!\left(VE^{2}+V\right)}}{\ensuremath{O\!\left(V+E\right)}}
\par\medskip

\Needspace{8\baselineskip}
\subsubsection{网络最大流（Dinic算法）}
\Code{0.0}{Dinic(G, s, t) \{}
\Hint{1.3}{只沿下一层送流。当前边尚可送流时保留游标，失效后才前进。}\nopagebreak[4]
\Code{1.3}{R = BuildResidualPairs(G);}
\Code{1.3}{flow = 0;}
\Code{1.3}{while (LevelBFS(R, s, t, level)) \{}
\Code{2.6}{ResetCursors(cur, R);}
\Code{2.6}{while ((d = Send(s, INF))>0) flow += d;}
\Code{1.3}{\}}
\Code{1.3}{return flow;}
\Code{0.0}{\}}
\Code{0.0}{Send(u, f) \{}
\Code{1.3}{if (u == t) return f;}
\Code{1.3}{while (cur[u] != END) \{}
\Code{2.6}{e = CurrentEdge(cur[u]);}
\Code{2.6}{if (e.cap>0 \&\& level[e.to] == level[u]+1) \{}
\Code{3.9}{d = Send(e.to, min(f, e.cap));}
\Code{3.9}{if (d>0) \{}
\Code{5.2}{e.cap -= d;}
\Code{5.2}{e.rev.cap += d;}
\Code{5.2}{return d;}
\Code{3.9}{\}}
\Code{2.6}{\}}
\Hint{2.6}{Send 的返回值是本次实际送出的流量；返回 0 表示此分层里从该处分支已无法继续。}\nopagebreak[4]
\Hint{2.6}{该边不能再为当前分层送流，游标才向后移动。}\nopagebreak[4]
\Code{2.6}{Advance(\&cur[u]);}
\Code{1.3}{\}}
\Code{1.3}{return 0;}
\Code{0.0}{\}}
\Code{0.0}{LevelBFS(R, s, t, level) \{}
\Code{1.3}{Fill(level, -1);}
\Code{1.3}{level[s] = 0;}
\Code{1.3}{EnQueue(\&Q, s);}
\Code{1.3}{while (!Empty(\&Q)) \{}
\Code{2.6}{u = DeQueue(\&Q);}
\Code{2.6}{for (e in R[u]) if (e.cap>0 \&\& level[e.to]<0) \{}
\Code{3.9}{level[e.to] = level[u]+1;}
\Code{3.9}{EnQueue(\&Q, e.to);}
\Code{2.6}{\}}
\Code{1.3}{\}}
\Code{1.3}{return level[t] >= 0;}
\Code{0.0}{\}}
\nopagebreak[4]
\Cost{\ensuremath{O\!\left(V^{2}E+V\right)}}{\ensuremath{O\!\left(V+E\right)}}
\par\medskip
\end{document}
